\documentclass[10pt,oneside,openany]{book}
\usepackage[T1]{fontenc}
\usepackage[utf8]{inputenc}
\usepackage{babel}
\babelprovide[import,main]{spanish}
\usepackage{lmodern,amsmath,amssymb,amsthm,mathrsfs}
\usepackage{iftex}
\ifLuaTeX\usepackage{luatex85}\fi
\usepackage[all]{xy}
\usepackage{tikz-cd}
\usepackage{xstring,xparse,xcolor,etoolbox}
\usepackage[a4paper,left=18mm,right=18mm,top=22mm,bottom=23mm,marginparwidth=14mm,marginparsep=2mm]{geometry}
\usepackage{array,longtable,multicol,booktabs}
\usepackage[shortlabels]{enumitem}
\usepackage{marginnote,microtype,graphicx,adjustbox}
\usepackage[unicode=true,pdfencoding=auto,colorlinks=true,linkcolor=blue,citecolor=blue,urlcolor=blue,hypertexnames=false,hyperfootnotes=false,linktocpage=true]{hyperref}
\usepackage{bookmark}
\renewcommand{\thesection}{\arabic{section}}
\renewcommand{\thesubsection}{\thesection.\arabic{subsection}}
% Keep PDF destinations unique when the cumulative reader resets table counters.
\renewcommand{\theHtable}{\arabic{part}.\arabic{table}}
\setcounter{tocdepth}{2}
\setlength{\emergencystretch}{12em}
\tolerance=3000
\binoppenalty=200
\relpenalty=200
\hfuzz=0.1pt
\vfuzz=0.1pt
\raggedbottom
\allowdisplaybreaks[3]
\newtheoremstyle{egaenv}{}{}{\normalfont}{}{\bfseries}{.}{ }{\thmnote{(#3)}}
\newtheoremstyle{egastatement}{}{}{\itshape}{}{\bfseries}{. ---}{ }{\thmname{#1}\thmnote{ (#3)}}
\newtheoremstyle{egadescription}{}{}{\normalfont}{}{\bfseries}{. ---}{ }{\thmname{#1}\thmnote{ (#3)}}
\theoremstyle{egaenv}\newtheorem*{env}{---}
\theoremstyle{egastatement}
\newtheorem*{theorem}{Teorema}
\newtheorem*{proposition}{Proposición}
\newtheorem*{lemma}{Lema}
\newtheorem*{corollary}{Corolario}
\theoremstyle{egadescription}
\newtheorem*{definition}{Definición}
\newtheorem*{remark}{Observación}
\newtheorem*{remarks}{Observaciones}
\newtheorem*{example}{Ejemplo}
\newtheorem*{examples}{Ejemplos}
\newtheorem*{scholium}{Escolio}
\newtheorem*{notation}{Notación}
\newenvironment{namedtheorem}[2][]{\par\medskip\noindent\setmanualcurrentlabel{#1}\textbf{#2\ifstrempty{#1}{}{ (#1)}. ---}\ifstrequal{#2}{Observación}{\normalfont}{\ifstrequal{#2}{Observaciones}{\normalfont}{\itshape}}}{\par\medskip}
\newenvironment{namedlemma}[2][]{\par\medskip\noindent\setmanualcurrentlabel{#1}\textbf{#2\ifstrempty{#1}{}{ (#1)}. ---}\itshape}{\par\medskip}
\renewcommand{\proofname}{Demostración}
\renewcommand{\contentsname}{Índice general}
\renewcommand{\bibname}{Bibliografía}
\providecommand{\refname}{Bibliografía}
\newcommand{\oldpage}[2][]{\marginnote{\normalfont\scriptsize\textbf{#1}\,|\,#2}\ignorespaces}
\newcommand{\sectionbreak}{\begin{center}***\end{center}}
\newcommand{\erratum}[2][]{#2\marginnote{\scriptsize\textbf{Err}\textsubscript{#1}}}
\providecommand{\shHom}{\mathcal{H}\!\mathit{om}}
\providecommand{\shProj}{\mathcal{P}\!\mathit{roj}}
% Literal macro definition from the French EGA II notation index.
\providecommand{\egaShProj}{\mathcal{P}\!\mathit{roj}}
\providecommand{\shExt}{\mathcal{E}\!\mathit{xt}}
\providecommand{\shTor}{\mathcal{T}\!\mathit{or}}
\providecommand{\shGr}{\mathcal{G}\!\mathit{r}}
\providecommand{\shDer}{\mathcal{D}\!\mathit{er}}
\providecommand{\shDiff}{\mathcal{D}\!\mathit{iff}}
\providecommand{\shHomcont}{\mathcal{H}\!\mathit{om.cont}}
\providecommand{\shAut}{\mathcal{A}\!\mathit{ut}}
\providecommand{\Frac}{\operatorname{Frac}}
\providecommand{\unsure}[1]{#1}
\providecommand{\completelyunsure}{\textbf{??}}
\providecommand{\og}{\guillemotleft{}}
\providecommand{\fg}{\guillemotright{}}
\makeatletter
\newcommand{\setmanualcurrentlabel}[1]{\def\@currentlabel{#1}}
\newcommand{\EGARefNumber}[1]{\begingroup\StrSubstitute{#1}{-fr}{}[\egarefkey]\StrSubstitute{\egarefkey}{-es}{}[\egarefkey]\StrSubstitute{\egarefkey}{.fr}{}[\egarefkey]\StrBehind{\egarefkey}{.}\endgroup}
\newcommand{\EGAMaybeRefLink}[2]{\@ifundefined{r@#1}{#2}{\hyperref[#1]{#2}}}
\newcommand{\EGAReferenceAlias}[2]{\ifcsdef{r@#2}{\csletcs{r@#1}{r@#2}}{}}
\NewDocumentCommand{\sref}{o m o}{\textup{(\IfNoValueF{#1}{\textbf{#1},~}\EGAMaybeRefLink{#2}{\EGARefNumber{#2}}\IfNoValueF{#3}{,~#3})}}
\newcommand{\eref}[1]{\textup{(\EGAMaybeRefLink{#1}{\EGARefNumber{#1}})}}
\newcommand{\EGAReaderTarget}[2]{\phantomsection\def\@currentlabel{#2}\label{#1}}
\newcommand{\EGAReadLink}[2]{\EGAMaybeRefLink{#1}{#2}}
\newcommand{\EGAUnavailableHyperref}[2]{\EGAMaybeRefLink{#1}{#2}}
\newcommand{\sgaStableTarget}[1]{\hypertarget{#1}{}}
\newcommand{\sgaStableOrOriginalLink}[2]{\EGAMaybeRefLink{#1}{#2}}
\newcommand{\EGAIVFourExternalLabel}[1]{}
\newcommand{\EGAIVFourExternalCitation}[2]{\expandafter\def\csname b@#1\endcsname{#2}}
\makeatother
\newcommand{\EGAResetVolumeCounters}{\setcounter{section}{0}\setcounter{subsection}{0}\setcounter{equation}{0}\setcounter{footnote}{0}\setcounter{figure}{0}\setcounter{table}{0}}
\newcommand{\EGAVolumeTitlePage}[3]{\clearpage\EGAResetVolumeCounters\chapter*{#2\\#3}\addcontentsline{toc}{chapter}{#2: #3}}
\newcommand{\MR}[1]{MR #1}
\newcommand{\MRhref}[2]{\href{#1}{#2}}
\providecommand{\bysame}{\leavevmode\hbox to3em{\hrulefill}\,}
\pagestyle{plain}

% BEGIN authoritative French CC0 preamble-base.tex
% Notation of the EGA editions: the abbreviations the transcription uses for alphabets, symbols, operators and
% arrows. Written for these editions and dedicated to the public domain under CC0 1.0.
\makeatletter
% a command that may or may not exist already is (re)defined in one step
\newcommand{\ega@define}[2]{\providecommand{#1}{}\renewcommand{#1}{#2}}
\newcommand{\ega@operator}[2]{\providecommand{#1}{}\renewcommand{#1}{\operatorname{#2}}}

% alphabets: sheaves, bold letters, categories
\ega@define{\sh}{\mathcal}
\ega@define{\bb}{\mathbf}
\ega@define{\cat}{\mathtt}
\ega@define{\C}{\cat{C}}
\ega@define{\Set}{\cat{Set}}
\ega@define{\CHom}{\cat{Hom}}
\ega@define{\op}{^{\cat{op}}}
\ega@define{\GL}{\bb{GL}}
\ega@define{\bbGamma}{\boldsymbol{\Gamma}}

% letters and symbols
\renewcommand{\leq}{\leqslant}
\renewcommand{\geq}{\geqslant}
\ega@define{\setmin}{-}
\ega@define{\emp}{\varnothing}
\ega@define{\vphi}{\varphi}
\ega@define{\kres}{k}
\ega@define{\rad}{\mathfrak{r}}
\ega@define{\nilrad}{\mathfrak{N}}
\ega@define{\HH}{\mathrm{H}}
\ega@define{\CHH}{\check{\mathrm{H}}}
\ega@define{\RR}{\mathrm{R}}
\ega@define{\LL}{\mathrm{L}}
\ega@define{\red}{\mathrm{red}}
\providecommand{\dual}{}\renewcommand{\dual}[1]{{#1}^{\vee}}
% the tilde of the sheaf associated with a module, written after it: M\supertilde
\ega@define{\supertilde}{{\,\widetilde{\,}\,}}

% operators
\ega@operator{\rg}{rg}
\ega@operator{\gr}{gr}
\ega@operator{\Gr}{Gr}
\ega@operator{\Hom}{Hom}
\ega@operator{\Proj}{Proj}
\ega@operator{\Tor}{Tor}
\ega@operator{\Ext}{Ext}
\ega@operator{\Supp}{Supp}
\ega@operator{\Ker}{Ker}
\ega@operator{\Im}{Im}
\ega@operator{\Coker}{Coker}
\ega@operator{\Spec}{Spec}
\ega@operator{\Spm}{Spm}
\ega@operator{\Spf}{Spf}
\ega@operator{\grad}{grad}
\ega@operator{\dimc}{dimc}
\ega@operator{\codim}{codim}
\ega@operator{\id}{id}
\ega@operator{\Der}{Der}
\ega@operator{\Diff}{Diff}
\ega@operator{\Hyp}{Hyp}
\ega@operator{\Tr}{Tr}

% arrows: long in displays, short in text, scripts and subscripts
\let\ega@shortmapsto\mapsto
\renewcommand{\to}{\mathchoice{\longrightarrow}{\rightarrow}{\rightarrow}{\rightarrow}}
\providecommand{\from}{}\renewcommand{\from}{\mathchoice{\longleftarrow}{\leftarrow}{\leftarrow}{\leftarrow}}
\renewcommand{\mapsto}{\mathchoice{\longmapsto}{\ega@shortmapsto}{\ega@shortmapsto}{\ega@shortmapsto}}
\ega@define{\isoto}{\simeq}
\ega@define{\simto}{\xrightarrow{\sim}}

% wide accents
\renewcommand{\hat}[1]{\widehat{#1}}
\renewcommand{\tilde}[1]{\widetilde{#1}}
\makeatother

% END authoritative preamble-base.tex

\hypersetup{pdftitle={EGA III, segunda parte},pdfauthor={Alexander Grothendieck y Jean Dieudonné}}
\begin{document}
\begin{titlepage}\centering\vspace*{20mm}
{\Huge\bfseries Elementos de geometría algebraica\par}\vspace{14mm}
{\LARGE EGA III, segunda parte\par}\vspace{10mm}
{\large Alexander Grothendieck\\con la colaboración de Jean Dieudonné\par}
\vfill{\Large Edición española actual\par}
\vspace{8mm}{\large 25 de noviembre de 2026\par}\end{titlepage}
\chapter*{Créditos y derechos}
\addcontentsline{toc}{chapter}{Créditos y derechos}
\noindent\textbf{Trabajo de esta ejecución.} Traducción y preparación por GPT-6.0 Astra (OpenAI), modo Pro de ChatGPT, noviembre de 2026. Traducción producida por IA; correcciones producidas por IA. Cotejo francés heredado: Claude Opus 5.5, esfuerzo max, noviembre de 2026.

\subsection*{Créditos previos conservados}

Las traducciones están producidas por IA, y también las correcciones, la edición y la composición. Tradujeron
modelos de OpenAI Codex (GPT-5.6 Sol y GPT-6.1 Sol, esfuerzo Ultra); GPT-6 Astra (esfuerzo Ultra) revisó, corrigió
y montó las ediciones. En la
preparación de algunas secciones de EGA se usaron fuentes inglesas que incorporan la traducción humana del proyecto
ryankeleti/ega; para FGA, la traducción inglesa de Tim Hosgood sirvió de referencia (véase «Fuentes y créditos»).

Los lectores de EGA I (traducción de OpenAI Codex, GPT-6.1 Sol, esfuerzo de razonamiento ultra, agosto–octubre
de 2026) se publicaron el 15 de noviembre de 2026; la composición (nada sobresale del margen más de 2 pt) y la
publicación las hizo Claude Opus 5.5 con esfuerzo high.

Las ediciones corregidas y las traducciones diplomáticas, que siguen el texto impreso, son documentos separados; las
notas matemáticas también.

\subsection*{Fuentes y créditos del repositorio}

\begin{itemize}
\item EGA: Alexander Grothendieck, con la colaboración de Jean Dieudonné, \emph{Éléments de géométrie algébrique},
  \emph{Publications mathématiques de l'IHÉS} 4, 8, 11, 17, 20, 24, 28 y 32 (1960–1967), en
  \href{https://www.numdam.org/journals/PMIHES/}{NUMDAM}.
\item FGA: Alexander Grothendieck, \emph{Fondements de la géométrie algébrique}, exposiciones 149, 182, 190, 195, 212, 221,
  232 y 236 del Seminario Bourbaki, en \href{https://www.numdam.org/}{NUMDAM}.
\item EGA: las fuentes inglesas usadas para preparar EGA III-1 (capítulo 0, §§8–10 y parte del §11; capítulo III, §§1 y
  3), EGA IV-1 (capítulo 0, §14) y EGA IV-4 (§16 y el comienzo de los §§17 y 20) incorporan la traducción humana de
  \href{https://github.com/ryankeleti/ega/tree/3d7afe28c99d1887ccdb85acf2ca1afd93fb1971}{ryankeleti/ega} (versión
  3d7afe2), cuyos colaboradores son Mateo Carmona, Tim Holzschuh, Tim Hosgood, Ryan Keleti, Thiago Solovera \& Nery y
  Wayne Yang (\href{https://github.com/ryankeleti/ega/blob/3d7afe28c99d1887ccdb85acf2ca1afd93fb1971/README.md}{créditos por sección}).
\item FGA: la traducción inglesa de \href{https://github.com/thosgood/fga/tree/9976dd5f8e2ddcb388f0a07f3aa334f120d1019f}{Tim Hosgood}
  (versión 9976dd5), que agradece las correcciones de Steve Hnizdur y el apoyo técnico de Jon Sterling.
\item La traducción española no se atribuye a estos traductores. \href{../../tree/repo/ATTRIBUTION.md}{ATTRIBUTION.md} da las versiones exactas
  de las fuentes inglesas.
\end{itemize}

\subsection*{Fuentes y autoría}

\subsection*{FGA: traducción inglesa de referencia}

La obra original, \emph{Fondements de la géométrie algébrique}, es de Alexandre Grothendieck. La traducción inglesa de \href{https://github.com/thosgood/fga/tree/9976dd5f8e2ddcb388f0a07f3aa334f120d1019f}{Tim Hosgood} es una fuente de referencia de esta edición; su proyecto agradece las correcciones de Steve Hnizdur y el apoyo técnico de Jon Sterling. La versión 9976dd5f8e2ddcb388f0a07f3aa334f120d1019f identifica el texto inglés usado en la comparación; no establece la versión exacta incorporada al principio. La traducción española y las correcciones y la composición posteriores están producidas por IA y no se atribuyen a Tim Hosgood. Esta atribución no constituye permiso ni aval.

\subsection*{EGA: fuentes inglesas}

Las fuentes inglesas utilizadas para preparar las traducciones españolas de EGA incorporan la traducción humana de \href{https://github.com/ryankeleti/ega}{ryankeleti/ega}, versión \href{https://github.com/ryankeleti/ega/tree/3d7afe28c99d1887ccdb85acf2ca1afd93fb1971}{3d7afe28c99d1887ccdb85acf2ca1afd93fb1971}. El proyecto nombra como colaboradores a Mateo Carmona, Tim Holzschuh, Tim Hosgood, Ryan Keleti, Thiago Solovera \& Nery y Wayne Yang; sus \href{https://github.com/ryankeleti/ega/blob/3d7afe28c99d1887ccdb85acf2ca1afd93fb1971/README.md}{créditos por sección} y su \href{https://github.com/ryankeleti/ega/commits/3d7afe28c99d1887ccdb85acf2ca1afd93fb1971}{historial de contribuciones} precisan la autoría. Las secciones afectadas son:

\begin{itemize}
\item EGA III-1: los §§ 8–10 y parte del § 11 del capítulo 0, y los §§ 1 y 3 del capítulo III.
\item EGA IV-1, capítulo 0: el § 14.
\item EGA IV-4: el § 16 y las partes iniciales de los §§ 17 y 20 del capítulo IV.
\end{itemize}

La traducción inglesa y la estructura TeX heredadas conservan esa procedencia. La traducción española y las adiciones y correcciones producidas por IA no se atribuyen a esos colaboradores. La atribución no implica autorización ni aval.

\subsection*{Traducción española}

La traducción al español, las correcciones, la edición y la composición están producidas por IA. Tradujeron modelos de OpenAI Codex (GPT-5.6 Sol y GPT-6.1 Sol, esfuerzo Ultra); GPT-6 Astra (esfuerzo Ultra) revisó, corrigió y montó las ediciones. No se afirma revisión humana. La autoría de las obras originales y de las contribuciones humanas citadas se conserva.

\begin{itemize}
\item \href{../../tree/repo/editions/II\_CURRENT/FUENTES.md}{EGA II: alcance e identidad de las fuentes}.
\item \href{../../tree/repo/editions/III1\_CURRENT/FUENTES.md}{EGA III-1: alcance e identidad de las fuentes}.
\item \href{../../tree/repo/editions/IV1\_CURRENT/FUENTES.md}{EGA IV-1, capítulo 0: alcance e identidad de las fuentes}.
\item \href{../../tree/repo/editions/IV4/FUENTES.md}{EGA IV-4: alcance e identidad de las fuentes}.
\end{itemize}
\subsection*{Fuentes y autoría}

A. Grothendieck, con la colaboración de J. Dieudonné, \emph{Éléments de géométrie algébrique III, seconde partie}, Publications Mathématiques de l’IHÉS 17 (1963), pp. 5–91. \href{https://www.numdam.org/item/PMIHES\_1963\_\_17\_\_5\_0/}{Original francés en NUMDAM}.

Esta edición española contiene el capítulo III, §§6–7, y el material final. La traducción, las correcciones y la composición están producidas por IA; la autoría de la obra original permanece con sus autores.
\subsection*{Declaración de derechos de esta edición}
Los derechos sobre las obras originales pertenecen a sus autores o a los representantes o herederos de estos. Nuestras traducciones, correcciones y composición tipográfica fueron producidas por IA. No tienen autor humano y, según nuestro entender, no están protegidas por derechos de autor; si aun así existiera algún derecho nuestro, se dedica al dominio público bajo CC0 1.0. Todo este trabajo reproduce o traduce las obras originales, por lo que nada de lo que aquí se publica concede licencia ni autoriza reproducción, traducción u otro uso de ellas. Estas ediciones se publican con fines educativos y sin ánimo de lucro, y no reclamamos ningún derecho que no poseamos. Si usted es titular de derechos sobre alguna de estas obras y se opone a su publicación aquí, abra una incidencia en GitHub y la retiraremos.
\clearpage\tableofcontents\clearpage
\EGAReferenceAlias{III.6}{section:III.6.fr}
\EGAReferenceAlias{III.6-fr}{section:III.6.fr}
\EGAReferenceAlias{III.6-es}{section:III.6.fr}
\EGAReferenceAlias{III.6.fr}{section:III.6.fr}
\EGAReferenceAlias{III.6.es}{section:III.6.fr}
\EGAReferenceAlias{III.6.eq}{section:III.6.fr}
\EGAReferenceAlias{eq:III.6}{section:III.6.fr}
\EGAReferenceAlias{section:III.6}{section:III.6.fr}
\EGAReferenceAlias{section:III.6-fr}{section:III.6.fr}
\EGAReferenceAlias{section:III.6-es}{section:III.6.fr}
\EGAReferenceAlias{section:III.6.es}{section:III.6.fr}
\EGAReferenceAlias{section:III.6.eq}{section:III.6.fr}
\EGAReferenceAlias{eq:section:III.6}{section:III.6.fr}
\EGAReferenceAlias{III.6.1}{subsection:III.6.1.fr}
\EGAReferenceAlias{III.6.1-fr}{subsection:III.6.1.fr}
\EGAReferenceAlias{III.6.1-es}{subsection:III.6.1.fr}
\EGAReferenceAlias{III.6.1.fr}{subsection:III.6.1.fr}
\EGAReferenceAlias{III.6.1.es}{subsection:III.6.1.fr}
\EGAReferenceAlias{III.6.1.eq}{subsection:III.6.1.fr}
\EGAReferenceAlias{eq:III.6.1}{subsection:III.6.1.fr}
\EGAReferenceAlias{subsection:III.6.1}{subsection:III.6.1.fr}
\EGAReferenceAlias{subsection:III.6.1-fr}{subsection:III.6.1.fr}
\EGAReferenceAlias{subsection:III.6.1-es}{subsection:III.6.1.fr}
\EGAReferenceAlias{subsection:III.6.1.es}{subsection:III.6.1.fr}
\EGAReferenceAlias{subsection:III.6.1.eq}{subsection:III.6.1.fr}
\EGAReferenceAlias{eq:subsection:III.6.1}{subsection:III.6.1.fr}
\EGAReferenceAlias{III.6.1.1-fr}{III.6.1.1}
\EGAReferenceAlias{III.6.1.1-es}{III.6.1.1}
\EGAReferenceAlias{III.6.1.1.fr}{III.6.1.1}
\EGAReferenceAlias{III.6.1.1.es}{III.6.1.1}
\EGAReferenceAlias{III.6.1.1.eq}{III.6.1.1}
\EGAReferenceAlias{eq:III.6.1.1}{III.6.1.1}
\EGAReferenceAlias{III.6.1.2-fr}{III.6.1.2}
\EGAReferenceAlias{III.6.1.2-es}{III.6.1.2}
\EGAReferenceAlias{III.6.1.2.fr}{III.6.1.2}
\EGAReferenceAlias{III.6.1.2.es}{III.6.1.2}
\EGAReferenceAlias{III.6.1.2.eq}{III.6.1.2}
\EGAReferenceAlias{eq:III.6.1.2}{III.6.1.2}
\EGAReferenceAlias{III.6.1.3-fr}{III.6.1.3}
\EGAReferenceAlias{III.6.1.3-es}{III.6.1.3}
\EGAReferenceAlias{III.6.1.3.fr}{III.6.1.3}
\EGAReferenceAlias{III.6.1.3.es}{III.6.1.3}
\EGAReferenceAlias{III.6.1.3.eq}{III.6.1.3}
\EGAReferenceAlias{eq:III.6.1.3}{III.6.1.3}
\EGAReferenceAlias{III.6.1.4-fr}{III.6.1.4}
\EGAReferenceAlias{III.6.1.4-es}{III.6.1.4}
\EGAReferenceAlias{III.6.1.4.fr}{III.6.1.4}
\EGAReferenceAlias{III.6.1.4.es}{III.6.1.4}
\EGAReferenceAlias{III.6.1.4.eq}{III.6.1.4}
\EGAReferenceAlias{eq:III.6.1.4}{III.6.1.4}
\EGAReferenceAlias{III.6.2}{subsection:III.6.2}
\EGAReferenceAlias{III.6.2-fr}{subsection:III.6.2}
\EGAReferenceAlias{III.6.2-es}{subsection:III.6.2}
\EGAReferenceAlias{III.6.2.fr}{subsection:III.6.2}
\EGAReferenceAlias{III.6.2.es}{subsection:III.6.2}
\EGAReferenceAlias{III.6.2.eq}{subsection:III.6.2}
\EGAReferenceAlias{eq:III.6.2}{subsection:III.6.2}
\EGAReferenceAlias{subsection:III.6.2-fr}{subsection:III.6.2}
\EGAReferenceAlias{subsection:III.6.2-es}{subsection:III.6.2}
\EGAReferenceAlias{subsection:III.6.2.es}{subsection:III.6.2}
\EGAReferenceAlias{subsection:III.6.2.eq}{subsection:III.6.2}
\EGAReferenceAlias{eq:subsection:III.6.2}{subsection:III.6.2}
\EGAReferenceAlias{III.6.2.1.1}{III.6.2.1.1.fr}
\EGAReferenceAlias{III.6.2.1.1-fr}{III.6.2.1.1.fr}
\EGAReferenceAlias{III.6.2.1.1-es}{III.6.2.1.1.fr}
\EGAReferenceAlias{III.6.2.1.1.es}{III.6.2.1.1.fr}
\EGAReferenceAlias{III.6.2.1.1.eq}{III.6.2.1.1.fr}
\EGAReferenceAlias{eq:III.6.2.1.1}{III.6.2.1.1.fr}
\EGAReferenceAlias{III.6.2.1.2}{III.6.2.1.2.fr}
\EGAReferenceAlias{III.6.2.1.2-fr}{III.6.2.1.2.fr}
\EGAReferenceAlias{III.6.2.1.2-es}{III.6.2.1.2.fr}
\EGAReferenceAlias{III.6.2.1.2.es}{III.6.2.1.2.fr}
\EGAReferenceAlias{III.6.2.1.2.eq}{III.6.2.1.2.fr}
\EGAReferenceAlias{eq:III.6.2.1.2}{III.6.2.1.2.fr}
\EGAReferenceAlias{III.6.2.4.1}{III.6.2.4.1.fr}
\EGAReferenceAlias{III.6.2.4.1-fr}{III.6.2.4.1.fr}
\EGAReferenceAlias{III.6.2.4.1-es}{III.6.2.4.1.fr}
\EGAReferenceAlias{III.6.2.4.1.es}{III.6.2.4.1.fr}
\EGAReferenceAlias{III.6.2.4.1.eq}{III.6.2.4.1.fr}
\EGAReferenceAlias{eq:III.6.2.4.1}{III.6.2.4.1.fr}
\EGAReferenceAlias{III.6.3}{subsection:III.6.3}
\EGAReferenceAlias{III.6.3-fr}{subsection:III.6.3}
\EGAReferenceAlias{III.6.3-es}{subsection:III.6.3}
\EGAReferenceAlias{III.6.3.fr}{subsection:III.6.3}
\EGAReferenceAlias{III.6.3.es}{subsection:III.6.3}
\EGAReferenceAlias{III.6.3.eq}{subsection:III.6.3}
\EGAReferenceAlias{eq:III.6.3}{subsection:III.6.3}
\EGAReferenceAlias{subsection:III.6.3-fr}{subsection:III.6.3}
\EGAReferenceAlias{subsection:III.6.3-es}{subsection:III.6.3}
\EGAReferenceAlias{subsection:III.6.3.es}{subsection:III.6.3}
\EGAReferenceAlias{subsection:III.6.3.eq}{subsection:III.6.3}
\EGAReferenceAlias{eq:subsection:III.6.3}{subsection:III.6.3}
\EGAReferenceAlias{III.6.3.1-fr}{III.6.3.1}
\EGAReferenceAlias{III.6.3.1-es}{III.6.3.1}
\EGAReferenceAlias{III.6.3.1.fr}{III.6.3.1}
\EGAReferenceAlias{III.6.3.1.es}{III.6.3.1}
\EGAReferenceAlias{III.6.3.1.eq}{III.6.3.1}
\EGAReferenceAlias{eq:III.6.3.1}{III.6.3.1}
\EGAReferenceAlias{III.6.3.1.1}{III.6.3.1.1.fr}
\EGAReferenceAlias{III.6.3.1.1-fr}{III.6.3.1.1.fr}
\EGAReferenceAlias{III.6.3.1.1-es}{III.6.3.1.1.fr}
\EGAReferenceAlias{III.6.3.1.1.es}{III.6.3.1.1.fr}
\EGAReferenceAlias{III.6.3.1.1.eq}{III.6.3.1.1.fr}
\EGAReferenceAlias{eq:III.6.3.1.1}{III.6.3.1.1.fr}
\EGAReferenceAlias{III.6.3.2-fr}{III.6.3.2}
\EGAReferenceAlias{III.6.3.2-es}{III.6.3.2}
\EGAReferenceAlias{III.6.3.2.fr}{III.6.3.2}
\EGAReferenceAlias{III.6.3.2.es}{III.6.3.2}
\EGAReferenceAlias{III.6.3.2.eq}{III.6.3.2}
\EGAReferenceAlias{eq:III.6.3.2}{III.6.3.2}
\EGAReferenceAlias{III.6.3.2.1}{III.6.3.2.1.fr}
\EGAReferenceAlias{III.6.3.2.1-fr}{III.6.3.2.1.fr}
\EGAReferenceAlias{III.6.3.2.1-es}{III.6.3.2.1.fr}
\EGAReferenceAlias{III.6.3.2.1.es}{III.6.3.2.1.fr}
\EGAReferenceAlias{III.6.3.2.1.eq}{III.6.3.2.1.fr}
\EGAReferenceAlias{eq:III.6.3.2.1}{III.6.3.2.1.fr}
\EGAReferenceAlias{III.6.3.2.2}{III.6.3.2.2.fr}
\EGAReferenceAlias{III.6.3.2.2-fr}{III.6.3.2.2.fr}
\EGAReferenceAlias{III.6.3.2.2-es}{III.6.3.2.2.fr}
\EGAReferenceAlias{III.6.3.2.2.es}{III.6.3.2.2.fr}
\EGAReferenceAlias{III.6.3.2.2.eq}{III.6.3.2.2.fr}
\EGAReferenceAlias{eq:III.6.3.2.2}{III.6.3.2.2.fr}
\EGAReferenceAlias{III.6.3.3-fr}{III.6.3.3}
\EGAReferenceAlias{III.6.3.3-es}{III.6.3.3}
\EGAReferenceAlias{III.6.3.3.fr}{III.6.3.3}
\EGAReferenceAlias{III.6.3.3.es}{III.6.3.3}
\EGAReferenceAlias{III.6.3.3.eq}{III.6.3.3}
\EGAReferenceAlias{eq:III.6.3.3}{III.6.3.3}
\EGAReferenceAlias{III.6.3.4-fr}{III.6.3.4}
\EGAReferenceAlias{III.6.3.4-es}{III.6.3.4}
\EGAReferenceAlias{III.6.3.4.fr}{III.6.3.4}
\EGAReferenceAlias{III.6.3.4.es}{III.6.3.4}
\EGAReferenceAlias{III.6.3.4.eq}{III.6.3.4}
\EGAReferenceAlias{eq:III.6.3.4}{III.6.3.4}
\EGAReferenceAlias{III.6.3.4.1}{III.6.3.4.1.fr}
\EGAReferenceAlias{III.6.3.4.1-fr}{III.6.3.4.1.fr}
\EGAReferenceAlias{III.6.3.4.1-es}{III.6.3.4.1.fr}
\EGAReferenceAlias{III.6.3.4.1.es}{III.6.3.4.1.fr}
\EGAReferenceAlias{III.6.3.4.1.eq}{III.6.3.4.1.fr}
\EGAReferenceAlias{eq:III.6.3.4.1}{III.6.3.4.1.fr}
\EGAReferenceAlias{III.6.3.5-fr}{III.6.3.5}
\EGAReferenceAlias{III.6.3.5-es}{III.6.3.5}
\EGAReferenceAlias{III.6.3.5.fr}{III.6.3.5}
\EGAReferenceAlias{III.6.3.5.es}{III.6.3.5}
\EGAReferenceAlias{III.6.3.5.eq}{III.6.3.5}
\EGAReferenceAlias{eq:III.6.3.5}{III.6.3.5}
\EGAReferenceAlias{III.6.3.6-fr}{III.6.3.6}
\EGAReferenceAlias{III.6.3.6-es}{III.6.3.6}
\EGAReferenceAlias{III.6.3.6.fr}{III.6.3.6}
\EGAReferenceAlias{III.6.3.6.es}{III.6.3.6}
\EGAReferenceAlias{III.6.3.6.eq}{III.6.3.6}
\EGAReferenceAlias{eq:III.6.3.6}{III.6.3.6}
\EGAReferenceAlias{III.6.3.6.1}{III.6.3.6.1.fr}
\EGAReferenceAlias{III.6.3.6.1-fr}{III.6.3.6.1.fr}
\EGAReferenceAlias{III.6.3.6.1-es}{III.6.3.6.1.fr}
\EGAReferenceAlias{III.6.3.6.1.es}{III.6.3.6.1.fr}
\EGAReferenceAlias{III.6.3.6.1.eq}{III.6.3.6.1.fr}
\EGAReferenceAlias{eq:III.6.3.6.1}{III.6.3.6.1.fr}
\EGAReferenceAlias{III.6.3.6.2}{III.6.3.6.2.fr}
\EGAReferenceAlias{III.6.3.6.2-fr}{III.6.3.6.2.fr}
\EGAReferenceAlias{III.6.3.6.2-es}{III.6.3.6.2.fr}
\EGAReferenceAlias{III.6.3.6.2.es}{III.6.3.6.2.fr}
\EGAReferenceAlias{III.6.3.6.2.eq}{III.6.3.6.2.fr}
\EGAReferenceAlias{eq:III.6.3.6.2}{III.6.3.6.2.fr}
\EGAReferenceAlias{III.6.3.6.3}{III.6.3.6.3.fr}
\EGAReferenceAlias{III.6.3.6.3-fr}{III.6.3.6.3.fr}
\EGAReferenceAlias{III.6.3.6.3-es}{III.6.3.6.3.fr}
\EGAReferenceAlias{III.6.3.6.3.es}{III.6.3.6.3.fr}
\EGAReferenceAlias{III.6.3.6.3.eq}{III.6.3.6.3.fr}
\EGAReferenceAlias{eq:III.6.3.6.3}{III.6.3.6.3.fr}
\EGAReferenceAlias{III.6.3.7-fr}{III.6.3.7}
\EGAReferenceAlias{III.6.3.7-es}{III.6.3.7}
\EGAReferenceAlias{III.6.3.7.fr}{III.6.3.7}
\EGAReferenceAlias{III.6.3.7.es}{III.6.3.7}
\EGAReferenceAlias{III.6.3.7.eq}{III.6.3.7}
\EGAReferenceAlias{eq:III.6.3.7}{III.6.3.7}
\EGAReferenceAlias{III.6.3.7.1}{III.6.3.7.1.fr}
\EGAReferenceAlias{III.6.3.7.1-fr}{III.6.3.7.1.fr}
\EGAReferenceAlias{III.6.3.7.1-es}{III.6.3.7.1.fr}
\EGAReferenceAlias{III.6.3.7.1.es}{III.6.3.7.1.fr}
\EGAReferenceAlias{III.6.3.7.1.eq}{III.6.3.7.1.fr}
\EGAReferenceAlias{eq:III.6.3.7.1}{III.6.3.7.1.fr}
\EGAReferenceAlias{III.6.3.8-fr}{III.6.3.8}
\EGAReferenceAlias{III.6.3.8-es}{III.6.3.8}
\EGAReferenceAlias{III.6.3.8.fr}{III.6.3.8}
\EGAReferenceAlias{III.6.3.8.es}{III.6.3.8}
\EGAReferenceAlias{III.6.3.8.eq}{III.6.3.8}
\EGAReferenceAlias{eq:III.6.3.8}{III.6.3.8}
\EGAReferenceAlias{III.6.3.8.1}{III.6.3.8.1.fr}
\EGAReferenceAlias{III.6.3.8.1-fr}{III.6.3.8.1.fr}
\EGAReferenceAlias{III.6.3.8.1-es}{III.6.3.8.1.fr}
\EGAReferenceAlias{III.6.3.8.1.es}{III.6.3.8.1.fr}
\EGAReferenceAlias{III.6.3.8.1.eq}{III.6.3.8.1.fr}
\EGAReferenceAlias{eq:III.6.3.8.1}{III.6.3.8.1.fr}
\EGAReferenceAlias{III.6.3.9-fr}{III.6.3.9}
\EGAReferenceAlias{III.6.3.9-es}{III.6.3.9}
\EGAReferenceAlias{III.6.3.9.fr}{III.6.3.9}
\EGAReferenceAlias{III.6.3.9.es}{III.6.3.9}
\EGAReferenceAlias{III.6.3.9.eq}{III.6.3.9}
\EGAReferenceAlias{eq:III.6.3.9}{III.6.3.9}
\EGAReferenceAlias{III.6.3.9.1}{III.6.3.9.1.fr}
\EGAReferenceAlias{III.6.3.9.1-fr}{III.6.3.9.1.fr}
\EGAReferenceAlias{III.6.3.9.1-es}{III.6.3.9.1.fr}
\EGAReferenceAlias{III.6.3.9.1.es}{III.6.3.9.1.fr}
\EGAReferenceAlias{III.6.3.9.1.eq}{III.6.3.9.1.fr}
\EGAReferenceAlias{eq:III.6.3.9.1}{III.6.3.9.1.fr}
\EGAReferenceAlias{III.6.3.9.2}{III.6.3.9.2.fr}
\EGAReferenceAlias{III.6.3.9.2-fr}{III.6.3.9.2.fr}
\EGAReferenceAlias{III.6.3.9.2-es}{III.6.3.9.2.fr}
\EGAReferenceAlias{III.6.3.9.2.es}{III.6.3.9.2.fr}
\EGAReferenceAlias{III.6.3.9.2.eq}{III.6.3.9.2.fr}
\EGAReferenceAlias{eq:III.6.3.9.2}{III.6.3.9.2.fr}
\EGAReferenceAlias{III.6.3.10-fr}{III.6.3.10}
\EGAReferenceAlias{III.6.3.10-es}{III.6.3.10}
\EGAReferenceAlias{III.6.3.10.fr}{III.6.3.10}
\EGAReferenceAlias{III.6.3.10.es}{III.6.3.10}
\EGAReferenceAlias{III.6.3.10.eq}{III.6.3.10}
\EGAReferenceAlias{eq:III.6.3.10}{III.6.3.10}
\EGAReferenceAlias{III.6.3.10.1}{III.6.3.10.1.fr}
\EGAReferenceAlias{III.6.3.10.1-fr}{III.6.3.10.1.fr}
\EGAReferenceAlias{III.6.3.10.1-es}{III.6.3.10.1.fr}
\EGAReferenceAlias{III.6.3.10.1.es}{III.6.3.10.1.fr}
\EGAReferenceAlias{III.6.3.10.1.eq}{III.6.3.10.1.fr}
\EGAReferenceAlias{eq:III.6.3.10.1}{III.6.3.10.1.fr}
\EGAReferenceAlias{III.6.3.11-fr}{III.6.3.11}
\EGAReferenceAlias{III.6.3.11-es}{III.6.3.11}
\EGAReferenceAlias{III.6.3.11.fr}{III.6.3.11}
\EGAReferenceAlias{III.6.3.11.es}{III.6.3.11}
\EGAReferenceAlias{III.6.3.11.eq}{III.6.3.11}
\EGAReferenceAlias{eq:III.6.3.11}{III.6.3.11}
\EGAReferenceAlias{III.6.3.11.1}{III.6.3.11.1.fr}
\EGAReferenceAlias{III.6.3.11.1-fr}{III.6.3.11.1.fr}
\EGAReferenceAlias{III.6.3.11.1-es}{III.6.3.11.1.fr}
\EGAReferenceAlias{III.6.3.11.1.es}{III.6.3.11.1.fr}
\EGAReferenceAlias{III.6.3.11.1.eq}{III.6.3.11.1.fr}
\EGAReferenceAlias{eq:III.6.3.11.1}{III.6.3.11.1.fr}
\EGAReferenceAlias{III.6.3.12-fr}{III.6.3.12}
\EGAReferenceAlias{III.6.3.12-es}{III.6.3.12}
\EGAReferenceAlias{III.6.3.12.fr}{III.6.3.12}
\EGAReferenceAlias{III.6.3.12.es}{III.6.3.12}
\EGAReferenceAlias{III.6.3.12.eq}{III.6.3.12}
\EGAReferenceAlias{eq:III.6.3.12}{III.6.3.12}
\EGAReferenceAlias{III.6.3.12.1}{III.6.3.12.1.fr}
\EGAReferenceAlias{III.6.3.12.1-fr}{III.6.3.12.1.fr}
\EGAReferenceAlias{III.6.3.12.1-es}{III.6.3.12.1.fr}
\EGAReferenceAlias{III.6.3.12.1.es}{III.6.3.12.1.fr}
\EGAReferenceAlias{III.6.3.12.1.eq}{III.6.3.12.1.fr}
\EGAReferenceAlias{eq:III.6.3.12.1}{III.6.3.12.1.fr}
\EGAReferenceAlias{III.6.3.13-fr}{III.6.3.13}
\EGAReferenceAlias{III.6.3.13-es}{III.6.3.13}
\EGAReferenceAlias{III.6.3.13.fr}{III.6.3.13}
\EGAReferenceAlias{III.6.3.13.es}{III.6.3.13}
\EGAReferenceAlias{III.6.3.13.eq}{III.6.3.13}
\EGAReferenceAlias{eq:III.6.3.13}{III.6.3.13}
\EGAReferenceAlias{III.6.4-fr}{subsection:III.6.4}
\EGAReferenceAlias{III.6.4-es}{subsection:III.6.4}
\EGAReferenceAlias{III.6.4.fr}{subsection:III.6.4}
\EGAReferenceAlias{III.6.4.es}{subsection:III.6.4}
\EGAReferenceAlias{III.6.4.eq}{subsection:III.6.4}
\EGAReferenceAlias{eq:III.6.4}{subsection:III.6.4}
\EGAReferenceAlias{subsection:III.6.4-fr}{subsection:III.6.4}
\EGAReferenceAlias{subsection:III.6.4-es}{subsection:III.6.4}
\EGAReferenceAlias{subsection:III.6.4.fr}{subsection:III.6.4}
\EGAReferenceAlias{subsection:III.6.4.es}{subsection:III.6.4}
\EGAReferenceAlias{subsection:III.6.4.eq}{subsection:III.6.4}
\EGAReferenceAlias{eq:subsection:III.6.4}{subsection:III.6.4}
\EGAReferenceAlias{III.6.4.1.1}{III.6.4.1.1.fr}
\EGAReferenceAlias{III.6.4.1.1-fr}{III.6.4.1.1.fr}
\EGAReferenceAlias{III.6.4.1.1-es}{III.6.4.1.1.fr}
\EGAReferenceAlias{III.6.4.1.1.es}{III.6.4.1.1.fr}
\EGAReferenceAlias{III.6.4.1.1.eq}{III.6.4.1.1.fr}
\EGAReferenceAlias{eq:III.6.4.1.1}{III.6.4.1.1.fr}
\EGAReferenceAlias{III.6.4.3-fr}{III.6.4.3}
\EGAReferenceAlias{III.6.4.3-es}{III.6.4.3}
\EGAReferenceAlias{III.6.4.3.fr}{III.6.4.3}
\EGAReferenceAlias{III.6.4.3.es}{III.6.4.3}
\EGAReferenceAlias{III.6.4.3.eq}{III.6.4.3}
\EGAReferenceAlias{eq:III.6.4.3}{III.6.4.3}
\EGAReferenceAlias{III.6.4.3.1}{III.6.4.3.1.fr}
\EGAReferenceAlias{III.6.4.3.1-fr}{III.6.4.3.1.fr}
\EGAReferenceAlias{III.6.4.3.1-es}{III.6.4.3.1.fr}
\EGAReferenceAlias{III.6.4.3.1.es}{III.6.4.3.1.fr}
\EGAReferenceAlias{III.6.4.3.1.eq}{III.6.4.3.1.fr}
\EGAReferenceAlias{eq:III.6.4.3.1}{III.6.4.3.1.fr}
\EGAReferenceAlias{III.6.4.3.2}{III.6.4.3.2.fr}
\EGAReferenceAlias{III.6.4.3.2-fr}{III.6.4.3.2.fr}
\EGAReferenceAlias{III.6.4.3.2-es}{III.6.4.3.2.fr}
\EGAReferenceAlias{III.6.4.3.2.es}{III.6.4.3.2.fr}
\EGAReferenceAlias{III.6.4.3.2.eq}{III.6.4.3.2.fr}
\EGAReferenceAlias{eq:III.6.4.3.2}{III.6.4.3.2.fr}
\EGAReferenceAlias{III.6.4.3.3}{III.6.4.3.3.fr}
\EGAReferenceAlias{III.6.4.3.3-fr}{III.6.4.3.3.fr}
\EGAReferenceAlias{III.6.4.3.3-es}{III.6.4.3.3.fr}
\EGAReferenceAlias{III.6.4.3.3.es}{III.6.4.3.3.fr}
\EGAReferenceAlias{III.6.4.3.3.eq}{III.6.4.3.3.fr}
\EGAReferenceAlias{eq:III.6.4.3.3}{III.6.4.3.3.fr}
\EGAReferenceAlias{III.6.4.4.1}{III.6.4.4.1.fr}
\EGAReferenceAlias{III.6.4.4.1-fr}{III.6.4.4.1.fr}
\EGAReferenceAlias{III.6.4.4.1-es}{III.6.4.4.1.fr}
\EGAReferenceAlias{III.6.4.4.1.es}{III.6.4.4.1.fr}
\EGAReferenceAlias{III.6.4.4.1.eq}{III.6.4.4.1.fr}
\EGAReferenceAlias{eq:III.6.4.4.1}{III.6.4.4.1.fr}
\EGAReferenceAlias{III.6.4.4.2}{III.6.4.4.2.fr}
\EGAReferenceAlias{III.6.4.4.2-fr}{III.6.4.4.2.fr}
\EGAReferenceAlias{III.6.4.4.2-es}{III.6.4.4.2.fr}
\EGAReferenceAlias{III.6.4.4.2.es}{III.6.4.4.2.fr}
\EGAReferenceAlias{III.6.4.4.2.eq}{III.6.4.4.2.fr}
\EGAReferenceAlias{eq:III.6.4.4.2}{III.6.4.4.2.fr}
\EGAReferenceAlias{III.6.4.5.1}{III.6.4.5.1.fr}
\EGAReferenceAlias{III.6.4.5.1-fr}{III.6.4.5.1.fr}
\EGAReferenceAlias{III.6.4.5.1-es}{III.6.4.5.1.fr}
\EGAReferenceAlias{III.6.4.5.1.es}{III.6.4.5.1.fr}
\EGAReferenceAlias{III.6.4.5.1.eq}{III.6.4.5.1.fr}
\EGAReferenceAlias{eq:III.6.4.5.1}{III.6.4.5.1.fr}
\EGAReferenceAlias{III.6.4.5.2}{III.6.4.5.2.fr}
\EGAReferenceAlias{III.6.4.5.2-fr}{III.6.4.5.2.fr}
\EGAReferenceAlias{III.6.4.5.2-es}{III.6.4.5.2.fr}
\EGAReferenceAlias{III.6.4.5.2.es}{III.6.4.5.2.fr}
\EGAReferenceAlias{III.6.4.5.2.eq}{III.6.4.5.2.fr}
\EGAReferenceAlias{eq:III.6.4.5.2}{III.6.4.5.2.fr}
\EGAReferenceAlias{III.6.4.5.3}{III.6.4.5.3.fr}
\EGAReferenceAlias{III.6.4.5.3-fr}{III.6.4.5.3.fr}
\EGAReferenceAlias{III.6.4.5.3-es}{III.6.4.5.3.fr}
\EGAReferenceAlias{III.6.4.5.3.es}{III.6.4.5.3.fr}
\EGAReferenceAlias{III.6.4.5.3.eq}{III.6.4.5.3.fr}
\EGAReferenceAlias{eq:III.6.4.5.3}{III.6.4.5.3.fr}
\EGAReferenceAlias{III.6.4.6-fr}{III.6.4.6}
\EGAReferenceAlias{III.6.4.6-es}{III.6.4.6}
\EGAReferenceAlias{III.6.4.6.fr}{III.6.4.6}
\EGAReferenceAlias{III.6.4.6.es}{III.6.4.6}
\EGAReferenceAlias{III.6.4.6.eq}{III.6.4.6}
\EGAReferenceAlias{eq:III.6.4.6}{III.6.4.6}
\EGAReferenceAlias{III.6.4.6.1}{III.6.4.6.1.fr}
\EGAReferenceAlias{III.6.4.6.1-fr}{III.6.4.6.1.fr}
\EGAReferenceAlias{III.6.4.6.1-es}{III.6.4.6.1.fr}
\EGAReferenceAlias{III.6.4.6.1.es}{III.6.4.6.1.fr}
\EGAReferenceAlias{III.6.4.6.1.eq}{III.6.4.6.1.fr}
\EGAReferenceAlias{eq:III.6.4.6.1}{III.6.4.6.1.fr}
\EGAReferenceAlias{III.6.5}{subsection:III.6.5}
\EGAReferenceAlias{III.6.5-fr}{subsection:III.6.5}
\EGAReferenceAlias{III.6.5-es}{subsection:III.6.5}
\EGAReferenceAlias{III.6.5.fr}{subsection:III.6.5}
\EGAReferenceAlias{III.6.5.es}{subsection:III.6.5}
\EGAReferenceAlias{III.6.5.eq}{subsection:III.6.5}
\EGAReferenceAlias{eq:III.6.5}{subsection:III.6.5}
\EGAReferenceAlias{subsection:III.6.5-fr}{subsection:III.6.5}
\EGAReferenceAlias{subsection:III.6.5-es}{subsection:III.6.5}
\EGAReferenceAlias{subsection:III.6.5.fr}{subsection:III.6.5}
\EGAReferenceAlias{subsection:III.6.5.es}{subsection:III.6.5}
\EGAReferenceAlias{subsection:III.6.5.eq}{subsection:III.6.5}
\EGAReferenceAlias{eq:subsection:III.6.5}{subsection:III.6.5}
\EGAReferenceAlias{III.6.5.4.1}{III.6.5.4.1.fr}
\EGAReferenceAlias{III.6.5.4.1-fr}{III.6.5.4.1.fr}
\EGAReferenceAlias{III.6.5.4.1-es}{III.6.5.4.1.fr}
\EGAReferenceAlias{III.6.5.4.1.es}{III.6.5.4.1.fr}
\EGAReferenceAlias{III.6.5.4.1.eq}{III.6.5.4.1.fr}
\EGAReferenceAlias{eq:III.6.5.4.1}{III.6.5.4.1.fr}
\EGAReferenceAlias{III.6.5.4.2}{III.6.5.4.2.fr}
\EGAReferenceAlias{III.6.5.4.2-fr}{III.6.5.4.2.fr}
\EGAReferenceAlias{III.6.5.4.2-es}{III.6.5.4.2.fr}
\EGAReferenceAlias{III.6.5.4.2.es}{III.6.5.4.2.fr}
\EGAReferenceAlias{III.6.5.4.2.eq}{III.6.5.4.2.fr}
\EGAReferenceAlias{eq:III.6.5.4.2}{III.6.5.4.2.fr}
\EGAReferenceAlias{III.6.5.6-fr}{III.6.5.6}
\EGAReferenceAlias{III.6.5.6-es}{III.6.5.6}
\EGAReferenceAlias{III.6.5.6.fr}{III.6.5.6}
\EGAReferenceAlias{III.6.5.6.es}{III.6.5.6}
\EGAReferenceAlias{III.6.5.6.eq}{III.6.5.6}
\EGAReferenceAlias{eq:III.6.5.6}{III.6.5.6}
\EGAReferenceAlias{III.6.5.6.1}{III.6.5.6.1.fr}
\EGAReferenceAlias{III.6.5.6.1-fr}{III.6.5.6.1.fr}
\EGAReferenceAlias{III.6.5.6.1-es}{III.6.5.6.1.fr}
\EGAReferenceAlias{III.6.5.6.1.es}{III.6.5.6.1.fr}
\EGAReferenceAlias{III.6.5.6.1.eq}{III.6.5.6.1.fr}
\EGAReferenceAlias{eq:III.6.5.6.1}{III.6.5.6.1.fr}
\EGAReferenceAlias{III.6.5.7.1}{III.6.5.7.1.fr}
\EGAReferenceAlias{III.6.5.7.1-fr}{III.6.5.7.1.fr}
\EGAReferenceAlias{III.6.5.7.1-es}{III.6.5.7.1.fr}
\EGAReferenceAlias{III.6.5.7.1.es}{III.6.5.7.1.fr}
\EGAReferenceAlias{III.6.5.7.1.eq}{III.6.5.7.1.fr}
\EGAReferenceAlias{eq:III.6.5.7.1}{III.6.5.7.1.fr}
\EGAReferenceAlias{III.6.5.9-fr}{III.6.5.9}
\EGAReferenceAlias{III.6.5.9-es}{III.6.5.9}
\EGAReferenceAlias{III.6.5.9.fr}{III.6.5.9}
\EGAReferenceAlias{III.6.5.9.es}{III.6.5.9}
\EGAReferenceAlias{III.6.5.9.eq}{III.6.5.9}
\EGAReferenceAlias{eq:III.6.5.9}{III.6.5.9}
\EGAReferenceAlias{III.6.5.9.1}{III.6.5.9.1.fr}
\EGAReferenceAlias{III.6.5.9.1-fr}{III.6.5.9.1.fr}
\EGAReferenceAlias{III.6.5.9.1-es}{III.6.5.9.1.fr}
\EGAReferenceAlias{III.6.5.9.1.es}{III.6.5.9.1.fr}
\EGAReferenceAlias{III.6.5.9.1.eq}{III.6.5.9.1.fr}
\EGAReferenceAlias{eq:III.6.5.9.1}{III.6.5.9.1.fr}
\EGAReferenceAlias{III.6.5.11-fr}{III.6.5.11}
\EGAReferenceAlias{III.6.5.11-es}{III.6.5.11}
\EGAReferenceAlias{III.6.5.11.fr}{III.6.5.11}
\EGAReferenceAlias{III.6.5.11.es}{III.6.5.11}
\EGAReferenceAlias{III.6.5.11.eq}{III.6.5.11}
\EGAReferenceAlias{eq:III.6.5.11}{III.6.5.11}
\EGAReferenceAlias{III.6.5.11.1}{III.6.5.11.1.fr}
\EGAReferenceAlias{III.6.5.11.1-fr}{III.6.5.11.1.fr}
\EGAReferenceAlias{III.6.5.11.1-es}{III.6.5.11.1.fr}
\EGAReferenceAlias{III.6.5.11.1.es}{III.6.5.11.1.fr}
\EGAReferenceAlias{III.6.5.11.1.eq}{III.6.5.11.1.fr}
\EGAReferenceAlias{eq:III.6.5.11.1}{III.6.5.11.1.fr}
\EGAReferenceAlias{III.6.5.12.1}{III.6.5.12.1.fr}
\EGAReferenceAlias{III.6.5.12.1-fr}{III.6.5.12.1.fr}
\EGAReferenceAlias{III.6.5.12.1-es}{III.6.5.12.1.fr}
\EGAReferenceAlias{III.6.5.12.1.es}{III.6.5.12.1.fr}
\EGAReferenceAlias{III.6.5.12.1.eq}{III.6.5.12.1.fr}
\EGAReferenceAlias{eq:III.6.5.12.1}{III.6.5.12.1.fr}
\EGAReferenceAlias{III.6.5.13-fr}{III.6.5.13}
\EGAReferenceAlias{III.6.5.13-es}{III.6.5.13}
\EGAReferenceAlias{III.6.5.13.fr}{III.6.5.13}
\EGAReferenceAlias{III.6.5.13.es}{III.6.5.13}
\EGAReferenceAlias{III.6.5.13.eq}{III.6.5.13}
\EGAReferenceAlias{eq:III.6.5.13}{III.6.5.13}
\EGAReferenceAlias{III.6.5.15.1}{III.6.5.15.1.fr}
\EGAReferenceAlias{III.6.5.15.1-fr}{III.6.5.15.1.fr}
\EGAReferenceAlias{III.6.5.15.1-es}{III.6.5.15.1.fr}
\EGAReferenceAlias{III.6.5.15.1.es}{III.6.5.15.1.fr}
\EGAReferenceAlias{III.6.5.15.1.eq}{III.6.5.15.1.fr}
\EGAReferenceAlias{eq:III.6.5.15.1}{III.6.5.15.1.fr}
\EGAReferenceAlias{III.6.6}{subsection:III.6.6}
\EGAReferenceAlias{III.6.6-fr}{subsection:III.6.6}
\EGAReferenceAlias{III.6.6-es}{subsection:III.6.6}
\EGAReferenceAlias{III.6.6.fr}{subsection:III.6.6}
\EGAReferenceAlias{III.6.6.es}{subsection:III.6.6}
\EGAReferenceAlias{III.6.6.eq}{subsection:III.6.6}
\EGAReferenceAlias{eq:III.6.6}{subsection:III.6.6}
\EGAReferenceAlias{subsection:III.6.6-fr}{subsection:III.6.6}
\EGAReferenceAlias{subsection:III.6.6-es}{subsection:III.6.6}
\EGAReferenceAlias{subsection:III.6.6.es}{subsection:III.6.6}
\EGAReferenceAlias{subsection:III.6.6.eq}{subsection:III.6.6}
\EGAReferenceAlias{eq:subsection:III.6.6}{subsection:III.6.6}
\EGAReferenceAlias{III.6.6.1-fr}{III.6.6.1}
\EGAReferenceAlias{III.6.6.1-es}{III.6.6.1}
\EGAReferenceAlias{III.6.6.1.fr}{III.6.6.1}
\EGAReferenceAlias{III.6.6.1.es}{III.6.6.1}
\EGAReferenceAlias{III.6.6.1.eq}{III.6.6.1}
\EGAReferenceAlias{eq:III.6.6.1}{III.6.6.1}
\EGAReferenceAlias{III.6.6.3-fr}{III.6.6.3}
\EGAReferenceAlias{III.6.6.3-es}{III.6.6.3}
\EGAReferenceAlias{III.6.6.3.fr}{III.6.6.3}
\EGAReferenceAlias{III.6.6.3.es}{III.6.6.3}
\EGAReferenceAlias{III.6.6.3.eq}{III.6.6.3}
\EGAReferenceAlias{eq:III.6.6.3}{III.6.6.3}
\EGAReferenceAlias{III.6.6.4-fr}{III.6.6.4}
\EGAReferenceAlias{III.6.6.4-es}{III.6.6.4}
\EGAReferenceAlias{III.6.6.4.fr}{III.6.6.4}
\EGAReferenceAlias{III.6.6.4.es}{III.6.6.4}
\EGAReferenceAlias{III.6.6.4.eq}{III.6.6.4}
\EGAReferenceAlias{eq:III.6.6.4}{III.6.6.4}
\EGAReferenceAlias{III.6.6.5-fr}{III.6.6.5}
\EGAReferenceAlias{III.6.6.5-es}{III.6.6.5}
\EGAReferenceAlias{III.6.6.5.fr}{III.6.6.5}
\EGAReferenceAlias{III.6.6.5.es}{III.6.6.5}
\EGAReferenceAlias{III.6.6.5.eq}{III.6.6.5}
\EGAReferenceAlias{eq:III.6.6.5}{III.6.6.5}
\EGAReferenceAlias{III.6.6.6-fr}{III.6.6.6}
\EGAReferenceAlias{III.6.6.6-es}{III.6.6.6}
\EGAReferenceAlias{III.6.6.6.fr}{III.6.6.6}
\EGAReferenceAlias{III.6.6.6.es}{III.6.6.6}
\EGAReferenceAlias{III.6.6.6.eq}{III.6.6.6}
\EGAReferenceAlias{eq:III.6.6.6}{III.6.6.6}
\EGAReferenceAlias{III.6.6.6.1-fr}{III.6.6.6.1}
\EGAReferenceAlias{III.6.6.6.1-es}{III.6.6.6.1}
\EGAReferenceAlias{III.6.6.6.1.es}{III.6.6.6.1}
\EGAReferenceAlias{III.6.6.6.1.eq}{III.6.6.6.1}
\EGAReferenceAlias{eq:III.6.6.6.1}{III.6.6.6.1}
\EGAReferenceAlias{III.6.6.7-fr}{III.6.6.7}
\EGAReferenceAlias{III.6.6.7-es}{III.6.6.7}
\EGAReferenceAlias{III.6.6.7.fr}{III.6.6.7}
\EGAReferenceAlias{III.6.6.7.es}{III.6.6.7}
\EGAReferenceAlias{III.6.6.7.eq}{III.6.6.7}
\EGAReferenceAlias{eq:III.6.6.7}{III.6.6.7}
\EGAReferenceAlias{III.6.6.7.t}{III.6.6.7.t.fr}
\EGAReferenceAlias{III.6.6.7.t-fr}{III.6.6.7.t.fr}
\EGAReferenceAlias{III.6.6.7.t-es}{III.6.6.7.t.fr}
\EGAReferenceAlias{III.6.6.7.t.es}{III.6.6.7.t.fr}
\EGAReferenceAlias{III.6.6.7.t.eq}{III.6.6.7.t.fr}
\EGAReferenceAlias{eq:III.6.6.7.t}{III.6.6.7.t.fr}
\EGAReferenceAlias{III.6.6.8-fr}{III.6.6.8}
\EGAReferenceAlias{III.6.6.8-es}{III.6.6.8}
\EGAReferenceAlias{III.6.6.8.fr}{III.6.6.8}
\EGAReferenceAlias{III.6.6.8.es}{III.6.6.8}
\EGAReferenceAlias{III.6.6.8.eq}{III.6.6.8}
\EGAReferenceAlias{eq:III.6.6.8}{III.6.6.8}
\EGAReferenceAlias{III.6.7}{subsection:III.6.7}
\EGAReferenceAlias{III.6.7-fr}{subsection:III.6.7}
\EGAReferenceAlias{III.6.7-es}{subsection:III.6.7}
\EGAReferenceAlias{III.6.7.fr}{subsection:III.6.7}
\EGAReferenceAlias{III.6.7.es}{subsection:III.6.7}
\EGAReferenceAlias{III.6.7.eq}{subsection:III.6.7}
\EGAReferenceAlias{eq:III.6.7}{subsection:III.6.7}
\EGAReferenceAlias{subsection:III.6.7-fr}{subsection:III.6.7}
\EGAReferenceAlias{subsection:III.6.7-es}{subsection:III.6.7}
\EGAReferenceAlias{subsection:III.6.7.fr}{subsection:III.6.7}
\EGAReferenceAlias{subsection:III.6.7.es}{subsection:III.6.7}
\EGAReferenceAlias{subsection:III.6.7.eq}{subsection:III.6.7}
\EGAReferenceAlias{eq:subsection:III.6.7}{subsection:III.6.7}
\EGAReferenceAlias{III.6.7.1-fr}{III.6.7.1}
\EGAReferenceAlias{III.6.7.1-es}{III.6.7.1}
\EGAReferenceAlias{III.6.7.1.fr}{III.6.7.1}
\EGAReferenceAlias{III.6.7.1.es}{III.6.7.1}
\EGAReferenceAlias{III.6.7.1.eq}{III.6.7.1}
\EGAReferenceAlias{eq:III.6.7.1}{III.6.7.1}
\EGAReferenceAlias{III.6.7.1.1-fr}{III.6.7.1.1}
\EGAReferenceAlias{III.6.7.1.1-es}{III.6.7.1.1}
\EGAReferenceAlias{III.6.7.1.1.es}{III.6.7.1.1}
\EGAReferenceAlias{III.6.7.1.1.eq}{III.6.7.1.1}
\EGAReferenceAlias{eq:III.6.7.1.1}{III.6.7.1.1}
\EGAReferenceAlias{III.6.7.2-fr}{III.6.7.2}
\EGAReferenceAlias{III.6.7.2-es}{III.6.7.2}
\EGAReferenceAlias{III.6.7.2.fr}{III.6.7.2}
\EGAReferenceAlias{III.6.7.2.es}{III.6.7.2}
\EGAReferenceAlias{III.6.7.2.eq}{III.6.7.2}
\EGAReferenceAlias{eq:III.6.7.2}{III.6.7.2}
\EGAReferenceAlias{III.6.7.2.1-fr}{III.6.7.2.1}
\EGAReferenceAlias{III.6.7.2.1-es}{III.6.7.2.1}
\EGAReferenceAlias{III.6.7.2.1.es}{III.6.7.2.1}
\EGAReferenceAlias{III.6.7.2.1.eq}{III.6.7.2.1}
\EGAReferenceAlias{eq:III.6.7.2.1}{III.6.7.2.1}
\EGAReferenceAlias{III.6.7.3-fr}{III.6.7.3}
\EGAReferenceAlias{III.6.7.3-es}{III.6.7.3}
\EGAReferenceAlias{III.6.7.3.fr}{III.6.7.3}
\EGAReferenceAlias{III.6.7.3.es}{III.6.7.3}
\EGAReferenceAlias{III.6.7.3.eq}{III.6.7.3}
\EGAReferenceAlias{eq:III.6.7.3}{III.6.7.3}
\EGAReferenceAlias{III.6.7.4-fr}{III.6.7.4}
\EGAReferenceAlias{III.6.7.4-es}{III.6.7.4}
\EGAReferenceAlias{III.6.7.4.fr}{III.6.7.4}
\EGAReferenceAlias{III.6.7.4.es}{III.6.7.4}
\EGAReferenceAlias{III.6.7.4.eq}{III.6.7.4}
\EGAReferenceAlias{eq:III.6.7.4}{III.6.7.4}
\EGAReferenceAlias{III.6.7.4.1-fr}{III.6.7.4.1}
\EGAReferenceAlias{III.6.7.4.1-es}{III.6.7.4.1}
\EGAReferenceAlias{III.6.7.4.1.es}{III.6.7.4.1}
\EGAReferenceAlias{III.6.7.4.1.eq}{III.6.7.4.1}
\EGAReferenceAlias{eq:III.6.7.4.1}{III.6.7.4.1}
\EGAReferenceAlias{III.6.7.4.2-fr}{III.6.7.4.2}
\EGAReferenceAlias{III.6.7.4.2-es}{III.6.7.4.2}
\EGAReferenceAlias{III.6.7.4.2.es}{III.6.7.4.2}
\EGAReferenceAlias{III.6.7.4.2.eq}{III.6.7.4.2}
\EGAReferenceAlias{eq:III.6.7.4.2}{III.6.7.4.2}
\EGAReferenceAlias{III.6.7.5-fr}{III.6.7.5}
\EGAReferenceAlias{III.6.7.5-es}{III.6.7.5}
\EGAReferenceAlias{III.6.7.5.fr}{III.6.7.5}
\EGAReferenceAlias{III.6.7.5.es}{III.6.7.5}
\EGAReferenceAlias{III.6.7.5.eq}{III.6.7.5}
\EGAReferenceAlias{eq:III.6.7.5}{III.6.7.5}
\EGAReferenceAlias{III.6.7.6-fr}{III.6.7.6}
\EGAReferenceAlias{III.6.7.6-es}{III.6.7.6}
\EGAReferenceAlias{III.6.7.6.fr}{III.6.7.6}
\EGAReferenceAlias{III.6.7.6.es}{III.6.7.6}
\EGAReferenceAlias{III.6.7.6.eq}{III.6.7.6}
\EGAReferenceAlias{eq:III.6.7.6}{III.6.7.6}
\EGAReferenceAlias{III.6.7.6.1-fr}{III.6.7.6.1}
\EGAReferenceAlias{III.6.7.6.1-es}{III.6.7.6.1}
\EGAReferenceAlias{III.6.7.6.1.es}{III.6.7.6.1}
\EGAReferenceAlias{III.6.7.6.1.eq}{III.6.7.6.1}
\EGAReferenceAlias{eq:III.6.7.6.1}{III.6.7.6.1}
\EGAReferenceAlias{III.6.7.6.2-fr}{III.6.7.6.2}
\EGAReferenceAlias{III.6.7.6.2-es}{III.6.7.6.2}
\EGAReferenceAlias{III.6.7.6.2.es}{III.6.7.6.2}
\EGAReferenceAlias{III.6.7.6.2.eq}{III.6.7.6.2}
\EGAReferenceAlias{eq:III.6.7.6.2}{III.6.7.6.2}
\EGAReferenceAlias{III.6.7.6.3-fr}{III.6.7.6.3}
\EGAReferenceAlias{III.6.7.6.3-es}{III.6.7.6.3}
\EGAReferenceAlias{III.6.7.6.3.es}{III.6.7.6.3}
\EGAReferenceAlias{III.6.7.6.3.eq}{III.6.7.6.3}
\EGAReferenceAlias{eq:III.6.7.6.3}{III.6.7.6.3}
\EGAReferenceAlias{III.6.7.7-fr}{III.6.7.7}
\EGAReferenceAlias{III.6.7.7-es}{III.6.7.7}
\EGAReferenceAlias{III.6.7.7.fr}{III.6.7.7}
\EGAReferenceAlias{III.6.7.7.es}{III.6.7.7}
\EGAReferenceAlias{III.6.7.7.eq}{III.6.7.7}
\EGAReferenceAlias{eq:III.6.7.7}{III.6.7.7}
\EGAReferenceAlias{III.6.7.7.1-fr}{III.6.7.7.1}
\EGAReferenceAlias{III.6.7.7.1-es}{III.6.7.7.1}
\EGAReferenceAlias{III.6.7.7.1.es}{III.6.7.7.1}
\EGAReferenceAlias{III.6.7.7.1.eq}{III.6.7.7.1}
\EGAReferenceAlias{eq:III.6.7.7.1}{III.6.7.7.1}
\EGAReferenceAlias{III.6.7.8-fr}{III.6.7.8}
\EGAReferenceAlias{III.6.7.8-es}{III.6.7.8}
\EGAReferenceAlias{III.6.7.8.fr}{III.6.7.8}
\EGAReferenceAlias{III.6.7.8.es}{III.6.7.8}
\EGAReferenceAlias{III.6.7.8.eq}{III.6.7.8}
\EGAReferenceAlias{eq:III.6.7.8}{III.6.7.8}
\EGAReferenceAlias{III.6.7.8.1-fr}{III.6.7.8.1}
\EGAReferenceAlias{III.6.7.8.1-es}{III.6.7.8.1}
\EGAReferenceAlias{III.6.7.8.1.es}{III.6.7.8.1}
\EGAReferenceAlias{III.6.7.8.1.eq}{III.6.7.8.1}
\EGAReferenceAlias{eq:III.6.7.8.1}{III.6.7.8.1}
\EGAReferenceAlias{III.6.7.9-fr}{III.6.7.9}
\EGAReferenceAlias{III.6.7.9-es}{III.6.7.9}
\EGAReferenceAlias{III.6.7.9.fr}{III.6.7.9}
\EGAReferenceAlias{III.6.7.9.es}{III.6.7.9}
\EGAReferenceAlias{III.6.7.9.eq}{III.6.7.9}
\EGAReferenceAlias{eq:III.6.7.9}{III.6.7.9}
\EGAReferenceAlias{III.6.7.10-fr}{III.6.7.10}
\EGAReferenceAlias{III.6.7.10-es}{III.6.7.10}
\EGAReferenceAlias{III.6.7.10.fr}{III.6.7.10}
\EGAReferenceAlias{III.6.7.10.es}{III.6.7.10}
\EGAReferenceAlias{III.6.7.10.eq}{III.6.7.10}
\EGAReferenceAlias{eq:III.6.7.10}{III.6.7.10}
\EGAReferenceAlias{III.6.7.10.1-fr}{III.6.7.10.1}
\EGAReferenceAlias{III.6.7.10.1-es}{III.6.7.10.1}
\EGAReferenceAlias{III.6.7.10.1.es}{III.6.7.10.1}
\EGAReferenceAlias{III.6.7.10.1.eq}{III.6.7.10.1}
\EGAReferenceAlias{eq:III.6.7.10.1}{III.6.7.10.1}
\EGAReferenceAlias{III.6.7.10.2-fr}{III.6.7.10.2}
\EGAReferenceAlias{III.6.7.10.2-es}{III.6.7.10.2}
\EGAReferenceAlias{III.6.7.10.2.es}{III.6.7.10.2}
\EGAReferenceAlias{III.6.7.10.2.eq}{III.6.7.10.2}
\EGAReferenceAlias{eq:III.6.7.10.2}{III.6.7.10.2}
\EGAReferenceAlias{III.6.7.10.3-fr}{III.6.7.10.3}
\EGAReferenceAlias{III.6.7.10.3-es}{III.6.7.10.3}
\EGAReferenceAlias{III.6.7.10.3.es}{III.6.7.10.3}
\EGAReferenceAlias{III.6.7.10.3.eq}{III.6.7.10.3}
\EGAReferenceAlias{eq:III.6.7.10.3}{III.6.7.10.3}
\EGAReferenceAlias{III.6.7.10.4-fr}{III.6.7.10.4}
\EGAReferenceAlias{III.6.7.10.4-es}{III.6.7.10.4}
\EGAReferenceAlias{III.6.7.10.4.es}{III.6.7.10.4}
\EGAReferenceAlias{III.6.7.10.4.eq}{III.6.7.10.4}
\EGAReferenceAlias{eq:III.6.7.10.4}{III.6.7.10.4}
\EGAReferenceAlias{III.6.7.10.5-fr}{III.6.7.10.5}
\EGAReferenceAlias{III.6.7.10.5-es}{III.6.7.10.5}
\EGAReferenceAlias{III.6.7.10.5.es}{III.6.7.10.5}
\EGAReferenceAlias{III.6.7.10.5.eq}{III.6.7.10.5}
\EGAReferenceAlias{eq:III.6.7.10.5}{III.6.7.10.5}
\EGAReferenceAlias{III.6.7.11-fr}{III.6.7.11}
\EGAReferenceAlias{III.6.7.11-es}{III.6.7.11}
\EGAReferenceAlias{III.6.7.11.fr}{III.6.7.11}
\EGAReferenceAlias{III.6.7.11.es}{III.6.7.11}
\EGAReferenceAlias{III.6.7.11.eq}{III.6.7.11}
\EGAReferenceAlias{eq:III.6.7.11}{III.6.7.11}
\EGAReferenceAlias{III.6.7.11.2-fr}{III.6.7.11.2}
\EGAReferenceAlias{III.6.7.11.2-es}{III.6.7.11.2}
\EGAReferenceAlias{III.6.7.11.2.fr}{III.6.7.11.2}
\EGAReferenceAlias{III.6.7.11.2.es}{III.6.7.11.2}
\EGAReferenceAlias{III.6.7.11.2.eq}{III.6.7.11.2}
\EGAReferenceAlias{eq:III.6.7.11.2}{III.6.7.11.2}
\EGAReferenceAlias{III.6.7.12-fr}{III.6.7.12}
\EGAReferenceAlias{III.6.7.12-es}{III.6.7.12}
\EGAReferenceAlias{III.6.7.12.fr}{III.6.7.12}
\EGAReferenceAlias{III.6.7.12.es}{III.6.7.12}
\EGAReferenceAlias{III.6.7.12.eq}{III.6.7.12}
\EGAReferenceAlias{eq:III.6.7.12}{III.6.7.12}
\EGAReferenceAlias{III.6.7.12.1-fr}{III.6.7.12.1}
\EGAReferenceAlias{III.6.7.12.1-es}{III.6.7.12.1}
\EGAReferenceAlias{III.6.7.12.1.es}{III.6.7.12.1}
\EGAReferenceAlias{III.6.7.12.1.eq}{III.6.7.12.1}
\EGAReferenceAlias{eq:III.6.7.12.1}{III.6.7.12.1}
\EGAReferenceAlias{III.6.7.13-fr}{III.6.7.13}
\EGAReferenceAlias{III.6.7.13-es}{III.6.7.13}
\EGAReferenceAlias{III.6.7.13.fr}{III.6.7.13}
\EGAReferenceAlias{III.6.7.13.es}{III.6.7.13}
\EGAReferenceAlias{III.6.7.13.eq}{III.6.7.13}
\EGAReferenceAlias{eq:III.6.7.13}{III.6.7.13}
\EGAReferenceAlias{III.6.9.1-fr}{III.6.9.1}
\EGAReferenceAlias{III.6.9.1-es}{III.6.9.1}
\EGAReferenceAlias{III.6.9.1.fr}{III.6.9.1}
\EGAReferenceAlias{III.6.9.1.es}{III.6.9.1}
\EGAReferenceAlias{III.6.9.1.eq}{III.6.9.1}
\EGAReferenceAlias{eq:III.6.9.1}{III.6.9.1}
\EGAReferenceAlias{III.6.9.2-fr}{III.6.9.2}
\EGAReferenceAlias{III.6.9.2-es}{III.6.9.2}
\EGAReferenceAlias{III.6.9.2.fr}{III.6.9.2}
\EGAReferenceAlias{III.6.9.2.es}{III.6.9.2}
\EGAReferenceAlias{III.6.9.2.eq}{III.6.9.2}
\EGAReferenceAlias{eq:III.6.9.2}{III.6.9.2}
\EGAReferenceAlias{III.6.9.2.1-fr}{III.6.9.2.1}
\EGAReferenceAlias{III.6.9.2.1-es}{III.6.9.2.1}
\EGAReferenceAlias{III.6.9.2.1.es}{III.6.9.2.1}
\EGAReferenceAlias{III.6.9.2.1.eq}{III.6.9.2.1}
\EGAReferenceAlias{eq:III.6.9.2.1}{III.6.9.2.1}
\EGAReferenceAlias{III.6.9.2.2-fr}{III.6.9.2.2}
\EGAReferenceAlias{III.6.9.2.2-es}{III.6.9.2.2}
\EGAReferenceAlias{III.6.9.2.2.fr}{III.6.9.2.2}
\EGAReferenceAlias{III.6.9.2.2.es}{III.6.9.2.2}
\EGAReferenceAlias{III.6.9.2.2.eq}{III.6.9.2.2}
\EGAReferenceAlias{eq:III.6.9.2.2}{III.6.9.2.2}
\EGAReferenceAlias{III.6.9.4-fr}{III.6.9.4}
\EGAReferenceAlias{III.6.9.4-es}{III.6.9.4}
\EGAReferenceAlias{III.6.9.4.fr}{III.6.9.4}
\EGAReferenceAlias{III.6.9.4.es}{III.6.9.4}
\EGAReferenceAlias{III.6.9.4.eq}{III.6.9.4}
\EGAReferenceAlias{eq:III.6.9.4}{III.6.9.4}
\EGAReferenceAlias{III.6.9.5-fr}{III.6.9.5}
\EGAReferenceAlias{III.6.9.5-es}{III.6.9.5}
\EGAReferenceAlias{III.6.9.5.fr}{III.6.9.5}
\EGAReferenceAlias{III.6.9.5.es}{III.6.9.5}
\EGAReferenceAlias{III.6.9.5.eq}{III.6.9.5}
\EGAReferenceAlias{eq:III.6.9.5}{III.6.9.5}
\EGAReferenceAlias{III.6.9.6-fr}{III.6.9.6}
\EGAReferenceAlias{III.6.9.6-es}{III.6.9.6}
\EGAReferenceAlias{III.6.9.6.fr}{III.6.9.6}
\EGAReferenceAlias{III.6.9.6.es}{III.6.9.6}
\EGAReferenceAlias{III.6.9.6.eq}{III.6.9.6}
\EGAReferenceAlias{eq:III.6.9.6}{III.6.9.6}
\EGAReferenceAlias{III.6.9.6.1-fr}{III.6.9.6.1}
\EGAReferenceAlias{III.6.9.6.1-es}{III.6.9.6.1}
\EGAReferenceAlias{III.6.9.6.1.es}{III.6.9.6.1}
\EGAReferenceAlias{III.6.9.6.1.eq}{III.6.9.6.1}
\EGAReferenceAlias{eq:III.6.9.6.1}{III.6.9.6.1}
\EGAReferenceAlias{III.6.9.6.2-fr}{III.6.9.6.2}
\EGAReferenceAlias{III.6.9.6.2-es}{III.6.9.6.2}
\EGAReferenceAlias{III.6.9.6.2.es}{III.6.9.6.2}
\EGAReferenceAlias{III.6.9.6.2.eq}{III.6.9.6.2}
\EGAReferenceAlias{eq:III.6.9.6.2}{III.6.9.6.2}
\EGAReferenceAlias{III.6.9.6.3-fr}{III.6.9.6.3}
\EGAReferenceAlias{III.6.9.6.3-es}{III.6.9.6.3}
\EGAReferenceAlias{III.6.9.6.3.es}{III.6.9.6.3}
\EGAReferenceAlias{III.6.9.6.3.eq}{III.6.9.6.3}
\EGAReferenceAlias{eq:III.6.9.6.3}{III.6.9.6.3}
\EGAReferenceAlias{III.6.9.6.4}{III.6.9.6.4.fr}
\EGAReferenceAlias{III.6.9.6.4-fr}{III.6.9.6.4.fr}
\EGAReferenceAlias{III.6.9.6.4-es}{III.6.9.6.4.fr}
\EGAReferenceAlias{III.6.9.6.4.es}{III.6.9.6.4.fr}
\EGAReferenceAlias{III.6.9.6.4.eq}{III.6.9.6.4.fr}
\EGAReferenceAlias{eq:III.6.9.6.4}{III.6.9.6.4.fr}
\EGAReferenceAlias{III.6.9.6.5}{III.6.9.6.5.fr}
\EGAReferenceAlias{III.6.9.6.5-fr}{III.6.9.6.5.fr}
\EGAReferenceAlias{III.6.9.6.5-es}{III.6.9.6.5.fr}
\EGAReferenceAlias{III.6.9.6.5.es}{III.6.9.6.5.fr}
\EGAReferenceAlias{III.6.9.6.5.eq}{III.6.9.6.5.fr}
\EGAReferenceAlias{eq:III.6.9.6.5}{III.6.9.6.5.fr}
\EGAReferenceAlias{III.6.9.7-fr}{III.6.9.7}
\EGAReferenceAlias{III.6.9.7-es}{III.6.9.7}
\EGAReferenceAlias{III.6.9.7.fr}{III.6.9.7}
\EGAReferenceAlias{III.6.9.7.es}{III.6.9.7}
\EGAReferenceAlias{III.6.9.7.eq}{III.6.9.7}
\EGAReferenceAlias{eq:III.6.9.7}{III.6.9.7}
\EGAReferenceAlias{III.6.9.7.1}{III.6.9.7.1.fr}
\EGAReferenceAlias{III.6.9.7.1-fr}{III.6.9.7.1.fr}
\EGAReferenceAlias{III.6.9.7.1-es}{III.6.9.7.1.fr}
\EGAReferenceAlias{III.6.9.7.1.es}{III.6.9.7.1.fr}
\EGAReferenceAlias{III.6.9.7.1.eq}{III.6.9.7.1.fr}
\EGAReferenceAlias{eq:III.6.9.7.1}{III.6.9.7.1.fr}
\EGAReferenceAlias{III.6.9.8-fr}{III.6.9.8}
\EGAReferenceAlias{III.6.9.8-es}{III.6.9.8}
\EGAReferenceAlias{III.6.9.8.fr}{III.6.9.8}
\EGAReferenceAlias{III.6.9.8.es}{III.6.9.8}
\EGAReferenceAlias{III.6.9.8.eq}{III.6.9.8}
\EGAReferenceAlias{eq:III.6.9.8}{III.6.9.8}
\EGAReferenceAlias{III.6.9.8.1-fr}{III.6.9.8.1}
\EGAReferenceAlias{III.6.9.8.1-es}{III.6.9.8.1}
\EGAReferenceAlias{III.6.9.8.1.es}{III.6.9.8.1}
\EGAReferenceAlias{III.6.9.8.1.eq}{III.6.9.8.1}
\EGAReferenceAlias{eq:III.6.9.8.1}{III.6.9.8.1}
\EGAReferenceAlias{III.6.9.8.2-fr}{III.6.9.8.2}
\EGAReferenceAlias{III.6.9.8.2-es}{III.6.9.8.2}
\EGAReferenceAlias{III.6.9.8.2.es}{III.6.9.8.2}
\EGAReferenceAlias{III.6.9.8.2.eq}{III.6.9.8.2}
\EGAReferenceAlias{eq:III.6.9.8.2}{III.6.9.8.2}
\EGAReferenceAlias{III.6.9.8.3}{III.6.9.8.3.fr}
\EGAReferenceAlias{III.6.9.8.3-fr}{III.6.9.8.3.fr}
\EGAReferenceAlias{III.6.9.8.3-es}{III.6.9.8.3.fr}
\EGAReferenceAlias{III.6.9.8.3.es}{III.6.9.8.3.fr}
\EGAReferenceAlias{III.6.9.8.3.eq}{III.6.9.8.3.fr}
\EGAReferenceAlias{eq:III.6.9.8.3}{III.6.9.8.3.fr}
\EGAReferenceAlias{III.6.9.8.4}{III.6.9.8.4.fr}
\EGAReferenceAlias{III.6.9.8.4-fr}{III.6.9.8.4.fr}
\EGAReferenceAlias{III.6.9.8.4-es}{III.6.9.8.4.fr}
\EGAReferenceAlias{III.6.9.8.4.es}{III.6.9.8.4.fr}
\EGAReferenceAlias{III.6.9.8.4.eq}{III.6.9.8.4.fr}
\EGAReferenceAlias{eq:III.6.9.8.4}{III.6.9.8.4.fr}
\EGAReferenceAlias{III.6.9.9-fr}{III.6.9.9}
\EGAReferenceAlias{III.6.9.9-es}{III.6.9.9}
\EGAReferenceAlias{III.6.9.9.fr}{III.6.9.9}
\EGAReferenceAlias{III.6.9.9.es}{III.6.9.9}
\EGAReferenceAlias{III.6.9.9.eq}{III.6.9.9}
\EGAReferenceAlias{eq:III.6.9.9}{III.6.9.9}
\EGAReferenceAlias{III.6.9.9.1-fr}{III.6.9.9.1}
\EGAReferenceAlias{III.6.9.9.1-es}{III.6.9.9.1}
\EGAReferenceAlias{III.6.9.9.1.es}{III.6.9.9.1}
\EGAReferenceAlias{III.6.9.9.1.eq}{III.6.9.9.1}
\EGAReferenceAlias{eq:III.6.9.9.1}{III.6.9.9.1}
\EGAReferenceAlias{III.6.9.9.2-fr}{III.6.9.9.2}
\EGAReferenceAlias{III.6.9.9.2-es}{III.6.9.9.2}
\EGAReferenceAlias{III.6.9.9.2.es}{III.6.9.9.2}
\EGAReferenceAlias{III.6.9.9.2.eq}{III.6.9.9.2}
\EGAReferenceAlias{eq:III.6.9.9.2}{III.6.9.9.2}
\EGAReferenceAlias{III.6.9.10-fr}{III.6.9.10}
\EGAReferenceAlias{III.6.9.10-es}{III.6.9.10}
\EGAReferenceAlias{III.6.9.10.fr}{III.6.9.10}
\EGAReferenceAlias{III.6.9.10.es}{III.6.9.10}
\EGAReferenceAlias{III.6.9.10.eq}{III.6.9.10}
\EGAReferenceAlias{eq:III.6.9.10}{III.6.9.10}
\EGAReferenceAlias{III.6.9.10.1-fr}{III.6.9.10.1}
\EGAReferenceAlias{III.6.9.10.1-es}{III.6.9.10.1}
\EGAReferenceAlias{III.6.9.10.1.es}{III.6.9.10.1}
\EGAReferenceAlias{III.6.9.10.1.eq}{III.6.9.10.1}
\EGAReferenceAlias{eq:III.6.9.10.1}{III.6.9.10.1}
\EGAReferenceAlias{III.6.10-fr}{subsection:III.6.10}
\EGAReferenceAlias{III.6.10-es}{subsection:III.6.10}
\EGAReferenceAlias{III.6.10.fr}{subsection:III.6.10}
\EGAReferenceAlias{III.6.10.es}{subsection:III.6.10}
\EGAReferenceAlias{III.6.10.eq}{subsection:III.6.10}
\EGAReferenceAlias{eq:III.6.10}{subsection:III.6.10}
\EGAReferenceAlias{subsection:III.6.10-fr}{subsection:III.6.10}
\EGAReferenceAlias{subsection:III.6.10-es}{subsection:III.6.10}
\EGAReferenceAlias{subsection:III.6.10.fr}{subsection:III.6.10}
\EGAReferenceAlias{subsection:III.6.10.es}{subsection:III.6.10}
\EGAReferenceAlias{subsection:III.6.10.eq}{subsection:III.6.10}
\EGAReferenceAlias{eq:subsection:III.6.10}{subsection:III.6.10}
\EGAReferenceAlias{III.6.10.1.1}{III.6.10.1.1.fr}
\EGAReferenceAlias{III.6.10.1.1-fr}{III.6.10.1.1.fr}
\EGAReferenceAlias{III.6.10.1.1-es}{III.6.10.1.1.fr}
\EGAReferenceAlias{III.6.10.1.1.es}{III.6.10.1.1.fr}
\EGAReferenceAlias{III.6.10.1.1.eq}{III.6.10.1.1.fr}
\EGAReferenceAlias{eq:III.6.10.1.1}{III.6.10.1.1.fr}
\EGAReferenceAlias{III.6.10.1.2}{III.6.10.1.2.fr}
\EGAReferenceAlias{III.6.10.1.2-fr}{III.6.10.1.2.fr}
\EGAReferenceAlias{III.6.10.1.2-es}{III.6.10.1.2.fr}
\EGAReferenceAlias{III.6.10.1.2.es}{III.6.10.1.2.fr}
\EGAReferenceAlias{III.6.10.1.2.eq}{III.6.10.1.2.fr}
\EGAReferenceAlias{eq:III.6.10.1.2}{III.6.10.1.2.fr}
\EGAReferenceAlias{III.6.10.2-fr}{III.6.10.2}
\EGAReferenceAlias{III.6.10.2-es}{III.6.10.2}
\EGAReferenceAlias{III.6.10.2.fr}{III.6.10.2}
\EGAReferenceAlias{III.6.10.2.es}{III.6.10.2}
\EGAReferenceAlias{III.6.10.2.eq}{III.6.10.2}
\EGAReferenceAlias{eq:III.6.10.2}{III.6.10.2}
\EGAReferenceAlias{III.6.10.3.1-fr}{III.6.10.3.1}
\EGAReferenceAlias{III.6.10.3.1-es}{III.6.10.3.1}
\EGAReferenceAlias{III.6.10.3.1.es}{III.6.10.3.1}
\EGAReferenceAlias{III.6.10.3.1.eq}{III.6.10.3.1}
\EGAReferenceAlias{eq:III.6.10.3.1}{III.6.10.3.1}
\EGAReferenceAlias{III.6.10.4-fr}{III.6.10.4}
\EGAReferenceAlias{III.6.10.4-es}{III.6.10.4}
\EGAReferenceAlias{III.6.10.4.fr}{III.6.10.4}
\EGAReferenceAlias{III.6.10.4.es}{III.6.10.4}
\EGAReferenceAlias{III.6.10.4.eq}{III.6.10.4}
\EGAReferenceAlias{eq:III.6.10.4}{III.6.10.4}
\EGAReferenceAlias{III.6.10.5.1}{III.6.10.5.1.fr}
\EGAReferenceAlias{III.6.10.5.1-fr}{III.6.10.5.1.fr}
\EGAReferenceAlias{III.6.10.5.1-es}{III.6.10.5.1.fr}
\EGAReferenceAlias{III.6.10.5.1.es}{III.6.10.5.1.fr}
\EGAReferenceAlias{III.6.10.5.1.eq}{III.6.10.5.1.fr}
\EGAReferenceAlias{eq:III.6.10.5.1}{III.6.10.5.1.fr}
\EGAReferenceAlias{III.6.10.5.2}{III.6.10.5.2.fr}
\EGAReferenceAlias{III.6.10.5.2-fr}{III.6.10.5.2.fr}
\EGAReferenceAlias{III.6.10.5.2-es}{III.6.10.5.2.fr}
\EGAReferenceAlias{III.6.10.5.2.es}{III.6.10.5.2.fr}
\EGAReferenceAlias{III.6.10.5.2.eq}{III.6.10.5.2.fr}
\EGAReferenceAlias{eq:III.6.10.5.2}{III.6.10.5.2.fr}
\EGAReferenceAlias{III.6.10.5.3}{III.6.10.5.3.fr}
\EGAReferenceAlias{III.6.10.5.3-fr}{III.6.10.5.3.fr}
\EGAReferenceAlias{III.6.10.5.3-es}{III.6.10.5.3.fr}
\EGAReferenceAlias{III.6.10.5.3.es}{III.6.10.5.3.fr}
\EGAReferenceAlias{III.6.10.5.3.eq}{III.6.10.5.3.fr}
\EGAReferenceAlias{eq:III.6.10.5.3}{III.6.10.5.3.fr}
\EGAReferenceAlias{III.6.10.5.4}{III.6.10.5.4.fr}
\EGAReferenceAlias{III.6.10.5.4-fr}{III.6.10.5.4.fr}
\EGAReferenceAlias{III.6.10.5.4-es}{III.6.10.5.4.fr}
\EGAReferenceAlias{III.6.10.5.4.es}{III.6.10.5.4.fr}
\EGAReferenceAlias{III.6.10.5.4.eq}{III.6.10.5.4.fr}
\EGAReferenceAlias{eq:III.6.10.5.4}{III.6.10.5.4.fr}
\EGAReferenceAlias{III.6.10.6-fr}{III.6.10.6}
\EGAReferenceAlias{III.6.10.6-es}{III.6.10.6}
\EGAReferenceAlias{III.6.10.6.fr}{III.6.10.6}
\EGAReferenceAlias{III.6.10.6.es}{III.6.10.6}
\EGAReferenceAlias{III.6.10.6.eq}{III.6.10.6}
\EGAReferenceAlias{eq:III.6.10.6}{III.6.10.6}
\EGAReferenceAlias{III.7}{section:III.7}
\EGAReferenceAlias{III.7-fr}{section:III.7}
\EGAReferenceAlias{III.7-es}{section:III.7}
\EGAReferenceAlias{III.7.fr}{section:III.7}
\EGAReferenceAlias{III.7.es}{section:III.7}
\EGAReferenceAlias{III.7.eq}{section:III.7}
\EGAReferenceAlias{eq:III.7}{section:III.7}
\EGAReferenceAlias{section:III.7-fr}{section:III.7}
\EGAReferenceAlias{section:III.7-es}{section:III.7}
\EGAReferenceAlias{section:III.7.es}{section:III.7}
\EGAReferenceAlias{section:III.7.eq}{section:III.7}
\EGAReferenceAlias{eq:section:III.7}{section:III.7}
\EGAReferenceAlias{III.7.1}{subsection:III.7.1}
\EGAReferenceAlias{III.7.1-fr}{subsection:III.7.1}
\EGAReferenceAlias{III.7.1-es}{subsection:III.7.1}
\EGAReferenceAlias{III.7.1.fr}{subsection:III.7.1}
\EGAReferenceAlias{III.7.1.es}{subsection:III.7.1}
\EGAReferenceAlias{III.7.1.eq}{subsection:III.7.1}
\EGAReferenceAlias{eq:III.7.1}{subsection:III.7.1}
\EGAReferenceAlias{subsection:III.7.1-fr}{subsection:III.7.1}
\EGAReferenceAlias{subsection:III.7.1-es}{subsection:III.7.1}
\EGAReferenceAlias{subsection:III.7.1.es}{subsection:III.7.1}
\EGAReferenceAlias{subsection:III.7.1.eq}{subsection:III.7.1}
\EGAReferenceAlias{eq:subsection:III.7.1}{subsection:III.7.1}
\EGAReferenceAlias{III.7.1.1-fr}{III.7.1.1}
\EGAReferenceAlias{III.7.1.1-es}{III.7.1.1}
\EGAReferenceAlias{III.7.1.1.fr}{III.7.1.1}
\EGAReferenceAlias{III.7.1.1.es}{III.7.1.1}
\EGAReferenceAlias{III.7.1.1.eq}{III.7.1.1}
\EGAReferenceAlias{eq:III.7.1.1}{III.7.1.1}
\EGAReferenceAlias{III.7.1.4-fr}{III.7.1.4}
\EGAReferenceAlias{III.7.1.4-es}{III.7.1.4}
\EGAReferenceAlias{III.7.1.4.fr}{III.7.1.4}
\EGAReferenceAlias{III.7.1.4.es}{III.7.1.4}
\EGAReferenceAlias{III.7.1.4.eq}{III.7.1.4}
\EGAReferenceAlias{eq:III.7.1.4}{III.7.1.4}
\EGAReferenceAlias{III.7.1.4.1-fr}{III.7.1.4.1}
\EGAReferenceAlias{III.7.1.4.1-es}{III.7.1.4.1}
\EGAReferenceAlias{III.7.1.4.1.es}{III.7.1.4.1}
\EGAReferenceAlias{III.7.1.4.1.eq}{III.7.1.4.1}
\EGAReferenceAlias{eq:III.7.1.4.1}{III.7.1.4.1}
\EGAReferenceAlias{III.7.1.4.2-fr}{III.7.1.4.2}
\EGAReferenceAlias{III.7.1.4.2-es}{III.7.1.4.2}
\EGAReferenceAlias{III.7.1.4.2.es}{III.7.1.4.2}
\EGAReferenceAlias{III.7.1.4.2.eq}{III.7.1.4.2}
\EGAReferenceAlias{eq:III.7.1.4.2}{III.7.1.4.2}
\EGAReferenceAlias{III.7.2}{subsection:III.7.2}
\EGAReferenceAlias{III.7.2-fr}{subsection:III.7.2}
\EGAReferenceAlias{III.7.2-es}{subsection:III.7.2}
\EGAReferenceAlias{III.7.2.fr}{subsection:III.7.2}
\EGAReferenceAlias{III.7.2.es}{subsection:III.7.2}
\EGAReferenceAlias{III.7.2.eq}{subsection:III.7.2}
\EGAReferenceAlias{eq:III.7.2}{subsection:III.7.2}
\EGAReferenceAlias{subsection:III.7.2-fr}{subsection:III.7.2}
\EGAReferenceAlias{subsection:III.7.2-es}{subsection:III.7.2}
\EGAReferenceAlias{subsection:III.7.2.es}{subsection:III.7.2}
\EGAReferenceAlias{subsection:III.7.2.eq}{subsection:III.7.2}
\EGAReferenceAlias{eq:subsection:III.7.2}{subsection:III.7.2}
\EGAReferenceAlias{III.7.2.1-fr}{III.7.2.1}
\EGAReferenceAlias{III.7.2.1-es}{III.7.2.1}
\EGAReferenceAlias{III.7.2.1.fr}{III.7.2.1}
\EGAReferenceAlias{III.7.2.1.es}{III.7.2.1}
\EGAReferenceAlias{III.7.2.1.eq}{III.7.2.1}
\EGAReferenceAlias{eq:III.7.2.1}{III.7.2.1}
\EGAReferenceAlias{III.7.2.2.1}{III.7.2.2.1.fr}
\EGAReferenceAlias{III.7.2.2.1-fr}{III.7.2.2.1.fr}
\EGAReferenceAlias{III.7.2.2.1-es}{III.7.2.2.1.fr}
\EGAReferenceAlias{III.7.2.2.1.es}{III.7.2.2.1.fr}
\EGAReferenceAlias{III.7.2.2.1.eq}{III.7.2.2.1.fr}
\EGAReferenceAlias{eq:III.7.2.2.1}{III.7.2.2.1.fr}
\EGAReferenceAlias{III.7.2.2.2}{III.7.2.2.2.fr}
\EGAReferenceAlias{III.7.2.2.2-fr}{III.7.2.2.2.fr}
\EGAReferenceAlias{III.7.2.2.2-es}{III.7.2.2.2.fr}
\EGAReferenceAlias{III.7.2.2.2.es}{III.7.2.2.2.fr}
\EGAReferenceAlias{III.7.2.2.2.eq}{III.7.2.2.2.fr}
\EGAReferenceAlias{eq:III.7.2.2.2}{III.7.2.2.2.fr}
\EGAReferenceAlias{III.7.2.3-fr}{III.7.2.3}
\EGAReferenceAlias{III.7.2.3-es}{III.7.2.3}
\EGAReferenceAlias{III.7.2.3.fr}{III.7.2.3}
\EGAReferenceAlias{III.7.2.3.es}{III.7.2.3}
\EGAReferenceAlias{III.7.2.3.eq}{III.7.2.3}
\EGAReferenceAlias{eq:III.7.2.3}{III.7.2.3}
\EGAReferenceAlias{III.7.2.3.1-fr}{III.7.2.3.1}
\EGAReferenceAlias{III.7.2.3.1-es}{III.7.2.3.1}
\EGAReferenceAlias{III.7.2.3.1.es}{III.7.2.3.1}
\EGAReferenceAlias{III.7.2.3.1.eq}{III.7.2.3.1}
\EGAReferenceAlias{eq:III.7.2.3.1}{III.7.2.3.1}
\EGAReferenceAlias{III.7.2.3.2-fr}{III.7.2.3.2}
\EGAReferenceAlias{III.7.2.3.2-es}{III.7.2.3.2}
\EGAReferenceAlias{III.7.2.3.2.es}{III.7.2.3.2}
\EGAReferenceAlias{III.7.2.3.2.eq}{III.7.2.3.2}
\EGAReferenceAlias{eq:III.7.2.3.2}{III.7.2.3.2}
\EGAReferenceAlias{III.7.2.3.3-fr}{III.7.2.3.3}
\EGAReferenceAlias{III.7.2.3.3-es}{III.7.2.3.3}
\EGAReferenceAlias{III.7.2.3.3.es}{III.7.2.3.3}
\EGAReferenceAlias{III.7.2.3.3.eq}{III.7.2.3.3}
\EGAReferenceAlias{eq:III.7.2.3.3}{III.7.2.3.3}
\EGAReferenceAlias{III.7.2.4-fr}{III.7.2.4}
\EGAReferenceAlias{III.7.2.4-es}{III.7.2.4}
\EGAReferenceAlias{III.7.2.4.fr}{III.7.2.4}
\EGAReferenceAlias{III.7.2.4.es}{III.7.2.4}
\EGAReferenceAlias{III.7.2.4.eq}{III.7.2.4}
\EGAReferenceAlias{eq:III.7.2.4}{III.7.2.4}
\EGAReferenceAlias{III.7.2.6-fr}{III.7.2.6}
\EGAReferenceAlias{III.7.2.6-es}{III.7.2.6}
\EGAReferenceAlias{III.7.2.6.fr}{III.7.2.6}
\EGAReferenceAlias{III.7.2.6.es}{III.7.2.6}
\EGAReferenceAlias{III.7.2.6.eq}{III.7.2.6}
\EGAReferenceAlias{eq:III.7.2.6}{III.7.2.6}
\EGAReferenceAlias{III.7.2.7-fr}{III.7.2.7}
\EGAReferenceAlias{III.7.2.7-es}{III.7.2.7}
\EGAReferenceAlias{III.7.2.7.fr}{III.7.2.7}
\EGAReferenceAlias{III.7.2.7.es}{III.7.2.7}
\EGAReferenceAlias{III.7.2.7.eq}{III.7.2.7}
\EGAReferenceAlias{eq:III.7.2.7}{III.7.2.7}
\EGAReferenceAlias{III.7.3}{subsection:III.7.3}
\EGAReferenceAlias{III.7.3-fr}{subsection:III.7.3}
\EGAReferenceAlias{III.7.3-es}{subsection:III.7.3}
\EGAReferenceAlias{III.7.3.fr}{subsection:III.7.3}
\EGAReferenceAlias{III.7.3.es}{subsection:III.7.3}
\EGAReferenceAlias{III.7.3.eq}{subsection:III.7.3}
\EGAReferenceAlias{eq:III.7.3}{subsection:III.7.3}
\EGAReferenceAlias{subsection:III.7.3-fr}{subsection:III.7.3}
\EGAReferenceAlias{subsection:III.7.3-es}{subsection:III.7.3}
\EGAReferenceAlias{subsection:III.7.3.es}{subsection:III.7.3}
\EGAReferenceAlias{subsection:III.7.3.eq}{subsection:III.7.3}
\EGAReferenceAlias{eq:subsection:III.7.3}{subsection:III.7.3}
\EGAReferenceAlias{III.7.3.1-fr}{III.7.3.1}
\EGAReferenceAlias{III.7.3.1-es}{III.7.3.1}
\EGAReferenceAlias{III.7.3.1.fr}{III.7.3.1}
\EGAReferenceAlias{III.7.3.1.es}{III.7.3.1}
\EGAReferenceAlias{III.7.3.1.eq}{III.7.3.1}
\EGAReferenceAlias{eq:III.7.3.1}{III.7.3.1}
\EGAReferenceAlias{III.7.3.1.1}{III.7.3.1.1.fr}
\EGAReferenceAlias{III.7.3.1.1-fr}{III.7.3.1.1.fr}
\EGAReferenceAlias{III.7.3.1.1-es}{III.7.3.1.1.fr}
\EGAReferenceAlias{III.7.3.1.1.es}{III.7.3.1.1.fr}
\EGAReferenceAlias{III.7.3.1.1.eq}{III.7.3.1.1.fr}
\EGAReferenceAlias{eq:III.7.3.1.1}{III.7.3.1.1.fr}
\EGAReferenceAlias{III.7.3.2-fr}{III.7.3.2}
\EGAReferenceAlias{III.7.3.2-es}{III.7.3.2}
\EGAReferenceAlias{III.7.3.2.fr}{III.7.3.2}
\EGAReferenceAlias{III.7.3.2.es}{III.7.3.2}
\EGAReferenceAlias{III.7.3.2.eq}{III.7.3.2}
\EGAReferenceAlias{eq:III.7.3.2}{III.7.3.2}
\EGAReferenceAlias{III.7.3.3-fr}{III.7.3.3}
\EGAReferenceAlias{III.7.3.3-es}{III.7.3.3}
\EGAReferenceAlias{III.7.3.3.fr}{III.7.3.3}
\EGAReferenceAlias{III.7.3.3.es}{III.7.3.3}
\EGAReferenceAlias{III.7.3.3.eq}{III.7.3.3}
\EGAReferenceAlias{eq:III.7.3.3}{III.7.3.3}
\EGAReferenceAlias{III.7.3.3.1}{III.7.3.3.1.fr}
\EGAReferenceAlias{III.7.3.3.1-fr}{III.7.3.3.1.fr}
\EGAReferenceAlias{III.7.3.3.1-es}{III.7.3.3.1.fr}
\EGAReferenceAlias{III.7.3.3.1.es}{III.7.3.3.1.fr}
\EGAReferenceAlias{III.7.3.3.1.eq}{III.7.3.3.1.fr}
\EGAReferenceAlias{eq:III.7.3.3.1}{III.7.3.3.1.fr}
\EGAReferenceAlias{III.7.3.4-fr}{III.7.3.4}
\EGAReferenceAlias{III.7.3.4-es}{III.7.3.4}
\EGAReferenceAlias{III.7.3.4.fr}{III.7.3.4}
\EGAReferenceAlias{III.7.3.4.es}{III.7.3.4}
\EGAReferenceAlias{III.7.3.4.eq}{III.7.3.4}
\EGAReferenceAlias{eq:III.7.3.4}{III.7.3.4}
\EGAReferenceAlias{III.7.3.5.1}{III.7.3.5.1.fr}
\EGAReferenceAlias{III.7.3.5.1-fr}{III.7.3.5.1.fr}
\EGAReferenceAlias{III.7.3.5.1-es}{III.7.3.5.1.fr}
\EGAReferenceAlias{III.7.3.5.1.es}{III.7.3.5.1.fr}
\EGAReferenceAlias{III.7.3.5.1.eq}{III.7.3.5.1.fr}
\EGAReferenceAlias{eq:III.7.3.5.1}{III.7.3.5.1.fr}
\EGAReferenceAlias{III.7.3.5.2}{III.7.3.5.2.fr}
\EGAReferenceAlias{III.7.3.5.2-fr}{III.7.3.5.2.fr}
\EGAReferenceAlias{III.7.3.5.2-es}{III.7.3.5.2.fr}
\EGAReferenceAlias{III.7.3.5.2.es}{III.7.3.5.2.fr}
\EGAReferenceAlias{III.7.3.5.2.eq}{III.7.3.5.2.fr}
\EGAReferenceAlias{eq:III.7.3.5.2}{III.7.3.5.2.fr}
\EGAReferenceAlias{III.7.3.5.4}{III.7.3.5.4.fr}
\EGAReferenceAlias{III.7.3.5.4-fr}{III.7.3.5.4.fr}
\EGAReferenceAlias{III.7.3.5.4-es}{III.7.3.5.4.fr}
\EGAReferenceAlias{III.7.3.5.4.es}{III.7.3.5.4.fr}
\EGAReferenceAlias{III.7.3.5.4.eq}{III.7.3.5.4.fr}
\EGAReferenceAlias{eq:III.7.3.5.4}{III.7.3.5.4.fr}
\EGAReferenceAlias{III.7.3.5.5}{III.7.3.5.5.fr}
\EGAReferenceAlias{III.7.3.5.5-fr}{III.7.3.5.5.fr}
\EGAReferenceAlias{III.7.3.5.5-es}{III.7.3.5.5.fr}
\EGAReferenceAlias{III.7.3.5.5.es}{III.7.3.5.5.fr}
\EGAReferenceAlias{III.7.3.5.5.eq}{III.7.3.5.5.fr}
\EGAReferenceAlias{eq:III.7.3.5.5}{III.7.3.5.5.fr}
\EGAReferenceAlias{III.7.3.5.6}{III.7.3.5.6.fr}
\EGAReferenceAlias{III.7.3.5.6-fr}{III.7.3.5.6.fr}
\EGAReferenceAlias{III.7.3.5.6-es}{III.7.3.5.6.fr}
\EGAReferenceAlias{III.7.3.5.6.es}{III.7.3.5.6.fr}
\EGAReferenceAlias{III.7.3.5.6.eq}{III.7.3.5.6.fr}
\EGAReferenceAlias{eq:III.7.3.5.6}{III.7.3.5.6.fr}
\EGAReferenceAlias{III.7.3.5.8}{III.7.3.5.8.fr}
\EGAReferenceAlias{III.7.3.5.8-fr}{III.7.3.5.8.fr}
\EGAReferenceAlias{III.7.3.5.8-es}{III.7.3.5.8.fr}
\EGAReferenceAlias{III.7.3.5.8.es}{III.7.3.5.8.fr}
\EGAReferenceAlias{III.7.3.5.8.eq}{III.7.3.5.8.fr}
\EGAReferenceAlias{eq:III.7.3.5.8}{III.7.3.5.8.fr}
\EGAReferenceAlias{III.7.3.6-fr}{III.7.3.6}
\EGAReferenceAlias{III.7.3.6-es}{III.7.3.6}
\EGAReferenceAlias{III.7.3.6.fr}{III.7.3.6}
\EGAReferenceAlias{III.7.3.6.es}{III.7.3.6}
\EGAReferenceAlias{III.7.3.6.eq}{III.7.3.6}
\EGAReferenceAlias{eq:III.7.3.6}{III.7.3.6}
\EGAReferenceAlias{III.7.3.7-fr}{III.7.3.7}
\EGAReferenceAlias{III.7.3.7-es}{III.7.3.7}
\EGAReferenceAlias{III.7.3.7.fr}{III.7.3.7}
\EGAReferenceAlias{III.7.3.7.es}{III.7.3.7}
\EGAReferenceAlias{III.7.3.7.eq}{III.7.3.7}
\EGAReferenceAlias{eq:III.7.3.7}{III.7.3.7}
\EGAReferenceAlias{III.7.3.8-fr}{III.7.3.8}
\EGAReferenceAlias{III.7.3.8-es}{III.7.3.8}
\EGAReferenceAlias{III.7.3.8.fr}{III.7.3.8}
\EGAReferenceAlias{III.7.3.8.es}{III.7.3.8}
\EGAReferenceAlias{III.7.3.8.eq}{III.7.3.8}
\EGAReferenceAlias{eq:III.7.3.8}{III.7.3.8}
\EGAReferenceAlias{III.7.3.9-fr}{III.7.3.9}
\EGAReferenceAlias{III.7.3.9-es}{III.7.3.9}
\EGAReferenceAlias{III.7.3.9.fr}{III.7.3.9}
\EGAReferenceAlias{III.7.3.9.es}{III.7.3.9}
\EGAReferenceAlias{III.7.3.9.eq}{III.7.3.9}
\EGAReferenceAlias{eq:III.7.3.9}{III.7.3.9}
\EGAReferenceAlias{III.7.4}{subsection:III.7.4.fr}
\EGAReferenceAlias{III.7.4-fr}{subsection:III.7.4.fr}
\EGAReferenceAlias{III.7.4-es}{subsection:III.7.4.fr}
\EGAReferenceAlias{III.7.4.fr}{subsection:III.7.4.fr}
\EGAReferenceAlias{III.7.4.es}{subsection:III.7.4.fr}
\EGAReferenceAlias{III.7.4.eq}{subsection:III.7.4.fr}
\EGAReferenceAlias{eq:III.7.4}{subsection:III.7.4.fr}
\EGAReferenceAlias{subsection:III.7.4}{subsection:III.7.4.fr}
\EGAReferenceAlias{subsection:III.7.4-fr}{subsection:III.7.4.fr}
\EGAReferenceAlias{subsection:III.7.4-es}{subsection:III.7.4.fr}
\EGAReferenceAlias{subsection:III.7.4.es}{subsection:III.7.4.fr}
\EGAReferenceAlias{subsection:III.7.4.eq}{subsection:III.7.4.fr}
\EGAReferenceAlias{eq:subsection:III.7.4}{subsection:III.7.4.fr}
\EGAReferenceAlias{III.7.4.1-fr}{III.7.4.1}
\EGAReferenceAlias{III.7.4.1-es}{III.7.4.1}
\EGAReferenceAlias{III.7.4.1.fr}{III.7.4.1}
\EGAReferenceAlias{III.7.4.1.es}{III.7.4.1}
\EGAReferenceAlias{III.7.4.1.eq}{III.7.4.1}
\EGAReferenceAlias{eq:III.7.4.1}{III.7.4.1}
\EGAReferenceAlias{III.7.4.1.1}{III.7.4.1.1.fr}
\EGAReferenceAlias{III.7.4.1.1-fr}{III.7.4.1.1.fr}
\EGAReferenceAlias{III.7.4.1.1-es}{III.7.4.1.1.fr}
\EGAReferenceAlias{III.7.4.1.1.es}{III.7.4.1.1.fr}
\EGAReferenceAlias{III.7.4.1.1.eq}{III.7.4.1.1.fr}
\EGAReferenceAlias{eq:III.7.4.1.1}{III.7.4.1.1.fr}
\EGAReferenceAlias{III.7.4.1.2}{III.7.4.1.2.fr}
\EGAReferenceAlias{III.7.4.1.2-fr}{III.7.4.1.2.fr}
\EGAReferenceAlias{III.7.4.1.2-es}{III.7.4.1.2.fr}
\EGAReferenceAlias{III.7.4.1.2.es}{III.7.4.1.2.fr}
\EGAReferenceAlias{III.7.4.1.2.eq}{III.7.4.1.2.fr}
\EGAReferenceAlias{eq:III.7.4.1.2}{III.7.4.1.2.fr}
\EGAReferenceAlias{III.7.4.1.3}{III.7.4.1.3.fr}
\EGAReferenceAlias{III.7.4.1.3-fr}{III.7.4.1.3.fr}
\EGAReferenceAlias{III.7.4.1.3-es}{III.7.4.1.3.fr}
\EGAReferenceAlias{III.7.4.1.3.es}{III.7.4.1.3.fr}
\EGAReferenceAlias{III.7.4.1.3.eq}{III.7.4.1.3.fr}
\EGAReferenceAlias{eq:III.7.4.1.3}{III.7.4.1.3.fr}
\EGAReferenceAlias{III.7.4.2.1}{III.7.4.2.1.fr}
\EGAReferenceAlias{III.7.4.2.1-fr}{III.7.4.2.1.fr}
\EGAReferenceAlias{III.7.4.2.1-es}{III.7.4.2.1.fr}
\EGAReferenceAlias{III.7.4.2.1.es}{III.7.4.2.1.fr}
\EGAReferenceAlias{III.7.4.2.1.eq}{III.7.4.2.1.fr}
\EGAReferenceAlias{eq:III.7.4.2.1}{III.7.4.2.1.fr}
\EGAReferenceAlias{III.7.4.3-fr}{III.7.4.3}
\EGAReferenceAlias{III.7.4.3-es}{III.7.4.3}
\EGAReferenceAlias{III.7.4.3.fr}{III.7.4.3}
\EGAReferenceAlias{III.7.4.3.es}{III.7.4.3}
\EGAReferenceAlias{III.7.4.3.eq}{III.7.4.3}
\EGAReferenceAlias{eq:III.7.4.3}{III.7.4.3}
\EGAReferenceAlias{III.7.4.4-fr}{III.7.4.4}
\EGAReferenceAlias{III.7.4.4-es}{III.7.4.4}
\EGAReferenceAlias{III.7.4.4.fr}{III.7.4.4}
\EGAReferenceAlias{III.7.4.4.es}{III.7.4.4}
\EGAReferenceAlias{III.7.4.4.eq}{III.7.4.4}
\EGAReferenceAlias{eq:III.7.4.4}{III.7.4.4}
\EGAReferenceAlias{III.7.4.5-fr}{III.7.4.5}
\EGAReferenceAlias{III.7.4.5-es}{III.7.4.5}
\EGAReferenceAlias{III.7.4.5.fr}{III.7.4.5}
\EGAReferenceAlias{III.7.4.5.es}{III.7.4.5}
\EGAReferenceAlias{III.7.4.5.eq}{III.7.4.5}
\EGAReferenceAlias{eq:III.7.4.5}{III.7.4.5}
\EGAReferenceAlias{III.7.4.6.1}{III.7.4.6.1.fr}
\EGAReferenceAlias{III.7.4.6.1-fr}{III.7.4.6.1.fr}
\EGAReferenceAlias{III.7.4.6.1-es}{III.7.4.6.1.fr}
\EGAReferenceAlias{III.7.4.6.1.es}{III.7.4.6.1.fr}
\EGAReferenceAlias{III.7.4.6.1.eq}{III.7.4.6.1.fr}
\EGAReferenceAlias{eq:III.7.4.6.1}{III.7.4.6.1.fr}
\EGAReferenceAlias{III.7.4.7-fr}{III.7.4.7}
\EGAReferenceAlias{III.7.4.7-es}{III.7.4.7}
\EGAReferenceAlias{III.7.4.7.fr}{III.7.4.7}
\EGAReferenceAlias{III.7.4.7.es}{III.7.4.7}
\EGAReferenceAlias{III.7.4.7.eq}{III.7.4.7}
\EGAReferenceAlias{eq:III.7.4.7}{III.7.4.7}
\EGAReferenceAlias{III.7.4.7.1}{III.7.4.7.1.fr}
\EGAReferenceAlias{III.7.4.7.1-fr}{III.7.4.7.1.fr}
\EGAReferenceAlias{III.7.4.7.1-es}{III.7.4.7.1.fr}
\EGAReferenceAlias{III.7.4.7.1.es}{III.7.4.7.1.fr}
\EGAReferenceAlias{III.7.4.7.1.eq}{III.7.4.7.1.fr}
\EGAReferenceAlias{eq:III.7.4.7.1}{III.7.4.7.1.fr}
\EGAReferenceAlias{III.7.4.7.2-fr}{III.7.4.7.2}
\EGAReferenceAlias{III.7.4.7.2-es}{III.7.4.7.2}
\EGAReferenceAlias{III.7.4.7.2.fr}{III.7.4.7.2}
\EGAReferenceAlias{III.7.4.7.2.es}{III.7.4.7.2}
\EGAReferenceAlias{III.7.4.7.2.eq}{III.7.4.7.2}
\EGAReferenceAlias{eq:III.7.4.7.2}{III.7.4.7.2}
\EGAReferenceAlias{III.7.4.7.3}{III.7.4.7.3.fr}
\EGAReferenceAlias{III.7.4.7.3-fr}{III.7.4.7.3.fr}
\EGAReferenceAlias{III.7.4.7.3-es}{III.7.4.7.3.fr}
\EGAReferenceAlias{III.7.4.7.3.es}{III.7.4.7.3.fr}
\EGAReferenceAlias{III.7.4.7.3.eq}{III.7.4.7.3.fr}
\EGAReferenceAlias{eq:III.7.4.7.3}{III.7.4.7.3.fr}
\EGAReferenceAlias{III.7.4.8-fr}{III.7.4.8}
\EGAReferenceAlias{III.7.4.8-es}{III.7.4.8}
\EGAReferenceAlias{III.7.4.8.fr}{III.7.4.8}
\EGAReferenceAlias{III.7.4.8.es}{III.7.4.8}
\EGAReferenceAlias{III.7.4.8.eq}{III.7.4.8}
\EGAReferenceAlias{eq:III.7.4.8}{III.7.4.8}
\EGAReferenceAlias{III.7.5}{subsection:III.7.5}
\EGAReferenceAlias{III.7.5-fr}{subsection:III.7.5}
\EGAReferenceAlias{III.7.5-es}{subsection:III.7.5}
\EGAReferenceAlias{III.7.5.fr}{subsection:III.7.5}
\EGAReferenceAlias{III.7.5.es}{subsection:III.7.5}
\EGAReferenceAlias{III.7.5.eq}{subsection:III.7.5}
\EGAReferenceAlias{eq:III.7.5}{subsection:III.7.5}
\EGAReferenceAlias{subsection:III.7.5-fr}{subsection:III.7.5}
\EGAReferenceAlias{subsection:III.7.5-es}{subsection:III.7.5}
\EGAReferenceAlias{subsection:III.7.5.es}{subsection:III.7.5}
\EGAReferenceAlias{subsection:III.7.5.eq}{subsection:III.7.5}
\EGAReferenceAlias{eq:subsection:III.7.5}{subsection:III.7.5}
\EGAReferenceAlias{III.7.5.1-fr}{III.7.5.1}
\EGAReferenceAlias{III.7.5.1-es}{III.7.5.1}
\EGAReferenceAlias{III.7.5.1.fr}{III.7.5.1}
\EGAReferenceAlias{III.7.5.1.es}{III.7.5.1}
\EGAReferenceAlias{III.7.5.1.eq}{III.7.5.1}
\EGAReferenceAlias{eq:III.7.5.1}{III.7.5.1}
\EGAReferenceAlias{III.7.5.1.1}{III.7.5.1.1.fr}
\EGAReferenceAlias{III.7.5.1.1-fr}{III.7.5.1.1.fr}
\EGAReferenceAlias{III.7.5.1.1-es}{III.7.5.1.1.fr}
\EGAReferenceAlias{III.7.5.1.1.es}{III.7.5.1.1.fr}
\EGAReferenceAlias{III.7.5.1.1.eq}{III.7.5.1.1.fr}
\EGAReferenceAlias{eq:III.7.5.1.1}{III.7.5.1.1.fr}
\EGAReferenceAlias{III.7.5.2-fr}{III.7.5.2}
\EGAReferenceAlias{III.7.5.2-es}{III.7.5.2}
\EGAReferenceAlias{III.7.5.2.fr}{III.7.5.2}
\EGAReferenceAlias{III.7.5.2.es}{III.7.5.2}
\EGAReferenceAlias{III.7.5.2.eq}{III.7.5.2}
\EGAReferenceAlias{eq:III.7.5.2}{III.7.5.2}
\EGAReferenceAlias{III.7.5.3-fr}{III.7.5.3}
\EGAReferenceAlias{III.7.5.3-es}{III.7.5.3}
\EGAReferenceAlias{III.7.5.3.fr}{III.7.5.3}
\EGAReferenceAlias{III.7.5.3.es}{III.7.5.3}
\EGAReferenceAlias{III.7.5.3.eq}{III.7.5.3}
\EGAReferenceAlias{eq:III.7.5.3}{III.7.5.3}
\EGAReferenceAlias{III.7.5.4-fr}{III.7.5.4}
\EGAReferenceAlias{III.7.5.4-es}{III.7.5.4}
\EGAReferenceAlias{III.7.5.4.fr}{III.7.5.4}
\EGAReferenceAlias{III.7.5.4.es}{III.7.5.4}
\EGAReferenceAlias{III.7.5.4.eq}{III.7.5.4}
\EGAReferenceAlias{eq:III.7.5.4}{III.7.5.4}
\EGAReferenceAlias{III.7.5.5-fr}{III.7.5.5}
\EGAReferenceAlias{III.7.5.5-es}{III.7.5.5}
\EGAReferenceAlias{III.7.5.5.fr}{III.7.5.5}
\EGAReferenceAlias{III.7.5.5.es}{III.7.5.5}
\EGAReferenceAlias{III.7.5.5.eq}{III.7.5.5}
\EGAReferenceAlias{eq:III.7.5.5}{III.7.5.5}
\EGAReferenceAlias{III.7.5.7-fr}{III.7.5.7}
\EGAReferenceAlias{III.7.5.7-es}{III.7.5.7}
\EGAReferenceAlias{III.7.5.7.fr}{III.7.5.7}
\EGAReferenceAlias{III.7.5.7.es}{III.7.5.7}
\EGAReferenceAlias{III.7.5.7.eq}{III.7.5.7}
\EGAReferenceAlias{eq:III.7.5.7}{III.7.5.7}
\EGAReferenceAlias{III.7.6-fr}{subsection:III.7.6}
\EGAReferenceAlias{III.7.6-es}{subsection:III.7.6}
\EGAReferenceAlias{III.7.6.fr}{subsection:III.7.6}
\EGAReferenceAlias{III.7.6.es}{subsection:III.7.6}
\EGAReferenceAlias{III.7.6.eq}{subsection:III.7.6}
\EGAReferenceAlias{eq:III.7.6}{subsection:III.7.6}
\EGAReferenceAlias{subsection:III.7.6-fr}{subsection:III.7.6}
\EGAReferenceAlias{subsection:III.7.6-es}{subsection:III.7.6}
\EGAReferenceAlias{subsection:III.7.6.es}{subsection:III.7.6}
\EGAReferenceAlias{subsection:III.7.6.eq}{subsection:III.7.6}
\EGAReferenceAlias{eq:subsection:III.7.6}{subsection:III.7.6}
\EGAReferenceAlias{III.7.6.1-fr}{III.7.6.1}
\EGAReferenceAlias{III.7.6.1-es}{III.7.6.1}
\EGAReferenceAlias{III.7.6.1.fr}{III.7.6.1}
\EGAReferenceAlias{III.7.6.1.es}{III.7.6.1}
\EGAReferenceAlias{III.7.6.1.eq}{III.7.6.1}
\EGAReferenceAlias{eq:III.7.6.1}{III.7.6.1}
\EGAReferenceAlias{III.7.6.1.1}{III.7.6.1.1.fr}
\EGAReferenceAlias{III.7.6.1.1-fr}{III.7.6.1.1.fr}
\EGAReferenceAlias{III.7.6.1.1-es}{III.7.6.1.1.fr}
\EGAReferenceAlias{III.7.6.1.1.es}{III.7.6.1.1.fr}
\EGAReferenceAlias{III.7.6.1.1.eq}{III.7.6.1.1.fr}
\EGAReferenceAlias{eq:III.7.6.1.1}{III.7.6.1.1.fr}
\EGAReferenceAlias{III.7.6.2-fr}{III.7.6.2}
\EGAReferenceAlias{III.7.6.2-es}{III.7.6.2}
\EGAReferenceAlias{III.7.6.2.fr}{III.7.6.2}
\EGAReferenceAlias{III.7.6.2.es}{III.7.6.2}
\EGAReferenceAlias{III.7.6.2.eq}{III.7.6.2}
\EGAReferenceAlias{eq:III.7.6.2}{III.7.6.2}
\EGAReferenceAlias{III.7.6.3.1}{III.7.6.3.1.fr}
\EGAReferenceAlias{III.7.6.3.1-fr}{III.7.6.3.1.fr}
\EGAReferenceAlias{III.7.6.3.1-es}{III.7.6.3.1.fr}
\EGAReferenceAlias{III.7.6.3.1.es}{III.7.6.3.1.fr}
\EGAReferenceAlias{III.7.6.3.1.eq}{III.7.6.3.1.fr}
\EGAReferenceAlias{eq:III.7.6.3.1}{III.7.6.3.1.fr}
\EGAReferenceAlias{III.7.6.3.2}{III.7.6.3.2.fr}
\EGAReferenceAlias{III.7.6.3.2-fr}{III.7.6.3.2.fr}
\EGAReferenceAlias{III.7.6.3.2-es}{III.7.6.3.2.fr}
\EGAReferenceAlias{III.7.6.3.2.es}{III.7.6.3.2.fr}
\EGAReferenceAlias{III.7.6.3.2.eq}{III.7.6.3.2.fr}
\EGAReferenceAlias{eq:III.7.6.3.2}{III.7.6.3.2.fr}
\EGAReferenceAlias{III.7.6.3.3}{III.7.6.3.3.fr}
\EGAReferenceAlias{III.7.6.3.3-fr}{III.7.6.3.3.fr}
\EGAReferenceAlias{III.7.6.3.3-es}{III.7.6.3.3.fr}
\EGAReferenceAlias{III.7.6.3.3.es}{III.7.6.3.3.fr}
\EGAReferenceAlias{III.7.6.3.3.eq}{III.7.6.3.3.fr}
\EGAReferenceAlias{eq:III.7.6.3.3}{III.7.6.3.3.fr}
\EGAReferenceAlias{III.7.6.3.4}{III.7.6.3.4.fr}
\EGAReferenceAlias{III.7.6.3.4-fr}{III.7.6.3.4.fr}
\EGAReferenceAlias{III.7.6.3.4-es}{III.7.6.3.4.fr}
\EGAReferenceAlias{III.7.6.3.4.es}{III.7.6.3.4.fr}
\EGAReferenceAlias{III.7.6.3.4.eq}{III.7.6.3.4.fr}
\EGAReferenceAlias{eq:III.7.6.3.4}{III.7.6.3.4.fr}
\EGAReferenceAlias{III.7.6.4-fr}{III.7.6.4}
\EGAReferenceAlias{III.7.6.4-es}{III.7.6.4}
\EGAReferenceAlias{III.7.6.4.fr}{III.7.6.4}
\EGAReferenceAlias{III.7.6.4.es}{III.7.6.4}
\EGAReferenceAlias{III.7.6.4.eq}{III.7.6.4}
\EGAReferenceAlias{eq:III.7.6.4}{III.7.6.4}
\EGAReferenceAlias{III.7.6.5-fr}{III.7.6.5}
\EGAReferenceAlias{III.7.6.5-es}{III.7.6.5}
\EGAReferenceAlias{III.7.6.5.fr}{III.7.6.5}
\EGAReferenceAlias{III.7.6.5.es}{III.7.6.5}
\EGAReferenceAlias{III.7.6.5.eq}{III.7.6.5}
\EGAReferenceAlias{eq:III.7.6.5}{III.7.6.5}
\EGAReferenceAlias{III.7.6.5.1}{III.7.6.5.1.fr}
\EGAReferenceAlias{III.7.6.5.1-fr}{III.7.6.5.1.fr}
\EGAReferenceAlias{III.7.6.5.1-es}{III.7.6.5.1.fr}
\EGAReferenceAlias{III.7.6.5.1.es}{III.7.6.5.1.fr}
\EGAReferenceAlias{III.7.6.5.1.eq}{III.7.6.5.1.fr}
\EGAReferenceAlias{eq:III.7.6.5.1}{III.7.6.5.1.fr}
\EGAReferenceAlias{III.7.6.5.2}{III.7.6.5.2.fr}
\EGAReferenceAlias{III.7.6.5.2-fr}{III.7.6.5.2.fr}
\EGAReferenceAlias{III.7.6.5.2-es}{III.7.6.5.2.fr}
\EGAReferenceAlias{III.7.6.5.2.es}{III.7.6.5.2.fr}
\EGAReferenceAlias{III.7.6.5.2.eq}{III.7.6.5.2.fr}
\EGAReferenceAlias{eq:III.7.6.5.2}{III.7.6.5.2.fr}
\EGAReferenceAlias{III.7.6.6.1}{III.7.6.6.1.fr}
\EGAReferenceAlias{III.7.6.6.1-fr}{III.7.6.6.1.fr}
\EGAReferenceAlias{III.7.6.6.1-es}{III.7.6.6.1.fr}
\EGAReferenceAlias{III.7.6.6.1.es}{III.7.6.6.1.fr}
\EGAReferenceAlias{III.7.6.6.1.eq}{III.7.6.6.1.fr}
\EGAReferenceAlias{eq:III.7.6.6.1}{III.7.6.6.1.fr}
\EGAReferenceAlias{III.7.6.7-fr}{III.7.6.7}
\EGAReferenceAlias{III.7.6.7-es}{III.7.6.7}
\EGAReferenceAlias{III.7.6.7.fr}{III.7.6.7}
\EGAReferenceAlias{III.7.6.7.es}{III.7.6.7}
\EGAReferenceAlias{III.7.6.7.eq}{III.7.6.7}
\EGAReferenceAlias{eq:III.7.6.7}{III.7.6.7}
\EGAReferenceAlias{III.7.6.7.1}{III.7.6.7.1.fr}
\EGAReferenceAlias{III.7.6.7.1-fr}{III.7.6.7.1.fr}
\EGAReferenceAlias{III.7.6.7.1-es}{III.7.6.7.1.fr}
\EGAReferenceAlias{III.7.6.7.1.es}{III.7.6.7.1.fr}
\EGAReferenceAlias{III.7.6.7.1.eq}{III.7.6.7.1.fr}
\EGAReferenceAlias{eq:III.7.6.7.1}{III.7.6.7.1.fr}
\EGAReferenceAlias{III.7.6.8-fr}{III.7.6.8}
\EGAReferenceAlias{III.7.6.8-es}{III.7.6.8}
\EGAReferenceAlias{III.7.6.8.fr}{III.7.6.8}
\EGAReferenceAlias{III.7.6.8.es}{III.7.6.8}
\EGAReferenceAlias{III.7.6.8.eq}{III.7.6.8}
\EGAReferenceAlias{eq:III.7.6.8}{III.7.6.8}
\EGAReferenceAlias{III.7.6.9-fr}{III.7.6.9}
\EGAReferenceAlias{III.7.6.9-es}{III.7.6.9}
\EGAReferenceAlias{III.7.6.9.fr}{III.7.6.9}
\EGAReferenceAlias{III.7.6.9.es}{III.7.6.9}
\EGAReferenceAlias{III.7.6.9.eq}{III.7.6.9}
\EGAReferenceAlias{eq:III.7.6.9}{III.7.6.9}
\EGAReferenceAlias{III.7.6.10-fr}{III.7.6.10}
\EGAReferenceAlias{III.7.6.10-es}{III.7.6.10}
\EGAReferenceAlias{III.7.6.10.fr}{III.7.6.10}
\EGAReferenceAlias{III.7.6.10.es}{III.7.6.10}
\EGAReferenceAlias{III.7.6.10.eq}{III.7.6.10}
\EGAReferenceAlias{eq:III.7.6.10}{III.7.6.10}
\EGAReferenceAlias{III.7.6.10.1}{III.7.6.10.1.fr}
\EGAReferenceAlias{III.7.6.10.1-fr}{III.7.6.10.1.fr}
\EGAReferenceAlias{III.7.6.10.1-es}{III.7.6.10.1.fr}
\EGAReferenceAlias{III.7.6.10.1.es}{III.7.6.10.1.fr}
\EGAReferenceAlias{III.7.6.10.1.eq}{III.7.6.10.1.fr}
\EGAReferenceAlias{eq:III.7.6.10.1}{III.7.6.10.1.fr}
\EGAReferenceAlias{III.7.6.10.2}{III.7.6.10.2.fr}
\EGAReferenceAlias{III.7.6.10.2-fr}{III.7.6.10.2.fr}
\EGAReferenceAlias{III.7.6.10.2-es}{III.7.6.10.2.fr}
\EGAReferenceAlias{III.7.6.10.2.es}{III.7.6.10.2.fr}
\EGAReferenceAlias{III.7.6.10.2.eq}{III.7.6.10.2.fr}
\EGAReferenceAlias{eq:III.7.6.10.2}{III.7.6.10.2.fr}
\EGAReferenceAlias{III.7.6.10.3}{III.7.6.10.3.fr}
\EGAReferenceAlias{III.7.6.10.3-fr}{III.7.6.10.3.fr}
\EGAReferenceAlias{III.7.6.10.3-es}{III.7.6.10.3.fr}
\EGAReferenceAlias{III.7.6.10.3.es}{III.7.6.10.3.fr}
\EGAReferenceAlias{III.7.6.10.3.eq}{III.7.6.10.3.fr}
\EGAReferenceAlias{eq:III.7.6.10.3}{III.7.6.10.3.fr}
\EGAReferenceAlias{III.7.6.11.1}{III.7.6.11.1.fr}
\EGAReferenceAlias{III.7.6.11.1-fr}{III.7.6.11.1.fr}
\EGAReferenceAlias{III.7.6.11.1-es}{III.7.6.11.1.fr}
\EGAReferenceAlias{III.7.6.11.1.es}{III.7.6.11.1.fr}
\EGAReferenceAlias{III.7.6.11.1.eq}{III.7.6.11.1.fr}
\EGAReferenceAlias{eq:III.7.6.11.1}{III.7.6.11.1.fr}
\EGAReferenceAlias{III.7.6.11.2}{III.7.6.11.2.fr}
\EGAReferenceAlias{III.7.6.11.2-fr}{III.7.6.11.2.fr}
\EGAReferenceAlias{III.7.6.11.2-es}{III.7.6.11.2.fr}
\EGAReferenceAlias{III.7.6.11.2.es}{III.7.6.11.2.fr}
\EGAReferenceAlias{III.7.6.11.2.eq}{III.7.6.11.2.fr}
\EGAReferenceAlias{eq:III.7.6.11.2}{III.7.6.11.2.fr}
\EGAReferenceAlias{III.7.6.11.3}{III.7.6.11.3.fr}
\EGAReferenceAlias{III.7.6.11.3-fr}{III.7.6.11.3.fr}
\EGAReferenceAlias{III.7.6.11.3-es}{III.7.6.11.3.fr}
\EGAReferenceAlias{III.7.6.11.3.es}{III.7.6.11.3.fr}
\EGAReferenceAlias{III.7.6.11.3.eq}{III.7.6.11.3.fr}
\EGAReferenceAlias{eq:III.7.6.11.3}{III.7.6.11.3.fr}
\EGAReferenceAlias{III.7.6.11.4}{III.7.6.11.4.fr}
\EGAReferenceAlias{III.7.6.11.4-fr}{III.7.6.11.4.fr}
\EGAReferenceAlias{III.7.6.11.4-es}{III.7.6.11.4.fr}
\EGAReferenceAlias{III.7.6.11.4.es}{III.7.6.11.4.fr}
\EGAReferenceAlias{III.7.6.11.4.eq}{III.7.6.11.4.fr}
\EGAReferenceAlias{eq:III.7.6.11.4}{III.7.6.11.4.fr}
\EGAReferenceAlias{III.7.6.11.5}{III.7.6.11.5.fr}
\EGAReferenceAlias{III.7.6.11.5-fr}{III.7.6.11.5.fr}
\EGAReferenceAlias{III.7.6.11.5-es}{III.7.6.11.5.fr}
\EGAReferenceAlias{III.7.6.11.5.es}{III.7.6.11.5.fr}
\EGAReferenceAlias{III.7.6.11.5.eq}{III.7.6.11.5.fr}
\EGAReferenceAlias{eq:III.7.6.11.5}{III.7.6.11.5.fr}
\EGAReferenceAlias{III.7.6.11.7}{III.7.6.11.7.fr}
\EGAReferenceAlias{III.7.6.11.7-fr}{III.7.6.11.7.fr}
\EGAReferenceAlias{III.7.6.11.7-es}{III.7.6.11.7.fr}
\EGAReferenceAlias{III.7.6.11.7.es}{III.7.6.11.7.fr}
\EGAReferenceAlias{III.7.6.11.7.eq}{III.7.6.11.7.fr}
\EGAReferenceAlias{eq:III.7.6.11.7}{III.7.6.11.7.fr}
\EGAReferenceAlias{III.7.7}{subsection:III.7.7.fr}
\EGAReferenceAlias{III.7.7-fr}{subsection:III.7.7.fr}
\EGAReferenceAlias{III.7.7-es}{subsection:III.7.7.fr}
\EGAReferenceAlias{III.7.7.fr}{subsection:III.7.7.fr}
\EGAReferenceAlias{III.7.7.es}{subsection:III.7.7.fr}
\EGAReferenceAlias{III.7.7.eq}{subsection:III.7.7.fr}
\EGAReferenceAlias{eq:III.7.7}{subsection:III.7.7.fr}
\EGAReferenceAlias{subsection:III.7.7}{subsection:III.7.7.fr}
\EGAReferenceAlias{subsection:III.7.7-fr}{subsection:III.7.7.fr}
\EGAReferenceAlias{subsection:III.7.7-es}{subsection:III.7.7.fr}
\EGAReferenceAlias{subsection:III.7.7.es}{subsection:III.7.7.fr}
\EGAReferenceAlias{subsection:III.7.7.eq}{subsection:III.7.7.fr}
\EGAReferenceAlias{eq:subsection:III.7.7}{subsection:III.7.7.fr}
\EGAReferenceAlias{III.7.7.1-fr}{III.7.7.1}
\EGAReferenceAlias{III.7.7.1-es}{III.7.7.1}
\EGAReferenceAlias{III.7.7.1.fr}{III.7.7.1}
\EGAReferenceAlias{III.7.7.1.es}{III.7.7.1}
\EGAReferenceAlias{III.7.7.1.eq}{III.7.7.1}
\EGAReferenceAlias{eq:III.7.7.1}{III.7.7.1}
\EGAReferenceAlias{III.7.7.1.1}{III.7.7.1.1.fr}
\EGAReferenceAlias{III.7.7.1.1-fr}{III.7.7.1.1.fr}
\EGAReferenceAlias{III.7.7.1.1-es}{III.7.7.1.1.fr}
\EGAReferenceAlias{III.7.7.1.1.es}{III.7.7.1.1.fr}
\EGAReferenceAlias{III.7.7.1.1.eq}{III.7.7.1.1.fr}
\EGAReferenceAlias{eq:III.7.7.1.1}{III.7.7.1.1.fr}
\EGAReferenceAlias{III.7.7.2-fr}{III.7.7.2}
\EGAReferenceAlias{III.7.7.2-es}{III.7.7.2}
\EGAReferenceAlias{III.7.7.2.fr}{III.7.7.2}
\EGAReferenceAlias{III.7.7.2.es}{III.7.7.2}
\EGAReferenceAlias{III.7.7.2.eq}{III.7.7.2}
\EGAReferenceAlias{eq:III.7.7.2}{III.7.7.2}
\EGAReferenceAlias{III.7.7.2.1}{III.7.7.2.1.fr}
\EGAReferenceAlias{III.7.7.2.1-fr}{III.7.7.2.1.fr}
\EGAReferenceAlias{III.7.7.2.1-es}{III.7.7.2.1.fr}
\EGAReferenceAlias{III.7.7.2.1.es}{III.7.7.2.1.fr}
\EGAReferenceAlias{III.7.7.2.1.eq}{III.7.7.2.1.fr}
\EGAReferenceAlias{eq:III.7.7.2.1}{III.7.7.2.1.fr}
\EGAReferenceAlias{III.7.7.2.2}{III.7.7.2.2.fr}
\EGAReferenceAlias{III.7.7.2.2-fr}{III.7.7.2.2.fr}
\EGAReferenceAlias{III.7.7.2.2-es}{III.7.7.2.2.fr}
\EGAReferenceAlias{III.7.7.2.2.es}{III.7.7.2.2.fr}
\EGAReferenceAlias{III.7.7.2.2.eq}{III.7.7.2.2.fr}
\EGAReferenceAlias{eq:III.7.7.2.2}{III.7.7.2.2.fr}
\EGAReferenceAlias{III.7.7.3-fr}{III.7.7.3}
\EGAReferenceAlias{III.7.7.3-es}{III.7.7.3}
\EGAReferenceAlias{III.7.7.3.fr}{III.7.7.3}
\EGAReferenceAlias{III.7.7.3.es}{III.7.7.3}
\EGAReferenceAlias{III.7.7.3.eq}{III.7.7.3}
\EGAReferenceAlias{eq:III.7.7.3}{III.7.7.3}
\EGAReferenceAlias{III.7.7.3.1}{III.7.7.3.1.fr}
\EGAReferenceAlias{III.7.7.3.1-fr}{III.7.7.3.1.fr}
\EGAReferenceAlias{III.7.7.3.1-es}{III.7.7.3.1.fr}
\EGAReferenceAlias{III.7.7.3.1.es}{III.7.7.3.1.fr}
\EGAReferenceAlias{III.7.7.3.1.eq}{III.7.7.3.1.fr}
\EGAReferenceAlias{eq:III.7.7.3.1}{III.7.7.3.1.fr}
\EGAReferenceAlias{III.7.7.3.2}{III.7.7.3.2.fr}
\EGAReferenceAlias{III.7.7.3.2-fr}{III.7.7.3.2.fr}
\EGAReferenceAlias{III.7.7.3.2-es}{III.7.7.3.2.fr}
\EGAReferenceAlias{III.7.7.3.2.es}{III.7.7.3.2.fr}
\EGAReferenceAlias{III.7.7.3.2.eq}{III.7.7.3.2.fr}
\EGAReferenceAlias{eq:III.7.7.3.2}{III.7.7.3.2.fr}
\EGAReferenceAlias{III.7.7.4-fr}{III.7.7.4}
\EGAReferenceAlias{III.7.7.4-es}{III.7.7.4}
\EGAReferenceAlias{III.7.7.4.fr}{III.7.7.4}
\EGAReferenceAlias{III.7.7.4.es}{III.7.7.4}
\EGAReferenceAlias{III.7.7.4.eq}{III.7.7.4}
\EGAReferenceAlias{eq:III.7.7.4}{III.7.7.4}
\EGAReferenceAlias{III.7.7.5-fr}{III.7.7.5}
\EGAReferenceAlias{III.7.7.5-es}{III.7.7.5}
\EGAReferenceAlias{III.7.7.5.fr}{III.7.7.5}
\EGAReferenceAlias{III.7.7.5.es}{III.7.7.5}
\EGAReferenceAlias{III.7.7.5.eq}{III.7.7.5}
\EGAReferenceAlias{eq:III.7.7.5}{III.7.7.5}
\EGAReferenceAlias{III.7.7.5.1}{III.7.7.5.1.fr}
\EGAReferenceAlias{III.7.7.5.1-fr}{III.7.7.5.1.fr}
\EGAReferenceAlias{III.7.7.5.1-es}{III.7.7.5.1.fr}
\EGAReferenceAlias{III.7.7.5.1.es}{III.7.7.5.1.fr}
\EGAReferenceAlias{III.7.7.5.1.eq}{III.7.7.5.1.fr}
\EGAReferenceAlias{eq:III.7.7.5.1}{III.7.7.5.1.fr}
\EGAReferenceAlias{III.7.7.5.2}{III.7.7.5.2.fr}
\EGAReferenceAlias{III.7.7.5.2-fr}{III.7.7.5.2.fr}
\EGAReferenceAlias{III.7.7.5.2-es}{III.7.7.5.2.fr}
\EGAReferenceAlias{III.7.7.5.2.es}{III.7.7.5.2.fr}
\EGAReferenceAlias{III.7.7.5.2.eq}{III.7.7.5.2.fr}
\EGAReferenceAlias{eq:III.7.7.5.2}{III.7.7.5.2.fr}
\EGAReferenceAlias{III.7.7.5.3}{III.7.7.5.3.fr}
\EGAReferenceAlias{III.7.7.5.3-fr}{III.7.7.5.3.fr}
\EGAReferenceAlias{III.7.7.5.3-es}{III.7.7.5.3.fr}
\EGAReferenceAlias{III.7.7.5.3.es}{III.7.7.5.3.fr}
\EGAReferenceAlias{III.7.7.5.3.eq}{III.7.7.5.3.fr}
\EGAReferenceAlias{eq:III.7.7.5.3}{III.7.7.5.3.fr}
\EGAReferenceAlias{III.7.7.6-fr}{III.7.7.6}
\EGAReferenceAlias{III.7.7.6-es}{III.7.7.6}
\EGAReferenceAlias{III.7.7.6.fr}{III.7.7.6}
\EGAReferenceAlias{III.7.7.6.es}{III.7.7.6}
\EGAReferenceAlias{III.7.7.6.eq}{III.7.7.6}
\EGAReferenceAlias{eq:III.7.7.6}{III.7.7.6}
\EGAReferenceAlias{III.7.7.6.1}{III.7.7.6.1.fr}
\EGAReferenceAlias{III.7.7.6.1-fr}{III.7.7.6.1.fr}
\EGAReferenceAlias{III.7.7.6.1-es}{III.7.7.6.1.fr}
\EGAReferenceAlias{III.7.7.6.1.es}{III.7.7.6.1.fr}
\EGAReferenceAlias{III.7.7.6.1.eq}{III.7.7.6.1.fr}
\EGAReferenceAlias{eq:III.7.7.6.1}{III.7.7.6.1.fr}
\EGAReferenceAlias{III.7.7.6.2}{III.7.7.6.2.fr}
\EGAReferenceAlias{III.7.7.6.2-fr}{III.7.7.6.2.fr}
\EGAReferenceAlias{III.7.7.6.2-es}{III.7.7.6.2.fr}
\EGAReferenceAlias{III.7.7.6.2.es}{III.7.7.6.2.fr}
\EGAReferenceAlias{III.7.7.6.2.eq}{III.7.7.6.2.fr}
\EGAReferenceAlias{eq:III.7.7.6.2}{III.7.7.6.2.fr}
\EGAReferenceAlias{III.7.7.7-fr}{III.7.7.7}
\EGAReferenceAlias{III.7.7.7-es}{III.7.7.7}
\EGAReferenceAlias{III.7.7.7.fr}{III.7.7.7}
\EGAReferenceAlias{III.7.7.7.es}{III.7.7.7}
\EGAReferenceAlias{III.7.7.7.eq}{III.7.7.7}
\EGAReferenceAlias{eq:III.7.7.7}{III.7.7.7}
\EGAReferenceAlias{III.7.7.7.1}{III.7.7.7.1.fr}
\EGAReferenceAlias{III.7.7.7.1-fr}{III.7.7.7.1.fr}
\EGAReferenceAlias{III.7.7.7.1-es}{III.7.7.7.1.fr}
\EGAReferenceAlias{III.7.7.7.1.es}{III.7.7.7.1.fr}
\EGAReferenceAlias{III.7.7.7.1.eq}{III.7.7.7.1.fr}
\EGAReferenceAlias{eq:III.7.7.7.1}{III.7.7.7.1.fr}
\EGAReferenceAlias{III.7.7.7.2}{III.7.7.7.2.fr}
\EGAReferenceAlias{III.7.7.7.2-fr}{III.7.7.7.2.fr}
\EGAReferenceAlias{III.7.7.7.2-es}{III.7.7.7.2.fr}
\EGAReferenceAlias{III.7.7.7.2.es}{III.7.7.7.2.fr}
\EGAReferenceAlias{III.7.7.7.2.eq}{III.7.7.7.2.fr}
\EGAReferenceAlias{eq:III.7.7.7.2}{III.7.7.7.2.fr}
\EGAReferenceAlias{III.7.7.8.1}{III.7.7.8.1.fr}
\EGAReferenceAlias{III.7.7.8.1-fr}{III.7.7.8.1.fr}
\EGAReferenceAlias{III.7.7.8.1-es}{III.7.7.8.1.fr}
\EGAReferenceAlias{III.7.7.8.1.es}{III.7.7.8.1.fr}
\EGAReferenceAlias{III.7.7.8.1.eq}{III.7.7.8.1.fr}
\EGAReferenceAlias{eq:III.7.7.8.1}{III.7.7.8.1.fr}
\EGAReferenceAlias{III.7.7.8.2}{III.7.7.8.2.fr}
\EGAReferenceAlias{III.7.7.8.2-fr}{III.7.7.8.2.fr}
\EGAReferenceAlias{III.7.7.8.2-es}{III.7.7.8.2.fr}
\EGAReferenceAlias{III.7.7.8.2.es}{III.7.7.8.2.fr}
\EGAReferenceAlias{III.7.7.8.2.eq}{III.7.7.8.2.fr}
\EGAReferenceAlias{eq:III.7.7.8.2}{III.7.7.8.2.fr}
\EGAReferenceAlias{III.7.7.9-fr}{III.7.7.9}
\EGAReferenceAlias{III.7.7.9-es}{III.7.7.9}
\EGAReferenceAlias{III.7.7.9.fr}{III.7.7.9}
\EGAReferenceAlias{III.7.7.9.es}{III.7.7.9}
\EGAReferenceAlias{III.7.7.9.eq}{III.7.7.9}
\EGAReferenceAlias{eq:III.7.7.9}{III.7.7.9}
\EGAReferenceAlias{III.7.7.10-fr}{III.7.7.10}
\EGAReferenceAlias{III.7.7.10-es}{III.7.7.10}
\EGAReferenceAlias{III.7.7.10.fr}{III.7.7.10}
\EGAReferenceAlias{III.7.7.10.es}{III.7.7.10}
\EGAReferenceAlias{III.7.7.10.eq}{III.7.7.10}
\EGAReferenceAlias{eq:III.7.7.10}{III.7.7.10}
\EGAReferenceAlias{III.7.7.11-fr}{III.7.7.11}
\EGAReferenceAlias{III.7.7.11-es}{III.7.7.11}
\EGAReferenceAlias{III.7.7.11.fr}{III.7.7.11}
\EGAReferenceAlias{III.7.7.11.es}{III.7.7.11}
\EGAReferenceAlias{III.7.7.11.eq}{III.7.7.11}
\EGAReferenceAlias{eq:III.7.7.11}{III.7.7.11}
\EGAReferenceAlias{III.7.7.12-fr}{III.7.7.12}
\EGAReferenceAlias{III.7.7.12-es}{III.7.7.12}
\EGAReferenceAlias{III.7.7.12.fr}{III.7.7.12}
\EGAReferenceAlias{III.7.7.12.es}{III.7.7.12}
\EGAReferenceAlias{III.7.7.12.eq}{III.7.7.12}
\EGAReferenceAlias{eq:III.7.7.12}{III.7.7.12}
\EGAReferenceAlias{III.7.7.12.1}{III.7.7.12.1.fr}
\EGAReferenceAlias{III.7.7.12.1-fr}{III.7.7.12.1.fr}
\EGAReferenceAlias{III.7.7.12.1-es}{III.7.7.12.1.fr}
\EGAReferenceAlias{III.7.7.12.1.es}{III.7.7.12.1.fr}
\EGAReferenceAlias{III.7.7.12.1.eq}{III.7.7.12.1.fr}
\EGAReferenceAlias{eq:III.7.7.12.1}{III.7.7.12.1.fr}
\EGAReferenceAlias{III.7.8-fr}{subsection:III.7.8}
\EGAReferenceAlias{III.7.8-es}{subsection:III.7.8}
\EGAReferenceAlias{III.7.8.fr}{subsection:III.7.8}
\EGAReferenceAlias{III.7.8.es}{subsection:III.7.8}
\EGAReferenceAlias{III.7.8.eq}{subsection:III.7.8}
\EGAReferenceAlias{eq:III.7.8}{subsection:III.7.8}
\EGAReferenceAlias{subsection:III.7.8-fr}{subsection:III.7.8}
\EGAReferenceAlias{subsection:III.7.8-es}{subsection:III.7.8}
\EGAReferenceAlias{subsection:III.7.8.es}{subsection:III.7.8}
\EGAReferenceAlias{subsection:III.7.8.eq}{subsection:III.7.8}
\EGAReferenceAlias{eq:subsection:III.7.8}{subsection:III.7.8}
\EGAReferenceAlias{III.7.8.1-fr}{III.7.8.1}
\EGAReferenceAlias{III.7.8.1-es}{III.7.8.1}
\EGAReferenceAlias{III.7.8.1.fr}{III.7.8.1}
\EGAReferenceAlias{III.7.8.1.es}{III.7.8.1}
\EGAReferenceAlias{III.7.8.1.eq}{III.7.8.1}
\EGAReferenceAlias{eq:III.7.8.1}{III.7.8.1}
\EGAReferenceAlias{III.7.8.3-fr}{III.7.8.3}
\EGAReferenceAlias{III.7.8.3-es}{III.7.8.3}
\EGAReferenceAlias{III.7.8.3.fr}{III.7.8.3}
\EGAReferenceAlias{III.7.8.3.es}{III.7.8.3}
\EGAReferenceAlias{III.7.8.3.eq}{III.7.8.3}
\EGAReferenceAlias{eq:III.7.8.3}{III.7.8.3}
\EGAReferenceAlias{III.7.8.4-fr}{III.7.8.4}
\EGAReferenceAlias{III.7.8.4-es}{III.7.8.4}
\EGAReferenceAlias{III.7.8.4.fr}{III.7.8.4}
\EGAReferenceAlias{III.7.8.4.es}{III.7.8.4}
\EGAReferenceAlias{III.7.8.4.eq}{III.7.8.4}
\EGAReferenceAlias{eq:III.7.8.4}{III.7.8.4}
\EGAReferenceAlias{III.7.8.4.1}{III.7.8.4.1.fr}
\EGAReferenceAlias{III.7.8.4.1-fr}{III.7.8.4.1.fr}
\EGAReferenceAlias{III.7.8.4.1-es}{III.7.8.4.1.fr}
\EGAReferenceAlias{III.7.8.4.1.es}{III.7.8.4.1.fr}
\EGAReferenceAlias{III.7.8.4.1.eq}{III.7.8.4.1.fr}
\EGAReferenceAlias{eq:III.7.8.4.1}{III.7.8.4.1.fr}
\EGAReferenceAlias{III.7.8.4.2}{III.7.8.4.2.fr}
\EGAReferenceAlias{III.7.8.4.2-fr}{III.7.8.4.2.fr}
\EGAReferenceAlias{III.7.8.4.2-es}{III.7.8.4.2.fr}
\EGAReferenceAlias{III.7.8.4.2.es}{III.7.8.4.2.fr}
\EGAReferenceAlias{III.7.8.4.2.eq}{III.7.8.4.2.fr}
\EGAReferenceAlias{eq:III.7.8.4.2}{III.7.8.4.2.fr}
\EGAReferenceAlias{III.7.8.5-fr}{III.7.8.5}
\EGAReferenceAlias{III.7.8.5-es}{III.7.8.5}
\EGAReferenceAlias{III.7.8.5.fr}{III.7.8.5}
\EGAReferenceAlias{III.7.8.5.es}{III.7.8.5}
\EGAReferenceAlias{III.7.8.5.eq}{III.7.8.5}
\EGAReferenceAlias{eq:III.7.8.5}{III.7.8.5}
\EGAReferenceAlias{III.7.8.7-fr}{III.7.8.7}
\EGAReferenceAlias{III.7.8.7-es}{III.7.8.7}
\EGAReferenceAlias{III.7.8.7.fr}{III.7.8.7}
\EGAReferenceAlias{III.7.8.7.es}{III.7.8.7}
\EGAReferenceAlias{III.7.8.7.eq}{III.7.8.7}
\EGAReferenceAlias{eq:III.7.8.7}{III.7.8.7}
\EGAReferenceAlias{III.7.8.8-fr}{III.7.8.8}
\EGAReferenceAlias{III.7.8.8-es}{III.7.8.8}
\EGAReferenceAlias{III.7.8.8.fr}{III.7.8.8}
\EGAReferenceAlias{III.7.8.8.es}{III.7.8.8}
\EGAReferenceAlias{III.7.8.8.eq}{III.7.8.8}
\EGAReferenceAlias{eq:III.7.8.8}{III.7.8.8}
\EGAReferenceAlias{III.7.8.9-fr}{III.7.8.9}
\EGAReferenceAlias{III.7.8.9-es}{III.7.8.9}
\EGAReferenceAlias{III.7.8.9.fr}{III.7.8.9}
\EGAReferenceAlias{III.7.8.9.es}{III.7.8.9}
\EGAReferenceAlias{III.7.8.9.eq}{III.7.8.9}
\EGAReferenceAlias{eq:III.7.8.9}{III.7.8.9}
\EGAReferenceAlias{III.7.8.9.1}{III.7.8.9.1.fr}
\EGAReferenceAlias{III.7.8.9.1-fr}{III.7.8.9.1.fr}
\EGAReferenceAlias{III.7.8.9.1-es}{III.7.8.9.1.fr}
\EGAReferenceAlias{III.7.8.9.1.es}{III.7.8.9.1.fr}
\EGAReferenceAlias{III.7.8.9.1.eq}{III.7.8.9.1.fr}
\EGAReferenceAlias{eq:III.7.8.9.1}{III.7.8.9.1.fr}
\EGAReferenceAlias{III.7.8.10-fr}{III.7.8.10}
\EGAReferenceAlias{III.7.8.10-es}{III.7.8.10}
\EGAReferenceAlias{III.7.8.10.fr}{III.7.8.10}
\EGAReferenceAlias{III.7.8.10.es}{III.7.8.10}
\EGAReferenceAlias{III.7.8.10.eq}{III.7.8.10}
\EGAReferenceAlias{eq:III.7.8.10}{III.7.8.10}
\EGAReferenceAlias{III.7.9-fr}{subsection:III.7.9}
\EGAReferenceAlias{III.7.9-es}{subsection:III.7.9}
\EGAReferenceAlias{III.7.9.es}{subsection:III.7.9}
\EGAReferenceAlias{III.7.9.eq}{subsection:III.7.9}
\EGAReferenceAlias{eq:III.7.9}{subsection:III.7.9}
\EGAReferenceAlias{subsection:III.7.9-fr}{subsection:III.7.9}
\EGAReferenceAlias{subsection:III.7.9-es}{subsection:III.7.9}
\EGAReferenceAlias{subsection:III.7.9.es}{subsection:III.7.9}
\EGAReferenceAlias{subsection:III.7.9.eq}{subsection:III.7.9}
\EGAReferenceAlias{eq:subsection:III.7.9}{subsection:III.7.9}
\EGAReferenceAlias{III.7.9.1-fr}{III.7.9.1}
\EGAReferenceAlias{III.7.9.1-es}{III.7.9.1}
\EGAReferenceAlias{III.7.9.1.es}{III.7.9.1}
\EGAReferenceAlias{III.7.9.1.eq}{III.7.9.1}
\EGAReferenceAlias{eq:III.7.9.1}{III.7.9.1}
\EGAReferenceAlias{III.7.9.1.1}{III.7.9.1.1.fr}
\EGAReferenceAlias{III.7.9.1.1-fr}{III.7.9.1.1.fr}
\EGAReferenceAlias{III.7.9.1.1-es}{III.7.9.1.1.fr}
\EGAReferenceAlias{III.7.9.1.1.es}{III.7.9.1.1.fr}
\EGAReferenceAlias{III.7.9.1.1.eq}{III.7.9.1.1.fr}
\EGAReferenceAlias{eq:III.7.9.1.1}{III.7.9.1.1.fr}
\EGAReferenceAlias{III.7.9.2-fr}{III.7.9.2}
\EGAReferenceAlias{III.7.9.2-es}{III.7.9.2}
\EGAReferenceAlias{III.7.9.2.es}{III.7.9.2}
\EGAReferenceAlias{III.7.9.2.eq}{III.7.9.2}
\EGAReferenceAlias{eq:III.7.9.2}{III.7.9.2}
\EGAReferenceAlias{III.7.9.2.1}{III.7.9.2.1.fr}
\EGAReferenceAlias{III.7.9.2.1-fr}{III.7.9.2.1.fr}
\EGAReferenceAlias{III.7.9.2.1-es}{III.7.9.2.1.fr}
\EGAReferenceAlias{III.7.9.2.1.es}{III.7.9.2.1.fr}
\EGAReferenceAlias{III.7.9.2.1.eq}{III.7.9.2.1.fr}
\EGAReferenceAlias{eq:III.7.9.2.1}{III.7.9.2.1.fr}
\EGAReferenceAlias{III.7.9.3-fr}{III.7.9.3}
\EGAReferenceAlias{III.7.9.3-es}{III.7.9.3}
\EGAReferenceAlias{III.7.9.3.es}{III.7.9.3}
\EGAReferenceAlias{III.7.9.3.eq}{III.7.9.3}
\EGAReferenceAlias{eq:III.7.9.3}{III.7.9.3}
\EGAReferenceAlias{III.7.9.4-fr}{III.7.9.4}
\EGAReferenceAlias{III.7.9.4-es}{III.7.9.4}
\EGAReferenceAlias{III.7.9.4.es}{III.7.9.4}
\EGAReferenceAlias{III.7.9.4.eq}{III.7.9.4}
\EGAReferenceAlias{eq:III.7.9.4}{III.7.9.4}
\EGAReferenceAlias{III.7.9.4.1}{III.7.9.4.1.fr}
\EGAReferenceAlias{III.7.9.4.1-fr}{III.7.9.4.1.fr}
\EGAReferenceAlias{III.7.9.4.1-es}{III.7.9.4.1.fr}
\EGAReferenceAlias{III.7.9.4.1.es}{III.7.9.4.1.fr}
\EGAReferenceAlias{III.7.9.4.1.eq}{III.7.9.4.1.fr}
\EGAReferenceAlias{eq:III.7.9.4.1}{III.7.9.4.1.fr}
\EGAReferenceAlias{III.7.9.5-fr}{III.7.9.5}
\EGAReferenceAlias{III.7.9.5-es}{III.7.9.5}
\EGAReferenceAlias{III.7.9.5.es}{III.7.9.5}
\EGAReferenceAlias{III.7.9.5.eq}{III.7.9.5}
\EGAReferenceAlias{eq:III.7.9.5}{III.7.9.5}
\EGAReferenceAlias{III.7.9.6-fr}{III.7.9.6}
\EGAReferenceAlias{III.7.9.6-es}{III.7.9.6}
\EGAReferenceAlias{III.7.9.6.es}{III.7.9.6}
\EGAReferenceAlias{III.7.9.6.eq}{III.7.9.6}
\EGAReferenceAlias{eq:III.7.9.6}{III.7.9.6}
\EGAReferenceAlias{III.7.9.6.1}{III.7.9.6.1.fr}
\EGAReferenceAlias{III.7.9.6.1-fr}{III.7.9.6.1.fr}
\EGAReferenceAlias{III.7.9.6.1-es}{III.7.9.6.1.fr}
\EGAReferenceAlias{III.7.9.6.1.es}{III.7.9.6.1.fr}
\EGAReferenceAlias{III.7.9.6.1.eq}{III.7.9.6.1.fr}
\EGAReferenceAlias{eq:III.7.9.6.1}{III.7.9.6.1.fr}
\EGAReferenceAlias{III.7.9.6.2}{III.7.9.6.2.fr}
\EGAReferenceAlias{III.7.9.6.2-fr}{III.7.9.6.2.fr}
\EGAReferenceAlias{III.7.9.6.2-es}{III.7.9.6.2.fr}
\EGAReferenceAlias{III.7.9.6.2.es}{III.7.9.6.2.fr}
\EGAReferenceAlias{III.7.9.6.2.eq}{III.7.9.6.2.fr}
\EGAReferenceAlias{eq:III.7.9.6.2}{III.7.9.6.2.fr}
\EGAReferenceAlias{III.7.9.7-fr}{III.7.9.7}
\EGAReferenceAlias{III.7.9.7-es}{III.7.9.7}
\EGAReferenceAlias{III.7.9.7.es}{III.7.9.7}
\EGAReferenceAlias{III.7.9.7.eq}{III.7.9.7}
\EGAReferenceAlias{eq:III.7.9.7}{III.7.9.7}
\EGAReferenceAlias{III.7.9.7.1}{III.7.9.7.1.fr}
\EGAReferenceAlias{III.7.9.7.1-fr}{III.7.9.7.1.fr}
\EGAReferenceAlias{III.7.9.7.1-es}{III.7.9.7.1.fr}
\EGAReferenceAlias{III.7.9.7.1.es}{III.7.9.7.1.fr}
\EGAReferenceAlias{III.7.9.7.1.eq}{III.7.9.7.1.fr}
\EGAReferenceAlias{eq:III.7.9.7.1}{III.7.9.7.1.fr}
\EGAReferenceAlias{III.7.9.8-fr}{III.7.9.8}
\EGAReferenceAlias{III.7.9.8-es}{III.7.9.8}
\EGAReferenceAlias{III.7.9.8.es}{III.7.9.8}
\EGAReferenceAlias{III.7.9.8.eq}{III.7.9.8}
\EGAReferenceAlias{eq:III.7.9.8}{III.7.9.8}
\EGAReferenceAlias{III.7.9.9-fr}{III.7.9.9}
\EGAReferenceAlias{III.7.9.9-es}{III.7.9.9}
\EGAReferenceAlias{III.7.9.9.es}{III.7.9.9}
\EGAReferenceAlias{III.7.9.9.eq}{III.7.9.9}
\EGAReferenceAlias{eq:III.7.9.9}{III.7.9.9}
\EGAReferenceAlias{III.7.9.10-fr}{III.7.9.10}
\EGAReferenceAlias{III.7.9.10-es}{III.7.9.10}
\EGAReferenceAlias{III.7.9.10.es}{III.7.9.10}
\EGAReferenceAlias{III.7.9.10.eq}{III.7.9.10}
\EGAReferenceAlias{eq:III.7.9.10}{III.7.9.10}
\EGAReferenceAlias{III.7.9.10.1-fr}{III.7.9.10.1}
\EGAReferenceAlias{III.7.9.10.1-es}{III.7.9.10.1}
\EGAReferenceAlias{III.7.9.10.1.es}{III.7.9.10.1}
\EGAReferenceAlias{III.7.9.10.1.eq}{III.7.9.10.1}
\EGAReferenceAlias{eq:III.7.9.10.1}{III.7.9.10.1}
\EGAReferenceAlias{III.7.9.11}{III.7.9.11.fr}
\EGAReferenceAlias{III.7.9.11-fr}{III.7.9.11.fr}
\EGAReferenceAlias{III.7.9.11-es}{III.7.9.11.fr}
\EGAReferenceAlias{III.7.9.11.es}{III.7.9.11.fr}
\EGAReferenceAlias{III.7.9.11.eq}{III.7.9.11.fr}
\EGAReferenceAlias{eq:III.7.9.11}{III.7.9.11.fr}
\EGAReferenceAlias{III.7.9.11.1}{III.7.9.11.1.fr}
\EGAReferenceAlias{III.7.9.11.1-fr}{III.7.9.11.1.fr}
\EGAReferenceAlias{III.7.9.11.1-es}{III.7.9.11.1.fr}
\EGAReferenceAlias{III.7.9.11.1.es}{III.7.9.11.1.fr}
\EGAReferenceAlias{III.7.9.11.1.eq}{III.7.9.11.1.fr}
\EGAReferenceAlias{eq:III.7.9.11.1}{III.7.9.11.1.fr}
\EGAReferenceAlias{III.7.9.12}{III.7.9.12.fr}
\EGAReferenceAlias{III.7.9.12-fr}{III.7.9.12.fr}
\EGAReferenceAlias{III.7.9.12-es}{III.7.9.12.fr}
\EGAReferenceAlias{III.7.9.12.es}{III.7.9.12.fr}
\EGAReferenceAlias{III.7.9.12.eq}{III.7.9.12.fr}
\EGAReferenceAlias{eq:III.7.9.12}{III.7.9.12.fr}
\EGAReferenceAlias{III.7.9.12.1}{III.7.9.12.1.fr}
\EGAReferenceAlias{III.7.9.12.1-fr}{III.7.9.12.1.fr}
\EGAReferenceAlias{III.7.9.12.1-es}{III.7.9.12.1.fr}
\EGAReferenceAlias{III.7.9.12.1.es}{III.7.9.12.1.fr}
\EGAReferenceAlias{III.7.9.12.1.eq}{III.7.9.12.1.fr}
\EGAReferenceAlias{eq:III.7.9.12.1}{III.7.9.12.1.fr}
\EGAReferenceAlias{III.7.9.12.2}{III.7.9.12.2.fr}
\EGAReferenceAlias{III.7.9.12.2-fr}{III.7.9.12.2.fr}
\EGAReferenceAlias{III.7.9.12.2-es}{III.7.9.12.2.fr}
\EGAReferenceAlias{III.7.9.12.2.es}{III.7.9.12.2.fr}
\EGAReferenceAlias{III.7.9.12.2.eq}{III.7.9.12.2.fr}
\EGAReferenceAlias{eq:III.7.9.12.2}{III.7.9.12.2.fr}
\EGAReferenceAlias{III.7.9.12.3}{III.7.9.12.3.fr}
\EGAReferenceAlias{III.7.9.12.3-fr}{III.7.9.12.3.fr}
\EGAReferenceAlias{III.7.9.12.3-es}{III.7.9.12.3.fr}
\EGAReferenceAlias{III.7.9.12.3.es}{III.7.9.12.3.fr}
\EGAReferenceAlias{III.7.9.12.3.eq}{III.7.9.12.3.fr}
\EGAReferenceAlias{eq:III.7.9.12.3}{III.7.9.12.3.fr}
\EGAReferenceAlias{III.7.9.13-fr}{III.7.9.13}
\EGAReferenceAlias{III.7.9.13-es}{III.7.9.13}
\EGAReferenceAlias{III.7.9.13.es}{III.7.9.13}
\EGAReferenceAlias{III.7.9.13.eq}{III.7.9.13}
\EGAReferenceAlias{eq:III.7.9.13}{III.7.9.13}
\EGAReferenceAlias{III.7.9.14-fr}{III.7.9.14}
\EGAReferenceAlias{III.7.9.14-es}{III.7.9.14}
\EGAReferenceAlias{III.7.9.14.es}{III.7.9.14}
\EGAReferenceAlias{III.7.9.14.eq}{III.7.9.14}
\EGAReferenceAlias{eq:III.7.9.14}{III.7.9.14}



\begin{titlepage}
  \centering
  {\Large A. Grothendieck y J. Dieudonné\par}
  \vfill
  {\Huge Elementos de geometría algebraica III\par}
  \vspace{1.5ex}
  {\LARGE Segunda parte\par}
  \vfill
  {\large Traducción española\par}
  \vspace{1ex}
  {\small Texto original: A. Grothendieck y J. Dieudonné, Publications Mathématiques de l’IHÉS, volumen 17 (1963)\par}
\end{titlepage}

\chapter*{CAPÍTULO III (continuación)\\[1.2ex]
\Large ESTUDIO COHOMOLÓGICO DE LOS HACES COHERENTES}

\setcounter{section}{5}
\setcounter{subsection}{0}
\setcounter{equation}{0}
\setcounter{footnote}{0}
\chapter*{CAPÍTULO III \emph{(continuación)}\\[1.2ex]
\Large ESTUDIO COHOMOLÓGICO DE LOS HACES COHERENTES}

\oldpage[III]{137}

\section{Funtores «Tor» locales y globales; fórmula de Künneth}
\label{section:III.6.fr}

\subsection{Introducción.}
\label{subsection:III.6.1.fr}

\begin{env}[6.1.1]
\label{III.6.1.1}
Sean $f:X\to Y$ un morfismo de preesquemas y $\mathscr F$ un
$\mathscr O_X$-módulo cuasicoherente. En el estudio de las «imágenes
directas superiores» $\mathrm R^n f_*(\mathscr F)$, se llega a considerar el
siguiente problema general: dado un morfismo de «cambio de base»
$g:Y'\to Y$, se pone
$X'=X_{(Y')}=X\times_Y Y'$, $\mathscr F'=\mathscr F\otimes_Y\mathscr O_{Y'}$, $f'=f_{(Y')}:X'\to Y'$,
y se propone obtener información sobre las imágenes directas
superiores $\mathrm R^n f'_*(\mathscr F')$ (suponiendo conocidas las
$\mathrm R^n f_*(\mathscr F)$). Se ve fácilmente (cf., por ejemplo, (7.7.2))
que el problema se reduce a estudiar las variaciones de
$\mathrm R^n f_*(\mathscr F\otimes_{\mathscr O_Y}\mathscr G)$ para un
$\mathscr O_Y$-\emph{módulo} cuasicoherente variable $\mathscr G$, es decir, el
funtor
$\mathscr G\leadsto\mathrm R^n f_*(\mathscr F\otimes_{\mathscr O_Y}\mathscr G)$.
Si $\mathscr F$ es \emph{plano} sobre $Y$, el funtor
$\mathscr G\leadsto\mathscr F\otimes_{\mathscr O_Y}\mathscr G$ es \emph{exacto}
(\textbf{0}\textsubscript{I}, 6.7.4) y, por consiguiente, el funtor compuesto
$\mathscr G\leadsto\mathrm R^\bullet f_*(\mathscr F\otimes_{\mathscr O_Y}\mathscr G)$
sigue siendo un funtor \emph{cohomológico}. Pero esto ya no sucede en el caso
general; para poder aplicar los métodos cohomológicos, se ve uno llevado a
sustituir
$\mathscr G\leadsto\mathrm R^n f_*(\mathscr F\otimes_{\mathscr O_Y}\mathscr G)$
por otros funtores, que esta vez son \emph{siempre} funtores cohomológicos. Estos
funtores, que generalizan los funtores «Tor» de la teoría de módulos, se
definen en los números 6.3 a 6.7; hay, por lo demás, dos generalizaciones de
este tipo, una «local» y la otra «global», relacionadas por sucesiones
espectrales que se discutirán en el número 6.7; como aplicación de estas
sucesiones espectrales, se obtiene en particular, bajo ciertas condiciones,
una «fórmula de Künneth» que expresa
$\mathrm R^n(f_1\times f_2)_*(\mathscr F_1\otimes_Y\mathscr F_2)$
por medio de las imágenes directas superiores
$\mathrm R^p f_{1*}(\mathscr F_1)$ y $\mathrm R^q f_{2*}(\mathscr F_2)$.
Otras sucesiones espectrales (6.8) generalizan las sucesiones espectrales de
asociatividad del funtor «Tor» de módulos; por último, el problema del cambio
de base conduce también a sucesiones espectrales (6.9).
\end{env}

\begin{env}[6.1.2]
\label{III.6.1.2}
Además, se comprueba (6.10) que los funtores cohomológicos
$\mathscr G\leadsto\mathscr T_\bullet(\mathscr G)$ así definidos son
\emph{localmente} (sobre $Y$) de la forma
$\mathscr G\leadsto\mathscr H_\bullet(\mathscr L_\bullet\otimes_Y\mathscr G)$,
donde $\mathscr L_\bullet$ es un complejo de $\mathscr O_Y$-módulos
\emph{localmente libres} (definido salvo homotopía) y $\mathscr H_\bullet$ denota la
homología. Conviene entonces dejar de lado la situación particular que dio
origen a $\mathscr T_\bullet$ y estudiar, de

\oldpage[III]{138}

manera general, los funtores de la forma precedente
$\mathscr G\leadsto\mathscr H_\bullet(\mathscr L_\bullet\otimes_Y\mathscr G)$
(imponiendo además, cuando proceda, hipótesis de finitud adecuadas sobre los
$\mathscr L_i$ o los $\mathscr H_i(\mathscr L_\bullet)$): este es el objeto
del § 7, cuya lectura es, en lo esencial, independiente del § 6. Las
propiedades más importantes de estos funtores conciernen a las propiedades de
\emph{exactitud} de una componente $\mathscr T_i$ de $\mathscr T_\bullet$; se darán
diversos criterios que permiten establecer tales propiedades; como
aplicación, se obtendrán condiciones que permiten afirmar (con las notaciones
de (6.1.1)) que el funtor
$\mathscr G\leadsto\mathrm R^n f_*(\mathscr F\otimes_{\mathscr O_Y}\mathscr G)$
es \emph{exacto} (lo que se expresará diciendo que $\mathscr F$ es
\emph{cohomológicamente plano} sobre $Y$ en dimensión $n$). Otra propiedad
importante de las componentes $\mathscr T_i$ de
$\mathscr T_\bullet(\mathscr G)=
\mathscr H_\bullet(\mathscr L_\bullet\otimes_Y\mathscr G)$ es una propiedad
de \emph{semicontinuidad} de la función
$y\leadsto\dim_{\boldsymbol{k}(y)}(T_i(\boldsymbol{k}(y)))$; cuando $\mathscr T_i$ es \emph{exacto}, esta
propiedad es sustituida por una propiedad de \emph{continuidad}, y el recíproco es
además cierto, por Grauert, cuando $Y$ es reducido (7.8.4).
\end{env}

\begin{env}[6.1.3]
\label{III.6.1.3}
En los §§ 6 y 7 hemos utilizado sistemáticamente la \emph{hipercohomología}, tomando
en todas partes como argumentos \emph{complejos} de haces en lugar de haces, aunque
la necesidad de este punto de vista solo se manifieste en capítulos
posteriores. El formalismo cohomológico desarrollado con este motivo se hará,
por lo demás, más transparente en el capítulo de este Tratado que estará
dedicado al desarrollo de un álgebra de funtores cohomológicos de haces
coherentes, incluido el formalismo de la dualidad. Pero esto requerirá
desarrollos que exceden el marco del presente capítulo.
\end{env}

\begin{env}[6.1.4]
\label{III.6.1.4}
Para abreviar, dados dos complejos $K^\bullet$, $K'^\bullet$ en una categoría
abeliana $\boldsymbol{C}$, diremos que un morfismo de complejos
$f:K^\bullet\to K'^\bullet$ es una \emph{equivalencia homotópica} si existe un
morfismo $g:K'^\bullet\to K^\bullet$ tal que los morfismos compuestos
$f\circ g$ y $g\circ f$ sean ambos \emph{homotópicos a la identidad} (por abuso de
lenguaje, cuando exista una equivalencia homotópica de este tipo, se dirá
asimismo que $K^\bullet$ y $K'^\bullet$ son \emph{homotópicamente equivalentes}).
Cuando se puede definir la \emph{hipercohomología} de un funtor covariante aditivo
$T$ de $\boldsymbol{C}$ en una categoría abeliana $\boldsymbol{C}'$, respecto de un complejo de $\boldsymbol{C}$
(\textbf 0, 11.4.3), es inmediato que una equivalencia homotópica
$K^\bullet\to K'^\bullet$ de complejos de $\boldsymbol{C}$ define canónicamente un
\emph{isomorfismo}
$\mathbf R^\bullet T(K^\bullet)\to\mathbf R^\bullet T(K'^\bullet)$
en hipercohomología \emph{(loc. cit.)}.
\end{env}

\subsection{Hipercohomología de los complejos de módulos sobre un preesquema}
\label{subsection:III.6.2}
\label{subsection:III.6.2.fr}

\begin{env}[6.2.1]
Sean $X$ un preesquema y
$\mathscr K^\bullet=(\mathscr K^i)_{i\in\mathbf Z}$ un complejo de
$\mathscr O_X$-módulos cuya diferencial es de grado $+1$.
Recordemos que, para todo morfismo $f:X\to Y$ de preesquemas, se han definido
(\textbf 0, 12.4.1) los $\mathscr O_Y$-\emph{módulos de hipercohomología}
$\mathscr H^n(f,\mathscr K^\bullet)$ (denotados también por
$\mathscr H_f^n(\mathscr K^\bullet)$ o
$\mathbf R^n f_*(\mathscr K^\bullet)$) para todo $n\in\mathbf Z$;
la hipercohomología $\mathscr H^\bullet(f,\mathscr K^\bullet)$ es el objeto de
convergencia de los dos funtores espectrales
$'\mathscr E(f,\mathscr K^\bullet)$ y
$''\mathscr E(f,\mathscr K^\bullet)$, cuyos términos $\mathscr E_2$ vienen
dados por
\begin{align}
  '\mathscr E_2^{pq}
    &=\mathscr H^p\bigl(\mathscr H^q(f,\mathscr K^\bullet)\bigr)
    \tag{6.2.1.1}\label{III.6.2.1.1.fr}\\
  ''\mathscr E_2^{pq}
    &=\mathscr H^p\bigl(f,\mathscr H^q(\mathscr K^\bullet)\bigr)
      =\mathrm R^p f_*\bigl(\mathscr H^q(\mathscr K^\bullet)\bigr)
    \tag{6.2.1.2}\label{III.6.2.1.2.fr}
\end{align}

\oldpage[III]{139}

donde $\mathscr H^q(f,\mathscr K^\bullet)$ es el \emph{complejo} cuyo componente de
grado $i$ es
$\mathscr H^q(f,\mathscr K^i)=\mathrm R^q f_*(\mathscr K^i)$ \emph{(loc. cit.)}.
Recordemos también que, cuando $Y$ se reduce a un punto, se denota la
hipercohomología correspondiente por $\mathbf H^\bullet(X,\mathscr K^\bullet)$
(la cual está formada por \emph{módulos} sobre $\Gamma(X,\mathscr O_X)$ independientes
del preesquema puntual $Y$ considerado); cuando $Y=X$ y $f=1_X$, se tiene
$\mathscr H^n(f,\mathscr K^\bullet)=\mathscr H^n(\mathscr K^\bullet)$
(cohomología del complejo $\mathscr K^\bullet$); cuando
$\mathscr K^i=0$, salvo para $i=i_0$, se tiene
\[
  \mathscr H^n(f,\mathscr K^\bullet)
    =\mathrm R^{n-i_0}f_*(\mathscr K^{i_0}).
\]
La sucesión espectral $'\mathscr E(f,\mathscr K^\bullet)$ es siempre
\emph{regular}; las dos sucesiones espectrales son \emph{birregulares} cuando
$\mathscr K^\bullet$ está acotado inferiormente (\textbf 0, 12.4.1).

Toda \emph{equivalencia homotópica} $h:\mathscr K^\bullet\to\mathscr K'^\bullet$ de complejos de
$\mathscr O_X$-módulos (6.1.4) da lugar a un \emph{isomorfismo}
$\mathscr H^\bullet(f,\mathscr K^\bullet)\xrightarrow{\sim}
\mathscr H^\bullet(f,\mathscr K'^\bullet)$ en hipercohomología. Lo mismo ocurre
si se supone solamente que
$\mathscr H^\bullet(h):\mathscr H^\bullet(\mathscr K^\bullet)\to
\mathscr H^\bullet(\mathscr K'^\bullet)$ es un \emph{isomorfismo} y que
$\mathscr K^\bullet$ y $\mathscr K'^\bullet$ están \emph{acotados inferiormente},
como se deduce inmediatamente de (\textbf 0, 11.1.5), aplicado a la sucesión espectral
(6.2.1.2) y a su análoga para $\mathscr K'^\bullet$. Por último, para todo
recubrimiento abierto $\mathfrak U=(U_\alpha)$ de $X$, se ha definido también
(\textbf 0, 12.4.5) la hipercohomología
$\mathbf H^\bullet(\mathfrak U,\mathscr K^\bullet)$ como la cohomología del
bicomplejo $C^\bullet(\mathfrak U,\mathscr K^\bullet)$ (cuyo componente de
índices $(i,j)$ es, por definición,
$C^i(\mathfrak U,\mathscr K^j)$); los
$\mathbf H^n(\mathfrak U,\mathscr K^\bullet)$ son también módulos sobre
$\Gamma(X,\mathscr O_X)$.
\end{env}

\begin{proposition}[6.2.2]
Sean $X$ un esquema y $\mathfrak U=(U_\alpha)$ un recubrimiento de $X$ por
abiertos afines. Para todo complejo $\mathscr K^\bullet$ de
$\mathscr O_X$-módulos cuasicoherentes, los módulos de hipercohomología
$\mathbf H^\bullet(X,\mathscr K^\bullet)$ y
$\mathbf H^\bullet(\mathfrak U,\mathscr K^\bullet)$ son canónicamente
isomorfos.
\end{proposition}

\begin{proof}
En efecto, toda intersección finita $V$ de abiertos del recubrimiento
$\mathfrak U$ es afín (\textbf I, 5.5.6), por lo que
$\mathrm H^q(V,\mathscr K^i)=0$ para todo $i$ y todo $q>0$ (1.3.1); la
proposición es, por tanto, un caso particular de (\textbf 0, 12.4.7).
\end{proof}

\begin{proposition}[6.2.3]
Sea $f:X\to Y$ un morfismo cuasicompacto y separado de preesquemas. Para todo
complejo $\mathscr K^\bullet$ de $\mathscr O_X$-módulos cuasicoherentes, los
$\mathscr O_Y$-módulos $\mathscr H^n(f,\mathscr K^\bullet)$ son
cuasicoherentes.
\end{proposition}

\begin{proof}
Como los $\mathscr H^q(f,\mathscr K^i)=\mathrm R^q f_*(\mathscr K^i)$ son
$\mathscr O_Y$-módulos cuasicoherentes (1.4.10), lo mismo ocurre con
$'\mathscr E_2^{pq}$, que, según (6.2.1.1), es cociente del núcleo de un
homomorfismo de módulos cuasicoherentes por la imagen de un homomorfismo de
esta clase (\textbf I, 4.1.1). Por la misma razón, todos los
$\mathscr O_Y$-módulos $'\mathscr E_r^{pq}$,
$B_k('\mathscr E_r^{pq})$, $Z_k('\mathscr E_r^{pq})$ de la primera sucesión
espectral son cuasicoherentes. La regularidad de la sucesión espectral
$'\mathscr E(f,\mathscr K^\bullet)$ implica que
$Z_\infty('\mathscr E_2^{pq})$ es igual a uno de los
$Z_k('\mathscr E_2^{pq})$, luego es cuasicoherente, y lo mismo ocurre con
$B_\infty('\mathscr E_2^{pq})=\varinjlim_k B_k('\mathscr E_2^{pq})$
(\textbf 0, 11.2.4 e \textbf I, 4.1.1); por tanto, los $'\mathscr E_\infty^{pq}$ son también
cuasicoherentes. Como la sucesión espectral precedente es regular, la
filtración de los $F^p(\mathscr H^n(f,\mathscr K^\bullet))$ es discreta y
exhaustiva; dicho de otro modo, el $\mathscr O_Y$-módulo
$\mathscr H^n(f,\mathscr K^\bullet)$ es la unión de una sucesión creciente
$(\mathscr G_k)_{k\geq0}$ de $\mathscr O_Y$-módulos tal que
$\mathscr G_0=0$ y cada $\mathscr G_k/\mathscr G_{k-1}$ es igual a uno de los
$\mathscr O_Y$-módulos $'\mathscr E_\infty^{pq}$ y, por tanto, es
cuasicoherente. Por inducción sobre $k$, se deduce que los $\mathscr G_k$ son
cuasicoherentes (1.4.17), y, como
$\mathscr H^n(f,\mathscr K^\bullet)=\varinjlim\mathscr G_k$, queda demostrada
la proposición (\textbf I, 4.1.1).
\end{proof}

\oldpage[III]{140}

\begin{corollary}[6.2.4]
Bajo las hipótesis de (6.2.3), para todo abierto afín $V$ de $Y$, el
homomorfismo canónico
\[
  \mathbf H^n(f^{-1}(V),\mathscr K^\bullet)
  \longrightarrow
  \Gamma(V,\mathscr H^n(f,\mathscr K^\bullet))
  \tag{6.2.4.1}\label{III.6.2.4.1.fr}
\]
es biyectivo para todo $n\in\mathbf Z$.
\end{corollary}

\begin{proof}
La demostración es la misma que la de (1.4.11), utilizando (6.2.2),
reemplazando $\mathscr F$ por $\mathscr K^\bullet$,
$\mathscr K^\bullet$ por
$f_*(\mathscr C^\bullet(\mathfrak U,\mathscr K^\bullet))$,
$\mathscr H^\bullet(\mathscr K^\bullet)$ por
$\mathscr H^\bullet(f,\mathscr K^\bullet)$, y observando que este último es un
$\mathscr O_Y$-módulo cuasicoherente por (6.2.3).
\end{proof}

\begin{proposition}[6.2.5]
Sean $Y$ un preesquema localmente noetheriano, $f:X\to Y$ un morfismo propio
y $\mathscr K^\bullet$ un complejo de $\mathscr O_X$-módulos tal que los
$\mathscr O_X$-módulos $\mathscr H^q(\mathscr K^\bullet)$ sean coherentes.
Entonces los $\mathscr O_Y$-módulos $\mathscr H^n(f,\mathscr K^\bullet)$ son
coherentes.
\end{proposition}

\begin{proof}
Como la cuestión es local sobre $Y$, se puede limitar al caso en que $Y$ es
noetheriano y afín; se trata entonces, en virtud de (6.2.4), de demostrar que
los $\mathbf H^n(X,\mathscr K^\bullet)$ son
$\Gamma(Y,\mathscr O_Y)$-módulos de tipo finito. Se tiene entonces
$\mathbf H^\bullet(X,\mathscr K^\bullet)=
\mathbf H^\bullet(\mathfrak U,\mathscr K^\bullet)$ (6.2.2), donde se puede
suponer que $\mathfrak U$ es \emph{finito}, puesto que $X$ es cuasicompacto. Como las
cocadenas de cada complejo $C^\bullet(\mathfrak U,\mathscr K^j)$ son
\emph{alternadas} por definición, existe un entero $r>0$ tal que
$C^i(\mathfrak U,\mathscr K^j)=0$ para $i<0$ y $i>r$; se concluye de ello
(\textbf 0, 11.3.3) que las dos sucesiones espectrales del bicomplejo
$C^\bullet(\mathfrak U,\mathscr K^\bullet)$ son \emph{birregulares}. Como las
intersecciones de los miembros de $\mathfrak U$ son abiertos afines
(\textbf I, 5.5.6), cada funtor
$\mathscr F\leadsto C^i(\mathfrak U,\mathscr F)$ es exacto en la categoría de
los $\mathscr O_X$-módulos cuasicoherentes; por consiguiente,
$\mathrm H_I^q(C^i(\mathfrak U,\mathscr K^\bullet))
=C^i(\mathfrak U,\mathscr H^q(\mathscr K^\bullet))$, y los términos
$E_2$ de la segunda sucesión espectral de
$C^\bullet(\mathfrak U,\mathscr K^\bullet)$ vienen dados (\textbf 0, 11.3.2) por
\[
  ''E_2^{pq}
  =\mathrm H^p(C^\bullet(\mathfrak U,\mathscr H^q(\mathscr K^\bullet)))
  =\mathrm H^p(\mathfrak U,\mathscr H^q(\mathscr K^\bullet))
  =\mathrm H^p(X,\mathscr H^q(\mathscr K^\bullet))
\]
en virtud de (1.4.1); dado que $f$ es propio, estos son
$\Gamma(Y,\mathscr O_Y)$-módulos \emph{de tipo finito} (3.2.1). Puesto que la sucesión
espectral $''E(C^\bullet(\mathfrak U,\mathscr K^\bullet))$ es birregular, se
deduce que los $\mathbf H^n(X,\mathscr K^\bullet)$ son
$\Gamma(Y,\mathscr O_Y)$-módulos de tipo finito (\textbf 0, 11.1.8).

\emph{A fortiori}, si $\mathscr K^\bullet$ es un complejo de
$\mathscr O_X$-módulos \emph{coherentes}, los $\mathscr O_Y$-módulos
$\mathscr H^n(f,\mathscr K^\bullet)$ son \emph{coherentes} bajo las hipótesis de
(6.2.5) relativas a $Y$ y $f$ (\textbf{0}\textsubscript{I}, 5.3.4).
\end{proof}

\begin{env}[6.2.6]
La hipercohomología $\mathscr H^\bullet(f,\mathscr K^\bullet)$ es un funtor
\emph{cohomológico} en la categoría de los complejos de $\mathscr O_X$-módulos
\emph{acotados inferiormente} (\textbf 0, 12.4.4). Es un funtor cohomológico en la categoría
de \emph{todos} los complejos de $\mathscr O_X$-módulos cuando el morfismo $f$ es
\emph{cuasicompacto} y el espacio subyacente a $X$ es \emph{localmente noetheriano}: en
efecto, de (G, II, 3.10.1) se deduce entonces que $f_*$ conmuta con los límites
inductivos (pues la cuestión es local sobre $Y$), y puede aplicarse
(\textbf 0, 11.5.2).

Por último, si $f$ es \emph{separado}, $\mathscr H^\bullet(f,\mathscr K^\bullet)$ es
un funtor cohomológico en la categoría de los complejos de
$\mathscr O_X$-módulos \emph{cuasicoherentes}. Esto es inmediato cuando $Y$ es
afín, pues entonces $X$ es un \emph{esquema}; por tanto, en virtud del isomorfismo
canónico (6.2.2), todo se reduce a ver que
$\mathscr K^\bullet\leadsto\mathbf H^\bullet(\mathfrak U,\mathscr K^\bullet)$
es un funtor cohomológico en la categoría de los complejos de
$\mathscr O_X$-módulos cuasicoherentes, lo cual es inmediato, ya que el
funtor $\mathscr K^\bullet\leadsto
C^\bullet(\mathfrak U,\mathscr K^\bullet)$ es exacto en esta categoría
(\textbf I, 1.3.7). En el caso general, para todo abierto afín $V$

\oldpage[III]{141}

de $Y$, $f^{-1}(V)$ es un esquema y, para aplicar lo que precede, basta
verificar que, para una sucesión exacta
$0\longrightarrow\mathscr K'^\bullet\longrightarrow\mathscr K^\bullet\longrightarrow\mathscr K''^\bullet\longrightarrow0$
de complejos de $\mathscr O_X$-módulos cuasicoherentes, el homomorfismo
$\partial:\mathscr H^n(f,\mathscr K''^\bullet)|V\longrightarrow\mathscr H^{n+1}(f,\mathscr K'^\bullet)|V$
no depende del recubrimiento abierto afín $\mathfrak U$ de $f^{-1}(V)$ utilizado
para definirlo. Pero esto se deduce de que, si $\mathfrak U'$ es un
recubrimiento abierto afín más fino que $\mathfrak U$, el diagrama
\[
  \xymatrix@C=34pt@R=11pt{
    \mathbf H^\bullet(\mathfrak U,\mathscr K^\bullet)
      \ar[rrd]^-{\sim}\ar[dd]^{\wr}
    && \\
    &&
    \mathbf H^\bullet(f^{-1}(V),\mathscr K^\bullet)
    \\
    \mathbf H^\bullet(\mathfrak U',\mathscr K^\bullet)
      \ar[rru]^-{\sim}
    &&
  }
\]
de isomorfismos canónicos es conmutativo, al igual que el diagrama
\[
  \xymatrix@C=46pt@R=25pt{
    \mathbf H^n(\mathfrak U,\mathscr K''^\bullet)
      \ar[r]^-{\partial}\ar[d]^{\wr}
    &
    \mathbf H^{n+1}(\mathfrak U,\mathscr K'^\bullet)
      \ar[d]^{\wr}
    \\
    \mathbf H^n(\mathfrak U',\mathscr K''^\bullet)
      \ar[r]_-{\partial}
    &
    \mathbf H^{n+1}(\mathfrak U',\mathscr K'^\bullet)
  }
\]

Cuando se cumple una de las condiciones precedentes y
$\mathscr K^\bullet\to\mathscr K'^\bullet$ es una equivalencia homotópica (6.1.4), el
isomorfismo correspondiente
$\mathscr H^\bullet(f,\mathscr K^\bullet)\xrightarrow{\sim}
\mathscr H^\bullet(f,\mathscr K'^\bullet)$ es entonces un isomorfismo de
$\partial$-funtores (\textbf 0, 11.4.4).
\end{env}

\begin{env}[6.2.7]
Todo lo que precede se aplica naturalmente, sin más cambio que el de la
notación, a un complejo $\mathscr K_\bullet$ de $\mathscr O_X$-módulos
cuasicoherentes cuya diferencial es de grado $-1$; basta considerar el
complejo $\mathscr K^\bullet=(\mathscr K^i)$, donde
$\mathscr K^i=\mathscr K_{-i}$ para todo $i\in\mathbf Z$.
\end{env}

\subsection{Hipertor de dos complejos de módulos}
\label{subsection:III.6.3}
\label{subsection:III.6.3.fr}

\begin{env}[6.3.1]
\label{III.6.3.1}
Sean $A$ un anillo conmutativo y $P_\bullet$, $Q_\bullet$ dos complejos de
$A$-módulos cuyos diferenciales son de grado $-1$; sea
$L_{\bullet\bullet}$ (resp. $M_{\bullet\bullet}$) una resolución proyectiva
de Cartan-Eilenberg de $P_\bullet$ (resp. $Q_\bullet$)
\sref[0]{0.11.6.1};
$L_{\bullet\bullet}\otimes_A M_{\bullet\bullet}$ es entonces (para la suma
de los primeros grados y la suma de los segundos grados) un bicomplejo (con
diferenciales de grado $-1$), cuya homología
$H_\bullet(L_{\bullet\bullet}\otimes_A M_{\bullet\bullet})$ no depende de las
resoluciones de Cartan-Eilenberg $L_{\bullet\bullet}$,
$M_{\bullet\bullet}$ elegidas y es, por definición, la hiperhomología del
bifuntor $P_\bullet\otimes_A Q_\bullet$ en $P_\bullet$ y $Q_\bullet$
\sref[0]{0.11.6.5}. Pondremos, por definición,
\[
  \mathbf{Tor}_n^A(P_\bullet,Q_\bullet)
  =H_n(L_{\bullet\bullet}\otimes_A M_{\bullet\bullet})
  \tag{6.3.1.1}\label{III.6.3.1.1.fr}
\]

\oldpage[III]{142}

y diremos que este $A$-módulo es el \emph{hipertor} de índice $n$ de los dos
complejos $P_\bullet$, $Q_\bullet$. Se sabe que, en la categoría de los
complejos de $A$-módulos \emph{acotados inferiormente}, los
$\mathbf{Tor}_n^A(P_\bullet,Q_\bullet)$ forman un \emph{bifuntor homológico} en
$P_\bullet$, $Q_\bullet$ \sref[0]{0.11.6.5}. Además:
\end{env}

\begin{proposition}[6.3.2]
\label{III.6.3.2}
El bifuntor $\mathbf{Tor}_\bullet^A(P_\bullet,Q_\bullet)$ es el término
límite común de dos bifuntores espectrales
$'E(P_\bullet,Q_\bullet)$, $''E(P_\bullet,Q_\bullet)$, cuyos términos $E_2$
son
\begin{align}
  'E^2_{pq}
    &=H_p\bigl(\operatorname{Tor}_q^A(P_\bullet,Q_\bullet)\bigr),
    \tag{6.3.2.1}\label{III.6.3.2.1.fr}\\
  ''E^2_{pq}
    &=\bigoplus_{q'+q''=q}
      \operatorname{Tor}_p^A\bigl(H_{q'}(P_\bullet),H_{q''}(Q_\bullet)\bigr).
    \tag{6.3.2.2}\label{III.6.3.2.2.fr}
\end{align}
donde, en \sref{III.6.3.2.1},
$\operatorname{Tor}_q^A(P_\bullet,Q_\bullet)$ designa el
bicomplejo formado por los $A$-módulos $\operatorname{Tor}_q^A(P_i,Q_j)$. La
sucesión espectral \sref{III.6.3.2.2} es siempre regular; si $P_\bullet$ y
$Q_\bullet$
están acotados inferiormente, o si $A$ es de dimensión cohomológica finita,
las dos sucesiones espectrales \sref{III.6.3.2.1} y
\sref{III.6.3.2.2} son birregulares.
\end{proposition}

\begin{proof}
Esto resulta de \sref[0]{0.11.6.5}, pues cuando $A$ es de dimensión
cohomológica
finita $n$, todo $A$-módulo admite una resolución proyectiva de longitud $n$
\emph{(M, VI, 2.1)}.
\end{proof}

\begin{corollary}[6.3.3]
\label{III.6.3.3}
Sean $P'_\bullet$, $Q'_\bullet$ dos complejos de $A$-módulos, y
$u:P_\bullet\to P'_\bullet$, $v:Q_\bullet\to Q'_\bullet$ dos homomorfismos
de complejos. Si los homomorfismos
$H_\bullet(u):H_\bullet(P_\bullet)\to H_\bullet(P'_\bullet)$,
$H_\bullet(v):H_\bullet(Q_\bullet)\to H_\bullet(Q'_\bullet)$ deducidos
respectivamente de $u$ y $v$ son biyectivos, entonces el homomorfismo
$\mathbf{Tor}_\bullet^A(P_\bullet,Q_\bullet)\to
\mathbf{Tor}_\bullet^A(P'_\bullet,Q'_\bullet)$ deducido de $u$ y $v$ es
biyectivo.
\end{corollary}

\begin{proof}
En efecto, el homomorfismo de sucesiones espectrales
$''E(P_\bullet,Q_\bullet)\to''E(P'_\bullet,Q'_\bullet)$ deducido de $u$ y $v$
es entonces un isomorfismo para los términos $E_2$, y la conclusión resulta
de que estas sucesiones son regulares en virtud de \sref{III.6.3.2} y
\sref[0]{0.11.1.5}.
\end{proof}

\begin{proposition}[6.3.4]
\label{III.6.3.4}
Sean $P_\bullet$, $Q_\bullet$ dos complejos de $A$-módulos acotados
inferiormente. Sea $L_{\bullet\bullet}$ \textup{(resp. $M_{\bullet\bullet}$)} un
bicomplejo formado por $A$-módulos planos, tal que, para todo $i$,
$L_{i,\bullet}$ \textup{(resp. $M_{i,\bullet}$)} sea una resolución de $P_i$
\textup{(resp. $Q_i$)}. Se tienen entonces isomorfismos canónicos
\[
  \mathbf{Tor}_\bullet^A(P_\bullet,Q_\bullet)
  \xrightarrow{\sim}H_\bullet(L_{\bullet\bullet}\otimes_AQ_\bullet)
  \xrightarrow{\sim}H_\bullet(P_\bullet\otimes_AM_{\bullet\bullet})
  \xrightarrow{\sim}H_\bullet(L_{\bullet\bullet}\otimes_AM_{\bullet\bullet})
  \tag{6.3.4.1}\label{III.6.3.4.1.fr}
\]
\end{proposition}

\begin{proof}
Esto resulta de \sref[0]{0.11.6.5}, partes (ii) y (iii), y de la definición
de los
$A$-módulos planos.
\end{proof}

\begin{namedtheorem}[6.3.5]{Observaciones}
\label{III.6.3.5}
(i) Con las notaciones de \sref{III.6.3.1}, los bicomplejos
$L_{\bullet\bullet}\otimes_A M_{\bullet\bullet}$ y
$M_{\bullet\bullet}\otimes_A L_{\bullet\bullet}$ son canónicamente
isomorfos, de donde se obtiene un isomorfismo canónico
$\mathbf{Tor}_\bullet^A(P_\bullet,Q_\bullet)\xrightarrow{\sim}
\mathbf{Tor}_\bullet^A(Q_\bullet,P_\bullet)$.

(ii) Si $F$ y $G$ son dos $A$-módulos, y $P_\bullet$ y $Q_\bullet$ son los
complejos de $A$-módulos reducidos a $F$ y $G$, respectivamente, en grado 0 y
nulos en los demás grados, entonces dos resoluciones proyectivas $L_\bullet$,
$M_\bullet$ de $F$ y $G$, respectivamente, pueden considerarse como
resoluciones de Cartan-Eilenberg de $P_\bullet$ y $Q_\bullet$ completándolas
con ceros. Por consiguiente, en este caso se tiene
$\mathbf{Tor}_\bullet^A(P_\bullet,Q_\bullet)
=\operatorname{Tor}_\bullet^A(F,G)$.
\end{namedtheorem}

\begin{proposition}[6.3.6]
\label{III.6.3.6}
Sean $(P_\bullet^\lambda)$, $(Q_\bullet^\mu)$ dos sistemas inductivos
filtrantes de complejos de $A$-módulos; se tiene un isomorfismo canónico
\[
  \varinjlim_{\lambda,\mu}
  \mathbf{Tor}_\bullet^A(P_\bullet^\lambda,Q_\bullet^\mu)
  \xrightarrow{\sim}
  \mathbf{Tor}_\bullet^A
  \left(\varinjlim_\lambda P_\bullet^\lambda,
        \varinjlim_\mu Q_\bullet^\mu\right)
  \tag{6.3.6.1}\label{III.6.3.6.1.fr}
\]
\end{proposition}

\begin{proof}
Pongamos $P_\bullet=\varinjlim P_\bullet^\lambda$,
$Q_\bullet=\varinjlim Q_\bullet^\mu$; por funtorialidad, está claro que
los $\mathbf{Tor}_\bullet^A(P_\bullet^\lambda,Q_\bullet^\mu)$ forman un
sistema inductivo y que las aplicaciones
$\mathbf{Tor}_\bullet^A(P_\bullet^\lambda,Q_\bullet^\mu)\to
\mathbf{Tor}_\bullet^A(P_\bullet,Q_\bullet)$ deducidas

\oldpage[III]{143}

de las aplicaciones canónicas $P_\bullet^\lambda\to P_\bullet$,
$Q_\bullet^\mu\to Q_\bullet$, forman un sistema inductivo de homomorfismos,
de donde se obtiene un homomorfismo canónico \sref{III.6.3.6.1} y, más
generalmente,
un homomorfismo canónico
$\varinjlim''E(P_\bullet^\lambda,Q_\bullet^\mu)\to''E(P_\bullet,Q_\bullet)$
del cual \sref{III.6.3.6.1} es el homomorfismo de los términos límite.
Además, la sucesión espectral $''E(P_\bullet,Q_\bullet)$ es regular
\sref{III.6.3.2}, y lo mismo
ocurre con la sucesión espectral
$\varinjlim''E(P_\bullet^\lambda,Q_\bullet^\mu)$, como resulta
de las definiciones \sref[0]{0.11.1.7} y de la demostración de
\sref[0]{0.11.3.3}; para demostrar que \sref{III.6.3.6.1} es biyectivo basta,
pues, por \sref[0]{0.11.1.5}, probar que el
homomorfismo
\[
  \varinjlim''E(P_\bullet^\lambda,Q_\bullet^\mu)
  \longrightarrow ''E(P_\bullet,Q_\bullet)
  \tag{6.3.6.2}\label{III.6.3.6.2.fr}
\]
es biyectivo para los términos $E^2$. Como el funtor $H_\bullet$ conmuta con
el límite inductivo de complejos de módulos, finalmente se reduce todo a
probar que, para dos sistemas inductivos filtrantes $(F^\lambda)$,
$(G^\mu)$ de $A$-módulos, el homomorfismo canónico
\[
  \varinjlim_{\lambda,\mu}(\operatorname{Tor}_\bullet^A(F^\lambda,G^\mu))
  \longrightarrow
  \operatorname{Tor}_\bullet^A
  \left(\varinjlim_\lambda F^\lambda,
        \varinjlim_\mu G^\mu\right)
\]
es biyectivo. Para ello, consideremos para cada $F^\lambda$ la \emph{resolución
libre canónica}
\[
  L_\bullet^\lambda:\quad
  \ldots\longrightarrow L_{i+1}^\lambda
  \longrightarrow L_i^\lambda\longrightarrow\ldots
  \longrightarrow L_1^\lambda\longrightarrow L_0^\lambda\longrightarrow0
\]
donde $L_0^\lambda$ es el $A$-módulo de las combinaciones lineales formales
de elementos de $F^\lambda$, y $L_{i+1}^\lambda$ es el $A$-módulo de las
combinaciones lineales formales de elementos de
$\operatorname{Ker}(L_i^\lambda\to L_{i-1}^\lambda)$; se comprueba
inmediatamente que los $L_\bullet^\lambda$ forman un sistema inductivo de
complejos, y si se pone
$F=\varinjlim_\lambda F^\lambda$, $L_i=\varinjlim_\lambda L_i^\lambda$, los
$L_i$ forman una \emph{resolución} $L_\bullet$ de $F$, pues el funtor
$\varinjlim$ es exacto; además, los $L_i$, límites inductivos de
$A$-módulos libres, son \emph{planos} \sref[0]{0.6.1.2}. Se considera de
igual modo, para cada $\mu$, la resolución libre canónica $M_\bullet^\mu$ de
$G^\mu$, y $M_\bullet=\varinjlim M_\bullet^\mu$ es una resolución plana
de $G=\varinjlim G^\mu$. Se tiene entonces
$\operatorname{Tor}_\bullet^A(\varinjlim F^\lambda,
\varinjlim G^\mu)=H_\bullet(L_\bullet\otimes_A M_\bullet)$ en virtud de
\sref{III.6.3.5} y \sref{III.6.3.4}; pero
$H_\bullet(L_\bullet\otimes_A M_\bullet)
=\varinjlim_{\lambda,\mu}
H_\bullet(L_\bullet^\lambda\otimes_A M_\bullet^\mu)$
puesto que $H_\bullet$ conmuta con los límites inductivos de complejos de
módulos; como $H_\bullet(L_\bullet^\lambda\otimes_A M_\bullet^\mu)
=\operatorname{Tor}_\bullet^A(F^\lambda,G^\mu)$, esto demuestra
\sref{III.6.3.6.1}.

Cuando se supone que existe $i_0$ tal que
$P_i^\lambda=Q_i^\mu=0$ para $i<i_0$, cualesquiera que sean $\lambda$ y
$\mu$, se demuestra del mismo modo que el homomorfismo canónico
\[
  \varinjlim'E(P_\bullet^\lambda,Q_\bullet^\mu)
  \longrightarrow'E(P_\bullet,Q_\bullet)
  \tag{6.3.6.3}\label{III.6.3.6.3.fr}
\]
es \emph{biyectivo}.
\end{proof}

\begin{proposition}[6.3.7]
\label{III.6.3.7}
Supongamos que $P_\bullet$ y $Q_\bullet$ están acotados inferiormente. Si el
complejo $P_\bullet$ está formado por $A$-módulos planos, se tiene un
isomorfismo $A$-lineal canónico de $\partial$-funtores en $Q_\bullet$
\[
  \mathbf{Tor}_\bullet^A(P_\bullet,Q_\bullet)
  \xrightarrow{\sim}H_\bullet(P_\bullet\otimes_A Q_\bullet).
  \tag{6.3.7.1}\label{III.6.3.7.1.fr}
\]
\end{proposition}

\begin{proof}
En efecto, la sucesión espectral \sref{III.6.3.2.1} es birregular y
\emph{degenerada}, y la existencia del isomorfismo \sref{III.6.3.7.1} resulta de
\sref[0]{0.11.1.6}. Además,
calculando el hipertor a partir de una resolución proyectiva de
Cartan-Eilenberg de $P_\bullet$ como en \sref{III.6.3.4}, se ve
inmediatamente que el
isomorfismo así definido es un isomorfismo de $\partial$-funtores en
$Q_\bullet$.
\end{proof}
\begin{env}[6.3.8]
\label{III.6.3.8}
\oldpage[III]{144}
Sea $\rho:A\to A'$ un homomorfismo de anillos. Nos proponemos definir un
homomorfismo $A$-lineal funtorial de grado 0 asociado canónicamente a $\rho$:
\[
  \rho_{P_\bullet,Q_\bullet}:
  \mathbf{Tor}_\bullet^A(P_\bullet,Q_\bullet)
  \longrightarrow
  \mathbf{Tor}_\bullet^{A'}
  (P_\bullet\otimes_A A',Q_\bullet\otimes_A A').
  \tag{6.3.8.1}\label{III.6.3.8.1.fr}
\]

Para ello, consideremos una resolución \emph{proyectiva} de Cartan-Eilenberg
$L_{\bullet\bullet}$ de $P_\bullet$; consideremos, por otra parte, una
resolución \emph{proyectiva} de Cartan-Eilenberg $L'_{\bullet\bullet}$ de
$P_\bullet\otimes_A A'$. Vamos a ver que se puede definir un $A'$-lineal de complejos
$L_{\bullet\bullet}\otimes_A A'\to L'_{\bullet\bullet}$, determinado salvo
homotopía. En efecto, la construcción de $L_{\bullet\bullet}$ queda
enteramente determinada cuando se elige (arbitrariamente), para cada $i$,
una resolución proyectiva $(X_{ij}^{\mathrm B})_{j\geqslant0}$ de
$\mathrm B_i(P_\bullet)$ y una resolución proyectiva
$(X_{ij}^{\mathrm H})_{j\geqslant0}$ de $\mathrm H_i(P_\bullet)$, que son
respectivamente iguales a $\mathrm B_i^{\mathrm I}(L_{\bullet\bullet})$ y
$\mathrm H_i^{\mathrm I}(L_{\bullet\bullet})$; se deduce sucesivamente
$\mathrm Z_i^{\mathrm I}(L_{\bullet\bullet})
=\mathrm H_i^{\mathrm I}(L_{\bullet\bullet})
\oplus\mathrm B_i^{\mathrm I}(L_{\bullet\bullet})$, después
$L_{i,\bullet}=\mathrm Z_i^{\mathrm I}(L_{\bullet\bullet})
\oplus\mathrm B_{i-1}^{\mathrm I}(L_{\bullet\bullet})$. Dicho esto, $X_{i,\bullet}^{\mathrm B}\otimes_A A'$ ya no es, en general,
una resolución de $\mathrm B_i(P_\bullet)\otimes_A A'$, pero sigue siendo un
complejo formado por $A'$-módulos \emph{proyectivos}, y por tanto existe un
homomorfismo $A'$-lineal
$X_{i,\bullet}^{\mathrm B}\otimes_A A'
\to\mathrm B_i^{\mathrm I}(L'_{\bullet\bullet})$
compatible con las aumentaciones y determinado salvo homotopía
\emph{(M, V, 1.1)}. Se tiene asimismo un homomorfismo $A'$-lineal
$X_{i,\bullet}^{\mathrm H}\otimes_A A'
\to\mathrm H_i^{\mathrm I}(L'_{\bullet\bullet})$
determinado salvo homotopía, de donde se deduce, mediante la construcción
recordada más arriba, un homomorfismo $A'$-lineal
$L_{i,\bullet}\otimes_A A'\to L'_{i,\bullet}$ para todo $i$; estos
homomorfismos (para $i\in\mathbf Z$) son compatibles con los diferenciales
$L_{i,\bullet}\to L_{i-1,\bullet}$ y con los análogos para
$L'_{\bullet\bullet}$, en virtud de la misma construcción, y constituyen por
tanto el homomorfismo $A'$-lineal
$L_{\bullet\bullet}\otimes_A A'\to L'_{\bullet\bullet}$ buscado.

Para definir \sref{III.6.3.8.1}, basta entonces considerar de igual modo una
resolución proyectiva de Cartan-Eilenberg $M_{\bullet\bullet}$ (resp.
$M'_{\bullet\bullet}$) de $Q_\bullet$ (resp.
$Q_\bullet\otimes_A A'$), y un homomorfismo $A'$-lineal
$M_{\bullet\bullet}\otimes_A A'\to M'_{\bullet\bullet}$. De estos
homomorfismos se deduce un homomorfismo $A'$-lineal
$(L_{\bullet\bullet}\otimes_A A')\otimes_{A'}
(M_{\bullet\bullet}\otimes_A A')
\to
L'_{\bullet\bullet}\otimes_{A'}M'_{\bullet\bullet}$,
y luego, por composición, un homomorfismo $A$-lineal de bicomplejos
$L_{\bullet\bullet}\otimes_A M_{\bullet\bullet}
\to
L'_{\bullet\bullet}\otimes_{A'}M'_{\bullet\bullet}$,
y al pasar a la homología se obtiene \sref{III.6.3.8.1}, que está bien
definido
puesto que proviene de un morfismo de complejos determinado salvo homotopía.

Si $\rho':A'\to A''$ es un segundo homomorfismo de anillos, y
$\rho'':A\to A''$ es el homomorfismo compuesto $\rho'\circ\rho$, está claro que
$\rho''_{P_\bullet,Q_\bullet}
=\rho'_{P'_\bullet,Q'_\bullet}\circ\rho_{P_\bullet,Q_\bullet}$, donde
\[
  P'_\bullet=P_\bullet\otimes_A A',
  \quad Q'_\bullet=Q_\bullet\otimes_A A'.
\]

Observemos también que el morfismo de bicomplejos
$L_{\bullet\bullet}\otimes_A M_{\bullet\bullet}
\to L'_{\bullet\bullet}\otimes_{A'}M'_{\bullet\bullet}$ considerado
más arriba define morfismos funtoriales (en $P_\bullet$ y $Q_\bullet$)
de sucesiones espectrales
\[
  'E_{pq}^r(P_\bullet,Q_\bullet)
  \longrightarrow
  'E_{pq}^r(P_\bullet\otimes_A A',Q_\bullet\otimes_A A')
  \quad\text{y}\quad
  ''E_{pq}^r(P_\bullet,Q_\bullet)
  \longrightarrow
  ''E_{pq}^r(P_\bullet\otimes_A A',Q_\bullet\otimes_A A'),
\]
independientes de las resoluciones de Cartan-Eilenberg consideradas, y que tienen también
la propiedad de transitividad anterior.
\end{env}

\begin{proposition}[6.3.9]
\label{III.6.3.9}
Sea $\rho:A\to A'$ un homomorfismo de anillos tal que $A'$ sea un
$A$-módulo \textup{plano}. Se tienen entonces isomorfismos canónicos funtoriales
\[
  \mathbf{Tor}_\bullet^{A'}
  (P_\bullet\otimes_A A',Q_\bullet\otimes_A A')
  \xrightarrow{\sim}
  \mathbf{Tor}_\bullet^A(P_\bullet,Q_\bullet)\otimes_A A'
  \tag{6.3.9.1}\label{III.6.3.9.1.fr}
\]
\[
  \left\{
  \begin{aligned}
  'E(P_\bullet\otimes_A A',Q_\bullet\otimes_A A')
    &\xrightarrow{\sim}'E(P_\bullet,Q_\bullet)\otimes_A A'\\
  ''E(P_\bullet\otimes_A A',Q_\bullet\otimes_A A')
    &\xrightarrow{\sim}''E(P_\bullet,Q_\bullet)\otimes_A A'.
  \end{aligned}
  \right.
  \tag{6.3.9.2}\label{III.6.3.9.2.fr}
\]
\end{proposition}

\begin{proof}
\oldpage[III]{145}
En efecto, dada la exactitud del funtor $M\otimes_A A'$ en $M$,
$L_{\bullet\bullet}\otimes_A A'$ y $M_{\bullet\bullet}\otimes_A A'$ son
entonces \emph{resoluciones} proyectivas de Cartan-Eilenberg de
$P_\bullet\otimes_A A'$ y $Q_\bullet\otimes_A A'$, respectivamente, de
donde resulta la conclusión.
\end{proof}

\begin{env}[6.3.10]
\label{III.6.3.10}
Sea $\rho:A\to A'$ un homomorfismo de anillos; para todo complejo
$P'_\bullet$ de $A'$-módulos, $P'_{\bullet[\rho]}$ es un complejo de
$A$-módulos; además, la aplicación identidad
$P'_{\bullet[\rho]}\to P'_\bullet$ puede considerarse como el compuesto de
las aplicaciones canónicas
\[
  P'_{\bullet[\rho]}\longrightarrow
  P'_{\bullet[\rho]}\otimes_A A'
  \xrightarrow{\mu}P'_\bullet,
\]
donde $\mu$ es el homomorfismo $A'$-lineal $\mu(x\otimes a')=a'x$. Si
$Q'_\bullet$
es un segundo complejo de $A'$-módulos, se tienen por tanto homomorfismos
canónicos funtoriales de grado 0
\[
  \mathbf{Tor}_\bullet^A(P'_{\bullet[\rho]},Q'_{\bullet[\rho]})
  \longrightarrow
  \mathbf{Tor}_\bullet^{A'}
  (P'_{\bullet[\rho]}\otimes_A A',Q'_{\bullet[\rho]}\otimes_A A')
  \longrightarrow
  \mathbf{Tor}_\bullet^{A'}(P'_\bullet,Q'_\bullet)
  \tag{6.3.10.1}\label{III.6.3.10.1.fr}
\]
donde la primera flecha es el homomorfismo $A$-lineal definido en
\sref{III.6.3.8}, y la
segunda se deduce de los homomorfismos $A'$-lineales
$P'_{\bullet[\rho]}\otimes_A A'\to P'_\bullet$ y
$Q'_{\bullet[\rho]}\otimes_A A'\to Q'_\bullet$ por funtorialidad. Se tienen
homomorfismos análogos para las sucesiones espectrales de
\sref{III.6.3.2}, y
propiedades evidentes de transitividad, cuyo enunciado dejamos al cuidado del
lector.
\end{env}

\begin{proposition}[6.3.11]
\label{III.6.3.11}
Sea $\rho:A\to A'$ un homomorfismo de anillos que hace de $A'$ un
$A$-módulo plano. Para todo complejo $P'_\bullet$ de $A'$-módulos y todo
complejo $Q_\bullet$ de $A$-módulos acotado inferiormente, se tiene un
isomorfismo canónico funtorial
\[
  \mathbf{Tor}_\bullet^A(P'_{\bullet[\rho]},Q_\bullet)
  \xrightarrow{\sim}
  \mathbf{Tor}_\bullet^{A'}(P'_\bullet,Q_\bullet\otimes_A A').
  \tag{6.3.11.1}\label{III.6.3.11.1.fr}
\]
\end{proposition}

\begin{proof}
En efecto, si $M_{\bullet\bullet}$ es una resolución proyectiva de
Cartan-Eilenberg de $Q_\bullet$, $M_{\bullet\bullet}\otimes_A A'$ es una
resolución proyectiva de Cartan-Eilenberg de $Q_\bullet\otimes_A A'$, y se
tiene, salvo isomorfismo canónico,
$P'_{\bullet[\rho]}\otimes_A M_{\bullet\bullet}
=P'_\bullet\otimes_{A'}(M_{\bullet\bullet}\otimes_A A')$;
la conclusión resulta de \sref{III.6.3.4}.
\end{proof}

\begin{namedtheorem}[6.3.12]{Observación}
\label{III.6.3.12}
Sea $(A^\lambda)$ un sistema inductivo filtrante de anillos, y sean
$(P_\bullet^\lambda)$, $(Q_\bullet^\lambda)$ dos sistemas inductivos de
complejos de $(A^\lambda)$-módulos; se tiene entonces un isomorfismo canónico
que generaliza \sref{III.6.3.6.1}
\[
  \varinjlim
  \mathbf{Tor}_\bullet^{A^\lambda}
  (P_\bullet^\lambda,Q_\bullet^\lambda)
  \xrightarrow{\sim}
  \mathbf{Tor}_\bullet^A(P_\bullet,Q_\bullet)
  \tag{6.3.12.1}\label{III.6.3.12.1.fr}
\]
donde $A=\varinjlim A^\lambda$, $P_\bullet=\varinjlim P_\bullet^\lambda$,
$Q_\bullet=\varinjlim Q_\bullet^\lambda$. Una vez definidos, con ayuda de
\sref{III.6.3.10}, los homomorfismos
\[
  \mathbf{Tor}_\bullet^{A^\lambda}
  (P_\bullet^\lambda,Q_\bullet^\lambda)
  \longrightarrow
  \mathbf{Tor}_\bullet^{A^\mu}
  (P_\bullet^\mu,Q_\bullet^\mu)
\]
para $\lambda\leqslant\mu$, la demostración es la de \sref{III.6.3.6}.
\end{namedtheorem}

\begin{proposition}[6.3.13]
\label{III.6.3.13}
Sean $S$ un subconjunto multiplicativo de $A$, y $P_\bullet$ y $Q_\bullet$
dos complejos de $A$-módulos en los que las multiplicaciones por los
elementos de $S$ sean biyectivas, de modo que, si $A'=S^{-1}A$,
$P_\bullet$ y $Q_\bullet$ estén formados por $A'$-módulos. Entonces se tiene
un isomorfismo canónico
$\mathbf{Tor}_\bullet^A(P_\bullet,Q_\bullet)
\xrightarrow{\sim}
\mathbf{Tor}_\bullet^{A'}(P_\bullet,Q_\bullet)$.
\end{proposition}

\begin{proof}
En efecto, la hipótesis implica que los homomorfismos canónicos
$P_\bullet\to P_\bullet\otimes_A A'$,
$Q_\bullet\to Q_\bullet\otimes_A A'$ son biyectivos. Por otra parte, la
funtorialidad del hipertor muestra que la multiplicación por todo $s\in S$
es biyectiva en $\mathbf{Tor}_\bullet^A(P_\bullet,Q_\bullet)$ y, por
consiguiente,
\[
  \mathbf{Tor}_\bullet^A(P_\bullet,Q_\bullet)
  \longrightarrow
  \mathbf{Tor}_\bullet^A(P_\bullet,Q_\bullet)\otimes_A A'
\]

\oldpage[III]{146}

es también un homomorfismo biyectivo. Como $A'$ es un $A$-módulo plano, la
conclusión resulta de \sref{III.6.3.9}, y se tienen asimismo isomorfismos
canónicos
para las sucesiones espectrales.
\end{proof}

\subsection{Funtores hipertor locales de complejos de módulos cuasicoherentes: caso de los esquemas afines}
\label{subsection:III.6.4}
\label{III.6.4}

\begin{env}[6.4.1]
Sean $S$ un esquema afín de anillo $A$, $X$, $Y$ dos $S$-esquemas afines
de anillos $B$, $C$ respectivamente, de modo que $B$ y $C$ son álgebras
sobre $A$. Todo complejo $\mathscr P_\bullet$ (resp. $\mathscr Q_\bullet$) de
$\mathscr O_X$-Módulos (resp. $\mathscr O_Y$-Módulos) cuasicoherentes es de
la forma $\widetilde{P}_\bullet$ (resp. $\widetilde{Q}_\bullet$), donde
$P_\bullet$ (resp. $Q_\bullet$) es un complejo de $B$-módulos (resp.
$C$-módulos) (\textbf{I}, 1.3.7 y 1.3.8). Se puede evidentemente considerar $P_\bullet$
y $Q_\bullet$ como complejos de $A$-módulos y formar los
$\mathbf{Tor}_n^A(P_\bullet,Q_\bullet)$; además, en virtud del carácter
bifuntorial de $\mathbf{Tor}_n^A(P_\bullet,Q_\bullet)$, las $A$-álgebras
$B$ y $C$ operan en este $A$-módulo, y estas operaciones hacen de él un
$(B,C)$-bimódulo, o, lo que equivale a lo mismo, un módulo sobre
$B\otimes_A C=A(X\times_S Y)$. Se ha definido pues de este modo un
$\mathscr O_{X\times_S Y}$-Módulo cuasicoherente
\[
  \boldsymbol{\mathscr{T}\!or}_n^{\mathscr O_S}(\mathscr P_\bullet,\mathscr Q_\bullet)
  =\bigl(\mathbf{Tor}_n^A(P_\bullet,Q_\bullet)\bigr)^{\sim}
  \tag{6.4.1.1}\label{III.6.4.1.1.fr}
\]
que se llama el \emph{hipertor local} de índice $n$ de los complejos
$\mathscr P_\bullet$ y $\mathscr Q_\bullet$ y que se denota también por
$\boldsymbol{\mathscr{T}\!or}_n^S(\mathscr P_\bullet,\mathscr Q_\bullet)$.
\end{env}

\begin{namedtheorem}[6.4.2]{Lema}
Con las notaciones de \textup{(6.4.1)}, supongamos que el anillo $A$ es de la forma
$R^{-1}A'$, donde $A'$ es un anillo y $R$ una parte
multiplicativa de $A'$. Sea $S'=\operatorname{Spec}(A')$. Entonces $X$ e $Y$
pueden considerarse como $S'$-preesquemas y
$X\times_{S'}Y=X\times_S Y$
\textup{(}\sref[I]{I.1.6.2} y \sref[I]{I.3.2.4}\textup{)}; además,
$\boldsymbol{\mathscr{T}\!or}_\bullet^{S'}(\mathscr P_\bullet,\mathscr Q_\bullet)
=\boldsymbol{\mathscr{T}\!or}_\bullet^S(\mathscr P_\bullet,\mathscr Q_\bullet)$.
\end{namedtheorem}

\begin{proof}
Esto se deduce de \sref{III.6.4.1.1} y \sref{III.6.3.13}.
\end{proof}

\begin{env}[6.4.3]
\label{III.6.4.3}
Con las notaciones e hipótesis de \sref{III.6.4.1}, sean
$\mathscr F=\widetilde F$ un $\mathscr O_X$-módulo cuasicoherente y
$\mathscr G=\widetilde G$ un $\mathscr O_Y$-módulo cuasicoherente.
Considerando $\mathscr F$ y $\mathscr G$ como complejos de Módulos, se denotará por $\mathscr{T}\!or_n^{\mathscr O_S}(\mathscr F,\mathscr G)$, o
por $\mathscr{T}\!or_n^S(\mathscr F,\mathscr G)$, su hipertor local de índice $n$.
Por \sref{III.6.3.5}(ii),
\[
  \mathscr{T}\!or_n^{\mathscr O_S}(\mathscr F,\mathscr G)
  =\bigl(\operatorname{Tor}_n^A(F,G)\bigr)^{\sim}.
  \tag{6.4.3.1}\label{III.6.4.3.1.fr}
\]

Volvamos ahora a dos complejos arbitrarios de módulos
cuasicoherentes $\mathscr P_\bullet$ y $\mathscr Q_\bullet$. Las fórmulas
\sref{III.6.4.1.1} y \sref{III.6.4.3.1}, junto con
\sref{III.6.3.2}, muestran que
$\boldsymbol{\mathscr{T}\!or}_\bullet^S(\mathscr P_\bullet,\mathscr Q_\bullet)$ es el objeto al
que convergen dos sucesiones espectrales
$'\mathscr E(\mathscr P_\bullet,\mathscr Q_\bullet)$ y
$''\mathscr E(\mathscr P_\bullet,\mathscr Q_\bullet)$, cuyos términos $E_2$ son
\[
  '\mathscr E_{pq}^2
  =\mathscr H_p\bigl(\mathscr{T}\!or_q^S
    (\mathscr P_\bullet,\mathscr Q_\bullet)\bigr)
  \tag{6.4.3.2}\label{III.6.4.3.2.fr}
\]
y
\[
  ''\mathscr E_{pq}^2
  =\bigoplus_{q'+q''=q}
  \mathscr{T}\!or_p^S
  \bigl(\mathscr H_{q'}(\mathscr P_\bullet),
        \mathscr H_{q''}(\mathscr Q_\bullet)\bigr)
  \tag{6.4.3.3}\label{III.6.4.3.3.fr}
\]
Aquí $\mathscr{T}\!or_q^S(\mathscr P_\bullet,\mathscr Q_\bullet)$ en
\sref{III.6.4.3.2} denota el bicomplejo de
$\mathscr O_{X\times_S Y}$-módulos cuasicoherentes
$\mathscr{T}\!or_q^S(\mathscr P_i,\mathscr Q_j)$.
\end{env}

\begin{env}[6.4.4]
Consideremos ahora otros dos esquemas afines
$X^{(1)}=\operatorname{Spec}(B^{(1)})$,
$Y^{(1)}=\operatorname{Spec}(C^{(1)})$, donde $B^{(1)}$ y $C^{(1)}$ son
$A$-álgebras, y supongamos dados dos
\oldpage[III]{147}
S-morfismos $u:X^{(1)}\to X$, $v:Y^{(1)}\to Y$, correspondientes a
$A$-homomorfismos
$\varphi:B\to B^{(1)}$ y
$\psi:C\to C^{(1)}$.
Consideremos los complejos
$u^*(\mathscr P_\bullet)=(P_\bullet\otimes_B B^{(1)})^{\sim}$ de
$\mathscr O_{X^{(1)}}$-Módulos,
$v^*(\mathscr Q_\bullet)=(Q_\bullet\otimes_C C^{(1)})^{\sim}$ de
$\mathscr O_{Y^{(1)}}$-Módulos. Los $A$-homomorfismos canónicos
\[
  P_\bullet\longrightarrow P_\bullet\otimes_B B^{(1)},
  \qquad
  Q_\bullet\longrightarrow Q_\bullet\otimes_C C^{(1)}
\]
dan, por funtorialidad, un $A$-homomorfismo
\[
  \mathbf{Tor}_\bullet^A(P_\bullet,Q_\bullet)
  \longrightarrow
  \mathbf{Tor}_\bullet^A
  (P_\bullet\otimes_B B^{(1)},Q_\bullet\otimes_C C^{(1)});
\]
De nuevo por funtorialidad, se trata en realidad de un homomorfismo
de $(B\otimes_A C)$-módulos. De donde se concluye que se ha definido así un $(u\times_S v)$-morfismo
\[
  \theta:\boldsymbol{\mathscr{T}\!or}_\bullet^S(\mathscr P_\bullet,\mathscr Q_\bullet)
  \longrightarrow
  \boldsymbol{\mathscr{T}\!or}_\bullet^S
  (u^*(\mathscr P_\bullet),v^*(\mathscr Q_\bullet))
  \tag{6.4.4.1}\label{III.6.4.4.1.fr}
\]
y, en consecuencia, un homomorfismo de
$\mathscr O_{X^{(1)}\times_S Y^{(1)}}$-módulos
\[
  \theta^\sharp:(u\times_S v)^*
  \bigl(\boldsymbol{\mathscr{T}\!or}_\bullet^S(\mathscr P_\bullet,\mathscr Q_\bullet)\bigr)
  \longrightarrow
  \boldsymbol{\mathscr{T}\!or}_\bullet^S
  (u^*(\mathscr P_\bullet),v^*(\mathscr Q_\bullet))
  \tag{6.4.4.2}\label{III.6.4.4.2.fr}
\]
Este es manifiestamente un morfismo de \emph{bi-$\partial$-funtores} en
las categorías de módulos cuasicoherentes acotados inferiormente.
El homomorfismo \sref{III.6.4.4.2} no es necesariamente biyectivo;
sin embargo:
\end{env}

\begin{namedtheorem}[6.4.5]{Lema}
Con las notaciones de \textup{(6.4.4)}, supongamos que $u$ y $v$ son inmersiones
abiertas; entonces el homomorfismo \textup{(6.4.4.2)} es biyectivo.
\end{namedtheorem}

\begin{proof}
Identifiquemos $X^{(1)}$ (resp. $Y^{(1)}$) con un abierto de $X$ (resp. $Y$);
$X^{(1)}$ (resp. $Y^{(1)}$) es entonces unión de abiertos de la forma
$\mathrm D(f)$ (resp. $\mathrm D(g)$), donde $f\in B$ (resp. $g\in C$), y los
preesquemas inducidos $\mathrm D(\varphi(f))$ y $\mathrm D(f)$ (resp.
$\mathrm D(\psi(g))$ y $\mathrm D(g)$) son isomorfos. Bastará probar
el lema cuando $X^{(1)}$ (resp. $Y^{(1)}$) es de la forma $\mathrm D(f)$
(resp. $\mathrm D(g)$); en efecto, si se establece este punto, y se vuelve al
caso general, bastará probar que la restricción de $\theta^\sharp$ a
cada abierto $\mathrm D(f)\times_S\mathrm D(g)$ es un isomorfismo; ahora bien, si
$u_1:\mathrm D(f)\to X^{(1)}$, $v_1:\mathrm D(g)\to Y^{(1)}$ son las
inclusiones canónicas, la restricción anterior no es otra que
$(u_1\times_S v_1)^*(\theta^\sharp)$; pero es inmediato, en virtud de las
definiciones (6.4.4) y de ($\mathbf{0}_{\mathrm I}$, 4.4.8), que al componerla con
el homomorfismo canónico
\[
  (u_1\times_S v_1)^*
  \bigl(\boldsymbol{\mathscr{T}\!or}_\bullet^S
    (u^*(\mathscr P_\bullet),v^*(\mathscr Q_\bullet))\bigr)
  \longrightarrow
  \boldsymbol{\mathscr{T}\!or}_\bullet^S
  (u'^*(\mathscr P_\bullet),v'^*(\mathscr Q_\bullet))
  \tag{6.4.5.1}\label{III.6.4.5.1.fr}
\]
donde $u'=u\circ u_1$ y $v'=v\circ v_1$, se obtiene el homomorfismo canónico
\[
  (u'\times_S v')^*
  \bigl(\boldsymbol{\mathscr{T}\!or}_\bullet^S(\mathscr P_\bullet,\mathscr Q_\bullet)\bigr)
  \longrightarrow
  \boldsymbol{\mathscr{T}\!or}_\bullet^S
  (u'^*(\mathscr P_\bullet),v'^*(\mathscr Q_\bullet))
  \tag{6.4.5.2}\label{III.6.4.5.2.fr}
\]
y si se sabe que (6.4.5.1) y (6.4.5.2) son isomorfismos,
resultará que lo mismo ocurre con $(u_1\times_S v_1)^*(\theta^\sharp)$.

Supongamos pues que $X^{(1)}=\mathrm D(f)$ y $Y^{(1)}=\mathrm D(g)$, de modo
que $B^{(1)}=B_f$ y $C^{(1)}=C_g$; $u^*(\mathscr P_\bullet)$ (resp.
$v^*(\mathscr Q_\bullet)$) se identifica entonces con
$(P_\bullet)_f^{\sim}$ (resp. $(Q_\bullet)_g^{\sim}$); por otra parte,
$X^{(1)}\times_S Y^{(1)}$ se identifica con el subesquema abierto
$\mathrm D(f\otimes g)$ de
$X\times_S Y=\operatorname{Spec}(B\otimes_A C)$ (\textbf{II}, 4.3.2.4); se trata de
probar que el homomorfismo
\[
  \bigl(\mathbf{Tor}_\bullet^A(P_\bullet,Q_\bullet)\bigr)_{f\otimes g}
  \longrightarrow
  \mathbf{Tor}_\bullet^A((P_\bullet)_f,(Q_\bullet)_g)
  \tag{6.4.5.3}\label{III.6.4.5.3.fr}
\]
\oldpage[III]{148}
deducido por funtorialidad de los homomorfismos canónicos
$P_\bullet\to(P_\bullet)_f$, $Q_\bullet\to(Q_\bullet)_g$, es biyectivo. Ahora bien
($\mathbf{0}_{\mathrm I}$, 1.6.1), se puede escribir
$(P_\bullet)_f=\varinjlim P_\bullet^{(n)}$, donde los $P_\bullet^{(n)}$ son todos
complejos de $B$-módulos idénticos a $P_\bullet$, siendo la aplicación
$P_\bullet^{(m)}\to P_\bullet^{(n)}$ para $m\leqslant n$ la
multiplicación por $f^{n-m}$; se tiene un resultado análogo para $Q_\bullet$
reemplazando $f$ por $g$; por otra parte, está claro que el homomorfismo
\[
  \mathbf{Tor}_\bullet^A(P_\bullet^{(m)},Q_\bullet^{(m)})
  \longrightarrow
  \mathbf{Tor}_\bullet^A(P_\bullet^{(n)},Q_\bullet^{(n)})
\]
correspondiente a los homomorfismos
$P_\bullet^{(m)}\to P_\bullet^{(n)}$ y
$Q_\bullet^{(m)}\to Q_\bullet^{(n)}$ es por definición la multiplicación
por $(f\otimes g)^{n-m}$. La conclusión resulta entonces de ($\mathbf{0}_{\mathrm I}$,
1.6.1) aplicado al primer miembro de (6.4.5.3) y de (6.3.6).
\end{proof}

\begin{env}[6.4.6]
\label{III.6.4.6}
Con las notaciones de \sref{III.6.4.4}, se definen del mismo modo
homomorfismos canónicos de funtores espectrales
\[
  \left\{
  \begin{aligned}
  (u\times_S v)^*\bigl('\mathscr E(\mathscr P_\bullet,\mathscr Q_\bullet)\bigr)
    &\longrightarrow
      '\mathscr E(u^*(\mathscr P_\bullet),v^*(\mathscr Q_\bullet))\\
  (u\times_S v)^*\bigl(''\mathscr E(\mathscr P_\bullet,\mathscr Q_\bullet)\bigr)
    &\longrightarrow
      ''\mathscr E(u^*(\mathscr P_\bullet),v^*(\mathscr Q_\bullet))
  \end{aligned}
  \right.
  \tag{6.4.6.1}\label{III.6.4.6.1.fr}
\]
El razonamiento de \sref{III.6.4.5} muestra que, cuando $u$ y $v$
son \emph{inmersiones abiertas}, estos homomorfismos son \emph{biyectivos}.
Teniendo en cuenta \sref{III.6.3.6.2} y \sref{III.6.3.6.3}, el
razonamiento demuestra que son isomorfismos en los términos $E^2$,
mientras que \sref{III.6.4.5} demuestra que son isomorfismos en los
objetos límite. La afirmación se deduce entonces de
\sref[0]{0.11.1.2} y \sref[0]{0.11.2.4}.
\end{env}

\subsection{Funtores hipertor locales de complejos de módulos cuasicoherentes: caso general}
\label{subsection:III.6.5}

\begin{env}[6.5.1]
Consideremos ahora un preesquema $S$ cualquiera y dos $S$-preesquemas
cualesquiera $X$, $Y$; sea $\mathscr P_\bullet$ (resp. $\mathscr Q_\bullet$)
un complejo de $\mathscr O_X$-Módulos (resp. $\mathscr O_Y$-Módulos)
cuasicoherentes. Pongamos $Z=X\times_S Y$; vamos a definir
$\mathscr O_Z$-Módulos cuasicoherentes
$\boldsymbol{\mathscr{T}\!or}_n^S(\mathscr P_\bullet,\mathscr Q_\bullet)$ llamados
\emph{hipertor locales} de $\mathscr P_\bullet$ y $\mathscr Q_\bullet$, que se
reducirán a los ya definidos en (6.4) cuando $S$, $X$ y $Y$ sean afines.

Cuando $\mathscr P_\bullet$ y $\mathscr Q_\bullet$ se reducen
respectivamente a sus términos de grado 0, $\mathscr F$ y $\mathscr G$, (siendo
nulos los demás), se escribirá
$\mathscr{T}\!or_n^S(\mathscr F,\mathscr G)$ en lugar de
$\boldsymbol{\mathscr{T}\!or}_n^S(\mathscr P_\bullet,\mathscr Q_\bullet)$.
\end{env}

\begin{env}[6.5.2]
Supongamos primero $S$ \emph{afín}, y sean $(X_\lambda)$, $(Y_\mu)$
recubrimientos de $X$ y $Y$ respectivamente por abiertos afines; entonces los
$Z_{\lambda\mu}=X_\lambda\times_S Y_\mu$ forman un recubrimiento abierto afín
de $Z$. Pongamos
$\mathscr P_{\lambda\bullet}=\mathscr P_\bullet|X_\lambda$,
$\mathscr Q_{\mu\bullet}=\mathscr Q_\bullet|Y_\mu$; tenemos pues para todo
par $(\lambda,\mu)$ un $\mathscr O_{Z_{\lambda\mu}}$-Módulo cuasicoherente
$\mathscr F_{\lambda\mu}
  =\boldsymbol{\mathscr{T}\!or}_n^S
   (\mathscr P_{\lambda\bullet},\mathscr Q_{\mu\bullet})$,
y hay que mostrar que los $\mathscr F_{\lambda\mu}$ verifican la condición de
pegado ($\mathbf{0}_{\mathrm I}$, 3.3.1). Para ello, basta verificar que para
todo abierto afín $U\subset X_\lambda\cap X_{\lambda'}$ (resp.
$V\subset Y_\mu\cap Y_{\mu'}$), las restricciones de
$\mathscr F_{\lambda\mu}$ y de $\mathscr F_{\lambda'\mu'}$ a $U\times_S V$
son canónicamente isomorfas; pero esto se deduce inmediatamente de la existencia
de isomorfismos canónicos de estas restricciones sobre
$\boldsymbol{\mathscr{T}\!or}_n^S(\mathscr P_\bullet|U,\mathscr Q_\bullet|V)$ (6.4.5).
Además, resulta inmediatamente de esta definición y de (6.4.5) que el
$\mathscr O_Z$-Módulo así definido no depende (salvo isomorfismo) de los
recubrimientos abiertos $(X_\lambda)$, $(Y_\mu)$ consi\oldpage[III]{149}derados;
lo denotaremos pues por
$\boldsymbol{\mathscr{T}\!or}_n^S(\mathscr P_\bullet,\mathscr Q_\bullet)$; resulta finalmente
de (6.4.5) que para \emph{toda} parte abierta $U$ (resp. $V$) de $X$ (resp. $Y$), la
restricción de $\boldsymbol{\mathscr{T}\!or}_n^S(\mathscr P_\bullet,\mathscr Q_\bullet)$ a
$U\times_S V$ es canónicamente isomorfa a
$\boldsymbol{\mathscr{T}\!or}_n^S(\mathscr P_\bullet|U,\mathscr Q_\bullet|V)$.
\end{env}

\begin{env}[6.5.3]
Pasemos ahora al caso general en que $S$ es arbitrario, y sea $(S_\alpha)$
un recubrimiento de $S$ formado por abiertos afines; designemos por $X_\alpha$
(resp. $Y_\alpha$) la imagen inversa de $S_\alpha$ en $X$ (resp. $Y$);
hay que probar también que los haces
$\boldsymbol{\mathscr{T}\!or}_n^{S_\alpha}
  (\mathscr P_\bullet|X_\alpha,\mathscr Q_\bullet|Y_\alpha)
  =\mathscr G_\alpha$
verifican la condición de pegado. Basta definir, para todo abierto
afín $T$ contenido en $S_\alpha\cap S_\beta$, isomorfismos canónicos de las
restricciones de $\mathscr G_\alpha$ y $\mathscr G_\beta$ a $U\times_S V$ (designando
por $U$ y $V$ las imágenes inversas de $T$ en $X$ y $Y$
respectivamente) sobre
$\boldsymbol{\mathscr{T}\!or}_n^T(\mathscr P_\bullet|U,\mathscr Q_\bullet|V)$; se puede,
además, limitar al caso en que $T$ se escribe a la vez $\mathrm D(f_\alpha)$ y
$\mathrm D(f_\beta)$, siendo $f_\alpha$ (resp. $f_\beta$) una sección de
$\mathscr O_S$ sobre $S_\alpha$ (resp. $S_\beta$); pero entonces
$\boldsymbol{\mathscr{T}\!or}_n^T(\mathscr P_\bullet|U,\mathscr Q_\bullet|V)$ es
canónicamente isomorfo a
$\boldsymbol{\mathscr{T}\!or}_n^{S_\alpha}(\mathscr P_\bullet|U,\mathscr Q_\bullet|V)$
por una parte, a
$\boldsymbol{\mathscr{T}\!or}_n^{S_\beta}(\mathscr P_\bullet|U,\mathscr Q_\bullet|V)$
por otra, en virtud de (6.4.2); como se acaban de definir isomorfismos
canónicos de $\mathscr G_\alpha$ sobre
$\boldsymbol{\mathscr{T}\!or}_n^{S_\alpha}(\mathscr P_\bullet|U,\mathscr Q_\bullet|V)$ y
de $\mathscr G_\beta$ sobre
$\boldsymbol{\mathscr{T}\!or}_n^{S_\beta}(\mathscr P_\bullet|U,\mathscr Q_\bullet|V)$
(6.5.2), esto termina de definir el $\mathscr O_Z$-Módulo
$\boldsymbol{\mathscr{T}\!or}_n^S(\mathscr P_\bullet,\mathscr Q_\bullet)$. Además, para
toda parte abierta $U$ (resp. $V$) de $X$ (resp. $Y$),
$\boldsymbol{\mathscr{T}\!or}_n^S(\mathscr P_\bullet|U,\mathscr Q_\bullet|V)$ es
canónicamente isomorfo a la restricción de
$\boldsymbol{\mathscr{T}\!or}_n^S(\mathscr P_\bullet,\mathscr Q_\bullet)$ a $U\times_S V$.

Es inmediato que se ha definido así (en las categorías de complejos de
Módulos cuasicoherentes acotados inferiormente) un \emph{bi-$\partial$-funtor}
$\boldsymbol{\mathscr{T}\!or}_\bullet^S(\mathscr P_\bullet,\mathscr Q_\bullet)$ con valores
en la categoría de los $\mathscr O_Z$-Módulos, pues está claro que la cuestión
es local sobre $X$, $Y$ y $S$, en virtud de (6.4.5) y de la observación de que
(6.4.4.2) es un morfismo de bi-$\partial$-funtores. Se observará que si
$\mathscr P_\bullet$ y $\mathscr Q_\bullet$ se reducen respectivamente a
sus términos de grado 0, $\mathscr F$ y $\mathscr G$,
$\mathscr{T}\!or_0^S(\mathscr F,\mathscr G)$ no es otro, en virtud de
(6.4.1.1), que el producto tensorial externo $\mathscr F\otimes_S\mathscr G$
definido en (\textbf{I}, 9.1.2); esto resulta en efecto de (\textbf{I}, 9.1.3).
\end{env}

\begin{env}[6.5.4]
Resulta de la construcción precedente y de las observaciones hechas en (6.4.6) que $\boldsymbol{\mathscr{T}\!or}_\bullet^S(\mathscr P_\bullet,\mathscr Q_\bullet)$ es el objeto al
que convergen dos funtores espectrales
$'\mathscr E(\mathscr P_\bullet,\mathscr Q_\bullet)$ y
$''\mathscr E(\mathscr P_\bullet,\mathscr Q_\bullet)$, cuyos términos $E_2$ son
\[
  '\mathscr E_{pq}^2
  =\mathscr H_p\bigl(\mathscr{T}\!or_q^S
    (\mathscr P_\bullet,\mathscr Q_\bullet)\bigr)
  \tag{6.5.4.1}\label{III.6.5.4.1.fr}
\]
y
\[
  ''\mathscr E_{pq}^2
  =\bigoplus_{q'+q''=q}
   \mathscr{T}\!or_p^S
   \bigl(\mathscr H_{q'}(\mathscr P_\bullet),
         \mathscr H_{q''}(\mathscr Q_\bullet)\bigr)
  \tag{6.5.4.2}\label{III.6.5.4.2.fr}
\]
La sucesión espectral \sref{III.6.5.4.2} es siempre regular.
Ambas sucesiones espectrales son birregulares cuando
$\mathscr P_\bullet$ y $\mathscr Q_\bullet$ están \emph{acotados inferiormente}.
A continuación se describe otro caso en que son birregulares.
\end{env}

\begin{env}[6.5.5]
Diremos que sobre un espacio topológico $T$ un haz de anillos
$\mathscr A$ es \emph{de dimensión cohomológica $\leqslant n$} si, para todo
$t\in T$, el anillo $\mathscr A_t$ es de dimensión cohomológica $\leqslant n$;
se dirá entonces también que el \emph{espacio anillado} $(T,\mathscr A)$ es
\emph{de dimensión cohomológica $\leqslant n$}. Se dirá que un haz
de anillos (resp. un espacio anillado) es de dimensión cohomológica \emph{finita}
si existe un entero $n$ tal que es de dimensión cohomológica
$\leqslant n$. Se observará que si los $\mathscr A_t$ son \emph{anillos locales
(conmutativos) noetherianos},
\oldpage[III]{150}
decir que son de dimensión cohomológica $\leqslant n$ significa que son
\emph{regulares} y de \emph{dimensión} (de Krull) $\leqslant n$ ($\mathbf{0}_{\mathrm{IV}}$, 17.3.1).
Con la terminología de la teoría de la dimensión que introduciremos en el
cap. IV, equivale a decir que un \emph{preesquema localmente noetheriano} $T$
es de \emph{dimensión cohomológica} $\leqslant n$, o a decir que es \emph{regular}
($\mathbf{0}_{\mathrm I}$, 4.1.4) y de \emph{dimensión} $\leqslant n$; esto significa que para
todo abierto afín $U$ de $T$, el anillo $\Gamma(U,\mathscr O_T)$ es \emph{de
dimensión cohomológica} $\leqslant n$ ($\mathbf{0}_{\mathrm{IV}}$, 17.2.6). Dicho esto,
esta última observación, junto con (6.3.2), prueba que si $S$ es \emph{localmente
noetheriano y de dimensión cohomológica finita}, las sucesiones espectrales
$'\mathscr E(\mathscr P_\bullet,\mathscr Q_\bullet)$ y
$''\mathscr E(\mathscr P_\bullet,\mathscr Q_\bullet)$ son \emph{birregulares}.

Está claro que
$\boldsymbol{\mathscr{T}\!or}_\bullet^S(\mathscr P_\bullet,\mathscr Q_\bullet)$ se
transforma en
$\boldsymbol{\mathscr{T}\!or}_\bullet^S(\mathscr Q_\bullet,\mathscr P_\bullet)$ (salvo
isomorfismo) por el isomorfismo canónico de $X\times_S Y$ sobre
$Y\times_S X$.
\end{env}

\begin{proposition}[6.5.6]
\label{III.6.5.6}
Sea $(\mathscr P_{\alpha\bullet})$ un sistema inductivo filtrante de
complejos de $\mathscr O_X$-módulos cuasicoherentes. Existe un
isomorfismo canónico
\[
  \varinjlim
  (\boldsymbol{\mathscr{T}\!or}_\bullet^S
  (\mathscr P_{\alpha\bullet},\mathscr Q_\bullet))
  \xrightarrow{\sim}
  \boldsymbol{\mathscr{T}\!or}_\bullet^S
  \left(\varinjlim\mathscr P_{\alpha\bullet},\mathscr Q_\bullet\right).
  \tag{6.5.6.1}\label{III.6.5.6.1.fr}
\]
\end{proposition}

\begin{proof}
La cuestión es local en $S$, $X$ e $Y$, se puede suponer $S$, $X$ e $Y$ son afines, en cuyo caso la proposición
se reduce a \sref{III.6.3.6}.
\end{proof}

\begin{namedtheorem}[6.5.7]{Observaciones}
\textnormal{(i)} Consideremos en particular el caso en que $S=X=Y$,
siendo pues $\mathscr P_\bullet$ y $\mathscr Q_\bullet$ dos complejos de
$\mathscr O_S$-Módulos cuasicoherentes; entonces los
$\boldsymbol{\mathscr{T}\!or}_n^S(\mathscr P_\bullet,\mathscr Q_\bullet)$ son
$\mathscr O_S$-Módulos cuasicoherentes; además, para todo punto $z\in S$,
resulta de (6.5.6) que se tiene un isomorfismo canónico
\[
  \bigl(\boldsymbol{\mathscr{T}\!or}_n^S
    (\mathscr P_\bullet,\mathscr Q_\bullet)\bigr)_z
  \xrightarrow{\sim}
  \mathbf{Tor}_n^{\mathscr O_z}
  \bigl((\mathscr P_\bullet)_z,(\mathscr Q_\bullet)_z\bigr)
  \tag{6.5.7.1}\label{III.6.5.7.1.fr}
\]
pues la cuestión es local y se está reducido al caso de los módulos, en virtud de
(6.4.1.1).

\textnormal{(ii)} Se puede generalizar la definición de los hipertor al caso de
\emph{dos complejos de $\mathscr O_X$-Módulos} $\mathscr P_\bullet$,
$\mathscr Q_\bullet$ sobre un mismo \emph{espacio anillado} $(X,\mathscr O_X)$; para todo
abierto $U$ de $X$, pongamos en efecto
$A(U)=\Gamma(U,\mathscr O_X)$,
$P_\bullet(U)=\Gamma(U,\mathscr P_\bullet)$,
$Q_\bullet(U)=\Gamma(U,\mathscr Q_\bullet)$; los $A(U)$-módulos
$\mathbf{Tor}_n^{A(U)}(P_\bullet(U),Q_\bullet(U))$
forman entonces un \emph{prehaz} sobre $X$, y se designa por
$\boldsymbol{\mathscr{T}\!or}_n^X(\mathscr P_\bullet,\mathscr Q_\bullet)$ el
$\mathscr O_X$-Módulo asociado a este prehaz. Cuando $X$ es un \emph{preesquema},
resulta de (6.3.12) que este $\mathscr O_X$-Módulo es canónicamente isomorfo
al hipertor definido antes. No desarrollaremos más esta
generalización.
\end{namedtheorem}

\begin{proposition}[6.5.8]
Sean $X$, $Y$ dos $S$-preesquemas, $\mathscr F$ \textup{(resp.} $\mathscr G$\textup{)} un
$\mathscr O_X$-Módulo \textup{(resp.} un $\mathscr O_Y$-Módulo\textup{)} cuasicoherente. Si
$\mathscr F$ o $\mathscr G$ es $S$-plano, se tiene
$\mathscr{T}\!or_n^S(\mathscr F,\mathscr G)=0$ para $n\neq0$.
\end{proposition}

\begin{proof}
Al ser la cuestión local sobre $X$ y $Y$, se puede suponer $X$, $Y$ y $S$
afines, de anillos respectivos $B$, $C$, $A$, y
$\mathscr F=\widetilde M$, $\mathscr G=\widetilde N$, siendo $M$ (resp. $N$)
un $B$-módulo (resp. un $C$-módulo). Supongamos por ejemplo que $\mathscr F$
es $S$-plano, lo que significa que para todo $s\in S$, $M_s$ es un
$A_s$-módulo plano ($\mathbf{0}_{\mathrm I}$, 6.7.1); por tanto $M$ es un $A$-módulo plano
($\mathbf{0}_{\mathrm I}$, 6.3.3), y se sabe que
$\operatorname{Tor}_n^A(M,N)=0$ para $n>0$ y para todo $C$-módulo $N$
($\mathbf{0}_{\mathrm I}$, 6.1.1), de donde la conclusión por (6.4.1.1).
\end{proof}

\begin{namedtheorem}[6.5.9]{Corolario}
\label{III.6.5.9}
Sean $X$ e $Y$ dos $S$-preesquemas, y sea $\mathscr P_\bullet$
\textup{(resp.} $\mathscr Q_\bullet$\textup{)} un complejo
\oldpage[III]{151}
acotado
inferiormente de $\mathscr O_X$-módulos cuasicoherentes
\textup{(resp.} de $\mathscr O_Y$-módulos cuasicoherentes\textup{)}.
Supongamos que cada $\mathscr P_i$ es plano sobre $S$. Entonces existe
un isomorfismo canónico de $\partial$-funtores en $\mathscr Q_\bullet$
\[
  \boldsymbol{\mathscr{T}\!or}_\bullet^S(\mathscr P_\bullet,\mathscr Q_\bullet)
  \xrightarrow{\sim}
  \mathscr H_\bullet(\mathscr P_\bullet\otimes_S\mathscr Q_\bullet).
  \tag{6.5.9.1}\label{III.6.5.9.1.fr}
\]
\end{namedtheorem}

\begin{proof}
Cuando $S$, $X$ e $Y$ son afines, esto es precisamente
\sref{III.6.3.7}; el caso general se deduce mediante los
razonamientos de \sref{III.6.5.2} y \sref{III.6.5.3}.
\end{proof}

\begin{namedtheorem}[6.5.10]{Corolario}
Supongamos que $X$ es plano sobre $S$ \textup{($\mathbf{0}_{\mathrm I}$, 6.7.1)}, que
$\mathscr P_\bullet$ y $\mathscr Q_\bullet$ están acotados inferiormente y
que todos los $\mathscr P_i$ son $\mathscr O_X$-Módulos localmente
libres (no necesariamente de tipo finito). Entonces el homomorfismo \textup{(6.5.9.1)} es
biyectivo.
\end{namedtheorem}

\begin{proof}
En efecto, se cumple la hipótesis de (6.5.9), al ser la planitud una propiedad
puntual sobre $X$ por definición y toda suma directa de módulos planos
un módulo plano ($\mathbf{0}_{\mathrm I}$, 6.1.2).
\end{proof}

\begin{proposition}[6.5.11]
\label{III.6.5.11}
Sean $X'$ e $Y'$ dos $S$-preesquemas, y sean
$f:X\to X'$ y $g:Y\to Y'$ dos $S$-morfismos afines. Para un
complejo $\mathscr P_\bullet$ \textup{(resp.}
$\mathscr Q_\bullet$\textup{)} un complejo de $\mathscr O_X$-Módulos \textup{(resp.} de
$\mathscr O_Y$-Módulos\textup{)} cuasicoherentes; se tiene entonces un isomorfismo canónico funtorial
\[
  (f\times_S g)_*
  \bigl(\boldsymbol{\mathscr{T}\!or}_\bullet^S
    (\mathscr P_\bullet,\mathscr Q_\bullet)\bigr)
  \xrightarrow{\sim}
  \boldsymbol{\mathscr{T}\!or}_\bullet^S
  (f_*(\mathscr P_\bullet),g_*(\mathscr Q_\bullet)).
  \tag{6.5.11.1}\label{III.6.5.11.1.fr}
\]
\end{proposition}

\begin{proof}
Como $f$ y $g$ son afines, $f_*(\mathscr P_\bullet)$ y
$g_*(\mathscr Q_\bullet)$ son complejos de Módulos cuasicoherentes (\textbf{II}, 1.2.6), y si se pone $Z'=X'\times_S Y'$, ambos miembros de
\sref{III.6.5.11.1} son $\mathscr O_{Z'}$-Módulos cuasicoherentes (6.5.1); se reduce fácilmente al caso en que $S$, $X'$ e $Y'$ son afines; pero entonces lo mismo ocurre por hipótesis con $X$ e $Y$ también son
afines, y la afirmación se deduce de \sref{III.6.4.1.1} y
\sref[I]{I.1.6.3}.
\end{proof}

\begin{namedtheorem}[6.5.12]{Observación}
Sean $X'$, $Y'$ dos $S$-preesquemas y supongamos que $X'$ es $S$-\emph{plano};
sean $X$ un subpreesquema cerrado de $X'$, $i:X\to X'$ la inclusión canónica,
$\mathscr P_\bullet$ (resp. $\mathscr Q_\bullet$) un complejo de
$\mathscr O_X$-Módulos (resp. de $\mathscr O_{Y'}$-Módulos) cuasicoherentes,
acotado inferiormente. Sea finalmente $\mathscr L'_{\bullet\bullet}$ una resolución
de $i_*(\mathscr P_\bullet)$ formada por $\mathscr O_{X'}$-Módulos localmente
libres, tal que todo punto de $X'$ tiene una vecindad abierta afín $U$ para
la cual $\mathscr L'_{j,\bullet}|U$ es una resolución libre de
$i_*(\mathscr P_j)|U$ para todo $j$. Se tiene entonces un isomorfismo canónico
\[
  (i\times_S 1)_*
  \bigl(\boldsymbol{\mathscr{T}\!or}_\bullet^S
    (\mathscr P_\bullet,\mathscr Q_\bullet)\bigr)
  \xrightarrow{\sim}
  \mathscr H_\bullet
  (\mathscr L'_{\bullet\bullet}\otimes_S\mathscr Q_\bullet).
  \tag{6.5.12.1}\label{III.6.5.12.1.fr}
\]

Si $S$, $X'$, $Y'$ son afines y si $\mathscr L'_{j,\bullet}$ es una
resolución libre de $i_*(\mathscr P_j)$ para todo $j$, se está reducido, en virtud
de (6.5.11), al caso en que $X'=X$ es $S$-plano, y basta aplicar (6.3.4).
En el caso general, se define localmente el isomorfismo (6.5.12.1), y se
trata de verificar que esta definición da efectivamente un isomorfismo global.
Para ello, hay que remitirse a la definición del primer isomorfismo
(6.3.4.1), que procede de un isomorfismo de sucesiones espectrales (\textbf{0}, 11.6.5 y
11.5.3), obtenido a su vez a partir de un morfismo de bicomplejos
$L_{\bullet\bullet}\to L''_{\bullet\bullet}$, donde $L_{\bullet\bullet}$ es la
resolución dada de $P_\bullet$, $L''_{\bullet\bullet}$ una resolución
proyectiva de $P_\bullet$ en la categoría de complejos de $A$-módulos
acotados
\oldpage[III]{152}
inferiormente (cf. (\textbf{0}, 11.5.2.2)); nuestra afirmación resulta de que
el isomorfismo (6.3.4.1) no depende de la resolución proyectiva
$L''_{\bullet\bullet}$ elegida, en virtud de la existencia de un homotopismo entre
dos de tales resoluciones (M, V, 1.2).
\end{namedtheorem}

\begin{proposition}[6.5.13]
\label{III.6.5.13}
Sean $X$ e $Y$ dos $S$-preesquemas, y supongamos que se cumple
una de las condiciones siguientes:
\begin{enumerate}
  \item[\textup{(i)}] $X$ y $Z=X\times_S Y$ son localmente noetherianos, y
    $X$ es plano sobre $S$;
  \item[\textup{(ii)}] $S$ y $X$ son localmente noetherianos, y $Y$ es de
    tipo finito sobre $S$.
\end{enumerate}
Sea $\mathscr P_\bullet$ \textup{(resp.} $\mathscr Q_\bullet$) un complejo
acotado inferiormente de $\mathscr O_X$-módulos cuasicoherentes
\textup{(resp.} de $\mathscr O_Y$-Módulos) cuasicoherentes, acotado inferiormente. Se supone además que, para todo $n$,
$\mathscr H_n(\mathscr P_\bullet)$ \textup{(resp.}
$\mathscr H_n(\mathscr Q_\bullet)$) es un $\mathscr O_X$-módulo
\textup{(resp.} un $\mathscr O_Y$-módulo) de tipo finito. Entonces
$\mathscr{T}\!or_n^S(\mathscr P_\bullet,\mathscr Q_\bullet)$
es un $\mathscr O_Z$-Módulos coherentes.
\end{proposition}

\begin{proof}
Como $\mathscr P_\bullet$ y $\mathscr Q_\bullet$ están acotados
inferiormente, la sucesión espectral
$''\mathscr E(\mathscr P_\bullet,\mathscr Q_\bullet)$ es birregular por
\sref{III.6.5.4}. Tanto en el caso (i) como en el caso (ii), $Z$
es localmente noetheriano, de modo que \sref[0]{0.11.1.8} reduce
la afirmación a la coherencia de los términos
$''\mathscr E_{pq}^2$ son coherentes. La hipótesis sobre los
$\mathscr H_n(\mathscr P_\bullet)$ y
$\mathscr H_n(\mathscr Q_\bullet)$ y la expresión (6.5.4.2) de los
$''\mathscr E_{pq}^2$ muestran pues que la proposición equivale a su
caso particular correspondiente a $\mathscr P_\bullet$ y $\mathscr Q_\bullet$ están
concentrados en grado cero, es decir, al corolario siguiente.
\end{proof}

\begin{corollary}[6.5.14]
Supongamos verificada una de las condiciones \textup{(i), (ii)} de \textup{(6.5.13)}, y sea
$\mathscr F$ \textup{(resp.} $\mathscr G$) un $\mathscr O_X$-Módulo \textup{(resp.} un
$\mathscr O_Y$-Módulo) cuasicoherente de tipo finito; entonces los
$\mathscr{T}\!or_n^S(\mathscr F,\mathscr G)$ son $\mathscr O_Z$-Módulos
coherentes.
\end{corollary}

\begin{proof}
Al ser la cuestión local sobre $X$ y $Y$, se puede suponer $S$, $X$, $Y$
afines.

\textnormal{(i)} Bajo las hipótesis de (i), $S$, $X$ y $Z$ son noetherianos.
Existe pues una resolución localmente libre $\mathscr L_\bullet$ de
$\mathscr F$ formada por $\mathscr O_X$-Módulos de tipo finito (\textbf{I}, 1.3.7); como
$X$ es plano sobre $S$, resulta de (6.3.4) que se tiene
$\mathscr{T}\!or_n^S(\mathscr F,\mathscr G)
=\mathscr H_n(\mathscr L_\bullet\otimes_S\mathscr G)$;
ahora bien, los $\mathscr L_i\otimes_S\mathscr G$ son $\mathscr O_Z$-Módulos
cuasicoherentes de tipo finito (\textbf{I}, 9.1.1), luego coherentes. Se concluye que
$\mathscr H_n(\mathscr L_\bullet\otimes_S\mathscr G)$ es coherente
($\mathbf{0}_{\mathrm I}$, 5.3.4).

\textnormal{(ii)} Supongamos ahora verificadas las condiciones (ii). Como
el anillo $A(Y)$ es cociente de una $A(S)$-álgebra de polinomios en un número finito
de indeterminadas (\textbf{I}, 6.3.3), $Y$ es un sub-$S$-preesquema cerrado de un
$S$-preesquema afín $Y'$, plano y de tipo finito sobre $S$; al ser $Y'$ noetheriano
(\textbf{I}, 6.3.7), existe una resolución localmente libre $\mathscr M_\bullet$ de
$j_*(\mathscr G)$ por $\mathscr O_{Y'}$-Módulos de tipo finito
(siendo $j:Y\to Y'$ la inclusión canónica); en virtud de (6.5.12),
$\mathscr{T}\!or_n^S(\mathscr F,\mathscr G)$ es la imagen inversa por
$1\times j$ del $\mathscr O_{Z'}$-Módulo
$\mathscr H_n(\mathscr F\otimes_S\mathscr M_\bullet)$ (donde
$Z'=X\times_S Y'$); se ve como en (i) que los
$\mathscr F\otimes_S\mathscr M_i$ son $\mathscr O_{Z'}$-Módulos
coherentes, y se obtiene también la conclusión por ($\mathbf{0}_{\mathrm I}$, 5.3.4).
\end{proof}

\begin{env}[6.5.15]
La teoría desarrollada antes para dos complejos $\mathscr P_\bullet$,
$\mathscr Q_\bullet$ de haces cuasicoherentes sobre dos $S$-preesquemas $X$ e $Y$ se extiende sin dificultad a una
familia finita arbitraria de $S$-preesquemas $X^{(i)}$ ($1\leqslant i\leqslant m$) y sobre
cada $X^{(i)}$ un complejo $\mathscr P_\bullet^{(i)}$ de $\mathscr O_{X^{(i)}}$-módulos
cuasicoherentes. Si
$Z=X^{(1)}\times_S X^{(2)}\times_S\ldots\times_S X^{(m)}$,
se define así un $\mathscr O_Z$-módulo cuasicoherente
$\mathscr{T}\!or_q^S(\mathscr P_\bullet^{(1)},\mathscr P_\bullet^{(2)},
\ldots,\mathscr P_\bullet^{(m)})$.
Dejamos al lector el desarrollo de la teoría en este caso general,
y registramos únicamente, para referencia posterior, el término
$E_2$ de la segunda sucesión espectral
(regular) cuyo resultado es
$\mathscr{T}\!or_\bullet^S(\mathscr P_\bullet^{(1)},
\mathscr P_\bullet^{(2)},\ldots,\mathscr P_\bullet^{(m)})$:
\oldpage[III]{153}
\[
  ''\mathscr E_{pq}^2
  =\bigoplus_{q_1+q_2+\ldots+q_m=q}
   \mathscr{T}\!or_p^S
   \bigl(\mathscr H_{q_1}(\mathscr P_\bullet^{(1)}),\ldots,
         \mathscr H_{q_m}(\mathscr P_\bullet^{(m)})\bigr).
  \tag{6.5.15.1}\label{III.6.5.15.1.fr}
\]
En \sref{III.6.8} estudiaremos las sucesiones espectrales de
asociatividad a las que dan lugar estos funtores hipertor de un
número finito arbitrario de complejos.
\end{env}

\subsection{Funtores de hipertor global para complejos de módulos
cuasicoherentes y sucesiones espectrales de Künneth: caso de base afín}
\label{subsection:III.6.6}
\label{subsection:III.6.6.fr}

\begin{env}[6.6.1]
\label{III.6.6.1}
Consideremos un esquema afín $S=\operatorname{Spec}(A)$ y dos $S$-\emph{esquemas}
cuasicompactos $X^{(i)}$ ($i=1,2$); sea $\mathscr P_\bullet^{(i)}$ un complejo
de $\mathscr O_{X^{(i)}}$-módulos cuasicoherentes, \emph{acotado inferiormente},
cuyo diferencial es de grado $-1$ ($i=1,2$). Consideremos, por otra parte, un
recubrimiento \emph{finito}
$\mathfrak U^{(i)}=(U_\alpha^{(i)})$ de $X^{(i)}$ por abiertos afines; sea
$X=X^{(1)}\times_S X^{(2)}$,
que es un $S$-esquema cuasicompacto (\textbf I, 5.5.1 y 6.6.4), y sea
$\mathfrak U=\mathfrak U^{(1)}\times_S\mathfrak U^{(2)}$ el recubrimiento de
$X$ formado por los abiertos afines
$U_\alpha^{(1)}\times_S U_\beta^{(2)}$. Para todo par de enteros
$p\leqslant0$, $q\in\mathbf Z$, el grupo de las $(-p)$-cocadenas alternadas
$C^{-p}(\mathfrak U^{(i)},\mathscr P_q^{(i)})$
del recubrimiento $\mathfrak U^{(i)}$ con coeficientes en el haz
$\mathscr P_q^{(i)}$ (G, II, 5.1) es un $A$-módulo; para $p>0$, se pondrá
$C^p(\mathfrak U^{(i)},\mathscr P_q^{(i)})=0$; queda así definido un
\emph{bicomplejo}
$C^\bullet(\mathfrak U^{(i)},\mathscr P_\bullet^{(i)})$
de $A$-módulos, cuyos \emph{dos} diferenciales son de grado $-1$. De las
definiciones (\textbf{0}, 12.1.2) resulta que el $A$-módulo de homología
$\mathrm H_n\bigl(C^\bullet
(\mathfrak U^{(i)},\mathscr P_\bullet^{(i)})\bigr)$
de este bicomplejo (considerado, como de costumbre, como un complejo simple
para el grado total) no es sino el $A$-módulo de \emph{hipercohomología}
$\mathbf H^{-n}
(\mathfrak U^{(i)},\mathscr P^{(i)\bullet})$
donde $\mathscr P^{(i)\bullet}$ es el complejo con diferencial de grado $+1$
obtenido tomando $\mathscr P_{-q}^{(i)}$ como componente de grado $q$; por
abuso de notación, lo escribiremos
$\mathbf H^{-n}(\mathfrak U^{(i)},\mathscr P_\bullet^{(i)})$. Resulta entonces
de (6.2.2) que este $A$-módulo es canónicamente isomorfo al $A$-módulo de
hipercohomología
$\mathbf H^{-n}(X^{(i)},\mathscr P^{(i)\bullet})$, que escribiremos asimismo
$\mathbf H^{-n}(X^{(i)},\mathscr P_\bullet^{(i)})$; por tanto, no depende del
recubrimiento finito $\mathfrak U^{(i)}$ elegido.
\end{env}

\begin{env}[6.6.2]
Vamos a aplicar a los dos bicomplejos de $A$-módulos
\[
  L_{\bullet\bullet}^{(i)}
  =C^\bullet(\mathfrak U^{(i)},\mathscr P_\bullet^{(i)})
  \qquad(i=1,2)
\] y al bifuntor covariante $L_{\bullet\bullet}^{(1)}\otimes_A L_{\bullet\bullet}^{(2)}$ a estos dos bicomplejos, la teoría general de la hiperhomología de los funtores respecto de los bicomplejos (\textbf{0}, 11.7.4). Como las cocadenas consideradas son \emph{alternadas} y los recubrimientos $\mathfrak U^{(i)}$ finitos, se observará que los módulos $C^{-p}(\mathfrak U^{(i)},\mathscr P_q^{(i)})$ solo son $\ne0$
para un número \emph{finito (independiente de $q$)} de valores de $p$ y, en
particular, las \emph{dos} graduaciones de cada $L_{\bullet\bullet}^{(i)}$ están
\emph{acotadas inferiormente}. Denotaremos por
$\mathbf{Tor}_n^A(L_{\bullet\bullet}^{(1)},L_{\bullet\bullet}^{(2)})$
o por
$\mathbf{Tor}_n^S
(\mathfrak U^{(1)},\mathfrak U^{(2)};
\mathscr P_\bullet^{(1)},\mathscr P_\bullet^{(2)})$
el $n$-ésimo módulo de hiperhomología de
$L_{\bullet\bullet}^{(1)}\otimes_A L_{\bullet\bullet}^{(2)}$, al que
llamaremos el \emph{hipertor} de índice $n$ de
$\mathscr P_\bullet^{(1)}$ y $\mathscr P_\bullet^{(2)}$, relativo a los
recubrimientos $\mathfrak U^{(1)}$ y $\mathfrak U^{(2)}$. Cuando
$\mathscr P_\bullet^{(1)}$ y $\mathscr P_\bullet^{(2)}$ están reducidos a sus
términos de grado 0, $\mathscr F^{(1)}$ y $\mathscr F^{(2)}$, se escribe
$\mathscr{T}\!or_n^S
(\mathfrak U^{(1)},\mathfrak U^{(2)};
\mathscr F^{(1)},\mathscr F^{(2)})$ su hipertor. Se designará por $\mathscr{T}\!or_n^S
(\mathfrak U^{(1)},\mathfrak U^{(2)};
\mathscr P_\bullet^{(1)},\mathscr P_\bullet^{(2)})$, según las convenciones generales, el \emph{bicomplejo} cuya componente de índices $(j,k)$ es
$\mathscr{T}\!or_n^S
(\mathfrak U^{(1)},\mathfrak U^{(2)};
\mathscr P_j^{(1)},\mathscr P_k^{(2)})$.

Como $L_{\bullet\bullet}^{(i)}$ es un funtor exacto respecto de
$\mathscr P_\bullet^{(i)}$, puesto que las intersecciones de los miembros
\oldpage[III]{154}
de $\mathfrak U^{(i)}$ son afines (\textbf I, 5.5.6 y 1.3.11),
$\mathbf{Tor}_\bullet^S
(\mathfrak U^{(1)},\mathfrak U^{(2)};
 \mathscr P_\bullet^{(1)},\mathscr P_\bullet^{(2)})$ es un
\emph{bi-$\partial$-funtor covariante} en $\mathscr P_\bullet^{(1)}$ y
$\mathscr P_\bullet^{(2)}$, con valores en la categoría de los $A$-módulos
(\textbf{0}, 11.7.3). Además, se sabe (\textbf{0}, 11.7.2) que este
bifuntor es el término límite común de seis funtores espectrales \emph{birregulares},
que denotaremos por
${}^{(t)}E^S
(\mathfrak U^{(1)},\mathfrak U^{(2)};
\mathscr P_\bullet^{(1)},\mathscr P_\bullet^{(2)})$
o por
${}^{(t)}E
(\mathfrak U^{(1)},\mathfrak U^{(2)};
\mathscr P_\bullet^{(1)},\mathscr P_\bullet^{(2)})$,
donde $t$ ha de sustituirse por una de las letras $a$, $b$, $a'$, $b'$, $c$,
$d$, y cuyos términos $E_2$ son los siguientes:
\[
\begin{aligned}
  {}^{(a)}E_{pq}^2
  &=\bigoplus_{q_1+q_2=q}
    \operatorname{Tor}_p^A
    \bigl(\mathrm H_{q_1}(L_{\bullet\bullet}^{(1)}),
          \mathrm H_{q_2}(L_{\bullet\bullet}^{(2)})\bigr)\\
  {}^{(b)}E_{pq}^2
  &=\mathrm H_p\bigl(
    \mathbf{Tor}_q^{A,\mathrm{II}}
    (L_{\bullet\bullet}^{(1)},L_{\bullet\bullet}^{(2)})\bigr)\\
  {}^{(a')}E_{pq}^2
  &=\bigoplus_{q_1+q_2=q}
    \mathbf{Tor}_p^A
    \bigl(\mathrm H_{q_1}^{\mathrm I}(L_{\bullet\bullet}^{(1)}),
          \mathrm H_{q_2}^{\mathrm I}(L_{\bullet\bullet}^{(2)})\bigr)\\
  {}^{(b')}E_{pq}^2
  &=\mathrm H_p\bigl(
    \operatorname{Tor}_q^A
    (L_{\bullet\bullet}^{(1)},L_{\bullet\bullet}^{(2)})\bigr)\\
  {}^{(c)}E_{pq}^2
  &=\bigoplus_{q_1+q_2=q}
    \mathbf{Tor}_p^A
    \bigl(\mathrm H_{q_1}^{\mathrm{II}}(L_{\bullet\bullet}^{(1)}),
          \mathrm H_{q_2}^{\mathrm{II}}(L_{\bullet\bullet}^{(2)})\bigr)\\
  {}^{(d)}E_{pq}^2
  &=\mathrm H_p\bigl(
    \mathbf{Tor}_q^{A,\mathrm I}
    (L_{\bullet\bullet}^{(1)},L_{\bullet\bullet}^{(2)})\bigr),
\end{aligned}
\]
donde las notaciones concuerdan con las de la teoría general de la
hiperhomología. A continuación explicitaremos con mayor detalle estos términos
iniciales.
\end{env}

\begin{env}[6.6.3]
\label{III.6.6.3}
\emph{Sucesiones espectrales $(a)$ y $(a')$.} Hemos visto en (6.6.1) que el
módulo de homología $\mathrm H_n(L_{\bullet\bullet}^{(i)})$ del bicomplejo
$L_{\bullet\bullet}^{(i)}$ era igual a
$\mathbf H^{-n}(X^{(i)},\mathscr P_\bullet^{(i)})$; por tanto,
\[
  {}^{(a)}E_{pq}^2
  =\bigoplus_{q_1+q_2=q}
   \operatorname{Tor}_p^A
   \bigl(\mathbf H^{-q_1}(X^{(1)},\mathscr P_\bullet^{(1)}),
         \mathbf H^{-q_2}(X^{(2)},\mathscr P_\bullet^{(2)})\bigr).
\]
Por definición, el complejo
$\mathrm H_n^{\mathrm I}(L_{\bullet\bullet}^{(i)})$ tiene por término de grado
$k$ el módulo de homología
$\mathrm H_n(C^\bullet(\mathfrak U^{(i)},\mathscr P_k^{(i)}))$, es decir,
por definición, el módulo de \emph{cohomología}
$\mathrm H^{-n}(\mathfrak U^{(i)},\mathscr P_k^{(i)})$; se sabe por (1.4.1)
que este módulo es canónicamente isomorfo a
$\mathrm H^{-n}(X^{(i)},\mathscr P_k^{(i)})$; por tanto,
\[
  {}^{(a')}E_{pq}^2
  =\bigoplus_{q_1+q_2=q}
   \mathbf{Tor}_p^A
   \bigl(\mathrm H^{-q_1}(X^{(1)},\mathscr P_\bullet^{(1)}),
         \mathrm H^{-q_2}(X^{(2)},\mathscr P_\bullet^{(2)})\bigr).
\]
\end{env}

\begin{env}[6.6.4]
\label{III.6.6.4}
\emph{Sucesiones espectrales $(b)$ y $(b')$.} Por definición,
$\mathbf{Tor}_q^{A,\mathrm{II}}
(L_{\bullet\bullet}^{(1)},L_{\bullet\bullet}^{(2)})$ es un bicomplejo cuya
componente de bigrado $(h,k)$ es el $A$-módulo
\[
  \mathbf{Tor}_q^A
  \bigl(C^{-h}(\mathfrak U^{(1)},\mathscr P_\bullet^{(1)}),
        C^{-k}(\mathfrak U^{(2)},\mathscr P_\bullet^{(2)})\bigr).
\]
Sea $\Phi^{(i)}$ el conjunto de índices de $\mathfrak U^{(i)}$; por
definición, el complejo de módulos
$C^r(\mathfrak U^{(i)},\mathscr P_\bullet^{(i)})$ ($r\geqslant0$) es suma
directa de los complejos
$\Gamma(U_\rho^{(i)},\mathscr P_\bullet^{(i)})$, donde $U_\rho^{(i)}$ es la
intersección de los $U_\xi^{(i)}$ para $\xi\in\rho$, y $\rho$ recorre
$\mathfrak P(\Phi^{(i)})$; por tanto, el $A$-módulo
\[
  \mathbf{Tor}_q^A
  \bigl(C^{-h}(\mathfrak U^{(1)},\mathscr P_\bullet^{(1)}),
        C^{-k}(\mathfrak U^{(2)},\mathscr P_\bullet^{(2)})\bigr)
\]
es suma directa de los $A$-módulos
$\mathbf{Tor}_q^A
\bigl(\Gamma(U_\sigma^{(1)},\mathscr P_\bullet^{(1)}),
\Gamma(U_\tau^{(2)},\mathscr P_\bullet^{(2)})\bigr)$,
donde $\sigma$ (resp. $\tau$) recorre los elementos de
$\mathfrak P(\Phi^{(1)})$ (resp. $\mathfrak P(\Phi^{(2)})$) tales
que $\operatorname{Card}(\sigma)=1-h$ (resp.
$\operatorname{Card}(\tau)=1-k$). Como $X^{(1)}$ y $X^{(2)}$ son
esquemas, los $U_\rho^{(i)}$ son afines y, por tanto, se tiene (6.4.1.1)
\[
  \mathbf{Tor}_q^A
  \bigl(\Gamma(U_\sigma^{(1)},\mathscr P_\bullet^{(1)}),
        \Gamma(U_\tau^{(2)},\mathscr P_\bullet^{(2)})\bigr)
  =\Gamma\bigl(U_\sigma^{(1)}\times_S U_\tau^{(2)},
    \mathscr{T}\!or_q^S
    (\mathscr P_\bullet^{(1)},\mathscr P_\bullet^{(2)})\bigr).
\]
\oldpage[III]{155}
Se ve, pues, que ${}^{(b)}E_{pq}^2$ es el $(-p)$-ésimo módulo de
cohomología del complejo
$L^\bullet(\Phi^{(1)},\Phi^{(2)};\mathscr S)$
de cocadenas \emph{bialternadas} sobre $\Phi^{(1)}$ y $\Phi^{(2)}$ con valores en el
\emph{sistema de coeficientes}
\[
  \mathscr S:(\sigma,\tau)\leadsto
  \Gamma\bigl(U_\sigma^{(1)}\times_S U_\tau^{(2)},
    \mathscr{T}\!or_q^S
    (\mathscr P_\bullet^{(1)},\mathscr P_\bullet^{(2)})\bigr)
\]
(\textbf{0}, 11.8.4). Se sabe entonces (\textbf{0}, 11.8.5 y
11.8.6) que la cohomología de este complejo es la misma que la del complejo
$C^\bullet(\Phi^{(1)},\Phi^{(2)};\mathscr S)$ de \emph{todas} las cocadenas sobre
$\Phi^{(1)}$ y $\Phi^{(2)}$ con valores en $\mathscr S$, y también la del
complejo $P^\bullet(\Phi^{(1)},\Phi^{(2)};\mathscr S)$, cuyos elementos son
las combinaciones lineales de los
\[
  \lambda(\sigma,\tau)\in
  \Gamma\bigl(U_\sigma^{(1)}\times_S U_\tau^{(2)},
    \mathscr{T}\!or_q^S
    (\mathscr P_\bullet^{(1)},\mathscr P_\bullet^{(2)})\bigr)
\]
donde $\sigma=(\alpha_0,\ldots,\alpha_h)$ y
$\tau=(\beta_0,\ldots,\beta_h)$ son sucesiones con el \emph{mismo número de
elementos}. Pero entonces se tiene
$U_\sigma^{(1)}\times_S U_\tau^{(2)}
=(U_{\alpha_0}^{(1)}\times_S U_{\beta_0}^{(2)})\cap\ldots\cap
(U_{\alpha_h}^{(1)}\times_S U_{\beta_h}^{(2)})$
Si se denota por $\mathfrak U$ el recubrimiento de
$Z=X^{(1)}\times_S X^{(2)}$ por los abiertos afines
$U_\alpha^{(1)}\times_S U_\beta^{(2)}$, se ve finalmente, teniendo en cuenta
que $X^{(1)}\times_S X^{(2)}$ es un \emph{esquema}, que se tiene, en virtud de
(1.3.1),
\[
  {}^{(b)}E_{pq}^2
  =\mathbf H^{-p}\bigl(X^{(1)}\times_S X^{(2)},
    \mathscr{T}\!or_q^S
    (\mathscr P_\bullet^{(1)},\mathscr P_\bullet^{(2)})\bigr).
\]

En segundo lugar,
$\operatorname{Tor}_q^A(L_{\bullet\bullet}^{(1)},L_{\bullet\bullet}^{(2)})$ es un
bicomplejo cuya componente de bigrado $(h,k)$ es la suma directa de los
$A$-módulos
\[
  \operatorname{Tor}_q^A
  \bigl(C^{-h_1}(\mathfrak U^{(1)},\mathscr P_{k_1}^{(1)}),
        C^{-h_2}(\mathfrak U^{(2)},\mathscr P_{k_2}^{(2)})\bigr)
\]
tales que $h_1+h_2=h$ y $k_1+k_2=k$; explicitando los módulos
$C^r(\mathfrak U^{(i)},\mathscr P_s^{(i)})$ como antes, se ve nuevamente que
esta componente es suma directa de los $A$-módulos
\[
  \Gamma\bigl(U_\sigma^{(1)}\times_S U_\tau^{(2)},
    \mathscr{T}\!or_q^S
    (\mathscr P_{k_1}^{(1)},\mathscr P_{k_2}^{(2)})\bigr)
\]
donde $k_1+k_2=k$, y $\sigma$ (resp. $\tau$) recorre los elementos
de $\mathfrak P(\Phi^{(1)})$ (resp.
$\mathfrak P(\Phi^{(2)})$) tales que
$\operatorname{Card}(\sigma)+\operatorname{Card}(\tau)=2-h$. El término
${}^{(b')}E_{pq}^2$ que estamos calculando es el $(-p)$-ésimo módulo
de cohomología de un bicomplejo $N^{\bullet\bullet}=(N^{hk})$, donde el
complejo simple $N^{\bullet k}$ es el complejo de cocadenas \emph{bialternadas} sobre
$\Phi^{(1)}$ y $\Phi^{(2)}$, con valores en el sistema de coeficientes
\[
  \mathscr S_k:(\sigma,\tau)\leadsto
  \Gamma\left(U_\sigma^{(1)}\times_S U_\tau^{(2)},
    \bigoplus_{k_1+k_2=k}
    \mathscr{T}\!or_q^S
    (\mathscr P_{-k_1}^{(1)},\mathscr P_{-k_2}^{(2)})\right)
\]
y estos sistemas de coeficientes forman un \emph{complejo} $\mathscr S^\bullet$ cuyo
diferencial proviene del diferencial del complejo simple asociado al bicomplejo
$\mathscr{T}\!or_q^S
(\mathscr P^{(1)\bullet},\mathscr P^{(2)\bullet})$. Se sabe que la
cohomología de $N^{\bullet\bullet}$ es la misma que la del bicomplejo
$C^\bullet(\Phi^{(1)},\Phi^{(2)};\mathscr S^\bullet)$
(\textbf{0}, 11.8.9), y también la misma que la del bicomplejo
$P^\bullet(\Phi^{(1)},\Phi^{(2)};\mathscr S^\bullet)$, donde los elementos de
bigrado $(h,k)$ son las combinaciones lineales de los
\[
  \lambda(\sigma,\tau)\in
  \Gamma\left(U_\sigma^{(1)}\times_S U_\tau^{(2)},
    \bigoplus_{k_1+k_2=k}
    \mathscr{T}\!or_q^S
    (\mathscr P_{-k_1}^{(1)},\mathscr P_{-k_2}^{(2)})\right)
\]
donde $\sigma=(\alpha_0,\ldots,\alpha_h)$ y
$\tau=(\beta_0,\ldots,\beta_h)$ son sucesiones con el \emph{mismo número de
elementos} (\textbf{0}, 11.8.10). Se ve entonces, como antes, que
${}^{(b')}E_{pq}^2$ es el $(-p)$-ésimo módulo de cohomología del
bicomplejo $C^\bullet(\mathfrak U,\mathscr Q^\bullet)$, donde
$\mathscr Q^\bullet$ es el complejo simple asociado al bicomplejo
$\mathscr{T}\!or_q^S
(\mathscr P^{(1)\bullet},\mathscr P^{(2)\bullet})$ de
$\mathscr O_Z$-módulos. Con las convenciones adoptadas en (6.6.1), se tiene,
por tanto,
\[
  {}^{(b')}E_{pq}^2
  =\mathbf H^{-p}\bigl(X^{(1)}\times_S X^{(2)},
    \mathscr{T}\!or_q^S
    (\mathscr P_\bullet^{(1)},\mathscr P_\bullet^{(2)})\bigr).
\]
\end{env}

\oldpage[III]{156}
\begin{env}[6.6.5]
\label{III.6.6.5}
\emph{Sucesiones espectrales $(c)$ y $(d)$.} Por definición, el complejo
$\mathrm H_n^{\mathrm{II}}(L_{\bullet\bullet}^{(i)})$ tiene por término de
grado $h$ el $A$-módulo
$C^{-h}(\mathfrak U^{(i)},\mathscr H_n(\mathscr P_\bullet^{(i)}))$, en virtud
de la exactitud del funtor $C^{-h}$. Por tanto, por definición del hipertor de
dos módulos relativo a dos recubrimientos (6.6.2),
\[
  {}^{(c)}E_{pq}^2
  =\bigoplus_{q_1+q_2=q}
   \operatorname{Tor}_p^S
   \bigl(\mathfrak U^{(1)},\mathfrak U^{(2)};
         \mathscr H_{q_1}(\mathscr P_\bullet^{(1)}),
         \mathscr H_{q_2}(\mathscr P_\bullet^{(2)})\bigr).
\]
Por último, por definición,
$\mathbf{Tor}_q^{A,\mathrm I}
(L_{\bullet\bullet}^{(1)},L_{\bullet\bullet}^{(2)})$ es un bicomplejo cuya
componente de bigrado $(h,k)$ es el $A$-módulo
$\operatorname{Tor}_q^S
(\mathfrak U^{(1)},\mathfrak U^{(2)};
 \mathscr P_h^{(1)},\mathscr P_k^{(2)})$. Por tanto,
\[
  {}^{(d)}E_{pq}^2
  =\mathrm H_p\bigl(
    \operatorname{Tor}_q^S
    (\mathfrak U^{(1)},\mathfrak U^{(2)};
     \mathscr P_\bullet^{(1)},\mathscr P_\bullet^{(2)})\bigr).
\]
\end{env}

\begin{env}[6.6.6]
\label{III.6.6.6}
La teoría de la hiperhomología de los funtores de bicomplejos
(\textbf{0}, 11.7.3) muestra, como en (6.3.4), que, para toda resolución
\emph{plana} de Cartan--Eilenberg $M_{\bullet\bullet\bullet}^{(i)}$ de
$L_{\bullet\bullet}^{(i)}$ (en la categoría de los complejos de módulos
acotados inferiormente) ($i=1,2$), se tienen isomorfismos canónicos de
bi-$\partial$-funtores

\label{III.6.6.6.1}
\[
\begin{aligned}
  \mathbf{Tor}_\bullet^S
  (\mathfrak U^{(1)},\mathfrak U^{(2)};
   \mathscr P_\bullet^{(1)},\mathscr P_\bullet^{(2)})
  &\xrightarrow{\sim}
   \mathrm H_\bullet
   (M_{\bullet\bullet\bullet}^{(1)}\otimes_A M_{\bullet\bullet\bullet}^{(2)})\\
  &\xrightarrow{\sim}
   \mathrm H_\bullet
   (M_{\bullet\bullet\bullet}^{(1)}\otimes_A L_{\bullet\bullet}^{(2)})
   \xrightarrow{\sim}
   \mathrm H_\bullet
   (L_{\bullet\bullet}^{(1)}\otimes_A M_{\bullet\bullet\bullet}^{(2)}).
\end{aligned}
\tag{6.6.6.1}\label{III.6.6.6.1.fr}
\]

\end{env}

\begin{env}[6.6.7]
\label{III.6.6.7}
Vamos a mostrar ahora que el hipertor global definido en (6.6.2), así como
las seis sucesiones espectrales correspondientes, no dependen de los
recubrimientos abiertos afines finitos $\mathfrak U^{(i)}$ que se utilizaron
para definirlos (salvo isomorfismos canónicos). Bastará para ello mostrar que,
si $\mathfrak V^{(i)}$ son otros dos recubrimientos de la misma naturaleza,
tales que $\mathfrak V^{(i)}$ sea \emph{más fino} que $\mathfrak U^{(i)}$
para $i=1,2$, entonces se tienen \emph{isomorfismos canónicos} de funtores espectrales
\[
  {}^{(t)}E
  (\mathfrak U^{(1)},\mathfrak U^{(2)};
   \mathscr P_\bullet^{(1)},\mathscr P_\bullet^{(2)})
  \xrightarrow{\sim}
  {}^{(t)}E
  (\mathfrak V^{(1)},\mathfrak V^{(2)};
   \mathscr P_\bullet^{(1)},\mathscr P_\bullet^{(2)})
  \tag{6.6.7, $t$)}\label{III.6.6.7.t.fr}
\]

donde $t$ se sustituye por $a$, $b$, $a'$, $b'$, $c$ o $d$.

Ahora bien, para $i=1,2$ se tienen homomorfismos de bicomplejos
\[
  C^\bullet(\mathfrak U^{(i)},\mathscr P_\bullet^{(i)})
  \longrightarrow
  C^\bullet(\mathfrak V^{(i)},\mathscr P_\bullet^{(i)})
\]
bien definidos salvo homotopía (G, II, 5.7.1); de ello resultan ya
homomorfismos (6.6.7, $t$)) definidos canónicamente y compatibles con los
operadores de borde en los términos límite (\textbf{0}, 11.3.2). Además,
el cálculo de los términos $E_2$ de las sucesiones espectrales $(a)$, $(b)$,
$(a')$, $(b')$ muestra que, para estas sucesiones espectrales, el homomorfismo
(6.6.7, $t$)) es un isomorfismo sobre los términos $E_2$; como estas sucesiones
espectrales son birregulares, se ve que (6.6.7, $t$)) es un isomorfismo para
estos funtores espectrales y, por consiguiente, induce un \emph{isomorfismo de
bi-$\partial$-funtores} en su término límite común
(\textbf{0}, 11.1.5).

En particular, para $\mathscr O_{X^{(i)}}$-módulos cuasicoherentes
$\mathscr F^{(i)}$ ($i=1,2$), el homomorfismo canónico
\[
  \operatorname{Tor}_\bullet^S
  (\mathfrak U^{(1)},\mathfrak U^{(2)};
   \mathscr F^{(1)},\mathscr F^{(2)})
  \longrightarrow
  \operatorname{Tor}_\bullet^S
  (\mathfrak V^{(1)},\mathfrak V^{(2)};
   \mathscr F^{(1)},\mathscr F^{(2)})
\]
es biyectivo; visto el cálculo de (6.6.5), se ve que (6.6.7, $t$)) es también
un isomorfismo de los términos $E_2$ para $t=c$ y $t=d$. Se concluye como
antes que (6.6.7, $t$)) es también un isomorfismo de sucesiones espectrales
para $t=c$ y $t=d$.
\oldpage[III]{157}
Se puede considerar que los isomorfismos (6.6.7, $t$)) definen sistemas
inductivos de funtores espectrales sobre el conjunto filtrante de pares
$(\mathfrak U^{(1)},\mathfrak U^{(2)})$ de recubrimientos abiertos afines
finitos de $X^{(1)}$ y $X^{(2)}$. Denotaremos por
\[
  {}^{(t)}E
  (X^{(1)},X^{(2)};
   \mathscr P_\bullet^{(1)},\mathscr P_\bullet^{(2)})
  \quad\text{o}\quad
  {}^{(t)}E^S
  (X^{(1)},X^{(2)};
   \mathscr P_\bullet^{(1)},\mathscr P_\bullet^{(2)})
\]
el límite inductivo de este sistema, y por
$\mathbf{Tor}_\bullet^S
(X^{(1)},X^{(2)};\mathscr P_\bullet^{(1)},\mathscr P_\bullet^{(2)})$
el término límite de este funtor espectral, al que llamaremos el
\emph{hipertor global} de los dos complejos $\mathscr P_\bullet^{(1)}$ y
$\mathscr P_\bullet^{(2)}$; si $\mathscr P_\bullet^{(1)}$ y
$\mathscr P_\bullet^{(2)}$ están reducidos a sus términos de grado 0,
$\mathscr F^{(1)}$ y $\mathscr F^{(2)}$, escribiremos
\[
  \operatorname{Tor}_\bullet^S
  (X^{(1)},X^{(2)};\mathscr F^{(1)},\mathscr F^{(2)}),
\]
y, conforme a las convenciones generales,
$
  \operatorname{Tor}_q^S
  (X^{(1)},X^{(2)};\mathscr P_\bullet^{(1)},\mathscr P_\bullet^{(2)})
$
será, por tanto, el bicomplejo cuyas componentes son los módulos
$
  \operatorname{Tor}_q^S
  (X^{(1)},X^{(2)};\mathscr P_h^{(1)},\mathscr P_k^{(2)})$.
\end{env}

\begin{env}[6.6.8]
\label{III.6.6.8}
Bajo las hipótesis de (6.6.1), consideremos ahora dos $S$-morfismos
$f_i:X^{(i)}\to Y^{(i)}$, donde
$Y^{(i)}=\operatorname{Spec}(B_i)$ es un $S$-\emph{esquema afín}, siendo
$B_i$ una $A$-álgebra ($i=1,2$); esto define, pues, un $A$-homomorfismo
$B_i\to\Gamma(X^{(i)},\mathscr O_{X^{(i)}})$ (\textbf I, 2.2.4) y, por consiguiente,
cada uno de los $L_{\bullet\bullet}^{(i)}$ definidos en (6.6.2) es un
bicomplejo de $B_i$-módulos; se concluye que
$L_{\bullet\bullet}^{(1)}\otimes_A L_{\bullet\bullet}^{(2)}$ es un
complejo cuádruple de $(B_1\otimes_A B_2)$-módulos, y sus seis funtores
espectrales de hiperhomología pueden considerarse, por tanto, con valores en
la categoría de las sucesiones espectrales de
$(B_1\otimes_A B_2)$-\emph{módulos}. Si se pone
$Y=Y^{(1)}\times_S Y^{(2)}=\operatorname{Spec}(B_1\otimes_A B_2)$, se pueden
considerar los $\mathscr O_Y$-módulos cuasicoherentes asociados a estos
módulos (\textbf I, 1.3.4); denotaremos por
${}^{(t)}\mathscr E
(f_1,f_2;\mathscr P_\bullet^{(1)},\mathscr P_\bullet^{(2)})$
(para $t=a$, $b$, $a'$, $b'$, $c$ o $d$) las seis sucesiones espectrales
$
  \bigl({}^{(t)}E
  (X^{(1)},X^{(2)};
   \mathscr P_\bullet^{(1)},\mathscr P_\bullet^{(2)})\bigr)^{\sim}$
de $\mathscr O_Y$-módulos, y por
$\boldsymbol{\mathscr{T}\!or}_\bullet^S
(f_1,f_2;\mathscr P_\bullet^{(1)},\mathscr P_\bullet^{(2)})$ su término
límite común
$
  \bigl(\mathbf{Tor}_\bullet^S
  (X^{(1)},X^{(2)};
   \mathscr P_\bullet^{(1)},\mathscr P_\bullet^{(2)})\bigr)^{\sim}$. Se denotará por $\mathscr{T}\!or_\bullet^S(f_1,f_2;\mathscr F^{(1)},\mathscr F^{(2)})$
cuando $\mathscr P_\bullet^{(i)}$ esté reducido a su término de grado 0,
$\mathscr F^{(i)}$ ($i=1,2$).
\end{env}

\subsection{Funtores de hipertor global de complejos de módulos
cuasicoherentes y sucesiones espectrales de Künneth: caso general}
\label{subsection:III.6.7}

\begin{env}[6.7.1]
\label{III.6.7.1}
Vamos a generalizar ahora las definiciones de (6.6.8) al caso en que $S$ es
un preesquema cualquiera, $X^{(i)}$ e $Y^{(i)}$ son $S$-preesquemas y
$f_i:X^{(i)}\to Y^{(i)}$ son morfismos \emph{separados} y \emph{cuasicompactos}. Se
trata entonces, para todo par de complejos $\mathscr P_\bullet^{(i)}$ de
$\mathscr O_{X^{(i)}}$-módulos cuasicoherentes, acotados inferiormente
($i=1,2$), de definir para todo $n$ un $\mathscr O_Y$-módulo cuasicoherente
$\boldsymbol{\mathscr{T}\!or}_n^S
(f_1,f_2;\mathscr P_\bullet^{(1)},\mathscr P_\bullet^{(2)})$, así como seis
funtores espectrales, que se reduzcan a las definiciones de (6.6.8) cuando
$S$, $Y^{(1)}$ e $Y^{(2)}$ son afines (se pone
$Y=Y^{(1)}\times_S Y^{(2)}$).

Supongamos primero que $S=\operatorname{Spec}(A)$ es \emph{afín}, pero que
$Y^{(1)}$ e $Y^{(2)}$ son cualesquiera; sea $W^{(i)}$ un abierto afín de
$Y^{(i)}$; entonces $f_i^{-1}(W^{(i)})$ es un $S$-esquema cuasicompacto y
$W=W^{(1)}\times_S W^{(2)}$ es un abierto afín de $Y$; sea
$f'_i:f_i^{-1}(W^{(i)})\to W^{(i)}$ la restricción de $f_i$, y sea
$\mathscr P'_\bullet{}^{(i)}$ la restricción
$\mathscr P_\bullet^{(i)}|f_i^{-1}(W^{(i)})$ ($i=1,2$). Se tienen entonces,
por (6.6.8), las sucesiones espectrales
${}^{(t)}\mathscr E
(f'_1,f'_2;\mathscr P'_\bullet{}^{(1)},\mathscr P'_\bullet{}^{(2)})$
de $(\mathscr O_Y|W)$-módulos cuasicoherentes, y hay que verificar que
satisfacen las condiciones de pegado ($\mathbf{0}_{\mathrm I}$, 3.3.1).
Se reduce inmediatamente al caso en que
$Y^{(i)}=\operatorname{Spec}(B_i)$ es afín y $W^{(i)}=D(g_i)$, donde
$g_i\in B_i$, de modo que
$W=D(g_1\otimes g_2)$
\oldpage[III]{158}
en $Y=\operatorname{Spec}(B_1\otimes_A B_2)$ (\textbf{II}, 4.3.2.1); si
$X'^{(i)}=f_i^{-1}(W^{(i)})$, se trata de establecer un isomorfismo canónico
de funtores espectrales

\label{III.6.7.1.1}
\[
  {}^{(t)}E
  (X'^{(1)},X'^{(2)};
   \mathscr P'_\bullet{}^{(1)},\mathscr P'_\bullet{}^{(2)})
  \xrightarrow{\sim}
  {}^{(t)}E
  (X^{(1)},X^{(2)};
   \mathscr P_\bullet^{(1)},\mathscr P_\bullet^{(2)})
  \otimes_B B_g
  \tag{6.7.1.1}\label{III.6.7.1.1.fr}
\]


donde se ha puesto $B=B_1\otimes_A B_2$ y $g=g_1\otimes g_2$. Para ello,
partamos de recubrimientos abiertos afines finitos $\mathfrak U^{(i)}$ de
$X^{(i)}$ ($i=1,2$), y sea $\mathfrak U'^{(i)}$ la traza de
$\mathfrak U^{(i)}$ sobre $X'^{(i)}$, que sigue estando formada por abiertos
afines (\textbf I, 5.5.10); más precisamente, se tiene
\[
  C^\bullet(\mathfrak U'^{(i)},\mathscr P'_\bullet{}^{(i)})
  =C^\bullet(\mathfrak U^{(i)},\mathscr P_\bullet^{(i)})
   \otimes_{B_i}(B_i)_{g_i}.
\]
Si se pone
$L'_{\bullet\bullet}{}^{(i)}
=C^\bullet(\mathfrak U'^{(i)},\mathscr P'_\bullet{}^{(i)})$, se tiene, por
tanto,
$
  L'_{\bullet\bullet}{}^{(1)}\otimes_A L'_{\bullet\bullet}{}^{(2)}
  =\bigl(L_{\bullet\bullet}^{(1)}\otimes_{B_1}(B_1)_{g_1}\bigr)
   \otimes_A
   \bigl(L_{\bullet\bullet}^{(2)}\otimes_{B_2}(B_2)_{g_2}\bigr)$;
como $B_g=(B_1)_{g_1}\otimes_A(B_2)_{g_2}$ salvo isomorfismo canónico, se
tiene, también salvo isomorfismo canónico,
$
  L'_{\bullet\bullet}{}^{(1)}\otimes_A L'_{\bullet\bullet}{}^{(2)}
  =(L_{\bullet\bullet}^{(1)}\otimes_A L_{\bullet\bullet}^{(2)})
    \otimes_B B_g$.
Si $M_{\bullet\bullet\bullet}^{(i)}$ es una resolución proyectiva de
Cartan--Eilenberg de $L_{\bullet\bullet}^{(i)}$, que puede suponerse formada
por $B_i$-módulos, del hecho de que $(B_i)_{g_i}$ es \emph{plano} sobre $B_i$
se deduce que
$M'_{\bullet\bullet\bullet}{}^{(i)}
=M_{\bullet\bullet\bullet}^{(i)}\otimes_{B_i}(B_i)_{g_i}$ es una resolución
proyectiva de Cartan--Eilenberg del bicomplejo
$L'_{\bullet\bullet}{}^{(i)}$; además, se tiene
\[
  M'_{\bullet\bullet\bullet}{}^{(1)}\otimes_A M'_{\bullet\bullet\bullet}{}^{(2)}
  =(M_{\bullet\bullet\bullet}^{(1)}\otimes_A M_{\bullet\bullet\bullet}^{(2)})
    \otimes_B B_g.
\]
El isomorfismo buscado (6.7.1.1) se deduce entonces inmediatamente de las
definiciones de la hiperhomología de un bicomplejo y de la exactitud del
funtor $G\otimes_B B_g$ respecto del $B$-módulo $G$.
\end{env}

\begin{env}[6.7.2]
\label{III.6.7.2}
Supongamos ahora que $S$ es cualquiera, y sean
$u_i:Y^{(i)}\to S$ los morfismos estructurales ($i=1,2$). Sea $(S_\alpha)$ un
recubrimiento abierto afín de $S$, pongamos
$Y_\alpha^{(i)}=u_i^{-1}(S_\alpha)$,
$X_\alpha^{(i)}=f_i^{-1}(Y_\alpha^{(i)})$, y sea
$f_{i\alpha}:X_\alpha^{(i)}\to Y_\alpha^{(i)}$ la restricción de $f_i$, que
es un morfismo separado y cuasicompacto. Los
$Y_\alpha=Y_\alpha^{(1)}\times_{S_\alpha}Y_\alpha^{(2)}$ forman un
recubrimiento abierto de $Y$, y sobre cada $Y_\alpha$ quedan definidos, por
(6.7.1), los funtores espectrales
\[
  {}^{(t)}\mathscr E^{S_\alpha}
  \bigl(f_{1\alpha},f_{2\alpha};
    \mathscr P_\bullet^{(1)}|X_\alpha^{(1)},
    \mathscr P_\bullet^{(2)}|X_\alpha^{(2)}\bigr);
\]
hay que demostrar de nuevo que estos funtores verifican las condiciones de
pegado.

Se reduce inmediatamente a la situación siguiente:
$S=\operatorname{Spec}(A)$ es afín, $S'=D(h)$, donde $h\in A$, y
$u_i(Y^{(i)})\subset S'$; puede suponerse además que
$Y^{(i)}=\operatorname{Spec}(B_i)$ es afín; se trata de definir isomorfismos
canónicos

\label{III.6.7.2.1}
\[
  {}^{(t)}E^S
  (X^{(1)},X^{(2)};
   \mathscr P_\bullet^{(1)},\mathscr P_\bullet^{(2)})
  \xrightarrow{\sim}
  {}^{(t)}E^{S'}
  (X^{(1)},X^{(2)};
   \mathscr P_\bullet^{(1)},\mathscr P_\bullet^{(2)}).
  \tag{6.7.2.1}\label{III.6.7.2.1.fr}
\]

Ahora bien, con las notaciones de (6.6.2), los
$L_{\bullet\bullet}^{(i)}$ están formados por $A_h$-módulos, y se tiene
$L_{\bullet\bullet}^{(1)}\otimes_{A_h}L_{\bullet\bullet}^{(2)}
=L_{\bullet\bullet}^{(1)}\otimes_A L_{\bullet\bullet}^{(2)}$ salvo
isomorfismo canónico; como puede tomarse una resolución proyectiva de
Cartan--Eilenberg $M_{\bullet\bullet\bullet}^{(i)}$ de
$L_{\bullet\bullet}^{(i)}$ formada por $A_h$-módulos, esto proporciona
inmediatamente el isomorfismo canónico buscado.
\end{env}

Hemos demostrado, en resumen, el siguiente

\begin{theorem}[6.7.3]
\label{III.6.7.3}
Sean $S$ un preesquema, $f_i:X_i\to Y_i$ un $S$-morfismo separado y
cuasicompacto de $S$-preesquemas y $\mathscr P_\bullet^{(i)}$ un complejo
acotado inferiormente de $\mathscr O_{X^{(i)}}$-módulos cuasicoherentes
($i=1,2$); se pone $Y=Y^{(1)}\times_S Y^{(2)}$. Existe un
bi-$\partial$-funtor
  $\boldsymbol{\mathscr{T}\!or}_\bullet^S
(f_1,f_2;\mathscr P_\bullet^{(1)},\mathscr P_\bullet^{(2)})$
con valores en la categoría de los $\mathscr O_Y$-módulos cuasicoherentes,
tal que, si $V^{(i)}$ es un abierto afín de $Y^{(i)}$ ($i=1,2$) y
$V=V^{(1)}\times_S V^{(2)}$, se tiene
\[
  \boldsymbol{\mathscr{T}\!or}_\bullet^S
  (f_1,f_2;\mathscr P_\bullet^{(1)},\mathscr P_\bullet^{(2)})|V
  =\Bigl(\mathbf{Tor}_\bullet^S
    \bigl(f_1^{-1}(V^{(1)}),f_2^{-1}(V^{(2)});
      \mathscr P_\bullet^{(1)}|f_1^{-1}(V^{(1)}),
      \mathscr P_\bullet^{(2)}|f_2^{-1}(V^{(2)})\bigr)\Bigr)^{\sim}.
\]
\oldpage[III]{159}
Este bifuntor es el término límite de seis funtores espectrales
birregulares
\[
  {}^{(t)}\mathscr E
  (f_1,f_2;\mathscr P_\bullet^{(1)},\mathscr P_\bullet^{(2)})
  \qquad(t=a,b,a',b',c,d)
\]
cuyos términos $E_2$ vienen dados por
\[
\begin{aligned}
  {}^{(a)}\mathscr E_{pq}^2
  &=\bigoplus_{q_1+q_2=q}
    \mathscr{T}\!or_p^S
    \bigl(\mathcal H^{-q_1}(f_1,\mathscr P_\bullet^{(1)}),
          \mathcal H^{-q_2}(f_2,\mathscr P_\bullet^{(2)})\bigr),\\
  {}^{(b)}\mathscr E_{pq}^2
  &=\mathscr H^{-p}\bigl(f_1\times_S f_2,
    \boldsymbol{\mathscr{T}\!or}_q^S
    (\mathscr P_\bullet^{(1)},\mathscr P_\bullet^{(2)})\bigr),\\
  {}^{(a')}\mathscr E_{pq}^2
  &=\bigoplus_{q_1+q_2=q}
    \boldsymbol{\mathscr{T}\!or}_p^S
    \bigl(\mathscr H^{-q_1}(f_1,\mathscr P_\bullet^{(1)}),
          \mathscr H^{-q_2}(f_2,\mathscr P_\bullet^{(2)})\bigr),\\
  {}^{(b')}\mathscr E_{pq}^2
  &=\mathcal H^{-p}\bigl(f_1\times_S f_2,
    \mathscr{T}\!or_q^S
    (\mathscr P_\bullet^{(1)},\mathscr P_\bullet^{(2)})\bigr),\\
  {}^{(c)}\mathscr E_{pq}^2
  &=\bigoplus_{q_1+q_2=q}
    \mathscr{T}\!or_p^S
    \bigl(f_1,f_2;
      \mathscr H_{q_1}(\mathscr P_\bullet^{(1)}),
      \mathscr H_{q_2}(\mathscr P_\bullet^{(2)})\bigr),\\
  {}^{(d)}\mathscr E_{pq}^2
  &=\mathscr H_p\bigl(\mathscr{T}\!or_q^S
    (f_1,f_2;\mathscr P_\bullet^{(1)},\mathscr P_\bullet^{(2)})\bigr).
\end{aligned}
\]
Se dice que las sucesiones espectrales $(a)$ y $(b)$ son las
\emph{sucesiones espectrales de Künneth}.

Se observará que las sucesiones espectrales $(a)$ y $(a')$ (resp.
$(b)$ y $(b')$) son idénticas cuando $\mathscr P_\bullet^{(1)}$ y
$\mathscr P_\bullet^{(2)}$ se reducen a sus términos de grado 0; en este caso,
las sucesiones $(c)$ y $(d)$ son degeneradas y, por tanto, no presentan
interés.
\end{theorem}

\begin{remark}[6.7.4]
\label{III.6.7.4}
Los hipertor globales que acabamos de definir comprenden como casos
particulares tanto los \emph{módulos de hipercohomología} definidos en
(6.2.1) como los \emph{hipertor locales} definidos en (6.5.3). Demostremos que
para todo morfismo $f:X\to Y$ \emph{cuasicompacto} y \emph{separado} y todo complejo
$\mathscr P_\bullet$ de $\mathscr O_X$-módulos cuasicoherentes, acotado
inferiormente, se tiene un isomorfismo canónico de $\partial$-funtores en
$\mathscr P_\bullet$

\label{III.6.7.4.1}
\[
  \boldsymbol{\mathscr{T}\!or}_n^Y(f,1_Y;\mathscr P_\bullet,\mathscr O_Y)
  \xrightarrow{\sim}
  \mathcal H^{-n}(f,\mathscr P_\bullet)
  \qquad\text{(para todo $n\in\mathbf Z$)}.
  \tag{6.7.4.1}\label{III.6.7.4.1.fr}
\]

En efecto, los métodos de pegado de (6.7.2) reducen inmediatamente al caso en
que $Y$ es afín; entonces, en virtud de (6.2.2), pueden calcularse los dos
miembros de (6.7.4.1) mediante un mismo recubrimiento finito $\mathfrak U$ de
$X$ por abiertos afines y, para el primer miembro, mediante el recubrimiento de
$Y$ formado por el propio $Y$; con las notaciones de (6.6.2), el bicomplejo
$L_{\bullet\bullet}^{(2)}$ se reduce entonces a su término de grados $(0,0)$,
igual a $A$, y la conclusión se deduce de (\textbf{0}, 11.7.5). Para una
generalización de este resultado, véase (6.7.7); pero se observará que, cuando en
el primer miembro de (6.7.4.1) se reemplaza $\mathscr O_Y$ por un
$\mathscr O_Y$-módulo cuasicoherente \emph{cualquiera} $\mathscr F$, ya no
se tiene en general un isomorfismo con
$\mathcal H^{-n}(f,\mathscr P_\bullet\otimes_{\mathscr O_Y}\mathscr F)$,
aunque, en el cálculo precedente, el bicomplejo
$L_{\bullet\bullet}^{(1)}\otimes_A L_{\bullet\bullet}^{(2)}$ todavía se
identifica con el bicomplejo
$C^\bullet(\mathfrak U,\mathscr P_\bullet\otimes_Y\mathscr F)$.

Por otra parte, se tiene un isomorfismo canónico de
bi-$\partial$-funtores

\label{III.6.7.4.2}
\[
  \boldsymbol{\mathscr{T}\!or}_\bullet^S
  (1_{X^{(1)}},1_{X^{(2)}};
   \mathscr P_\bullet^{(1)},\mathscr P_\bullet^{(2)})
  \xrightarrow{\sim}
  \boldsymbol{\mathscr{T}\!or}_\bullet^S
  (\mathscr P_\bullet^{(1)},\mathscr P_\bullet^{(2)}).
  \tag{6.7.4.2}\label{III.6.7.4.2.fr}
\]

En efecto, por (6.7.1) y (6.7.2), se reduce de nuevo al caso en que $S$ y
los $X^{(i)}$ son afines; en el cálculo del primer miembro de (6.7.4.2), puede
tomarse entonces como
\oldpage[III]{160}
recubrimiento $\mathfrak U^{(i)}$ la familia reducida al
único elemento $X^{(i)}$, de modo que, con las notaciones de (6.6.2),
$L_{\bullet\bullet}^{(i)}$ se reduce a
$\Gamma(X^{(i)},\mathscr P_\bullet^{(i)})$ (considerado como un bicomplejo
cuyos términos de primer grado $\ne0$ son nulos), y la igualdad de los dos
miembros de (6.7.4.2) se deduce de (6.4.1.1) y (6.3.1).
\end{remark}

\begin{proposition}[6.7.5]
\label{III.6.7.5}
Sea $u:\mathscr P_\bullet^{(1)}\to\mathscr Q_\bullet^{(1)}$ un
homomorfismo de complejos acotados inferiormente de
$\mathscr O_{X^{(1)}}$-módulos cuasicoherentes tal que el homomorfismo
\[
  \mathscr H_\bullet(u):
  \mathscr H_\bullet(\mathscr P_\bullet^{(1)})
  \longrightarrow
  \mathscr H_\bullet(\mathscr Q_\bullet^{(1)})
\]
inducido por $u$ sea un isomorfismo. Entonces los homomorfismos
\[
  {}^{(t)}\mathscr E
  (f_1,f_2;\mathscr P_\bullet^{(1)},\mathscr P_\bullet^{(2)})
  \longrightarrow
  {}^{(t)}\mathscr E
  (f_1,f_2;\mathscr Q_\bullet^{(1)},\mathscr P_\bullet^{(2)})
\]
inducidos por $u$ son isomorfismos para $t=a$, $t=b$ y $t=c$.
\end{proposition}

\begin{proof}
La afirmación relativa a la sucesión espectral $(c)$ se deduce de que esta
sucesión es birregular y de que, por hipótesis, el homomorfismo considerado es
un isomorfismo sobre los términos $E_2$ (\textbf{0}, 11.1.5). Esto
demuestra ya que
$
  \boldsymbol{\mathscr{T}\!or}_\bullet^S
  (f_1,f_2;\mathscr P_\bullet^{(1)},\mathscr P_\bullet^{(2)})
  \longrightarrow
  \boldsymbol{\mathscr{T}\!or}_\bullet^S
  (f_1,f_2;\mathscr Q_\bullet^{(1)},\mathscr P_\bullet^{(2)})
$
es un isomorfismo. Aplicando las relaciones (6.7.4.1) y (6.7.4.2), se ve
primero que los homomorfismos
$
  \mathcal H^{-n}(f_1,\mathscr P_\bullet^{(1)})
  \longrightarrow
  \mathcal H^{-n}(f_1,\mathscr Q_\bullet^{(1)})
  $ y $
  \boldsymbol{\mathscr{T}\!or}_n^S
  (\mathscr P_\bullet^{(1)},\mathscr P_\bullet^{(2)})
  \longrightarrow
  \boldsymbol{\mathscr{T}\!or}_n^S
  (\mathscr Q_\bullet^{(1)},\mathscr P_\bullet^{(2)})
$
deducidos de $u$ son isomorfismos. La afirmación relativa a las sucesiones
$(a)$ y $(b)$ se deduce entonces de que estas sucesiones son birregulares
(6.7.3) y de que los homomorfismos considerados son biyectivos sobre los
términos $E_2$ (\textbf{0}, 11.1.5).

Observemos, por otra parte, que, si
$u:\mathscr P_\bullet^{(1)}\to\mathscr Q_\bullet^{(1)}$ es una
\emph{equivalencia homotópica}, se deducen de ella isomorfismos canónicos
$
  {}^{(t)}\mathscr E
  (f_1,f_2;\mathscr P_\bullet^{(1)},\mathscr P_\bullet^{(2)})
  \longrightarrow
  {}^{(t)}\mathscr E
  (f_1,f_2;\mathscr Q_\bullet^{(1)},\mathscr P_\bullet^{(2)})
$
para las seis sucesiones espectrales. En efecto, si $S$ y los $Y^{(i)}$ son
afines, se deduce de $u$ una equivalencia homotópica de bicomplejos
$
  C^\bullet(\mathfrak U^{(1)},\mathscr P_\bullet^{(1)})
  \longrightarrow
  C^\bullet(\mathfrak U^{(1)},\mathscr Q_\bullet^{(1)})$,
y la proposición se deduce de la teoría general de la hiperhomología
(\textbf{0}, 11.3.2); el paso al caso general se efectúa por pegado,
utilizando el hecho de que una equivalencia homotópica de complejos induce
una equivalencia homotópica entre resoluciones proyectivas de
Cartan--Eilenberg de dichos complejos (M, XVII, 1.2).
\end{proof}

\begin{proposition}[6.7.6]
\label{III.6.7.6}
Supongamos que el complejo $\mathscr P_\bullet^{(1)}$ o el complejo
$\mathscr P_\bullet^{(2)}$ está formado por módulos $S$-planos (estando ambos
complejos acotados inferiormente). Se tiene entonces un isomorfismo canónico
de bi-$\partial$-funtores

\label{III.6.7.6.1}
\[
  \boldsymbol{\mathscr{T}\!or}_n^S
  (f_1,f_2;\mathscr P_\bullet^{(1)},\mathscr P_\bullet^{(2)})
  \xrightarrow{\sim}
  \mathcal H^{-n}\bigl(f_1\times_S f_2,
    \mathscr P_\bullet^{(1)}\otimes_S\mathscr P_\bullet^{(2)}\bigr).
  \tag{6.7.6.1}\label{III.6.7.6.1.fr}
\]

\end{proposition}

\begin{proof}
Supongamos primero que $S$, $Y^{(1)}$ e $Y^{(2)}$ son afines, de modo que
estamos en la situación de (6.6.2), cuyas notaciones conservamos. Supongamos,
por ejemplo, que $\mathscr P_\bullet^{(1)}$ está formado por módulos
$S$-planos, y calculemos el hipertor utilizando la observación (6.6.6): es,
por tanto, la homología de
$L_{\bullet\bullet}^{(1)}\otimes_A M_{\bullet\bullet\bullet}^{(2)}$, donde
$M_{\bullet\bullet\bullet}^{(2)}$ es una resolución proyectiva de
Cartan--Eilenberg de $L_{\bullet\bullet}^{(2)}$, en el sentido de
(\textbf{0}, 11.7.1). Por otra parte, los módulos
$L_{ij}^{(1)}$ son \emph{planos} sobre $A$ en virtud de la hipótesis
(1.4.15.1); de (\textbf{0}, 11.7.5) se deduce entonces un isomorfismo
canónico

\label{III.6.7.6.2}
\[
  \mathbf{Tor}_\bullet^S
  (\mathfrak U^{(1)},\mathfrak U^{(2)};
   \mathscr P_\bullet^{(1)},\mathscr P_\bullet^{(2)})
  \xrightarrow{\sim}
  \mathscr H_\bullet
  (L_{\bullet\bullet}^{(1)}\otimes_A L_{\bullet\bullet}^{(2)}).
  \tag{6.7.6.2}\label{III.6.7.6.2.fr}
\]

Por otra parte, se tiene un homomorfismo natural de bicomplejos de
$L_{\bullet\bullet}^{(1)}\otimes_A L_{\bullet\bullet}^{(2)}$ en
$C^\bullet(\mathfrak U,\mathscr Q_\bullet)$, donde $\mathfrak U$ es el
recubrimiento de $Z=X^{(1)}\times_S X^{(2)}$ por los abiertos afines
\oldpage[III]{161}
$U_\alpha^{(1)}\times_S U_\beta^{(2)}$ y
$\mathscr Q_\bullet
=\mathscr P_\bullet^{(1)}\otimes_S\mathscr P_\bullet^{(2)}$ (considerado como
complejo simple para el grado total); en efecto, la definición de este
homomorfismo se dio en sustancia durante el cálculo de la sucesión $(b')$ en
(6.6.4), para $q=0$; basta simplemente (conservando las notaciones de (6.6.4))
tener en cuenta que hay, por una parte, un homomorfismo natural del complejo
$N^{\bullet k}$ en el complejo
$C^\bullet(\Phi^{(1)},\Phi^{(2)};\mathscr S_k)$
(\textbf{0}, 11.8.5), por otra parte, un homomorfismo natural de este
último complejo en el complejo
$P^\bullet(\Phi^{(1)},\Phi^{(2)};\mathscr S_k)$
(\textbf{0}, 11.8.6), y, por último, un homomorfismo natural de este
último complejo de cocadenas en el subcomplejo de las cocadenas alternadas
(\textbf{0}, 11.8.7). Además, el homomorfismo de bicomplejos así definido
induce un isomorfismo en homología, como se vio en (6.6.4); componiendo con
(6.7.6.2), se obtiene, por tanto, un isomorfismo

\label{III.6.7.6.3}
\[
  \mathbf{Tor}_n^S
  (\mathfrak U^{(1)},\mathfrak U^{(2)};
   \mathscr P_\bullet^{(1)},\mathscr P_\bullet^{(2)})
  \xrightarrow{\sim}
  \mathbf H^{-n}
  (\mathfrak U,\mathscr P_\bullet^{(1)}\otimes_S\mathscr P_\bullet^{(2)}).
  \tag{6.7.6.3}\label{III.6.7.6.3.fr}
\]

Hay que demostrar a continuación que el isomorfismo así definido no depende
de los recubrimientos abiertos elegidos (el segundo miembro de (6.7.6.3) es
canónicamente isomorfo a
$\mathbf H^{-n}
(X^{(1)}\times_S X^{(2)},
 \mathscr P_\bullet^{(1)}\otimes_S\mathscr P_\bullet^{(2)})$ por (6.2.2));
esto se hace mediante (6.6.7), observando (con las notaciones de (6.6.7)) que
se tiene un diagrama conmutativo salvo homotopía
\[
\xymatrix@C=18pt@R=28pt{
  C^\bullet(\mathfrak U^{(1)},\mathscr P_\bullet^{(1)})
  \otimes_A
  C^\bullet(\mathfrak U^{(2)},\mathscr P_\bullet^{(2)})
  \ar[r]\ar[d]
  & C^\bullet(\mathfrak U,\mathscr Q_\bullet)\ar[d]
  \\
  C^\bullet(\mathfrak V^{(1)},\mathscr P_\bullet^{(1)})
  \otimes_A
  C^\bullet(\mathfrak V^{(2)},\mathscr P_\bullet^{(2)})
  \ar[r]
  & C^\bullet(\mathfrak V,\mathscr Q_\bullet)
}
\]
donde las flechas horizontales son los homomorfismos definidos anteriormente.
Por último, hay que pasar al caso general por pegado, lo que se hace sin
dificultad como en (6.7.1) y (6.7.2); dejamos los detalles al lector.
\end{proof}

\begin{proposition}[6.7.7]
\label{III.6.7.7}
Supongamos que $\mathscr P_\bullet^{(1)}$ y
$\mathscr P_\bullet^{(2)}$ están acotados inferiormente, y que todos los
módulos $\mathcal H^{-n}(f_1,\mathscr P_\bullet^{(1)})$ \textup{o} todos los módulos
$\mathcal H^{-n}(f_2,\mathscr P_\bullet^{(2)})$ son $S$-planos. Se tiene
entonces un isomorfismo canónico de bi-$\partial$-funtores ($n$ recorre
$\mathbf Z$)

\label{III.6.7.7.1}
\[
  \boldsymbol{\mathscr{T}\!or}_n^S
  (f_1,f_2;\mathscr P_\bullet^{(1)},\mathscr P_\bullet^{(2)})
  \xrightarrow{\sim}
  \bigoplus_{q_1+q_2=n}
  \mathcal H^{-q_1}(f_1,\mathscr P_\bullet^{(1)})
  \otimes_S
  \mathcal H^{-q_2}(f_2,\mathscr P_\bullet^{(2)}).
  \tag{6.7.7.1}\label{III.6.7.7.1.fr}
\]

\end{proposition}

\begin{proof}
En efecto, en vista de (6.5.8), la sucesión espectral $(a)$ de (6.7.3) es
degenerada, y la proposición se deduce inmediatamente de
(\textbf{0}, 11.1.6), pues esta sucesión es birregular (6.7.3).
\end{proof}

\begin{theorem}[6.7.8]
\label{III.6.7.8}
Supongamos que se cumplen las condiciones siguientes:

  $1^\circ$
    Los complejos $\mathscr P_\bullet^{(1)}$ y
    $\mathscr P_\bullet^{(2)}$ están acotados inferiormente.
  $2^\circ$
    El complejo $\mathscr P_\bullet^{(1)}$ \textup{o} el complejo
    $\mathscr P_\bullet^{(2)}$ está formado por módulos $S$-planos.

  $3^\circ$
    \oldpage[III]{162}
    Todos los módulos $\mathscr H^{-n}(f_1,\mathscr P_\bullet^{(1)})$ \textup{o} todos
    los módulos $\mathscr H^{-n}(f_2,\mathscr P_\bullet^{(2)})$ son
    $S$-planos.

Se tiene entonces un isomorfismo canónico de bi-$\partial$-funtores
($n$ recorre $\mathbf Z$)

\label{III.6.7.8.1}
\[
  \mathscr H^n\bigl(f_1\times_S f_2,
    \mathscr P_\bullet^{(1)}\otimes_S\mathscr P_\bullet^{(2)}\bigr)
  \xrightarrow{\sim}
  \bigoplus_{n_1+n_2=n}
  \mathscr H^{n_1}(f_1,\mathscr P_\bullet^{(1)})
  \otimes_S
  \mathscr H^{n_2}(f_2,\mathscr P_\bullet^{(2)})
  \tag{6.7.8.1}\label{III.6.7.8.1.fr}
\]

(«fórmula de Künneth»).
\end{theorem}

\begin{proof}
Esto se deduce de (6.7.6) y (6.7.7).

Cuando $S$, $Y^{(1)}$ e $Y^{(2)}$ son afines, el isomorfismo recíproco de
(6.7.8.1) se deduce (con las notaciones de (6.7.6)) del homomorfismo de
bicomplejos
\[
  C^\bullet(\mathfrak U^{(1)},\mathscr P_\bullet^{(1)})
  \otimes_A
  C^\bullet(\mathfrak U^{(2)},\mathscr P_\bullet^{(2)})
  \longrightarrow
  C^\bullet(\mathfrak U,\mathscr Q_\bullet)
\]
mediante el procedimiento definido en (G, I, 2.7), como se deduce de
(G, I, 5.5).
\end{proof}

\begin{proposition}[6.7.9]
\label{III.6.7.9}
Supongamos que se cumplen las tres condiciones siguientes:

  $1^\circ$
    $S$, $Y^{(1)}$ e $Y^{(2)}$ son localmente noetherianos, $f_1$ y $f_2$
    son propios, y $Y^{(1)}$ o $Y^{(2)}$ es de tipo finito sobre $S$.
  $2^\circ$
    $\mathscr P_\bullet^{(1)}$ y $\mathscr P_\bullet^{(2)}$ están acotados
    inferiormente.
  $3^\circ$
    Para todo $n\in\mathbf Z$,
    $\mathscr H_n(\mathscr P_\bullet^{(i)})$ es un módulo coherente
    ($i=1,2$).

En estas condiciones,
$\mathfrak{Tor}_n^S
(f_1,f_2;\mathscr P_\bullet^{(1)},\mathscr P_\bullet^{(2)})$ es un
$\mathscr O_Y$-módulo coherente (con
$Y=Y^{(1)}\times_S Y^{(2)}$).
\end{proposition}

\begin{proof}
De (6.5.13) se deduce que los hipertor locales
$\mathfrak{Tor}_n^S(\mathscr P_\bullet^{(1)},\mathscr P_\bullet^{(2)})$ son
$\mathscr O_X$-módulos coherentes
($X=X^{(1)}\times_S X^{(2)}$ es localmente noetheriano, pues uno de los
$X^{(i)}$ es, por hipótesis, de tipo finito sobre $S$ (\textbf I, 6.3.4 y 6.3.8)).
Como $Y$ es localmente noetheriano y $f_1\times_S f_2$ es propio
(\textbf{II}, 5.4.2), de (6.2.5) se deduce que los términos
${}^{(b)}\mathscr E_{pq}^2$ de (6.7.3) son $\mathscr O_Y$-módulos coherentes.
Como todas las sucesiones espectrales de (6.7.3) son birregulares en virtud
de la hipótesis $2^\circ$, se concluye por (\textbf{0}, 11.1.8).
\end{proof}

\begin{env}[6.7.10]
\label{III.6.7.10}
Sean ahora $Y'^{(i)}$ dos $S$-preesquemas,
$v_i:Y'^{(i)}\to Y^{(i)}$ dos $S$-morfismos ($i=1,2$) y
$v=v_1\times_S v_2$ su producto, que es un $S$-morfismo $Y'\to Y$, donde se
pone $Y'=Y'^{(1)}\times_S Y'^{(2)}$. Consideremos, por otra parte, para
$i=1,2$, un $S$-preesquema $X'^{(i)}$ y dos $S$-morfismos
$u_i:X'^{(i)}\to X^{(i)}$, $f'_i:X'^{(i)}\to Y'^{(i)}$, de modo que los
diagramas

\label{III.6.7.10.1}
\[
\xymatrix@C=36pt@R=28pt{
  X'^{(i)}\ar[r]^{u_i}\ar[d]_{f'_i}
  & X^{(i)}\ar[d]^{f_i}
  \\
  Y'^{(i)}\ar[r]_{v_i}
  & Y^{(i)}
}
\tag{6.7.10.1}\label{III.6.7.10.1.fr}
\]

sean conmutativos, siendo los morfismos $f'_i$ \emph{separados} y
\emph{cuasicompactos}. Se tienen entonces homomorfismos canónicos
$\mathscr O_{Y'}$-lineales de funtores espectrales

\label{III.6.7.10.2}
\[
  v^*\bigl({}^{(t)}\mathscr E
  (f_1,f_2;\mathscr P_\bullet^{(1)},\mathscr P_\bullet^{(2)})\bigr)
  \longrightarrow
  {}^{(t)}\mathscr E
  (f'_1,f'_2;u_1^*(\mathscr P_\bullet^{(1)}),
                 u_2^*(\mathscr P_\bullet^{(2)}))
  \tag{6.7.10.2}\label{III.6.7.10.2.fr}
\]

para $t=a$, $a'$, $b$, $b'$, $c$, $d$. Para definirlos, supongamos primero
que $S=\operatorname{Spec}(A)$,
$Y^{(i)}=\operatorname{Spec}(B_i)$ e
$Y'^{(i)}=\operatorname{Spec}(B'_i)$ son afines; entonces $X^{(i)}$ y
$X'^{(i)}$ son esquemas cuasicompactos. Para calcular las sucesiones
espectrales
${}^{(t)}\mathscr E
(f_1,f_2;\mathscr P_\bullet^{(1)},\mathscr P_\bullet^{(2)})$, consideraremos,
como en (6.6.1), recubrimientos finitos $\mathfrak U^{(i)}$ por abiertos afines
de $X^{(i)}$ ($i=1,2$); para calcular
${}^{(t)}\mathscr E
(f'_1,f'_2;u_1^*(\mathscr P_\bullet^{(1)}),
u_2^*(\mathscr P_\bullet^{(2)}))$, consideraremos recubrimientos finitos
$\mathfrak U'^{(i)}$

\oldpage[III]{163}
de $X'^{(i)}$ por abiertos afines, \emph{más finos}, respectivamente, que los
recubrimientos $u_i^{-1}(\mathfrak U^{(i)})$ ($i=1,2$). Está claro que el
bicomplejo
$C^\bullet(\mathfrak U^{(i)},\mathscr P_\bullet^{(i)})
=L_{\bullet\bullet}^{(i)}$ puede considerarse canónicamente como un
subbicomplejo de
$C^\bullet(u_i^{-1}(\mathfrak U^{(i)}),u_i^*(\mathscr P_\bullet^{(i)}))$
($\mathbf{0}_{\mathrm I}$, 4.4.3.2); además, eligiendo un morfismo simplicial
(G, II, 5.7) de $\mathfrak U'^{(i)}$ en
$u_i^{-1}(\mathfrak U^{(i)})$, se define un homomorfismo de bicomplejos
$C^\bullet(u_i^{-1}(\mathfrak U^{(i)}),u_i^*(\mathscr P_\bullet^{(i)}))
\to C^\bullet(\mathfrak U'^{(i)},u_i^*(\mathscr P_\bullet^{(i)}))$
del cual se obtiene, por composición, un homomorfismo de bicomplejos
$L_{\bullet\bullet}^{(i)}\to L_{\bullet\bullet}^{\prime(i)}
=C^\bullet(\mathfrak U'^{(i)},u_i^*(\mathscr P_\bullet^{(i)}))$.
Además, este homomorfismo se sustituye por uno homotópico cuando se cambia el
morfismo simplicial (G, II, 5.7.1); así se obtiene un homomorfismo bien
definido de funtores espectrales:

\label{III.6.7.10.3}
\[
  {}^{(t)}\mathscr E
  (\mathfrak U^{(1)},\mathfrak U^{(2)};
   \mathscr P_\bullet^{(1)},\mathscr P_\bullet^{(2)})
  \longrightarrow
  {}^{(t)}\mathscr E
  (\mathfrak U'^{(1)},\mathfrak U'^{(2)};
   u_1^*(\mathscr P_\bullet^{(1)}),u_2^*(\mathscr P_\bullet^{(2)})).
  \tag{6.7.10.3}\label{III.6.7.10.3.fr}
\]

Se verifica inmediatamente que, si $\mathfrak V^{(i)}$ es un recubrimiento
afín finito de $X^{(i)}$ más fino que $\mathfrak U^{(i)}$, y
$\mathfrak V'^{(i)}$ un recubrimiento afín finito de $X'^{(i)}$, más fino que
$u_i^{-1}(\mathfrak V^{(i)})$ y que $\mathfrak U'^{(i)}$, el diagrama
\[
\xymatrix@C=34pt@R=30pt{
  C^\bullet(\mathfrak U^{(i)},\mathscr P_\bullet^{(i)})
  \ar[r]\ar[d]
  & C^\bullet(\mathfrak V^{(i)},\mathscr P_\bullet^{(i)})\ar[d]
  \\
  C^\bullet(\mathfrak U'^{(i)},u_i^*(\mathscr P_\bullet^{(i)}))
  \ar[r]
  & C^\bullet(\mathfrak V'^{(i)},u_i^*(\mathscr P_\bullet^{(i)}))
}
\]
es conmutativo, lo que implica que el homomorfismo (6.7.10.3) no depende
esencialmente de los recubrimientos $\mathfrak U^{(i)}$ y
$\mathfrak U'^{(i)}$ considerados. Por tanto, se ha definido en realidad un
homomorfismo de $A$-\emph{módulos}

\label{III.6.7.10.4}
\[
  {}^{(t)}E
  (X^{(1)},X^{(2)};\mathscr P_\bullet^{(1)},\mathscr P_\bullet^{(2)})
  \longrightarrow
  {}^{(t)}E
  (X'^{(1)},X'^{(2)};
   u_1^*(\mathscr P_\bullet^{(1)}),u_2^*(\mathscr P_\bullet^{(2)}))
  \tag{6.7.10.4}\label{III.6.7.10.4.fr}
\]

pero, por la definición de los $u_i^*(\mathscr P_\bullet^{(i)})$ y en virtud
de la conmutatividad de (6.7.10.1), está claro que este homomorfismo es también
un homomorfismo de $(B_1\otimes_A B_2)$-módulos; como el segundo miembro de
(6.7.10.4) está formado por $(B'_1\otimes_A B'_2)$-módulos, se deduce
canónicamente de (6.7.10.4) un homomorfismo de
$(B'_1\otimes_A B'_2)$-módulos

\label{III.6.7.10.5}
\[
  {}^{(t)}E
  (X^{(1)},X^{(2)};\mathscr P_\bullet^{(1)},\mathscr P_\bullet^{(2)})
  \otimes_{B_1\otimes_A B_2}(B'_1\otimes_A B'_2)
  \longrightarrow
  {}^{(t)}E
  (X'^{(1)},X'^{(2)};
   u_1^*(\mathscr P_\bullet^{(1)}),u_2^*(\mathscr P_\bullet^{(2)}))
  \tag{6.7.10.5}\label{III.6.7.10.5.fr}
\]

lo que, teniendo en cuenta (\textbf I, 1.6.5), no es sino el homomorfismo buscado
(6.7.10.2) en el caso particular considerado.

Queda pasar al caso general siguiendo los procedimientos de pegado de (6.7.1)
y (6.7.2); el
segundo paso es inmediato; por lo que respecta al primero, se consideran, como
en (6.7.1), elementos $g_i\in B_i$ y sus imágenes $g'_i\in B'_i$, el producto
tensorial $g=g_1\otimes g_2$ en $B=B_1\otimes_A B_2$ y su imagen
$g'=g'_1\otimes g'_2$ en $B'=B'_1\otimes_A B'_2$, y todo se reduce a utilizar

\oldpage[III]{164}
el isomorfismo canónico
$(M\otimes_B B')_{g'}\xrightarrow{\sim}M_g\otimes_{B_g}B'_{g'}$
($\mathbf{0}_{\mathrm I}$, 1.5.4); dejamos los detalles al lector.
\end{env}

\begin{env}[6.7.11]
\label{III.6.7.11}
La teoría de los hipertor globales, desarrollada anteriormente para dos
$S$-morfismos $X^{(i)}\to Y^{(i)}$ y dos complejos
$\mathscr P_\bullet^{(i)}$ de módulos cuasicoherentes acotados inferiormente,
se extiende inmediatamente al caso general siguiente: se tiene un preesquema
$S$, una familia finita de $S$-preesquemas $Y^{(i)}$ ($i\in I$), una familia
finita de $S$-morfismos \emph{separados} y \emph{cuasicompactos}
$f_i:X^{(i)}\to Y^{(i)}$, y, para cada $i$, un complejo acotado inferiormente
de $\mathscr O_{X^{(i)}}$-módulos cuasicoherentes
$\mathscr P_\bullet^{(i)}$. Si $Y$ es el producto de los $S$-preesquemas
$Y^{(i)}$, se define entonces, para cada entero $n\in\mathbf Z$, un
$\mathscr O_Y$-módulo cuasicoherente
$\mathfrak{Tor}_n^S((f_i)_{i\in I};
(\mathscr P_\bullet^{(i)})_{i\in I})$;
estos módulos forman un $\partial$-funtor covariante en cada uno de los
complejos $\mathscr P_\bullet^{(i)}$; además, este funtor es el término límite
común de seis funtores espectrales
${}^{(t)}\mathscr E((f_i)_{i\in I};(\mathscr P_\bullet^{(i)})_{i\in I})$.
Dejamos al lector la tarea de repetir, para este caso general, las definiciones
y los razonamientos hechos anteriormente para $I=\{1,2\}$. Observemos
simplemente que, cuando $I$ se reduce a un solo elemento, se recupera la
hipercohomología $\mathscr H^\bullet(f,\mathscr P_\bullet)$ definida en
(6.2.7) (como ya se observó en (6.7.4)). Cuando $I$ es el intervalo
$1\leq i\leq m$ de $\mathbf N$, escribiremos
\phantomsection\label{III.6.7.11.2}
\[
  \mathfrak{Tor}_n^S
  (f_1,\ldots,f_m;
   \mathscr P_\bullet^{(1)},\ldots,\mathscr P_\bullet^{(m)})
  \quad\text{para}\quad
  \mathfrak{Tor}_n^S
  ((f_i)_{i\in I};(\mathscr P_\bullet^{(i)})_{i\in I}).
\]
\end{env}

\begin{proposition}[6.7.12]
\label{III.6.7.12}
Con las notaciones de (6.7.11), sea $J$ un subconjunto de $I$ tal que, para
$i\in I\setmin J$, se tenga $X^{(i)}=Y^{(i)}=S$, que $f_i$ sea la identidad y que
$\mathscr P_\bullet^{(i)}$ sea el complejo reducido al término de grado $0$
igual a $\mathscr O_S$. Hay entonces un isomorfismo canónico de
$\partial$-funtores

\label{III.6.7.12.1}
\[
  \mathfrak{Tor}_\bullet^S
  ((f_i)_{i\in I};(\mathscr P_\bullet^{(i)})_{i\in I})
  \xrightarrow{\sim}
  \mathfrak{Tor}_\bullet^S
  ((f_i)_{i\in J};(\mathscr P_\bullet^{(i)})_{i\in J}).
  \tag{6.7.12.1}\label{III.6.7.12.1.fr}
\]

\end{proposition}

\begin{proof}
Se puede limitar a definir este isomorfismo cuando $S$ y los $Y^{(i)}$ son
afines, pues el pegado se efectúa como de costumbre. Para $i\in I\setmin J$, puede
tomarse el recubrimiento $\mathfrak U^{(i)}$ formado por el único conjunto $S$,
y entonces
$L_{\bullet\bullet}^{(i)}=C^\bullet(\mathfrak U^{(i)},
\mathscr P_\bullet^{(i)})$ se reduce a su único término de grados $(0,0)$,
igual a $\Gamma(S,\mathscr O_S)=A(S)$; el isomorfismo (6.7.12.1) es entonces
evidente.
\end{proof}

\begin{remark}[6.7.13]
\label{III.6.7.13}
Con las notaciones de (6.7.3), consideremos el $S$-isomorfismo canónico
$Y^{(1)}\times_S Y^{(2)}\to Y^{(2)}\times_S Y^{(1)}$ (\textbf I, 3.3.5); entonces la
imagen por este isomorfismo de
$\mathfrak{Tor}_\bullet^S(f_1,f_2;
\mathscr P_\bullet^{(1)},\mathscr P_\bullet^{(2)})$ es
$\mathfrak{Tor}_\bullet^S(f_2,f_1;
\mathscr P_\bullet^{(2)},\mathscr P_\bullet^{(1)})$; como la cuestión es
local, se reduce al caso considerado en (6.6.2), y, si se designa por
$M_{\bullet\bullet\bullet}^{(i)}$ una resolución proyectiva de
Cartan--Eilenberg de $L_{\bullet\bullet}^{(i)}$ ($i=1,2$), el isomorfismo
considerado transforma
$M_{\bullet\bullet\bullet}^{(1)}\otimes_A
M_{\bullet\bullet\bullet}^{(2)}$ en
$M_{\bullet\bullet\bullet}^{(2)}\otimes_A
M_{\bullet\bullet\bullet}^{(1)}$, de donde nuestra aserción al considerar la homología de los complejos simples asociados a estos tricomplejos.
\end{remark}

\subsection{Las sucesiones espectrales de asociatividad de los hipertor globales.}

\begin{env}[6.8.1]
Siendo las hipótesis y notaciones las de (6.7.11) (y en particular los
$\mathscr P_\bullet^{(i)}$ supuestos \emph{acotados inferiormente}), supongamos
dada una partición $(I_j)_{j\in J}$ del conjunto de índices $I$; nos
proponemos dar una relación de «asociatividad» entre los hipertor
$\mathfrak{Tor}_n^S((f_i)_{i\in I};(\mathscr P_\bullet^{(i)})_{i\in I})$ y
cada uno de los hipertor «parciales»
\[
  \mathscr T_{\bullet j}
  =\mathfrak{Tor}_\bullet^S
  ((f_i)_{i\in I_j};(\mathscr P_\bullet^{(i)})_{i\in I_j}).
\]

\oldpage[III]{165}
Para simplificar la escritura, nos limitaremos al caso en que $I$ es
el intervalo $1\leq i\leq m$, y en que la partición $(I_j)$ se compone de los dos
intervalos $\{1,2,\ldots,r\}$ y $\{r+1,\ldots,m\}$.
\end{env}

\begin{proposition}[6.8.2]
Existe un funtor espectral canónico birregular \textup{(llamado «funtor espectral
de asociatividad»)} denotado por
\[
  {}^{(e)}\mathscr E^S
  (f_1,\ldots,f_m;
   \mathscr P_\bullet^{(1)},\ldots,\mathscr P_\bullet^{(m)})
  \qquad\textup{(o simplemente ${}^{(e)}\mathscr E(f_1,\ldots,f_m;
\mathscr P_\bullet^{(1)},\ldots,\mathscr P_\bullet^{(m)})$)}
\]
cuyo
resultado es
$\mathfrak{Tor}_\bullet^S(f_1,\ldots,f_m;
\mathscr P_\bullet^{(1)},\ldots,\mathscr P_\bullet^{(m)})$
y cuyo término $E_2$ es
\[
\begin{split}
  {}^{(e)}\mathscr E_{pq}^2
  =\bigoplus_{q_1+q_2=q}
  \mathscr{T}\!or_p^S\Bigl(
  &\mathfrak{Tor}_{q_1}^S
  (f_1,\ldots,f_r;
   \mathscr P_\bullet^{(1)},\ldots,\mathscr P_\bullet^{(r)}),\\
  &\mathfrak{Tor}_{q_2}^S
  (f_{r+1},\ldots,f_m;
   \mathscr P_\bullet^{(r+1)},\ldots,\mathscr P_\bullet^{(m)})
  \Bigr).
\end{split}
\]
\end{proposition}

\begin{proof}
Aquí se identifica canónicamente $Y$ con
$Z^{(1)}\times_S Z^{(2)}$, donde
$Z^{(1)}=Y^{(1)}\times_S Y^{(2)}\times_S\ldots\times_S Y^{(r)}$, y
$Z^{(2)}=Y^{(r+1)}\times_S\ldots\times_S Y^{(m)}$. Nos limitaremos al
caso en que $S$ y todos los $Y^{(i)}$ son \emph{afines}.
El paso de este caso al caso general utiliza los métodos de
\sref{III.6.7.1} y \sref{III.6.7.2}; dejamos los detalles al
lector. Basta, pues, demostrar el siguiente enunciado afín.
\end{proof}

\begin{corollary}[6.8.3]
Sean $A$ un anillo, $S=\operatorname{Spec}(A)$, $X^{(i)}$
($1\leq i\leq m$) unos $S$-esquemas cuasicompactos y, para cada $i$, sea
$\mathscr P_\bullet^{(i)}$ un complejo de
$\mathscr O_{X^{(i)}}$-Módulos cuasicoherentes acotado inferiormente. Existe
un funtor espectral canónico birregular que tiene por término límite
\[
  \mathbf{Tor}_\bullet^S
  (X^{(1)},\ldots,X^{(m)};
   \mathscr P_\bullet^{(1)},\ldots,\mathscr P_\bullet^{(m)})
\]
y cuyo término $E_2$ está dado por
\[
\begin{split}
  {}^{(e)}E_{pq}^2
  =\bigoplus_{q_1+q_2=q}
  \operatorname{Tor}_p^A\Bigl(
  &\mathbf{Tor}_{q_1}^S
  (X^{(1)},\ldots,X^{(r)};
   \mathscr P_\bullet^{(1)},\ldots,\mathscr P_\bullet^{(r)}),\\
  &\mathbf{Tor}_{q_2}^S
  (X^{(r+1)},\ldots,X^{(m)};
   \mathscr P_\bullet^{(r+1)},\ldots,\mathscr P_\bullet^{(m)})
  \Bigr)
\end{split}
\]
\end{corollary}

\begin{proof}
Según la definición dada en (6.6.2), el cálculo del hipertor considerado se
hace tomando para cada $i$ un recubrimiento abierto afín finito
$\mathfrak U^{(i)}$ de $X^{(i)}$, considerando los bicomplejos
$L_{\bullet\bullet}^{(i)}=C^\bullet(\mathfrak U^{(i)},
\mathscr P_\bullet^{(i)})$, una resolución proyectiva
$M_{\bullet\bullet\bullet}^{(i)}$ de Cartan-Eilenberg de cada uno de estos
bicomplejos (en el sentido de (\textbf{0}, 11.7.1)), el producto tensorial
$M_{\bullet\bullet\bullet}
=\bigotimes_{i=1}^m M_{\bullet\bullet\bullet}^{(i)}$
de estos tricomplejos, y tomando la homología de
$M_{\bullet\bullet\bullet}$. Consideremos $M_{\bullet\bullet\bullet}$ como
un complejo simple $N_\bullet$, producto tensorial de los dos complejos simples
\[
  N'_\bullet=\bigotimes_{i=1}^r M_{\bullet\bullet\bullet}^{(i)},
  \qquad
  N''_\bullet=\bigotimes_{i=r+1}^m M_{\bullet\bullet\bullet}^{(i)},
\]
donde $N'_\bullet$ y $N''_\bullet$ están graduados por la suma de los grados totales
de los $M_{\bullet\bullet\bullet}^{(i)}$. Además, los $A$-módulos de los
complejos $N'_\bullet$ y $N''_\bullet$ son \emph{proyectivos}, luego resulta de
(6.5.9) que se tiene
$H_\bullet(M_{\bullet\bullet\bullet})
=\mathbf{Tor}_\bullet^A(N'_\bullet,N''_\bullet)$;
la sucesión espectral buscada no es entonces sino la sucesión (6.3.2.2)
aplicada a los complejos $N'_\bullet$ y $N''_\bullet$, teniendo en cuenta
la interpretación de los módulos de homología de estos complejos que resulta de lo
que precede (cuando se aplican las observaciones del comienzo a cada uno de los productos
parciales
$X^{(1)}\times_S\ldots\times_S X^{(r)}$ y
$X^{(r+1)}\times_S\ldots\times_S X^{(m)}$). Finalmente, las propiedades de
regularidad resultan de (6.3.2) y del hecho de que, al estar los
$\mathscr P_\bullet^{(i)}$ acotados inferiormente, lo mismo ocurre
con los $M_{\bullet\bullet\bullet}^{(i)}$; por consiguiente, $N'_\bullet$ y
$N''_\bullet$ están acotados inferiormente.
\end{proof}

\oldpage[III]{166}
\subsection{Las sucesiones espectrales de cambio de base en los hipertor
globales.}

\begin{env}[6.9.1]
\label{III.6.9.1}
Se mantienen las hipótesis y notaciones de \sref{III.6.7.11} (y, en
particular, se supone que los $\mathscr P_\bullet^{(i)}$ están \emph{acotados
inferiormente}). Consideremos un morfismo $g:S'\to S$ de preesquemas, y
pongamos
$Y'^{(i)}=Y_{(S')}^{(i)}$, $X'^{(i)}=X_{(S')}^{(i)}$ y
$\mathscr P_\bullet'^{(i)}=
\mathscr P_\bullet^{(i)}\otimes_{\mathscr O_S}\mathscr O_{S'}$,
de modo que $\mathscr P_\bullet'^{(i)}$ es un complejo de
$\mathscr O_{X'^{(i)}}$-módulos cuasicoherentes; sea
$f'_i=(f_i)_{(S')}:X'^{(i)}\to Y'^{(i)}$, que es un morfismo separado y
cuasicompacto (\sref[I]{I.5.5.1} y \sref{III.6.6.4}). Nos proponemos
estudiar las relaciones entre los $\mathscr O_{Y'}$-módulos cuasicoherentes
$\mathfrak{Tor}_n^{S'}((f'_i)_{i\in I};
(\mathscr P_\bullet'^{(i)})_{i\in I})$ y
$\mathfrak{Tor}_n^S((f_i)_{i\in I};
(\mathscr P_\bullet^{(i)})_{i\in I})
\otimes_{\mathscr O_S}\mathscr O_{S'}$, donde
$Y'=Y\times_S S'=Y_{(S')}$. Un caso particularmente sencillo es el
siguiente, que se reduce a \sref[I]{I.1.4.15} cuando $I$ se reduce a un solo
elemento y $\mathscr P_\bullet$ a un solo módulo:
\end{env}

\begin{proposition}[6.9.2]
\label{III.6.9.2}
Si el morfismo $g:S'\to S$ es plano, se tiene un isomorfismo canónico de
$\partial$-funtores (en los $\mathscr P_\bullet^{(i)}$):

\label{III.6.9.2.1}
\[
  \mathfrak{Tor}_\bullet^S
  ((f_i)_{i\in I};(\mathscr P_\bullet^{(i)})_{i\in I})
  \otimes_{\mathscr O_S}\mathscr O_{S'}
  \xrightarrow{\sim}
  \mathfrak{Tor}_\bullet^{S'}
  ((f'_i)_{i\in I};(\mathscr P_\bullet'^{(i)})_{i\in I}).
  \tag{6.9.2.1}\label{III.6.9.2.1.fr}
\]

\end{proposition}

\begin{proof}
Todavía se puede restringir al caso en que $S$, $S'$ y los $Y^{(i)}$ son
afines, efectuándose el pegado mediante los métodos de \sref{III.6.7.1} y
\sref{III.6.7.2}. Sean $S=\operatorname{Spec}(A)$,
$S'=\operatorname{Spec}(A')$, y tomemos para cada $i$ un recubrimiento abierto
afín $\mathfrak U^{(i)}$ de $X^{(i)}$; si
$u_i:X'^{(i)}\to X^{(i)}$ es la proyección canónica,
$u_i^{-1}(\mathfrak U^{(i)})$ es un recubrimiento abierto afín de
$X'^{(i)}$ (\sref[II]{II.1.5.5}), que denotaremos por
$\mathfrak U'^{(i)}$; está claro entonces que
$C^\bullet(\mathfrak U'^{(i)},\mathscr P_\bullet'^{(i)})
=C^\bullet(\mathfrak U^{(i)},\mathscr P_\bullet^{(i)})\otimes_A A'$,
y la existencia del isomorfismo \eqref{III.6.9.2.1.eq} es inmediata, pues si
$M_{\bullet\bullet\bullet}^{(i)}$ es una resolución proyectiva de
Cartan-Eilenberg de
$L_{\bullet\bullet}^{(i)}=C^\bullet(\mathfrak U^{(i)},
\mathscr P_\bullet^{(i)})$ en el sentido de \sref[0]{0.11.7.1}, formada por
$A$-módulos, $M_{\bullet\bullet\bullet}^{(i)}\otimes_A A'$ es una resolución
proyectiva de Cartan-Eilenberg (en el mismo sentido) de
$L_{\bullet\bullet}^{(i)}\otimes_A A'$ formada por $A'$-módulos, en virtud
de la hipótesis de que $A'$ es un $A$-módulo plano; esta misma hipótesis
muestra además que
$H_\bullet(\bigotimes_{i=1}^m(M_{\bullet\bullet\bullet}^{(i)}\otimes_A A'))
=H_\bullet(\bigotimes_{i=1}^m M_{\bullet\bullet\bullet}^{(i)})\otimes_A A'$.
\end{proof}


\label{III.6.9.2.2}
Se observará que, cuando $I$ se reduce al único elemento $1$, la fórmula
\eqref{III.6.9.2.1.eq} se deduce directamente de \sref{III.6.7.7}, aplicado
tomando $Y_2=X_2=S'$, $f_2=1_{S'}$, y el complejo
$\mathscr P_\bullet^{(2)}$ reducido a su término de grado $0$, igual a
$\mathscr O_{S'}$; se sabe entonces que la hipercohomología
$\mathscr H^n(1_{S'},\mathscr O_{S'})$ es nula para todo $n\ne0$ y se reduce
a $\mathscr O_{S'}$ para $n=0$ (\sref{III.6.2.1}).


En el caso general, introduciremos, en lugar de $\mathscr O_{S'}$, un
complejo $\mathscr Q'_\bullet$ de $\mathscr O_{S'}$-módulos cuasicoherentes
\emph{acotado inferiormente}, de modo que, si para simplificar se toma
$I=\{1,2,\ldots,m\}$, se puede considerar el $\partial$-funtor
\[
  \mathfrak{Tor}_\bullet^S
  (f_1,\ldots,f_m,1_{S'};
   \mathscr P_\bullet^{(1)},\ldots,\mathscr P_\bullet^{(m)},
   \mathscr Q'_\bullet).
\]

\begin{proposition}[6.9.3]
Existen tres funtores espectrales canónicos birregulares denotados por
${}^{(t)}\mathscr E(f_1,\ldots,f_m;
\mathscr P_\bullet^{(1)},\ldots,\mathscr P_\bullet^{(m)},
\mathscr Q'_\bullet)$ (con $t=e$, $f$ o $f'$) cuyo resultado
común es
$\mathfrak{Tor}_\bullet^S(f_1,\ldots,f_m,1_{S'};
\mathscr P_\bullet^{(1)},\ldots,\mathscr P_\bullet^{(m)},
\mathscr Q'_\bullet)$
y cuyos términos $E_2$ son respectivamente

\oldpage[III]{167}
\[
\begin{split}
  {}^{(e)}\mathscr E_{pq}^2
  =\bigoplus_{q'+q''=q}
  \mathscr{T}\!or_p^S\Bigl(
  &\mathfrak{Tor}_{q'}^S
  (f_1,\ldots,f_m;
   \mathscr P_\bullet^{(1)},\ldots,\mathscr P_\bullet^{(m)}),\\
  &\mathscr H_{q''}(\mathscr Q'_\bullet)
  \Bigr)
\end{split}
\]
\[
\begin{split}
  {}^{(f)}\mathscr E_{pq}^2
  =\bigoplus_{q_1+q_2+\ldots+q_{m+1}=q}
  \mathscr{T}\!or_p^{S'}\Bigl(
  &\mathfrak{Tor}_{q_1}^S
  (f_1,1_{S'};\mathscr P_\bullet^{(1)},\mathscr O_{S'}),\ldots,\\
  &\ldots,\mathfrak{Tor}_{q_m}^S
  (f_m,1_{S'};\mathscr P_\bullet^{(m)},\mathscr O_{S'}),
  \mathscr H_{q_{m+1}}(\mathscr Q'_\bullet)
  \Bigr)
\end{split}
\]
\[
\begin{split}
  {}^{(f')}\mathscr E_{pq}^2
  =\bigoplus_{q_1+\ldots+q_m=q}
  \mathfrak{Tor}_p^{S'}\Bigl(
  &f'_1,\ldots,f'_m,1_{S'};\\
  &\mathscr{T}\!or_{q_1}^S
  (\mathscr P_\bullet^{(1)},\mathscr O_{S'}),\ldots,
  \mathscr{T}\!or_{q_m}^S
  (\mathscr P_\bullet^{(m)},\mathscr O_{S'}),
  \mathscr Q'_\bullet
  \Bigr)
\end{split}
\]
\end{proposition}

\begin{proof}
La sucesión $(e)$ no es otra que la sucesión de asociatividad de
\sref{III.6.8.2} para $r=m$. Para definir las otras dos sucesiones espectrales, se va a limitar también al caso en que $S$, $S'$ y los
$Y^{(i)}$ son afines, haciéndose el paso al caso general por los métodos de (6.7.1) y (6.7.2) y dejándose al lector. Demostraremos pues el
\end{proof}

\begin{corollary}[6.9.4]
\label{III.6.9.4}
Sean $A$ un anillo, $A'$ una $A$-álgebra,
$S=\operatorname{Spec}(A)$, $S'=\operatorname{Spec}(A')$,
$X^{(i)}$ ($1\leq i\leq m$) unos $S$-esquemas cuasicompactos y, para cada
$i$, sea $X'^{(i)}=X_{(S')}^{(i)}$, que es un $S'$-esquema cuasicompacto.
Para cada $i$, sea $\mathscr P_\bullet^{(i)}$ un complejo de
$\mathscr O_{X^{(i)}}$-módulos cuasicoherentes; sea finalmente
$Q'_\bullet$ un complejo de $A'$-módulos, estando estos complejos acotados
inferiormente. Existen tres funtores espectrales birregulares en los
$\mathscr P_\bullet^{(i)}$ y en $Q'_\bullet$, que tienen por término límite
común
\[
  \mathbf{Tor}_\bullet^S
  (X^{(1)},\ldots,X^{(m)},S';
   \mathscr P_\bullet^{(1)},\ldots,\mathscr P_\bullet^{(m)},
   \widetilde{Q}'_\bullet)
\]
y cuyos términos $E_2$ son respectivamente
\[
\begin{split}
  {}^{(e)}E_{pq}^2
  =\bigoplus_{q'+q''=q}
  \operatorname{Tor}_p^A\Bigl(
  &\mathbf{Tor}_{q'}^S
  (X^{(1)},\ldots,X^{(m)};
   \mathscr P_\bullet^{(1)},\ldots,\mathscr P_\bullet^{(m)}),\\
  &H_{q''}(Q'_\bullet)
  \Bigr)
\end{split}
\]
\[
\begin{split}
  {}^{(f)}E_{pq}^2
  =\bigoplus_{q_1+q_2+\ldots+q_{m+1}=q}
  \operatorname{Tor}_p^{A'}\Bigl(
  &\mathbf{Tor}_{q_1}^S
  (X^{(1)},S';\mathscr P_\bullet^{(1)},\mathscr O_{S'}),\ldots,\\
  &\ldots,\mathbf{Tor}_{q_m}^S
  (X^{(m)},S';\mathscr P_\bullet^{(m)},\mathscr O_{S'}),
  H_{q_{m+1}}(Q'_\bullet)
  \Bigr)
\end{split}
\]
\[
\begin{split}
  {}^{(f')}E_{pq}^2
  =\bigoplus_{q_1+\ldots+q_m=q}
  \mathbf{Tor}_p^{S'}\Bigl(
  &X'^{(1)},\ldots,X'^{(m)},S';\\
  &\mathscr{T}\!or_{q_1}^S
  (\mathscr P_\bullet^{(1)},\mathscr O_{S'}),\ldots,
  \mathscr{T}\!or_{q_m}^S
  (\mathscr P_\bullet^{(m)},\mathscr O_{S'}),
  \widetilde{Q}'_\bullet
  \Bigr).
\end{split}
\]
\end{corollary}

\begin{proof}
No volveremos sobre el primero de estos funtores espectrales, que fue tratado
en \sref{III.6.8.3} y solo se incluye aquí como recordatorio. Para definir
los otros, consideremos para cada $i$ un recubrimiento abierto afín finito
$\mathfrak U^{(i)}$ de $X^{(i)}$, y, si
$u_i:X'^{(i)}\to X^{(i)}$ es la proyección canónica, el recubrimiento abierto
afín finito correspondiente
$\mathfrak U'^{(i)}=u_i^{-1}(\mathfrak U^{(i)})$.

En virtud de \sref{III.6.6.6},
$\mathbf{Tor}_\bullet^S(X^{(1)},\ldots,X^{(m)},S';
\mathscr P_\bullet^{(1)},\ldots,\mathscr P_\bullet^{(m)},
\widetilde{Q}'_\bullet)$ se obtiene tomando, para $1\leq i\leq m$, una
resolución proyectiva de Cartan-Eilenberg
$M_{\bullet\bullet\bullet}^{(i)}$ de
$L_{\bullet\bullet}^{(i)}=C^\bullet(\mathfrak U^{(i)},
\mathscr P_\bullet^{(i)})$ (en el sentido de \sref[0]{0.11.7.1}),
considerando el tricomplejo
$M_{\bullet\bullet\bullet}
=M_{\bullet\bullet\bullet}^{(1)}\otimes_A
M_{\bullet\bullet\bullet}^{(2)}\otimes\ldots\otimes_A
M_{\bullet\bullet\bullet}^{(m)}\otimes_A Q'_\bullet$
(donde $Q'_\bullet$ se considera como un tricomplejo cuyos dos últimos
grados se reducen a $0$), y tomando su homología. Si se pone
$M_{\bullet\bullet\bullet}'^{(i)}=
M_{\bullet\bullet\bullet}^{(i)}\otimes_A A'$, se tiene (recordando que
$Q'_\bullet$ es un complejo de $A'$-módulos)
$M_{\bullet\bullet\bullet}
=M_{\bullet\bullet\bullet}'^{(1)}\otimes_{A'}
M_{\bullet\bullet\bullet}'^{(2)}\otimes\ldots\otimes_{A'}
M_{\bullet\bullet\bullet}'^{(m)}\otimes_{A'} Q'_\bullet$.
Ahora bien, consideremos cada uno de los complejos
$M_{\bullet\bullet\bullet}'^{(i)}$ como un complejo simple (para su grado
total), y observemos que este complejo está formado por $A'$-módulos
\emph{proyectivos}; de \sref{III.6.3.7} (extendido a un número arbitrario de
complejos) resulta que $H_\bullet(M_{\bullet\bullet\bullet})$ también es
igual a
$\mathbf{Tor}_\bullet^{A'}
(M_{\bullet\bullet\bullet}'^{(1)},\ldots,
M_{\bullet\bullet\bullet}'^{(m)},Q'_\bullet)$;
es, por tanto, por \sref{III.6.5.15}, el término límite de una sucesión
espectral que tiene las propiedades de regularidad requeridas (los tres
grados de $M_{\bullet\bullet\bullet}^{(i)}$ están acotados inferiormente
cuando $\mathscr P_\bullet^{(i)}$ está acotado inferiormente) y cuyo término
$E_2$ viene dado por
\[
  E_{pq}^2
  =\bigoplus_{q_1+\ldots+q_{m+1}=q}
  \operatorname{Tor}_p^{A'}
  (H_{q_1}(M_{\bullet\bullet\bullet}'^{(1)}),\ldots,
   H_{q_m}(M_{\bullet\bullet\bullet}'^{(m)}),H_{q_{m+1}}(Q'_\bullet)).
\]

\oldpage[III]{168}
Se tiene además
$H_{q_i}(M_{\bullet\bullet\bullet}'^{(i)})
=H_{q_i}(M_{\bullet\bullet\bullet}^{(i)}\otimes_A A')
=\mathbf{Tor}_{q_i}^S
(X^{(i)},S';\mathscr P_\bullet^{(i)},\mathscr O_{S'})$
en virtud de la definición de los hipertor globales, lo que da la sucesión
$(f)$ buscada. Por otra parte, se puede considerar
$M_{\bullet\bullet\bullet}'^{(i)}$ como un bicomplejo en el cual el primer
grado es la suma de los grados primero y segundo del tricomplejo
$M_{\bullet\bullet\bullet}^{(i)}$, y el segundo grado es el tercer grado de
este tricomplejo; como los módulos que forman los
$M_{\bullet\bullet\bullet}'^{(i)}$ son $A'$-módulos \emph{proyectivos}, la teoría
general de la hiperhomología muestra que la homología del bicomplejo
$M_{\bullet\bullet\bullet}'^{(1)}\otimes_{A'}
M_{\bullet\bullet\bullet}'^{(2)}\otimes\ldots\otimes_{A'}
M_{\bullet\bullet\bullet}'^{(m)}\otimes_{A'} Q'_\bullet$
es canónicamente isomorfa a su \emph{hiperhomología} \sref[0]{0.11.6.5}; por tanto,
es el término límite de una sucesión espectral cuyo término $E_2$ es igual a
\[
  E_{pq}^2
  =\bigoplus_{q_1+\ldots+q_{m+1}=q}
  \mathbf{Tor}_p^{A'}
  (H_{q_1}^{\mathrm{II}}(M_{\bullet\bullet\bullet}'^{(1)}),\ldots,
   H_{q_m}^{\mathrm{II}}(M_{\bullet\bullet\bullet}'^{(m)}),
   H_{q_{m+1}}^{\mathrm{II}}(Q'_\bullet)).
\]
Ahora bien, como el segundo grado de $Q'_\bullet$ se reduce a $0$, se tiene
$H_n^{\mathrm{II}}(Q'_\bullet)=0$ para $n\ne0$ y
$H_0^{\mathrm{II}}(Q'_\bullet)=Q'_\bullet$; la fórmula precedente también
se escribe
\[
  E_{pq}^2
  =\bigoplus_{q_1+\ldots+q_m=q}
  \mathbf{Tor}_p^{A'}
  (H_{q_1}^{\mathrm{II}}(M_{\bullet\bullet\bullet}'^{(1)}),\ldots,
   H_{q_m}^{\mathrm{II}}(M_{\bullet\bullet\bullet}'^{(m)}),Q'_\bullet).
\]
Además, se tiene
$H_{q_i}^{\mathrm{II}}(M_{\bullet\bullet\bullet}'^{(i)})
=H_{q_i}^{\mathrm{II}}(M_{\bullet\bullet\bullet}^{(i)}\otimes_A A')
=\operatorname{Tor}_{q_i}^A(L_{\bullet\bullet}^{(i)},A')$
en virtud de \sref{III.6.3.4}; pero
$L_{-j,k}^{(i)}=C^j(\mathfrak U^{(i)},\mathscr P_k^{(i)})$, suma directa de
los $\Gamma(V,\mathscr P_k^{(i)})$, donde $V$ recorre las intersecciones
(afines) de $j+1$ conjuntos del recubrimiento $\mathfrak U^{(i)}$; si
$V'=u_i^{-1}(V)$, $V'$ es afín en $X'^{(i)}$, y de
\sref{III.6.4.1.1} resulta que se tiene
\[
  \Gamma\bigl(V',\mathscr{T}\!or_{q_i}^S
  (\mathscr P_k^{(i)},\mathscr O_{S'})\bigr)
  =\operatorname{Tor}_{q_i}^A(\Gamma(V,\mathscr P_k^{(i)}),A')
\]
de donde, para el bicomplejo
$H_{q_i}^{\mathrm{II}}(M_{\bullet\bullet\bullet}'^{(i)})$, se obtiene la
expresión
\[
  C^\bullet\bigl(\mathfrak U'^{(i)},
  \mathscr{T}\!or_{q_i}^S(\mathscr P_\bullet^{(i)},\mathscr O_{S'})\bigr)
\]
lo que da finalmente la expresión buscada para el término $E_2$ de la
sucesión $(f')$. El hecho de que esta sucesión sea birregular bajo las
condiciones indicadas se comprueba como de costumbre, teniendo en cuenta que,
si $\mathscr P_\bullet^{(i)}$ está acotado inferiormente, todos los grados de
$M_{\bullet\bullet\bullet}'^{(i)}$ están acotados inferiormente.
\end{proof}

\begin{remark}[6.9.5]
\label{III.6.9.5}
Se ve, como en \sref{III.6.7.6}, que sustituir los
$\mathscr P_\bullet^{(i)}$ y $\mathscr Q'_\bullet$ por complejos
respectivamente \emph{homotópicos} a ellos no cambia las sucesiones $(e)$, $(f)$ y
$(f')$, salvo isomorfismo canónico. Además, para la sucesión $(f)$, unos
homomorfismos
$\mathscr P_\bullet^{(i)}\to\mathscr R_\bullet^{(i)}$,
$\mathscr Q'_\bullet\to\mathscr T'_\bullet$ de complejos que inducen
isomorfismos en homología
$\mathscr H_\bullet(\mathscr P_\bullet^{(i)})\xrightarrow{\sim}
\mathscr H_\bullet(\mathscr R_\bullet^{(i)})$,
$\mathscr H_\bullet(\mathscr Q'_\bullet)\xrightarrow{\sim}
\mathscr H_\bullet(\mathscr T'_\bullet)$ proporcionan un isomorfismo de
sucesiones espectrales
${}^{(f)}\mathscr E
(f_1,\ldots,f_m;
\mathscr P_\bullet^{(1)},\ldots,\mathscr P_\bullet^{(m)},
\mathscr Q'_\bullet)
\xrightarrow{\sim}
{}^{(f)}\mathscr E
(f_1,\ldots,f_m;
\mathscr R_\bullet^{(1)},\ldots,\mathscr R_\bullet^{(m)},
\mathscr T'_\bullet)$
la demostración es la misma que la de \sref{III.6.7.6}, teniendo en cuenta
el resultado de esa sección y la regularidad de la sucesión $(f)$.
\end{remark}

\begin{corollary}[6.9.6]
\label{III.6.9.6}
Bajo las condiciones de \sref{III.6.9.1}, supongamos que:
\begin{enumerate}
  \item[$1^\circ$] Los complejos $\mathscr P_\bullet^{(i)}$ están formados
  por módulos planos sobre $S$, y los $\mathscr O_Y$-módulos
  \[
    \mathfrak{Tor}_n^S
    (f_1,\ldots,f_m;
     \mathscr P_\bullet^{(1)},\ldots,\mathscr P_\bullet^{(m)})
  \]
  son planos sobre $S$.
  \item[$2^\circ$] Los $\mathscr P_\bullet^{(i)}$ y
  $\mathscr Q'_\bullet$ están acotados inferiormente.
\end{enumerate}

\oldpage[III]{169}
Se tienen entonces, poniendo
$\mathscr P_\bullet'^{(i)}=
\mathscr P_\bullet^{(i)}\otimes_{\mathscr O_S}\mathscr O_{S'}$,
isomorfismos canónicos funtoriales

\label{III.6.9.6.1}
\[
\begin{split}
  &\mathfrak{Tor}_n^{S'}
  (f'_1,\ldots,f'_m,1_{S'};
   \mathscr P_\bullet'^{(1)},\ldots,
   \mathscr P_\bullet'^{(m)},\mathscr Q'_\bullet)
  \xrightarrow{\sim}\\
  &\qquad\xrightarrow{\sim}
  \bigoplus_{n'+n''=n}
  \mathfrak{Tor}_{n'}^S
  (f_1,\ldots,f_m;
   \mathscr P_\bullet^{(1)},\ldots,
   \mathscr P_\bullet^{(m)})
  \otimes_S\mathscr H_{n''}(\mathscr Q'_\bullet).
\end{split}
\tag{6.9.6.1}\label{III.6.9.6.1.fr}
\]

En particular, para $\mathscr Q'_\bullet$ reducido a un solo término
$\mathscr F'$ de grado $0$, se tienen isomorfismos canónicos funtoriales

\label{III.6.9.6.2}
\[
\begin{split}
  &\mathfrak{Tor}_n^{S'}
  (f'_1,\ldots,f'_m,1_{S'};
   \mathscr P_\bullet'^{(1)},\ldots,
   \mathscr P_\bullet'^{(m)},\mathscr F')
  \xrightarrow{\sim}\\
  &\qquad
  \mathfrak{Tor}_n^S
  (f_1,\ldots,f_m;
   \mathscr P_\bullet^{(1)},\ldots,
   \mathscr P_\bullet^{(m)})
  \otimes_{\mathscr O_S}\mathscr F'
\end{split}
\tag{6.9.6.2}\label{III.6.9.6.2.fr}
\]

y, más particularmente, para $\mathscr F'=\mathscr O_{S'}$,

\label{III.6.9.6.3}
\[
\begin{split}
  &\mathfrak{Tor}_n^{S'}
  (f'_1,\ldots,f'_m;
   \mathscr P_\bullet'^{(1)},\ldots,
   \mathscr P_\bullet'^{(m)})
  \xrightarrow{\sim}\\
  &\qquad
  \mathfrak{Tor}_n^S
  (f_1,\ldots,f_m;
   \mathscr P_\bullet^{(1)},\ldots,
   \mathscr P_\bullet^{(m)})
  \otimes_{\mathscr O_S}\mathscr O_{S'}.
\end{split}
\tag{6.9.6.3}\label{III.6.9.6.3.fr}
\]

\end{corollary}

\begin{proof}
La hipótesis de planitud sobre los módulos que componen los
$\mathscr P_\bullet^{(i)}$ implica que los complejos
$\mathscr{T}\!or_{q_i}^S(\mathscr P_\bullet^{(i)},\mathscr O_{S'})$ son nulos
para $q_i\ne0$ por \sref{III.6.5.8}. La sucesión $(f')$ es, por tanto,
degenerada; la hipótesis $2^\circ$ implica además que es birregular
\sref{III.6.9.3}, de modo que el homomorfismo de borde
\[
\begin{split}
  &\mathfrak{Tor}_n^S
  (f_1,\ldots,f_m,1_{S'};
   \mathscr P_\bullet^{(1)},\ldots,
   \mathscr P_\bullet^{(m)},\mathscr Q'_\bullet)
  \longrightarrow {}^{(f')}\mathscr E_{n0}^2=\\
  &\qquad=
  \mathfrak{Tor}_n^{S'}
  (f'_1,\ldots,f'_m,1_{S'};
   \mathscr P_\bullet'^{(1)},\ldots,
   \mathscr P_\bullet'^{(m)},\mathscr Q'_\bullet)
\end{split}
\tag{6.9.6.4}\label{III.6.9.6.4.fr}
\]
es biyectivo por \sref[0]{0.11.1.6}. La hipótesis de planitud sobre los
módulos
\[
  \mathfrak{Tor}_n^S
  (f_1,\ldots,f_m;
   \mathscr P_\bullet^{(1)},\ldots,\mathscr P_\bullet^{(m)})
\]
implica que ${}^{(e)}\mathscr E_{pq}^2=0$ para $p\ne0$ por
\sref{III.6.5.8}. La sucesión $(e)$ es, por tanto, también degenerada y, como
es birregular, el homomorfismo de borde
\[
\begin{split}
  &\mathfrak{Tor}_n^S
  (f_1,\ldots,f_m,1_{S'};
   \mathscr P_\bullet^{(1)},\ldots,
   \mathscr P_\bullet^{(m)},\mathscr Q'_\bullet)
  \longrightarrow {}^{(e)}\mathscr E_{0n}^2=\\
  &\qquad=
  \bigoplus_{n'+n''=n}
  \mathfrak{Tor}_{n'}^S
  (f_1,\ldots,f_m;
   \mathscr P_\bullet^{(1)},\ldots,
   \mathscr P_\bullet^{(m)})
  \otimes_{\mathscr O_S}\mathscr H_{n''}(\mathscr Q'_\bullet)
\end{split}
\tag{6.9.6.5}\label{III.6.9.6.5.fr}
\]
es biyectivo; combinando los dos isomorfismos precedentes se obtiene el
isomorfismo \sref{III.6.9.6.1}. El isomorfismo \sref{III.6.9.6.2} se deduce
de este trivialmente, puesto que entonces se tiene
$\mathscr H_q(\mathscr Q'_\bullet)=0$ si $q\ne0$ y
$\mathscr H_0(\mathscr Q'_\bullet)=\mathscr F'$. Finalmente, el caso
$\mathscr F'=\mathscr O_{S'}$ en \sref{III.6.9.6.2}, teniendo en cuenta
\sref{III.6.7.12}, da el isomorfismo \sref{III.6.9.6.3}.
\end{proof}

\begin{corollary}[6.9.7]
\label{III.6.9.7}
Bajo las condiciones de \sref{III.6.9.1}, supongamos que $S$ y $S'$ afines, y supongamos dado para cada $i$ se ha dado un entero $d_i$
($1\leq i\leq m$). Existe entonces un entero $N$ que solo depende de $S$, de
los $X^{(i)}$ y de los $d_i$, y que tiene la propiedad siguiente: para todo
entero $n_0$, se tienen los isomorfismos canónicos \sref{III.6.9.6.3} para
$n\leq n_0$ y para todo sistema de complejos
$\mathscr P_\bullet^{(i)}$ que satisfaga las condiciones siguientes:

  $1^\circ$ $\mathscr P_k^{(i)}=0$ para $k<d_i$;
  $2^\circ$ $\mathscr P_k^{(i)}$ es plano sobre $S$ para
  $k<n_0+N$;
  $3^\circ$
  $\mathfrak{Tor}_q^S(f_1,\ldots,f_m;
\mathscr P_\bullet^{(1)},\ldots,\mathscr P_\bullet^{(m)})$ es plano sobre
  $S$ para $q<n_0+N$.

\end{corollary}

\begin{proof}
Supongamos que $\mathscr P_k^{(i)}$ es plano sobre $S$ para $k<r$; entonces
$\mathscr{T}\!or_{q_i}^S(\mathscr P_k^{(i)},\mathscr O_{S'})=0$ para $k<r$
y $q_i\ne0$; calculemos
\[
\begin{split}
  \mathfrak{Tor}_p^{S'}\Bigl(
  &f'_1,\ldots,f'_m;
  \mathscr{T}\!or_{q_1}^S(\mathscr P_\bullet^{(1)},\mathscr O_{S'}),
  \ldots,\\
  &\mathscr{T}\!or_{q_m}^S(\mathscr P_\bullet^{(m)},\mathscr O_{S'})
  \Bigr)
\end{split}
\tag{6.9.7.1}\label{III.6.9.7.1.fr}
\]

\oldpage[III]{170}
por el método de \sref{III.6.6.2}, con ayuda de la imagen inversa de un
recubrimiento afín fijo $\mathfrak U^{(i)}$ de $X^{(i)}$
($1\leq i\leq m$) (independiente de $S'$ y de los
$\mathscr P_\bullet^{(i)}$); los términos $L_{jk}^{(i)}$ son nulos para
$j<\setmin N_i$ (que solo depende de $\mathfrak U^{(i)}$); si alguno de los $q_i$
no es nulo, el complejo simple cuya homología de grado $p$ es
\eqref{III.6.9.7.1.eq} tiene términos nulos en todos los grados
$<r+\sum_{j\ne i}d_j\setmin\sum_{i=1}^m N_i$,
por lo que \eqref{III.6.9.7.1.eq} es nulo para $p<r\setmin N$, designando por $N$
el mayor de los números
$\sum_{i=1}^m N_i\setmin\sum_{j\ne i}d_j$.
Se concluye que ${}^{(f')}\mathscr E_{pq}^2=0$ para $q\ne0$ y $p<r\setmin N$;
como, por otra parte, ${}^{(f')}\mathscr E_{pq}^2=0$ para $q<0$, se ve que el
homomorfismo de borde \eqref{III.6.9.6.4.eq} es biyectivo para $n<r\setmin N$
\emph{(M, XV, 5.6)} (para $\mathscr Q'_\bullet=\mathscr O_{S'}$). En segundo
lugar, si
$\mathfrak{Tor}_q^S
(f_1,\ldots,f_m;
\mathscr P_\bullet^{(1)},\ldots,\mathscr P_\bullet^{(m)})$
es plano sobre $S$ para $q<r$, se tiene
${}^{(e)}\mathscr E_{pq}^2=0$ para $p\ne0$ y $q<r$; por otra parte,
${}^{(e)}\mathscr E_{pq}^2=0$ para $p<0$, por lo que el homomorfismo de borde
\eqref{III.6.9.6.5.eq} es biyectivo para $n<r$, lo que termina la
demostración.
\end{proof}

El caso de \sref{III.6.9.3} más importante en las aplicaciones es aquel en
que $m=1$, estando $\mathscr Q'_\bullet$ reducido a un solo término
$\mathscr F'$ de grado $0$; volveremos a enunciarlo en este caso con vistas
a referencias ulteriores\footnote{El caso tratado en \sref{III.6.9.8}, y en
particular las sucesiones espectrales \sref{III.6.9.8.3} y
\sref{III.6.9.8.4}, nos había sido señalado en 1957 por J.-P.~Serre.}:

\begin{proposition}[6.9.8]
\label{III.6.9.8}
Sean $S$ un preesquema, $g:S'\to S$ un morfismo,
$f:X\to Y$ un $S$-morfismo separado y cuasicompacto de $S$-preesquemas,
$\mathscr P_\bullet$ un complejo de $\mathscr O_X$-módulos cuasicoherentes
acotado inferiormente, y $\mathscr F'$ un $\mathscr O_{S'}$-módulo
cuasicoherente. Existen dos funtores espectrales birregulares en
$\mathscr P_\bullet$ y $\mathscr F'$, con valores en la categoría de los
$\mathscr O_{Y_{(S')}}$-módulos cuasicoherentes, que tienen el mismo término
límite $\mathfrak{Tor}_\bullet^S(f,1_{S'};\mathscr P_\bullet,\mathscr F')$, y cuyos términos $E_2$ son

\label{III.6.9.8.1}
\[
  {}'\mathscr E_{pq}^2
  =\mathscr{T}\!or_p^S
  (\mathscr H^{-q}(f,\mathscr P_\bullet),\mathscr F')
  \tag{6.9.8.1}\label{III.6.9.8.1.fr}
\]

y

\label{III.6.9.8.2}
\[
  {}''\mathscr E_{pq}^2
  =\mathscr H^{-p}
  (f',\mathscr{T}\!or_q^S(\mathscr P_\bullet,\mathscr F')),
  \tag{6.9.8.2}\label{III.6.9.8.2.fr}
\]

donde $f'=f_{(S')}:X_{(S')}\to Y_{(S')}$.
\end{proposition}

\begin{proof}
Las sucesiones en cuestión también pueden obtenerse, no partiendo de
\sref{III.6.9.3}, sino de las sucesiones $(a)$ y $(b')$ de
\sref{III.6.7.3} para $X^{(1)}=X$, $Y^{(1)}=Y$,
$X^{(2)}=Y^{(2)}=S'$, $f_1=f$, $f_2=1_{S'}$. Cuando $S=S'=Y$, siendo $Y$
afín, se obtienen dos sucesiones espectrales cuyos términos $E_2$ son iguales
a
\[
  {}'\mathscr E_{pq}^2
  =\mathscr{T}\!or_p^Y
  (\mathscr H^{-q}(f,\mathscr P_\bullet),\mathscr F)
  \tag{6.9.8.3}\label{III.6.9.8.3.fr}
\]
\[
  {}''\mathscr E_{pq}^2
  =\mathscr H^{-p}
  (f,\mathscr{T}\!or_q^Y(\mathscr P_\bullet,\mathscr F))
  \tag{6.9.8.4}\label{III.6.9.8.4.fr}
\]
que convergen (en virtud de \sref{III.6.7.6}) a la hipercohomología
$\mathscr H^\bullet(f,\mathscr P_\bullet\otimes_Y\mathscr F)$ del funtor
$f_*$ con respecto al complejo
$\mathscr P_\bullet\otimes_Y\mathscr F$ de $\mathscr O_X$-módulos, para todo
$\mathscr O_Y$-módulo cuasicoherente y \emph{plano} sobre $Y$, $\mathscr F$ (o para
todo $\mathscr O_Y$-módulo cuasicoherente $\mathscr F$ cuando
$\mathscr P_\bullet$ está formado por $\mathscr O_X$-módulos \emph{planos} sobre
$Y$), las cuales son distintas de las de \sref{III.6.2.1}.
\end{proof}

\begin{corollary}[6.9.9]
\label{III.6.9.9}
Bajo las condiciones de \sref{III.6.9.8}, supongamos que el complejo
$\mathscr P_\bullet$ está acotado inferiormente, formado por módulos planos
sobre $S$, y que los $\mathscr O_Y$-módulos
$\mathscr H^n(f,\mathscr P_\bullet)$ son planos sobre $S$.

\oldpage[III]{171}
Se tienen entonces isomorfismos canónicos funtoriales

\label{III.6.9.9.1}
\[
  \mathfrak{Tor}_n^{S'}
  (f',1_{S'};\mathscr P'_\bullet,\mathscr F')
  \xrightarrow{\sim}
  \mathscr H^{-n}(f,\mathscr P_\bullet)
  \otimes_{\mathscr O_S}\mathscr F'
  \tag{6.9.9.1}\label{III.6.9.9.1.fr}
\]

donde
$\mathscr P'_\bullet=\mathscr P_\bullet\otimes_{\mathscr O_S}
\mathscr O_{S'}$; en particular, para $\mathscr F'=\mathscr O_{S'}$, se
tienen isomorfismos canónicos funtoriales

\label{III.6.9.9.2}
\[
  \mathscr H^n(f',\mathscr P'_\bullet)
  \xrightarrow{\sim}
  \mathscr H^n(f,\mathscr P_\bullet)
  \otimes_{\mathscr O_S}\mathscr O_{S'}.
  \tag{6.9.9.2}\label{III.6.9.9.2.fr}
\]

\textup{Este es el caso particular $m=1$ de \sref{III.6.9.6}. Más particularmente:}
\end{corollary}

\begin{corollary}[6.9.10]
\label{III.6.9.10}
Sean $S$ un preesquema, $f:X\to Y$ un $S$-morfismo separado y cuasicompacto
de $S$-preesquemas, y $\mathscr P_\bullet$ un complejo acotado
inferiormente, formado por $\mathscr O_X$-módulos cuasicoherentes y planos
sobre $S$. Se supone además que los $\mathscr O_Y$-módulos
$\mathscr H^n(f,\mathscr P_\bullet)$ son planos sobre $S$. Para todo $s\in S$,
denotemos por $X_s$ y $Y_s$ las fibras
$X\otimes_S \boldsymbol{k}(s)$, $Y\otimes_S \boldsymbol{k}(s)$,
por $f_s:X_s\to Y_s$ el morfismo $f\times_S1$, y por
$\mathscr P_\bullet^s$ el complejo
$\mathscr P_\bullet\otimes_{\mathscr O_S}\boldsymbol{k}(s)$ de
$\mathscr O_{X_s}$-módulos. Se tienen entonces isomorfismos canónicos
funtoriales

\label{III.6.9.10.1}
\[
  \mathscr H^n(f_s,\mathscr P_\bullet^s)
  \xrightarrow{\sim}
  \mathscr H^n(f,\mathscr P_\bullet)\otimes_{\mathscr O_S}\boldsymbol{k}(s).
  \tag{6.9.10.1}\label{III.6.9.10.1.fr}
\]

\end{corollary}

Se tiene así, bajo hipótesis de planitud convenientes, un caso en el que la
formación de los funtores derivados $\mathrm R^nf_*(\mathscr F)$ «conmuta con el paso
a las fibras», que volveremos a encontrar por otro método en el §~7.

\subsection{Estructura local de ciertos funtores cohomológicos}
\label{subsection:III.6.10}
\label{III.6.10}

\begin{proposition}[6.10.1]
Sean $S=\operatorname{Spec}(A)$ un esquema afín,
$Y^{(i)}$ ($1\leq i\leq n$) una familia finita de $S$-esquemas afines, \textup{planos}
sobre $S$; para cada $i$, sea $f_i:X^{(i)}\to Y^{(i)}$ un $S$-morfismo
separado y cuasicompacto, y sea $\mathscr P_\bullet^{(i)}$ un complejo de
$\mathscr O_{X^{(i)}}$-Módulos cuasicoherentes acotado inferiormente. Sea
$Y$ el producto de los $S$-esquemas $Y^{(i)}$. Existe un complejo
$\mathscr R_\bullet$ de $\mathscr O_Y$-Módulos cuasicoherentes y \textup{planos} sobre
$S$, con la propiedad siguiente: para todo $S$-esquema afín $S'$ y todo
complejo de $\mathscr O_{S'}$-Módulos cuasicoherentes
$\mathscr Q'_\bullet$ acotado inferiormente, hay un isomorfismo
\[
  \mathfrak{Tor}_\bullet^S
  (f_1,\ldots,f_n,1_{S'};
   \mathscr P_\bullet^{(1)},\ldots,\mathscr P_\bullet^{(n)},
   \mathscr Q'_\bullet)
  \xrightarrow{\sim}
  \mathscr H_\bullet
  (\mathscr R_\bullet\otimes_{\mathscr O_S}\mathscr Q'_\bullet)
  \tag{6.10.1.1}\label{III.6.10.1.1.fr}
\]
que es un isomorfismo de $\partial$-funtores en $\mathscr Q'_\bullet$.
Además, para todo $S$-morfismo $u:S''\to S'$ de $S$-esquemas afines, el
diagrama
\[
\xymatrix@C=12pt@R=30pt{
  {\begin{gathered}
    \mathfrak{Tor}_\bullet^S
    (f_1,\ldots,f_n,1_{S'};\\
    \mathscr P_\bullet^{(1)},\ldots,\mathscr P_\bullet^{(n)},
    \mathscr Q'_\bullet)
  \end{gathered}}
  \ar[r]^-{\sim}\ar[d]
  &
  {\mathscr H_\bullet
    (\mathscr R_\bullet\otimes_{\mathscr O_S}\mathscr Q'_\bullet)}
  \ar[d]
  \\
  {\begin{gathered}
    \mathfrak{Tor}_\bullet^S
    (f_1,\ldots,f_n,1_{S''};\\
    \mathscr P_\bullet^{(1)},\ldots,\mathscr P_\bullet^{(n)},
    u^*(\mathscr Q'_\bullet))
  \end{gathered}}
  \ar[r]^-{\sim}
  &
  {\mathscr H_\bullet
    (\mathscr R_\bullet\otimes_{\mathscr O_S}u^*(\mathscr Q'_\bullet))}
}
\tag{6.10.1.2}\label{III.6.10.1.2.fr}
\]
(donde las flechas verticales son los $(1_Y\times_Su)$-morfismos canónicos
definidos en \textup{(6.7.10)}) es conmutativo.
\end{proposition}

\begin{proof}
\oldpage[III]{172}
Calculemos los hipertor por el método de (6.6.2), teniendo en cuenta la observación
(6.6.6); con las notaciones de (6.6.2), cada
$L_{\bullet\bullet}^{(i)}$ es un bicomplejo de $A_i$-módulos, designando por
$A_i$ el anillo de $Y^{(i)}$; admite pues una resolución de Cartan-Eilenberg
proyectiva $M_{\bullet\bullet\bullet}^{(i)}$ (en el sentido de (\textbf{0}, 11.7.1)) formada por
$A_i$-\emph{módulos}, y en virtud de (6.6.6), el primer miembro de (6.10.1.1) es
canónicamente isomorfo a
$\mathrm H_\bullet(M_{\bullet\bullet\bullet}\otimes_A Q'_\bullet)$, donde
$M_{\bullet\bullet\bullet}=M_{\bullet\bullet\bullet}^{(1)}\otimes_A M_{\bullet\bullet\bullet}^{(2)}\otimes\ldots\otimes_A M_{\bullet\bullet\bullet}^{(n)}$
y $\mathscr Q'_\bullet=\widetilde{Q}'_\bullet$. Como, por hipótesis, los
anillos $A_i$ son $A$-módulos \emph{planos}, los
$M_{\bullet\bullet\bullet}^{(i)}$ son tricomplejos de $A$-módulos planos
($\mathbf{0}_{\mathrm I}$, 6.2.1), y lo mismo ocurre con
$M_{\bullet\bullet\bullet}$; además, si $B$ es el anillo de $Y$, producto
tensorial de los $A_i$, $M_{\bullet\bullet\bullet}$ es un tricomplejo de
$B$-\emph{módulos}; el complejo
$\mathscr R_\bullet=(M_{\bullet\bullet\bullet})^\sim$ de $\mathscr O_Y$-Módulos (donde $M_{\bullet\bullet\bullet}$ se considera como
complejo simple) responde pues a la cuestión, como resulta fácilmente de
(6.7.10).
\end{proof}

\begin{corollary}[6.10.2]
\label{III.6.10.2}
En la Proposición \sref{III.6.10.1}, el complejo
$\mathscr R_\bullet$ puede tomarse acotado inferiormente. Si los
$\mathscr P_\bullet^{(i)}$ están acotados superiormente y los
$Y^{(i)}$ tienen dimensión cohomológica finita, entonces
$\mathscr R_\bullet$ puede tomarse acotado superiormente.
\end{corollary}

\begin{proof}
La primera afirmación se deduce de que los tres grados de cada
$M_{\bullet\bullet\bullet}^{(i)}$ están acotados inferiormente.
Por otra parte, si los anillos $A_i$ tienen dimensión
cohomológica finita, el tercer grado de cada
$M_{\bullet\bullet\bullet}^{(i)}$ solo toma un número finito de
valores, y por su construcción lo mismo ocurre con su primer grado
\sref{III.6.6.2}. Su segundo grado está acotado superiormente
exactamente cuando lo está el grado de $\mathscr P_\bullet^{(i)}$
\sref{III.6.6.2}, lo que demuestra la segunda afirmación.
\end{proof}

\begin{remarks}[6.10.3]
\textnormal{(i)} Con las notaciones de (6.10.1),
$\mathscr H_\bullet(\mathscr R_\bullet\otimes_S\mathscr Q'_\bullet)$ es
isomorfo a
$\mathscr{T}\!or_\bullet(\mathscr R_\bullet,\mathscr Q'_\bullet)$ puesto que
$\mathscr R_\bullet$ consta de $\mathscr O_Y$-módulos planos sobre
$S$ \sref{III.6.5.9}. Es, por tanto, según \sref{III.6.5.4}, el
objeto al que converge una sucesión espectral regular cuyo término
$E_2$ es

\label{III.6.10.3.1}
\[
  \mathscr E_{pq}^2
  =
  \bigoplus_{q'+q''=q}
  \mathscr{T}\!or_p^S
  \bigl(
    \mathscr H_{q'}(\mathscr R_\bullet),
    \mathscr H_{q''}(\mathscr Q'_\bullet)
  \bigr)
  \tag{6.10.3.1}\label{III.6.10.3.1.fr}
\]
que no es otra que la sucesión espectral $(e)$ del cambio de base (6.9.3).

\medskip
\textnormal{(ii)} Sea $\mathscr R'_\bullet$ otro complejo de $\mathscr O_Y$-módulos
cuasicoherentes, planos sobre $S$, y sea
$g:\mathscr R'_\bullet\to\mathscr R_\bullet$
un homomorfismo de complejos para el cual
$\mathscr H_\bullet(g):\mathscr H_\bullet(\mathscr R'_\bullet)\longrightarrow\mathscr H_\bullet(\mathscr R_\bullet)$ es un isomorfismo. Entonces, en virtud de (6.3.3) y (6.5.9), se deduce de $g$ induce un isomorfismo de
$\partial$-funtores en $\mathscr Q'_\bullet$,
$\mathscr H_\bullet(\mathscr R'_\bullet\otimes_S\mathscr Q'_\bullet)\xrightarrow{\sim}\mathscr H_\bullet(\mathscr R_\bullet\otimes_S\mathscr Q'_\bullet)$
tal que el diagrama
\[
\xymatrix@C=28pt@R=30pt{
  \mathscr H_\bullet
  (\mathscr R'_\bullet\otimes_S\mathscr Q'_\bullet)
  \ar[r]^-{\sim}\ar[d]
  &
  \mathscr H_\bullet
  (\mathscr R_\bullet\otimes_S\mathscr Q'_\bullet)
  \ar[d]
  \\
  \mathscr H_\bullet
  (\mathscr R'_\bullet\otimes_Su^*(\mathscr Q'_\bullet))
  \ar[r]^-{\sim}
  &
  \mathscr H_\bullet
  (\mathscr R_\bullet\otimes_Su^*(\mathscr Q'_\bullet))
}
\]
es conmutativo. Así, las propiedades de \sref{III.6.10.1} no
determinan por completo el complejo $\mathscr R_\bullet$ no está enteramente determinado por las propiedades de (6.10.1).

\medskip
\oldpage[III]{173}
\textnormal{(iii)} En la demostración de (6.10.1), se pueden suponer los $M_{\bullet\bullet\bullet}^{(i)}$ pueden elegirse formados por
$A_i$-módulos \emph{libres}, como se deduce inmediatamente dualizando la
demostración de \sref[0]{0.11.5.2.1}. Entonces los
$M_{\bullet\bullet\bullet}^{(i)}\otimes_{A_i}B$
constan de $B$-módulos libres; como
$M_{\bullet\bullet\bullet}$ es su producto tensorial sobre $B$,
se puede exigir además en \sref{III.6.10.1} que
$\mathscr R_\bullet$ esté asociado a un complejo de $B$-módulos
\emph{libres}.

Además, por (M, XVII, 1.2), el tricomplejo
$M_{\bullet\bullet\bullet}$ depende \emph{funtorialmente} de cada uno de
los bicomplejos $L_{\bullet\bullet}^{(i)}$ y, por tanto, de cada
$\mathscr P_\bullet^{(i)}$, una vez fijado un recubrimiento finito de
cada $X^{(i)}$. Aquí se entienden por «morfismos» de tricomplejos
las \emph{clases de homomorfismos} módulo homotopía. Sustituir un
recubrimiento de $X^{(i)}$ por otro más fino proporciona, para
$L_{\bullet\bullet}^{(i)}$ a homomorfismos definidos precisamente salvo homotopía (6.6.8), se ve finalmente que, con la convención precedente para los morfismos, el tricomplejo $M_{\bullet\bullet\bullet}$ es un \emph{funtor} en
cada $\mathscr P_\bullet^{(i)}$.

Precisaremos esta dependencia funtorial, en particular el
comportamiento de $\mathscr R_\bullet$ respecto de sucesiones exactas
cortas
$0\longrightarrow\mathscr P_\bullet^{\prime(i)}\longrightarrow\mathscr P_\bullet^{(i)}\longrightarrow\mathscr P_\bullet^{\prime\prime(i)}\longrightarrow0$
en el capítulo dedicado a un álgebra general de funtores
cohomológicos mencionado en \sref{III.6.1.3}.
\end{remarks}

\begin{scholium}[6.10.4]
\label{III.6.10.4}
El hecho de que $\mathscr R_\bullet$ conste de $\mathscr O_Y$-módulos
\emph{planos} sobre $S$ implica inmediatamente que
$\mathscr H_\bullet(\mathscr R_\bullet\otimes_S\mathscr Q'_\bullet)$
es un funtor \emph{homológico} en $\mathscr Q'_\bullet$ (véase el razonamiento de (7.7.1)). Es esta propiedad la que, como se ha mencionado en (6.1.1), motiva la introducción del hipertor. Pongamos en efecto \[
  X=X^{(1)}\times_SX^{(2)}\times_S\ldots\times_SX^{(n)},
  \qquad
  f=f_1\times_Sf_2\times_S\ldots\times_Sf_n,
\]
\[
  \mathscr P_\bullet
  =
  \mathscr P_\bullet^{(1)}
  \otimes_S\mathscr P_\bullet^{(2)}
  \otimes\ldots\otimes_S\mathscr P_\bullet^{(n)},
\]
y
\[
  X'=X\times_SS',
  \qquad
  Y'=Y\times_SS',
  \qquad
  f'=f\times_S1_{S'};
\]
Los problemas de cambio de base llevan a estudiar la
hipercohomología
$\mathscr H^\bullet_{f'}(\mathscr P_\bullet\otimes_S\mathscr N')$
como funtor del $\mathscr O_{S'}$-módulo cuasicoherente $\mathscr N'$, o,
equivalentemente, la hipercohomología
$\mathscr H^\bullet_f(\mathscr P_\bullet\otimes_S\mathscr N)$
como funtor del $\mathscr O_S$-módulo cuasicoherente $\mathscr N$. Cuando
los $\mathscr P_\bullet^{(i)}$, y por tanto $\mathscr P_\bullet$, son
\emph{planos} sobre $S$, los resultados precedentes y
\sref{III.6.7.6} muestran que este es efectivamente un funtor
\emph{cohomológico} en $\mathscr N$. Esto deja de ser cierto cuando no se
supone la planitud de los $\mathscr P_\bullet^{(i)}$, y ya no se puede entonces abordar el estudio de $\mathscr H^\bullet_f(\mathscr P_\bullet\otimes_S\mathscr N)$
ya no puede abordarse mediante los métodos usuales del álgebra
homológica.

Utilizaremos principalmente el caso $n=1$, $Y=S$, en el que
$\mathscr P_\bullet$ consta de $\mathscr O_X$-módulos \emph{planos} sobre $Y$.
En este caso se tiene el teorema siguiente.
\end{scholium}

\begin{theorem}[6.10.5]
Sean $Y=\operatorname{Spec}(A)$ un esquema afín noetheriano,
$f:X\to Y$ un morfismo propio, $\mathscr P_\bullet$ un complejo de
$\mathscr O_X$-Módulos coherentes, \textup{planos} sobre $Y$, acotado inferiormente.
Existe entonces un complejo $\mathscr L_\bullet$ de $\mathscr O_Y$-Módulos,
acotado inferiormente, cuyos términos $\mathscr L_i$ son
$\mathscr O_Y$-Módulos de la forma $\mathscr O_Y^{m_i}$, y un isomorfismo
\[
  \mathscr H^\bullet
  (f,\mathscr P_\bullet\otimes_Y\mathscr Q_\bullet)
  \xrightarrow{\sim}
  \mathscr H_\bullet
  (\mathscr L_\bullet\otimes_Y\mathscr Q_\bullet)
  \tag{6.10.5.1}\label{III.6.10.5.1.fr}
\]
de $\partial$-funtores en el complejo $\mathscr Q_\bullet$ de
$\mathscr O_Y$-Módulos cuasicoherentes, acotado inferiormente. Además, para
todo morfismo $u:Y'\to Y$, se tiene, poniendo
\[
  X'=X_{(Y')},
  \qquad
  f'=f_{(Y')},
  \qquad
  \mathscr P'_\bullet
  =
  \mathscr P_\bullet\otimes_Y\mathscr O_{Y'},
  \qquad
  \mathscr L'_\bullet=u^*(\mathscr L_\bullet)
\]
\oldpage[III]{174}
(que es un complejo de $\mathscr O_{Y'}$-Módulos localmente libres de tipo
finito), un isomorfismo
\[
  \mathscr H^\bullet
  (f',\mathscr P'_\bullet\otimes_{Y'}\mathscr Q'_\bullet)
  \xrightarrow{\sim}
  \mathscr H_\bullet
  (\mathscr L'_\bullet\otimes_{Y'}\mathscr Q'_\bullet)
  \tag{6.10.5.2}\label{III.6.10.5.2.fr}
\]
de $\partial$-funtores en el complejo $\mathscr Q'_\bullet$ de
$\mathscr O_{Y'}$-Módulos cuasicoherentes, acotado inferiormente, de modo que
el diagrama
\[
\xymatrix@C=25pt@R=30pt{
  \mathscr H^\bullet
  (f,\mathscr P_\bullet\otimes_Y\mathscr Q_\bullet)
  \ar[r]^-{\sim}\ar[d]
  &
  \mathscr H_\bullet
  (\mathscr L_\bullet\otimes_Y\mathscr Q_\bullet)
  \ar[d]
  \\
  \mathscr H^\bullet
  (f',\mathscr P'_\bullet\otimes_{Y'}u^*(\mathscr Q_\bullet))
  \ar[r]^-{\sim}
  &
  \mathscr H_\bullet
  (\mathscr L'_\bullet\otimes_{Y'}u^*(\mathscr Q_\bullet))
}
\tag{6.10.5.3}\label{III.6.10.5.3.fr}
\]
es conmutativo.
\end{theorem}

\begin{proof}
La aplicación de (6.10.1) da primero un complejo acotado inferiormente
(6.10.2) $\mathscr R_\bullet$ de $\mathscr O_Y$-Módulos cuasicoherentes y
$Y$-libres (6.10.3, (iii)) y un isomorfismo
\[
  \mathscr H^\bullet
  (f,\mathscr P_\bullet\otimes_Y\mathscr Q_\bullet)
  \xrightarrow{\sim}
  \mathscr H_\bullet
  (\mathscr R_\bullet\otimes_Y\mathscr Q_\bullet)
  \tag{6.10.5.4}\label{III.6.10.5.4.fr}
\]
de $\partial$-funtores en $\mathscr Q_\bullet$, pero \emph{a priori} los
términos de $\mathscr R_\bullet$ no son necesariamente
$\mathscr O_Y$-Módulos \emph{de tipo finito}. Pero si se aplica (6.10.5.4) al caso
en que $\mathscr Q_\bullet$ es un complejo reducido a un solo término
$\mathscr O_Y$, se ve que $\mathscr H_\bullet(\mathscr R_\bullet)$ es
isomorfo a $\mathscr H^\bullet(f,\mathscr P_\bullet)$, y está por tanto formado
por $\mathscr O_Y$-Módulos \emph{coherentes} (6.2.5). Se sabe entonces (\textbf{0}, 11.9.2) que
existe un complejo $\mathscr L_\bullet$ acotado inferiormente, formado por
$\mathscr O_Y$-Módulos asociados a $A$-módulos libres \emph{de tipo finito}, y un
homomorfismo $\mathscr L_\bullet\longrightarrow\mathscr R_\bullet$,
tales que el homomorfismo correspondiente para la homología,
$\mathscr H_\bullet(\mathscr L_\bullet)\longrightarrow\mathscr H_\bullet(\mathscr R_\bullet)$,
es biyectivo; de donde el isomorfismo (6.10.5.1), en virtud de (6.10.3, (ii)).
Las demás afirmaciones de (6.10.5) resultan de (6.10.1) y (6.10.3, (ii))
cuando $Y'$ es \emph{afín}; en el caso general, basta verificar que
cuando se considera un recubrimiento $(V_\alpha)$ de $Y'$ por abiertos
afines, y el isomorfismo correspondiente (6.10.5.2) relativo a cada uno de los
$V_\alpha$, las restricciones a un abierto afín
$W\subset V_\alpha\cap V_\beta$ de los isomorfismos correspondientes a
$V_\alpha$ y a $V_\beta$ coinciden con el isomorfismo correspondiente a
$W$, lo que resulta de la conmutatividad del diagrama (6.10.1.2) aplicado a las
inclusiones canónicas $W\to V_\alpha$ y $W\to V_\beta$.
\end{proof}

\begin{remark}[6.10.6]
\label{III.6.10.6}
En los capítulos siguientes aplicaremos \sref{III.6.10.5}
principalmente cuando $\mathscr P_\bullet$ está concentrado en un grado
y su único término es un $\mathscr O_X$-módulo coherente $\mathscr F$ plano
sobre $Y$. Como se tiene entonces
$\mathscr H_n(\mathscr L_\bullet)=\mathscr H^{-n}(f,\mathscr F)=\mathrm R^{-n}f_*(\mathscr F)$ (6.2.1),
se ve que los $\mathscr H_n(\mathscr L_\bullet)$ se anulan para $n>0$; veremos más adelante (7.7.12, (i)) que se puede entonces suponer que $\mathscr L_\bullet$ puede
elegirse \emph{con términos solo en grados} $\leqslant0$ (luego en
número \emph{finito}), a condición de reemplazar la hipótesis de que los $\mathscr L_i$
están asociados a $A$-módulos libres de tipo finito por la de que los
$\mathscr L_i$ son \emph{localmente libres} de tipo finito.

El complejo $\mathscr L_\bullet$ correspondiente a tal
$\mathscr O_X$-módulo $\mathscr F$ no parece tener ninguna propiedad
especial
\oldpage[III]{175}
aparte de la restricción precedente sobre sus grados. Cabe, pues,
preguntarse recíprocamente si, dado un complejo acotado
inferiormente $\mathscr L_\bullet$ cuyos términos son
$\mathscr O_Y$-módulos asociados a $A$-módulos proyectivos de tipo
finito y cuyos términos de grado $>0$ se anulan, existen un
$Y$-esquema $X$, proyectivo y plano sobre $Y$, y un
$\mathscr O_X$-módulo localmente libre $\mathscr F$ tales que exista un
isomorfismo
$\mathscr H^\bullet(f,\mathscr F\otimes_Y\mathscr Q_\bullet)\xrightarrow{\sim}\mathscr H_\bullet(\mathscr L_\bullet\otimes_Y\mathscr Q_\bullet)$
funtorial en $\mathscr Q_\bullet$. Un resultado de este tipo reduciría
por completo la teoría cohomológica de los módulos coherentes
planos sobre $Y$ en $Y$-esquemas propios a la teoría, «salvo
homotopía», de los complejos de $A$-módulos proyectivos de tipo
finito sobre un anillo noetheriano $A$.
\end{remark}

\section{Estudio del cambio de base en los funtores homológicos covariantes de módulos}
\label{section:III.7}
\label{section:III.7.fr}

\subsection{Funtores de $A$-módulos}
\label{subsection:III.7.1}
\label{subsection:III.7.1.fr}

\begin{env}[7.1.1]
\label{III.7.1.1}
Dado un anillo $A$, no necesariamente conmutativo, denotemos por
$\boldsymbol{Ab}_A$ la categoría de los $A$-\emph{módulos izquierdos}, y denotaremos simplemente por $\boldsymbol{Ab}$ la categoría de los $\mathbf Z$-módulos,
que es la categoría de los grupos abelianos. Sea
$T:\boldsymbol{Ab}_A\longrightarrow\boldsymbol{Ab}$
un \emph{funtor covariante aditivo}, y sea $M$ un $(A,A)$-\emph{bimódulo}.
Entonces $T(M)$ está dotado de manera natural de una estructura de
$A$-\emph{módulo a derecha}.

En efecto, para cada $a\in A$, denotemos por $h_{a,M}$, o
simplemente por $h_a$, el endomorfismo
$x\to xa$
del $A$-módulo \emph{a izquierda} $M$. Por hipótesis, $T(h_a)$ es un
endomorfismo del $\mathbf Z$-\emph{módulo} $T(M)$. Como $T$ es un funtor covariante aditivo, se tiene, para $a\in A$, $b\in A$,
\[
  T(h_{ab})
  =
  T(h_b\circ h_a)
  =
  T(h_b)\circ T(h_a)
  \qquad\text{y}\qquad
  T(h_{a+b})
  =
  T(h_a+h_b)
  =
  T(h_a)+T(h_b);
\]
esto prueba que la aplicación $(a,y)\to T(h_a)(y)$
es una ley externa de $A$-módulo a derecha sobre
$T(M)$. En particular, $T(A_s)$ es un $A$-\emph{módulo a derecha}.
\end{env}

\begin{env}[7.1.2]
Cuando $A$ es \emph{conmutativo}. De \sref{III.7.1.1} se deduce
que, para todo $A$-módulo $M$, el grupo $T(M)$ es naturalmente un
$A$-módulo. Además, si $u:M\to N$ es un homomorfismo
$A$-lineal, entonces, para todo $a\in A$,
$u\circ h_{a,M}=h_{a,N}\circ u$,
y, por consiguiente,
$T(u)\circ T(h_{a,M})=T(h_{a,N})\circ T(u)$,
Así, $T(u):T(M)\to T(N)$ es $A$-\emph{lineal}, y $T$ puede considerarse
como un funtor covariante aditivo de $\boldsymbol{Ab}_A$ en sí misma.

Más precisamente, esto define una \emph{equivalencia} canónica entre la
categoría de los funtores covariantes aditivos
$\boldsymbol{Ab}_A\to\boldsymbol{Ab}$
y la categoría de los funtores covariantes $A$-\emph{lineales}
$T:\boldsymbol{Ab}_A\longrightarrow\boldsymbol{Ab}_A$,
es decir, tales que
$T(h_{a,M})=h_{a,T(M)}$ para todo $a\in A$.
Como el funtor de inclusión $I:\boldsymbol{Ab}_A\longrightarrow\boldsymbol{Ab}$,
que asocia a un $A$-módulo su $\mathbf Z$-módulo subyacente, es exacto
y fiel. Por tanto, los dos funtores correspondientes bajo la
equivalencia precedente tienen las mismas propiedades de exactitud.
\end{env}

\begin{env}[7.1.3]
Suponiendo siempre conmutativo el anillo $A$, sea $B$ una $A$-\emph{álgebra} (no
necesariamente conmutativa), y sea $\rho:A\to B$ el homomorfismo de anillos
correspondiente a esta estructura de álgebra; este homomorfismo define un
funtor covariante
\oldpage[III]{176}
aditivo $\rho_*:M\leadsto M_{[\rho]}$
de la categoría $\boldsymbol{Ab}_B$ de los $B$-módulos izquierdos en la categoría
$\boldsymbol{Ab}_A$ de los $A$-módulos. Por composición, se deduce un funtor $T_{(B)}:\boldsymbol{Ab}_B\xrightarrow{\rho_*}\boldsymbol{Ab}_A\xrightarrow{T}\boldsymbol{Ab}$,
evidentemente covariante y aditivo, que denotaremos también por $T^{(B)}$ (por
razones tipográficas) o $T\otimes_A B$, y que diremos obtenido a partir
de $T$ por \emph{extensión de escalares} de $A$ a $B$. Por supuesto, si $B$ es
conmutativa, se puede considerar $T_{(B)}$ como un funtor de
$\boldsymbol{Ab}_B$ en sí misma (7.1.2). Cuando $B$ es conmutativa, y $C$
una $B$-álgebra, se ve inmediatamente que $T_{(C)}=(T_{(B)})_{(C)}$;
es inmediato que la extensión de escalares es \emph{funtorial} y \emph{aditiva} en
$T$; además, cuando $T$ conmuta con los límites inductivos o con las sumas
directas (resp. es exacto por la izquierda, exacto por la derecha, exacto), lo mismo ocurre
con $T_{(B)}$: en efecto, $\rho_*$ es exacto y conmuta con los límites inductivos
y las sumas directas.
\end{env}

\begin{env}[7.1.4]
\label{III.7.1.4}
Supongamos de nuevo que $A$ conmutativo, y sea
$T$ un funtor covariante aditivo
$A$-lineal $\boldsymbol{Ab}_A\to\boldsymbol{Ab}_A$, \emph{que conmuta con los límites inductivos}. Entonces, para toda parte
multiplicativa $S$ de $A$ y
todo $A$-módulo $M$, existe un isomorfismo canónico funtorial de
$A$-módulos

\label{III.7.1.4.1}
\[
  T(S^{-1}M)\xrightarrow{\sim}S^{-1}T(M).
  \tag{7.1.4.1}\label{III.7.1.4.1.fr}
\]


Supongamos primero que $S$ es el conjunto de las potencias $f^n$,
$(n\geqslant0)$, de un elemento $f\in A$. Entonces
$M_f=\varinjlim M_n$,
donde $(M_n,\varphi_{nm})$ es el sistema inductivo de $A$-módulos
con $M_n=M$,
con $\varphi_{nm}:z\leadsto f^{\,n-m}z$ ($\mathbf{0}_{\mathrm I}$, 1.6.1);
de donde, en este caso, el isomorfismo (7.1.4.1) en virtud de la hipótesis sobre $T$. Para $S$ arbitraria, el módulo $S^{-1}M$ es
el límite inductivo de los $M_f$, para $f\in S$ ($\mathbf{0}_{\mathrm I}$, 1.4.5) y se concluye del mismo modo. Además,
la funtorialidad del isomorfismo (7.1.4.1) muestra que es un isomorfismo
de $S^{-1}A$-\emph{módulos}. Así, salvo
isomorfismo canónico, se puede escribir

\label{III.7.1.4.2}
\[
  T_{(S^{-1}A)}(S^{-1}M)
  =
  S^{-1}T(M)
  =
  T(S^{-1}M).
  \tag{7.1.4.2}\label{III.7.1.4.2.fr}
\]

Si $S=A\setmin\mathfrak p$ es el complementario de un ideal primo
$\mathfrak p$ de $A$, se escribe $T_{\mathfrak p}$ en lugar de
$T_{(A_{\mathfrak p})}$.
\end{env}

\begin{proposition}[7.1.5]
Bajo las hipótesis de (7.1.4), si $T_{\mathfrak m}$ es exacto por la izquierda, resp.
exacto por la derecha o exacto, para todo ideal maximal
$\mathfrak m$ de $A$. Entonces $T$ es exacto por la izquierda,
resp. exacto por la derecha o exacto.
\end{proposition}

\begin{proof}
Si dos submódulos $N$ y $P$ de un $A$-módulo $M$ satisfacen
$N_{\mathfrak m}=P_{\mathfrak m}$
para todo ideal maximal $\mathfrak m$ de $A$, entonces $N=P$
(Bourbaki, \emph{Algèbre commutative}, cap.~II, \S~3,
n.º~3, teo.~1).
\end{proof}

\subsection{Caracterizaciones del funtor producto tensorial}
\label{subsection:III.7.2}
\label{subsection:III.7.2.fr}

\begin{env}[7.2.1]
\label{III.7.2.1}
Sea $A$ un anillo, no necesariamente conmutativo; sea $M$
(resp. $N$) un $A$-módulo a izquierda
(resp. a derecha), y sea $P$ un $\mathbf Z$-módulo.
Recordemos que dar un $\mathbf Z$-homomorfismo
$v:N\otimes_A M\longrightarrow P$
equivale a dar una aplicación $\mathbf Z$-bilineal
$u:N\times M\longrightarrow P$
tal que $u(ta,x)=u(t,ax)$ para $a\in A$, $t\in N$, $x\in M$,
estando relacionadas las dos aplicaciones por $v(t\otimes x)=u(t,x)$.
Por otra parte, dar $u$ equivale a dar un
\oldpage[III]{177}
$\mathbf Z$-\emph{homomorfismo}
$x\to f_x$
de $M$ en $\operatorname{Hom}_{\mathbf Z}(N,P)$ tal que
$f_{ax}(t)=f_x(ta)$ para $a\in A$, $t\in N$, $x\in M$,
estando relacionadas las dos aplicaciones por $u(t,x)=f_x(t)$.
\end{env}

\begin{env}[7.2.2]
Sea $T:\boldsymbol{Ab}_A\to\boldsymbol{Ab}$ un funtor covariante aditivo. Vamos
a definir, para todo $A$-módulo por la izquierda $M$, un \emph{homomorfismo canónico
funtorial en} $M$, \emph{de} $\mathbf Z$-\emph{módulos}
\[
  t_M:T(A_s)\otimes_A M\longrightarrow T(M).
  \tag{7.2.2.1}\label{III.7.2.2.1.fr}
\]

Para ello bastará, en virtud de (7.2.1), definir un
$\mathbf Z$-homomorfismo $x\leadsto t'_M(x)$ de $M$ en
$\operatorname{Hom}_{\mathbf Z}(T(A_s),T(M))$, tal que se tenga
$t'_M(ax)(y)=t'_M(x)(ya)$ para $a\in A$, $x\in M$ y $y\in T(A_s)$.
Observemos para ello que
$\operatorname{Hom}_{\mathbf Z}(T(A_s),T(M))$ está canónicamente provisto de una
estructura de $A$-módulo por la izquierda procedente de la estructura de $A$-módulo por la
derecha de $T(A_s)$, siendo la ley externa tal que si
$a\in A$, $v\in\operatorname{Hom}_{\mathbf Z}(T(A_s),T(M))$,
$(a.v)(y)=v(ya)$ para $y\in T(A_s)$.
Dicho esto, definiremos $t'_M$ como compuesto de los dos homomorfismos
canónicos
\[
  M\xrightarrow{\sim}\operatorname{Hom}_A(A_s,M)
  \xrightarrow{T}
  \operatorname{Hom}_{\mathbf Z}(T(A_s),T(M))
\]
siendo la segunda flecha la aplicación $u\leadsto T(u)$, la primera
el isomorfismo canónico de $A$-módulos $x\leadsto\theta_x$ tal que
$\theta_x(\xi)=\xi x$ para $\xi\in A$, $x\in M$. Se tiene
$\theta_{ax}=\theta_x\circ h_a$, luego
$T(\theta_{ax})=T(\theta_x\circ h_a)=T(\theta_x)\circ T(h_a)$
y por tanto, para $y\in T(A_s)$,
\[
  T(\theta_{ax})(y)
  =
  T(\theta_x)(T(h_a)(y))
  =
  T(\theta_x)(ya)
\]
por definición de la ley externa sobre $T(A_s)$, lo que prueba la existencia de
$t_M$; es inmediato verificar que este homomorfismo es funtorial en
$M$, es decir, que para todo homomorfismo $w:M\to M'$ de $A$-módulos por la
izquierda, el diagrama
\[
\xymatrix@C=30pt@R=30pt{
  T(A_s)\otimes_A M
  \ar[r]^-{t_M}\ar[d]_{1\otimes w}
  &
  T(M)\ar[d]^{T(w)}
  \\
  T(A_s)\otimes_A M'
  \ar[r]_-{t_{M'}}
  &
  T(M')
}
\tag{7.2.2.2}\label{III.7.2.2.2.fr}
\]
es conmutativo.

La funtorialidad del homomorfismo (7.2.2.1) muestra que cuando $A$ es
\emph{conmutativo}, es un homomorfismo de $A$-\emph{módulos} (cf. (7.1.2)).
\end{env}

\begin{env}[7.2.3]
\label{III.7.2.3}
Cuando $A$ es conmutativo, se puede definir más generalmente un
homomorfismo canónico de $A$-módulos

\label{III.7.2.3.1}
\[
  T(N)\otimes_A M\longrightarrow T(N\otimes_A M)
  \tag{7.2.3.1}\label{III.7.2.3.1.fr}
\]

para todo $A$-módulo $N$; basta, en la construcción de (7.2.2), reemplazar el homomorfismo $\theta_x$ por el homomorfismo de
$A$-módulos
$N\longrightarrow N\otimes_A M$
que a todo $y\in N$ hace corresponder $y\otimes x$. Es inmediato que este
homomorfismo es \emph{funtorial en} $M$ \emph{como
en} $N$.

\oldpage[III]{178}
En particular, si $B$ es una $A$-álgebra, no necesariamente
conmutativa, existe un homomorfismo, funtorial en $M$,

\label{III.7.2.3.2}
\[
  (T(M))_{(B)}
  =
  T(M)\otimes_A B
  \longrightarrow
  T(M\otimes_A B)
  =
  T_{(B)}(M_{(B)})
  \tag{7.2.3.2}\label{III.7.2.3.2.fr}
\]

Por la funtorialidad de \eqref{III.7.2.3.1.eq} en $M$, este es un
homomorfismo de $B$-\emph{módulos}.

Se tiene además un diagrama conmutativo

\label{III.7.2.3.3}
\[
\xymatrix@C=30pt@R=30pt{
  T(A)\otimes_A M
  \ar[r]^-{t_M}\ar[d]
  &
  T(M)\ar[d]
  \\
  T_{(B)}(B_s)\otimes_B M_{(B)}
  \ar[r]_-{t_{M_{(B)}}}
  &
  T_{(B)}(M_{(B)})
}
\tag{7.2.3.3}\label{III.7.2.3.3.fr}
\]

La flecha vertical derecha es el compuesto
\[
  T(M)
  \longrightarrow
  T(M)\otimes_A B
  \longrightarrow
  T(M\otimes_A B)
  =
  T_{(B)}(M_{(B)})
\]
del homomorfismo canónico y \eqref{III.7.2.3.2.eq}. La flecha
vertical izquierda de \eqref{III.7.2.3.3.eq} es el homomorfismo
$T(A)\otimes_A M\longrightarrow T_{(B)}(B_s)\otimes_B(B\otimes_A M)=T_{(B)}(B_s)\otimes_A M$
donde
$T(A)\longrightarrow T_{(B)}(B_s)=T(B)$
es $T(\rho)$, considerando $\rho$ como homomorfismo de
$A$-módulos $A\to B_{[\rho]}$.
\end{env}

\begin{lemma}[7.2.4]
\label{III.7.2.4}
Si $T$ es un funtor covariante aditivo de $\boldsymbol{Ab}_A$ en
$\boldsymbol{Ab}$ que conmuta con las sumas directas, entonces el
homomorfismo canónico $t_L$ de \eqref{III.7.2.2.1.eq} es un
isomorfismo para todo $A$-módulo libre $L$.
\end{lemma}

\begin{proof}
En efecto, se tiene $L=\bigoplus_{\alpha\in I}L_\alpha$
donde cada $L_\alpha$ es isomorfo a $A_s$ para todo $\alpha\in I$; la definición de
$t_M$
en \sref{III.7.2.2} muestra que
$t_L=\bigoplus_{\alpha\in I}t_{L_\alpha}$,
porque
\[
  T:
  \operatorname{Hom}_A(A_s,L)
  \longrightarrow
  \operatorname{Hom}_{\mathbf Z}(T(A_s),T(L))
\]
es la suma directa de las aplicaciones $\mathbf Z$-lineales
$T_\alpha:\operatorname{Hom}_A(A_s,L_\alpha)\longrightarrow\operatorname{Hom}_{\mathbf Z}(T(A_s),T(L_\alpha))$
por la hipótesis sobre $T$. Solo queda, por tanto, demostrar el
lema para $L=A_s$; pero entonces $t_L$ es precisamente el
isomorfismo canónico
$T(A_s)\otimes_A A_s\xrightarrow{\sim}T(A_s)$
válido para todo $A$-módulo a derecha.
\end{proof}

\begin{proposition}[7.2.5]
Sea $T$ un funtor covariante aditivo de $\boldsymbol{Ab}_A$ en
$\boldsymbol{Ab}$, que conmuta con las sumas directas. Las condiciones siguientes son
equivalentes:
\begin{enumerate}
  \item[a)] $T$ es exacto por la derecha.
  \item[b)] El homomorfismo canónico $t_M$ \textup{(7.2.2.1)} es un isomorfismo para
    todo $A$-módulo por la izquierda $M$.
  \item[$b'$)] $T$ es semiexacto y el homomorfismo $t_M$ es sobreyectivo para
    todo $A$-módulo por la izquierda $M$.
  \item[c)] $T$ es isomorfo a un funtor en $M$ de la forma
    $N\otimes_A M$, donde $N$ es un $A$-módulo por la derecha.
\end{enumerate}
\end{proposition}

\begin{proof}
Está claro que $b)$ implica $c)$ y que $c)$ implica $a)$; mostremos que
$a)$ implica $b)$.
\oldpage[III]{179}
Pongamos
$T'(M)=T(A_s)\otimes_A M$
para todo $A$-módulo por la izquierda $M$. Existe una sucesión exacta
$L'\longrightarrow L\longrightarrow M\longrightarrow0$,
donde $L$ y $L'$ son dos $A$-módulos por la izquierda \emph{libres}; como $T$ y $T'$ son
exactos por la derecha, se tiene pues el diagrama conmutativo
\[
\xymatrix@C=22pt@R=25pt{
  T'(L')\ar[r]\ar[d]_{t_{L'}}
  &
  T'(L)\ar[r]\ar[d]_{t_L}
  &
  T'(M)\ar[r]\ar[d]_{t_M}
  &
  0
  \\
  T(L')\ar[r]
  &
  T(L)\ar[r]
  &
  T(M)\ar[r]
  &
  0
}
\]
donde las dos filas son exactas; como $t_L$ y $t_{L'}$ son
isomorfismos en virtud de (7.2.4), lo mismo ocurre con $t_M$ por el lema de los
cinco. Finalmente, está claro que $b)$ implica $b')$. Para mostrar que $b')$
implica $a)$, basta probar el
\end{proof}

\begin{lemma}[7.2.5.1]
Sean $\boldsymbol{K}$, $\boldsymbol{K}'$ dos categorías abelianas, $F$, $G$ dos
funtores covariantes aditivos de $\boldsymbol{K}$ en $\boldsymbol{K}'$,
$f:F\to G$ un morfismo funtorial \textup{(T, I, 1.2)} tal que, para todo objeto $E$
de la categoría $\boldsymbol{K}$,
$f_E:F(E)\longrightarrow G(E)$
es un epimorfismo. Entonces, si $F$ es exacto por la derecha y $G$ semiexacto,
$G$ es exacto por la derecha.
\end{lemma}

\begin{proof}
Todo se reduce en efecto a demostrar que para todo epimorfismo
$v:E'\to E$ en $\boldsymbol{K}$,
$G(v):G(E')\longrightarrow G(E)$
es un epimorfismo; ahora bien, se tiene el diagrama conmutativo
\[
\xymatrix@C=35pt@R=30pt{
  F(E')\ar[r]^{F(v)}\ar[d]_{f_{E'}}
  &
  F(E)\ar[d]^{f_E}
  \\
  G(E')\ar[r]_{G(v)}
  &
  G(E)
}
\]
en el cual $F(v)$, $f_{E'}$ y $f_E$ son epimorfismos; lo mismo ocurre pues
con $G(v)$.
\end{proof}

\begin{remark}[7.2.6]
\label{III.7.2.6}
Para todo $A$-módulo a derecha $N$, pongamos
$T_N(M)=N\otimes_A M$
para todo $A$-módulo a izquierda $M$. Así, $T_N$ es un funtor
covariante aditivo de $\boldsymbol{Ab}_A$ en $\boldsymbol{Ab}$, es exacto por
la derecha y conmuta con las sumas directas. Si se identifica
canónicamente $T_N(A_s)$ con $N$, se comprueba inmediatamente que
el homomorfismo correspondiente \eqref{III.7.2.2.1.eq} es la
identidad. Se sigue que el $A$-módulo a derecha $N$ del enunciado de
(7.2.5, $c)$) está determinado salvo isomorfismo único y es
canónicamente isomorfo a $T(A_s)$. Se puede decir también que los morfismos
funtoriales
$T\leadsto T(A_s)$ y $N\leadsto T_N$
constituyen una \emph{equivalencia} (T, I, 1.2) de la categoría de los $A$-módulos a derecha y la categoría de los funtores
covariantes aditivos
$\boldsymbol{Ab}_A\to\boldsymbol{Ab}$
que son exactos por la derecha y conmutan con las sumas directas.
\end{remark}

\begin{proposition}[7.2.7]
\label{III.7.2.7}
Sea $A$ un anillo artiniano a izquierda cuyo cociente por su
radical $\mathfrak m$ es un cuerpo $k$, y sea $T$ un funtor
covariante aditivo de $\boldsymbol{Ab}_A$ en $\boldsymbol{Ab}$ que conmuta con
las sumas directas. Las condiciones de \sref{III.7.2.5} son
también equivalentes a la condición siguiente:
\begin{enumerate}
  \item[(d)] $T$ es semiexacto y el homomorfismo
    $T(\varepsilon):T(A_s)\longrightarrow T(k)$
    inducido por el homomorfismo canónico
    $\varepsilon:A_s\to k$ es sobreyectivo.
\end{enumerate}
\end{proposition}

\begin{proof}
\oldpage[III]{180}
Es claro que la condición ($b')$ de (7.2.5) implica $d)$; probemos que
$d)$ implica $b')$). Existe un entero $n$ tal que
$\mathfrak m^n=0$; pongamos, para todo $A$-módulo $M$, $M_h=\mathfrak m^hM$;
Demostraremos, por inducción descendente sobre $h$, que $t_{M_h}$
es sobreyectivo. La afirmación es clara para $h=n$. Si $h<n$,
existe una sucesión exacta
\[
  0\longrightarrow M_{h+1}
  \longrightarrow M_h
  \longrightarrow M_h/M_{h+1}
  \longrightarrow 0
\]
y la hipótesis de inducción dice que $t_{M_{h+1}}$ es
sobreyectivo. Además, $M_h/M_{h+1}$ es anulado por $\mathfrak m$
y, por tanto, es un $(A/\mathfrak m)$-módulo; en otras palabras,
es una suma directa de $A$-módulos isomorfos a $k$. Para probar que $t_{M_h/M_{h+1}}$ es sobreyectivo basta demostrar que $t_k$ lo es, puesto que $T$ conmuta con las sumas directas. Ahora bien,
en virtud de la conmutatividad del diagrama
\[
\xymatrix@C=35pt@R=30pt{
  T(A_s)\otimes_A A_s
  \ar[r]^-{t_{A_s}}\ar[d]_{1\otimes\varepsilon}
  &
  T(A_s)\ar[d]^{T(\varepsilon)}
  \\
  T(A_s)\otimes_A k
  \ar[r]_-{t_k}
  &
  T(k)
}
\]
y de (7.2.4), la hipótesis $d)$ implica que $t_k$ es sobreyectivo. Para terminar la demostración, bastará probar que si se tiene una sucesión exacta $0\longrightarrow M'\longrightarrow M\xrightarrow{v}M''\longrightarrow 0$
es una sucesión exacta de $A$-módulos para la que $t_{M'}$ y
$t_{M''}$ son sobreyectivos, entonces $t_M$ es sobreyectivo.
Existe un diagrama conmutativo
\[
\xymatrix@C=27pt@R=30pt{
  T'(M')\ar[r]\ar[d]_{t_{M'}}
  &
  T'(M)\ar[r]\ar[d]_{t_M}
  &
  T'(M'')\ar[r]\ar[d]_{t_{M''}}
  &
  0\ar[d]
  \\
  T(M')\ar[r]
  &
  T(M)\ar[r]_-{T(v)}
  &
  T(M'')\ar[r]
  &
  \operatorname{Coker}(T(v))
}
\]
cuyas dos filas son exactas, por la hipótesis de que $T$ es
semiexacto. La hipótesis de inducción dice que $t_{M'}$ y
$t_{M''}$ son epimorfismos, mientras que la última flecha vertical
es un monomorfismo. El lema de los cinco \emph{(M, I, 1.1)}
muestra entonces que $t_M$ es un epimorfismo.
\end{proof}

\subsection{Criterios de exactitud de los funtores homológicos de módulos}
\label{subsection:III.7.3}
\label{subsection:III.7.3.fr}

\begin{proposition}[7.3.1]
\label{III.7.3.1}
Sea $A$ un anillo, no necesariamente conmutativo, y sea
$T_\bullet$ un funtor homológico covariante
\emph{(T, II, 2.1)} de $\boldsymbol{Ab}_A$ en $\boldsymbol{Ab}$ que
conmuta con las sumas directas. Sea $p$ un entero tal que
$T_p$ y $T_{p-1}$ estén definidos. Las condiciones siguientes son
equivalentes:
\begin{enumerate}
  \item[(a)] $T_p$ es exacto por la derecha;
  \item[(b)] $T_{p-1}$ es exacto por la izquierda;

\oldpage[III]{181}

  \item[(c)] para todo $A$-módulo a izquierda $M$, el homomorfismo
    canónico y funtorial \eqref{III.7.2.2.1.eq}
    \[
      T_p(A_s)\otimes_A M\longrightarrow T_p(M)
      \tag{7.3.1.1}\label{III.7.3.1.1.fr}
    \]
    es un isomorfismo;
  \item[(d)] para todo $A$-módulo a izquierda $M$, el homomorfismo
    \eqref{III.7.3.1.1.eq} es un epimorfismo;
  \item[(e)] $T_p$ es isomorfo a un funtor de la forma
    $M\leadsto N\otimes_A M$,
    donde $N$ es un $A$-módulo a derecha.
\end{enumerate}
Si, además, se verifican las condiciones de \sref{III.7.2.7} sobre
$A$ y $\mathfrak m$, las condiciones precedentes son también
equivalentes a
\begin{enumerate}
  \item[(f)] el homomorfismo canónico
    $T_p(\varepsilon):T_p(A_s)\longrightarrow T_p(k)$
    es un epimorfismo.
\end{enumerate}
\end{proposition}

\begin{proof}
Por definición de funtor homológico, $T_i$ es semiexacto para todo
$i$ tales que $T_i$ está definido, y que, para toda sucesión exacta
$0\longrightarrow M'\xrightarrow{u}M\xrightarrow{v}M''\longrightarrow0$,
se tiene
$\operatorname{Ker}(T_{i-1}(u))=\operatorname{Coker}(T_i(v))$,
está claro que $a)$ y $b)$ son equivalentes y las demás afirmaciones
resultan trivialmente de (7.2.5) y (7.2.7).
\end{proof}

\begin{corollary}[7.3.2]
\label{III.7.3.2}
Sea $A$ un anillo conmutativo. Con las notaciones de
\sref{III.7.3.1}, supongamos que $T_p$ es exacto por la derecha. Si
$f\in A$ no pertenece al anulador de ningún elemento
$\ne0$ de un $A$-módulo $M$, entonces $f$ no pertenece al anulador de ningún
elemento $\ne0$ de $T_{p-1}(M)$. En particular, si $A$ es un
dominio entero, el $A$-módulo $T_{p-1}(A)$ es sin torsión.
\end{corollary}

\begin{proof}
Si $h_f$ designa la homotecia $x\to fx$ de $M$, la hipótesis
significa que $h_f$ es inyectiva. La condición (b) de
\sref{III.7.3.1} muestra entonces que $T_{p-1}(h_f)$ por
la condición $b)$ de (7.3.1).
\end{proof}

\begin{proposition}[7.3.3]
\label{III.7.3.3}
Sea $A$ un anillo, y sea $T_\bullet$ un funtor homológico
covariante de $\boldsymbol{Ab}_A$ en $\boldsymbol{Ab}$ que conmuta con las
sumas directas. Sea $p$ un entero tal que $T_{p-1}$, $T_p$ y
$T_{p+1}$ estén definidos. Las condiciones siguientes son
equivalentes:
\begin{enumerate}
  \item[(a)] $T_p$ es exacto;
  \item[(b)] $T_{p+1}$ y $T_p$ son exactos por la derecha;
  \item[(c)] $T_p$ y $T_{p-1}$ son exactos por la izquierda;
  \item[(d)] $T_{p+1}$ es exacto por la derecha y $T_{p-1}$ es
    exacto por la izquierda;
  \item[(e)] para todo $A$-módulo $M$, los homomorfismos canónicos
    \[
      T_i(A_s)\otimes_A M\longrightarrow T_i(M)
      \tag{7.3.3.1}\label{III.7.3.3.1.fr}
    \]
    son isomorfismos para $i=p$ e $i=p+1$;
  \item[($e'$)] para todo $A$-módulo $M$, los
    homomorfismos \eqref{III.7.3.3.1.eq} son epimorfismos para
    $i=p$ e $i=p+1$;
  \item[(f)] para todo $A$-módulo $M$, el homomorfismo
    \eqref{III.7.3.3.1.eq} es un isomorfismo para $i=p$, y
    $T_p(A_s)$ es plano como $A$-módulo a derecha;
  \item[($f'$)] para todo $A$-módulo $M$, el
    homomorfismo \eqref{III.7.3.3.1.eq} es un epimorfismo para
    $i=p$, y $T_p(A_s)$ es plano como $A$-módulo por la derecha plano.
\end{enumerate}
\end{proposition}

\begin{proof}
La equivalencia de las condiciones $a)$, $b)$, $c)$, $d)$ resulta de la equivalencia
de las condiciones $a)$ y $b)$ de (7.3.1). La equivalencia de $b)$, $e)$ y $e')$
resulta de la equivalencia de $a)$, $c)$ y $d)$ en (7.3.1). Finalmente, decir que
$T_p(A_s)$
es plano significa que el funtor
$M\leadsto T_p(A_s)\otimes_A M$
es exacto por la izquierda; la equivalencia de $a)$, $f)$ y $f')$ resulta también de
la equivalencia de $a)$, $c)$, $d)$ en (7.3.1).
\end{proof}

\oldpage[III]{182}

\begin{corollary}[7.3.4]
\label{III.7.3.4}
Supongamos que $A$ es conmutativo, que $T_p$ exacto y supongamos además que $T_p(A)$ es un $A$-módulo de presentación finita. Entonces la
función
$x\leadsto\operatorname{rang}_{\mathbf{k}(x)}(T_p(\boldsymbol{k}(x)))$
es localmente constante en $X=\operatorname{Spec}(A)$ y, por tanto, es constante
si $\operatorname{Spec}(A)$ es conexo.
\end{corollary}

\begin{proof}
Por la condición (f) de \sref{III.7.3.3}, $T_p(A)$ es un
$A$-módulo plano en virtud de (7.3.3, $f)$), es
proyectivo de tipo finito, y $\bigl(T_p(A)\bigr)^{\sim}$ es pues un
$\mathscr O_X$-módulo localmente
libre (Bourbaki, \emph{Alg\`ebre commutative}, chap.~II, \S~5,
no.~2, th.~1). Además,
$T_p(\boldsymbol{k}(x))=T_p(A)\otimes_A\boldsymbol{k}(x)$
(7.3.3, $e)$) y se sabe que el rango en el punto $x$ del $A$-módulo $T_p(A)$ es localmente constante
\emph{(loc.\ cit.)}. Esto prueba el corolario.
\end{proof}

\begin{proposition}[7.3.5]
Supongamos que $A$ es un anillo artiniano por la izquierda cuyo cociente por su
radical $\mathfrak m$ es un cuerpo $k$. Entonces las condiciones de \textup{(7.3.3)} son
también equivalentes a cada una de las siguientes:
\begin{enumerate}
  \item[g)] El homomorfismo canónico
    $T_i(\varepsilon):T_i(A_s)\longrightarrow T_i(k)$
    es un epimorfismo para $i=p$ y $i=p+1$.
  \item[h)] $T_p(\varepsilon)$ es un epimorfismo y $T_p(A_s)$ es un
    $A$-módulo por la derecha plano \textup{(o, lo que equivale a lo mismo (Bourbaki,
    \emph{Alg. comm.}, cap. II, \S 3, n\textsuperscript{o} 2, cor. 2 de la
    prop. 5), un $A$-módulo} libre).
\end{enumerate}
Supongamos además que $A$ es conmutativo y el $A$-módulo $T_p(k)$ de
longitud finita $d$. Entonces las condiciones anteriores equivalen también a
cada una de las siguientes:
\begin{enumerate}
  \item[i)] Para todo $A$-módulo $M$ de longitud finita, se tiene
    \[
      \operatorname{long}(T_p(M))=d\mathbin{.}\operatorname{long}(M).
      \tag{7.3.5.1}\label{III.7.3.5.1.fr}
    \]
  \item[j)] Se tiene
    \[
      \operatorname{long}(T_p(A))=d\mathbin{.}\operatorname{long}(A).
      \tag{7.3.5.2}\label{III.7.3.5.2.fr}
    \]
\end{enumerate}
\end{proposition}

La equivalencia de $g)$ y $h)$ con las condiciones de (7.3.3) se deduce inmediatamente
de (7.2.7). Para demostrar las demás afirmaciones, utilizaremos el lema
siguiente:

\begin{lemma}[7.3.5.3]
Sean $\boldsymbol{K}$, $\boldsymbol{K}'$ dos categorías abelianas,
$F:\boldsymbol{K}\to\boldsymbol{K}'$ un funtor
covariante aditivo; se supone que $F$ es semiexacto, y que, para todo objeto
simple $S$ de $\boldsymbol{K}$, $F(S)$ es un objeto de longitud finita en
$\boldsymbol{K}'$. Entonces, para todo objeto $E$ de longitud finita en $\boldsymbol{K}$,
$F(E)$ es de longitud finita en $\boldsymbol{K}'$. Para toda sucesión exacta
$0\longrightarrow E'\xrightarrow{u}E\xrightarrow{v}E''\longrightarrow0$
de objetos de longitud finita en $\boldsymbol{K}$, se tiene
\[
  \operatorname{long}F(E)
  \leq
  \operatorname{long}F(E')+\operatorname{long}F(E'')
  \tag{7.3.5.4}\label{III.7.3.5.4.fr}
\]
y para que los dos miembros de \textup{(7.3.5.4)} sean iguales, es necesario y suficiente
que la sucesión
\[
  0\longrightarrow F(E')\longrightarrow F(E)\longrightarrow F(E'')
  \longrightarrow0
\]
sea exacta.
\end{lemma}

\begin{proof}
En efecto, la sucesión
$F(E')\xrightarrow{F(u)}F(E)\xrightarrow{F(v)}F(E'')$
es exacta por hipótesis; si se supone $F(E')$ y $F(E'')$ de longitud
finita, lo mismo ocurre con $\operatorname{Im}(F(u))$ y con
$\operatorname{Im}(F(v))$ y como
$\operatorname{Ker}(F(v))=\operatorname{Im}(F(u))$, $F(E)$ es de longitud
finita y se tiene
\[
  \operatorname{long}F(E)
  =
  \operatorname{long}\operatorname{Im}(F(u))
  +
  \operatorname{long}\operatorname{Im}(F(v))
  \leq
  \operatorname{long}F(E')+\operatorname{long}F(E'').
  \tag{7.3.5.5}\label{III.7.3.5.5.fr}
\]

\oldpage[III]{183}
Por inducción sobre la longitud de $E$, esto prueba ya la primera afirmación;
además, los dos miembros de (7.3.5.5) solo pueden ser iguales si
$\operatorname{long}\operatorname{Im}(F(u))=\operatorname{long}F(E')$
(lo que equivale a
$\operatorname{long}\operatorname{Ker}(F(u))=0$, o
$\operatorname{Ker}(F(u))=0$) y
$\operatorname{long}\operatorname{Im}(F(v))=\operatorname{long}F(E'')$
(lo que equivale a
$\operatorname{long}\operatorname{Coker}(F(v))=0$, o
$\operatorname{Coker}(F(v))=0$).
\end{proof}

Observemos ahora que si $M$ es un $A$-módulo de longitud finita (siendo $A$
conmutativo), los cocientes de una sucesión de Jordan-Hölder de $M$ son
necesariamente isomorfos al $A$-módulo $k$, luego se deduce de (7.3.5.4), por
inducción sobre la longitud de $M$,
\[
  \operatorname{long}T_p(M)\leq d\mathbin{.}\operatorname{long}(M).
  \tag{7.3.5.6}\label{III.7.3.5.6.fr}
\]

Además, resulta de (7.3.5.3) que si $T_p$ es exacto, se tiene la igualdad
(7.3.5.1); luego la condición $a)$ de (7.3.3) implica $i)$; está claro que
$i)$ implica $j)$, y queda por probar el

\begin{lemma}[7.3.5.7]
La relación
$\operatorname{long}T_p(A)=d\mathbin{.}\operatorname{long}A$
implica que $T_p(\varepsilon)$ es un epimorfismo y $T_p(A)$ es un
$A$-módulo plano.
\end{lemma}

\begin{proof}
En efecto, partiendo de la sucesión exacta
$0\longrightarrow\mathfrak m\longrightarrow A\longrightarrow k\longrightarrow0$,
resulta de (7.3.5.4) y (7.3.5.6) que se tiene
\[
  \operatorname{long}T_p(A)
  \leq
  \operatorname{long}T_p(\mathfrak m)+\operatorname{long}T_p(k)
  \leq
  d\bigl(\operatorname{long}\mathfrak m+\operatorname{long}k\bigr)
  =d\mathbin{.}\operatorname{long}A
\]
y que la igualdad solo puede darse (7.3.5.3) si la sucesión
\[
  0\longrightarrow T_p(\mathfrak m)\longrightarrow T_p(A)
  \longrightarrow T_p(k)\longrightarrow0
  \tag{7.3.5.8}\label{III.7.3.5.8.fr}
\]
es exacta. En virtud de (7.2.7) y (7.2.5), $T_p$ es isomorfo a un funtor
$M\leadsto N\otimes_A M$, y la exactitud de la sucesión (7.3.5.8) muestra, en virtud
de la sucesión exacta de los $\operatorname{Tor}$, que se tiene
$\operatorname{Tor}_1^A(N,k)=0$. Se concluye que $N=T_p(A)$ es un
$A$-módulo plano (\textbf{0}, 10.1.3).
\end{proof}

\begin{lemma}[7.3.6]
\label{III.7.3.6}
Sea $A$ un anillo, y sea $T_\bullet$ un funtor homológico
covariante de $\boldsymbol{Ab}_A$ en $\boldsymbol{Ab}$ que conmuta con las
sumas directas. Supongamos que $T_p$ y $T_{p+1}$ están definidos y
que $T_p$ es exacto por la izquierda. Entonces $T_{p+1}$ es exacto
si y solo si $T_{p+1}(A_s)$ es plano como $A$-módulo derecho plano.
\end{lemma}

\begin{proof}
En efecto, se sabe según (7.3.1) que el homomorfismo canónico \[
  T_{p+1}(A_s)\otimes_A M\longrightarrow T_{p+1}(M)
\]
es un isomorfismo de funtores. Basta ahora aplicar la definición de
$A$-módulo plano.
\end{proof}

\begin{proposition}[7.3.7]
\label{III.7.3.7}
Sea $A$ un anillo, y sea $T_\bullet$ un funtor homológico
covariante de $\boldsymbol{Ab}_A$ en $\boldsymbol{Ab}$ que conmuta con las
sumas directas. Se supone que existe un $i_0$ tal que $T_i$ es
exacto para $i\leq i_0$. Entonces, para todo entero $p>i_0$, las
condiciones siguientes son equivalentes:
\begin{enumerate}
  \item[(a)] $T_q$ es exacto para $q\leq p$;
  \item[(b)] $T_q(A_s)$ es plano como $A$-módulo a derecha para
    $q\leq p$;
  \item[(c)] para todo $A$-módulo $M$, el homomorfismo canónico
    $T_q(A_s)\otimes_A M\longrightarrow T_q(M)$
    es sobreyectivo para $q\leq p+1$.
\end{enumerate}
\end{proposition}

\begin{proof}
La equivalencia de $a)$ y $b)$ resulta de (7.3.6) por inducción sobre $q$ de
\sref{III.7.3.6}, pues $T_{i_0}$ es exacto por hipótesis; la equivalencia de $a)$ y $c)$
resulta de la equivalencia de las condiciones $a)$ y $e')$ en (7.3.3).
\end{proof}

\oldpage[III]{184}

\begin{env}[7.3.8]
\label{III.7.3.8}
Sea $A$ un anillo conmutativo, sea $B$ una $A$-álgebra, no
necesariamente conmutativa, y sea $T_\bullet$ un funtor homológico
covariante de $\boldsymbol{Ab}_A$ en $\boldsymbol{Ab}$. Se sigue de las
definiciones de \sref{III.7.1.3} que el funtor de $\boldsymbol{Ab}_B$ en
$\boldsymbol{Ab}$ obtenido por extensión de escalares de $A$ a $B$, y que denotaremos por $T_\bullet^{(B)}=(T_i^{(B)})$,
es también un funtor \emph{homológico}.
\end{env}

\begin{corollary}[7.3.9]
\label{III.7.3.9}
Supongamos que $T_\bullet$ satisface las condiciones generales de \textup{(7.3.7)} y conmuta con los límites inductivos, y además que $A$ es un dominio entero y que cada $T_n(A)$ es un
$A$-módulo de presentación finita. Entonces, para todo entero $N$,
existe un $f\in A\setmin\{0\}$ tal que
$T_p^{(A_f)}:\boldsymbol{Ab}_{A_f}\longrightarrow\boldsymbol{Ab}$
es exacto para $p\leq N$.
\end{corollary}

\begin{proof}
Por hipótesis, $T_i$ es exacto para $i\leq i_0$, de modo que
$T_i(A)$ es plano para estos valores de $i$. En virtud de (7.3.7, $b)$), basta tomar $f$ de modo que
$T_p^{(A_f)}(A_f)=T_p(A_f)$
sea un $A_f$-módulo \emph{libre} para $i_0<p\leq N$. Ahora bien,
$T_p(A_f)=(T_p(A))_f$
pues $T_p$ conmuta con los límites inductivos \sref{III.7.1.4}. Si
$x$ es el punto genérico de $\operatorname{Spec}(A)$, entonces
$(T_p(A))_x$ es un espacio vectorial de dimensión finita
sobre el cuerpo de fracciones de $A$. Como cada $T_p(A)$ es de
presentación finita, existe un $f$ con la propiedad deseada
(Bourbaki, \emph{Alg\`ebre commutative}, chap.~II, \S~5,
no.~1, cor.\ de la prop.~2).

Si solo hay un número \emph{finito} de índices $i$ tales que $T_i\ne0$,
se puede elegir $f\in A\setmin\{0\}$ de modo que \emph{todos} los
$T_p^{(A_f)}$ sean exactos.
\end{proof}

\begin{corollary}[7.3.10]
Supongamos que $T_\bullet$ verifica las condiciones generales de \textup{(7.3.7)} y
conmuta con los límites inductivos, que $A$ es conmutativo y noetheriano y los
$T_n(A)$ son $A$-módulos de tipo finito. Entonces, para todo entero $N$, existe
un abierto denso $U$ de $\operatorname{Spec}(A)$ tal que, para todo $p\leq N$,
la función
$x\leadsto\operatorname{rang}_{\mathbf{k}(x)}(T_p(\boldsymbol{k}(x)))$
es constante en $U$.
\end{corollary}

\begin{proof}
Sea $\mathfrak p$ un ideal primo minimal de $A$; por hipótesis, el anillo
$B=A/\mathfrak p$ es íntegro y $\operatorname{Spec}(B)$ se identifica con una
componente irreducible del espacio topológico $\operatorname{Spec}(A)$. Vamos
a mostrar por inducción sobre $p\leq N$ que existe
$f_p\in B\setmin\{0\}$ tal que si se pone $B'=B_{f_p}$, $T_i^{(B')}$ es \emph{exacto} y
los $T_i(B')$ son $B'$-módulos \emph{de tipo finito} para $i\leq p$. La proposición es
verdadera para $p\leq i_0$ en virtud de la hipótesis, tomando $f_p=1$ (luego
$B'=B=A/\mathfrak p$), pues al ser entonces $T_p$ exacto, $T_p(B)$ es isomorfo a
$T_p(A)/T_p(\mathfrak p)$, luego es un $A$-módulo (y \emph{a fortiori} un
$B$-módulo) de tipo finito. Razonemos por inducción sobre $p$; $f_p$ es la imagen
canónica en $B$ de un elemento $g_p\in A$, y si se pone $A'=A_{g_p}$, se tiene
$B'=A'/\mathfrak p'$ donde $\mathfrak p'$ es un ideal primo minimal de $A'$,
igual a $\mathfrak p_{g_p}$. Como se tiene
$T_i(A_{g_p})=(T_i(A))_{g_p}$,
los $T_i(A')$ son $A'$-módulos de tipo finito, luego el funtor
$T_\bullet^{(A')}$ verifica las mismas hipótesis que $T_\bullet$, pero
reemplazando $i_0$ por $p$. Se puede pues limitar al caso en que $A'=A$, y $T_p$
es exacto; la sucesión exacta
$0\longrightarrow\mathfrak p\longrightarrow A\longrightarrow A/\mathfrak p\longrightarrow0$
da entonces la sucesión exacta
$T_{p+1}(A)\longrightarrow T_{p+1}(A/\mathfrak p)\xrightarrow{\partial}T_p(\mathfrak p)\longrightarrow T_p(A)$,
y como $T_p$ es exacto, la última flecha de la derecha es inyectiva, luego
$T_{p+1}(A/\mathfrak p)$ es un cociente de $T_{p+1}(A)$ y es por tanto de
tipo finito. Observemos ahora que el razonamiento de (7.3.9) solo utilizó el hecho
de que los $T_p(A)$ son de tipo finito para $p\leq N$; se puede pues aplicar
al anillo íntegro $B$, al funtor $T_\bullet^{(B)}$ y a $N=p+1$, lo que
termina el razonamiento por inducción. Dicho esto, hay un
$f_N\in B\setmin\{0\}$ tal que $\operatorname{Spec}(B_{f_N})$ es un abierto $V$
denso en $\operatorname{Spec}(B)$ que no interseca ninguna otra
componente irreducible de $\operatorname{Spec}(A)$. Si la proposición queda
\oldpage[III]{185}
demostrada para $B_{f_N}$, se tendrá un abierto $W$ denso en $V$ en
el que las funciones del enunciado serán constantes, puesto que
$A_x=(B_{f_N})_x$ para todo $x\in W$. Haciendo el mismo razonamiento para
toda componente irreducible de $\operatorname{Spec}(A)$, el corolario quedará
demostrado. Se puede pues limitar al caso en que $A$ es \emph{íntegro}; el razonamiento de
(7.3.9) prueba entonces la existencia de un $f\in A\setmin\{0\}$ tal que los $T_p(A_f)$
son $A_f$-módulos \emph{libres} de tipo finito para $p\leq N$, lo que implica
la conclusión de (7.3.10) en virtud de (7.3.4).
\end{proof}

\begin{proposition}[7.3.11]
Sean $A$ un anillo conmutativo local, $k$ su cuerpo residual, $T_\bullet$ un
funtor homológico covariante de $\boldsymbol{Ab}_A$ en $\boldsymbol{Ab}$,
que conmuta con las sumas directas. Se supone que existe $i_0$ tal que $T_i$ es
exacto para $i\leq i_0$, y que todos los $T_n(A)$ son $A$-módulos de
presentación finita. Entonces las condiciones equivalentes $a)$, $b)$, $c)$ de
\textup{(7.3.7)} implican las dos siguientes, y son equivalentes a ellas cuando
el anillo es además reducido:
\begin{enumerate}
  \item[d)] Para todo $x\in\operatorname{Spec}(A)$, se tiene
    $\operatorname{rang}_{\mathbf{k}(x)}T_q(\boldsymbol{k}(x))=\operatorname{rang}_kT_q(k)$
    para $q\leq p$.
  \item[$d'$)] Para todo punto genérico $x_j$ de una componente irreducible
    de $\operatorname{Spec}(A)$, se tiene
    \[
      \operatorname{rang}_{\mathbf{k}(x_j)}T_q(\boldsymbol{k}(x_j))
      =
      \operatorname{rang}_kT_q(k)
      \quad\text{para }q\leq p.
    \]
\end{enumerate}
\end{proposition}

\begin{proof}
Como $T_q(A)$ es un $A$-módulo de presentación finita, la condición $b)$ de
(7.3.7) equivale a decir que $T_q(A)$ es un $A$-módulo \emph{libre} para $q\leq p$
(Bourbaki, \emph{Alg. comm.}, cap. II, \S 3, n\textsuperscript{o} 2, cor. 2
de la prop. 5); la condición $c)$ implica que
$T_q(\boldsymbol{k}(x))=T_q(A)\otimes_A\boldsymbol{k}(x)$ para $q\leq p$,
luego las condiciones equivalentes de (7.3.7) implican $d)$, y es trivial
que $d)$ implica $d')$. Queda por probar que $d')$ implica $a)$ cuando $A$
es \emph{reducido}. Razonemos por inducción sobre $q\leq p$, puesto que $T_q$ es exacto
para $q\leq i_0$. Supongamos pues $T_k$ exacto para $k\leq q<p$ y mostremos que
$T_{q+1}(A)$ es un $A$-módulo \emph{libre}. En virtud de la hipótesis de inducción,
$T_{q+1}(A)\otimes_A M$ es isomorfo a $T_{q+1}(M)$ para todo $A$-módulo $M$,
por la condición $c)$ de (7.3.7) y (7.3.3); aplicando esta propiedad a
$M=\boldsymbol{k}(x_j)$ y $M=k$, se obtiene, en virtud de la hipótesis $d')$, que
\[
  \operatorname{rang}_{\mathbf{k}(x_j)}
  \bigl(T_{q+1}(A)\otimes_A\boldsymbol{k}(x_j)\bigr)
  =
  \operatorname{rang}_kT_{q+1}(k)
\]
para todo $i$; pero esto implica que $T_{q+1}(A)$ es libre (Bourbaki,
\emph{Alg. comm.}, cap. II, \S 3, n\textsuperscript{o} 2, prop. 7), lo que
termina la demostración.
\end{proof}

Los resultados anteriores se mejorarán considerablemente para los funtores
homológicos del tipo particular que vamos a estudiar en (7.4); se
obtendrán en efecto criterios de exactitud en los que solo interviene \emph{uno solo}
de los $T_p$.


\subsection{Criterios de exactitud para los funtores
$H_\bullet(P_\bullet\otimes_A M)$.}
\label{subsection:III.7.4.fr}

\begin{env}[7.4.1]
\label{III.7.4.1}
Sean $A$ un anillo (no necesariamente conmutativo) y $P_\bullet$ un complejo
de $A$-módulos \emph{planos} a derecha. Como el funtor
$M\leadsto P_k\otimes_A M$ es entonces exacto en $\boldsymbol{Ab}_A$ para todo $k$,
el $\partial$-funtor
\[
  T_\bullet(M)=H_\bullet(P_\bullet\otimes_A M)
  \tag{7.4.1.1}\label{III.7.4.1.1.fr}
\]
\oldpage[III]{186}
es un \emph{funtor homológico} de $\boldsymbol{Ab}_A$ en $\boldsymbol{Ab}$,
evidentemente $A$-lineal cuando $A$ es conmutativo
\sref{III.7.1.2}, y conmuta con los límites inductivos.

Si $A$ es conmutativo, entonces, para toda $A$-álgebra $B$, el funtor
homológico $T_\bullet^{(B)}$ \sref{III.7.3.8} viene dado, por definición, por
\[
  T_\bullet^{(B)}(N)
  =
  H_\bullet(P_\bullet\otimes_A N_{[\rho]})
  \tag{7.4.1.2}\label{III.7.4.1.2.fr}
\]
donde $\rho:A\to B$ es el homomorfismo que define la estructura de álgebra de
$B$; como también se puede escribir
$P_\bullet\otimes_A N_{[\rho]}=P_\bullet\otimes_A(B\otimes_B N)_{[\rho]}=(P_\bullet\otimes_A B)\otimes_B N$,
se ve que
\[
  T_\bullet^{(B)}(N)=H_\bullet(P'_\bullet\otimes_B N)
  \tag{7.4.1.3}\label{III.7.4.1.3.fr}
\]
para todo $B$-módulo $N$, siendo $P'_\bullet$ el complejo
$P_\bullet\otimes_A B$ de $B$-módulos \emph{planos}
\sref[0]{0.6.2.1}.
\end{env}

\begin{proposition}[7.4.2]
Bajo las condiciones generales de \textup{(7.4.1)}, y para un entero $p\in\mathbf Z$
dado, las propiedades siguientes son equivalentes:
\begin{enumerate}
  \item[a)] $T_p$ es exacto por la izquierda \textup{(o, lo que equivale a lo mismo,} $T_{p+1}$
    es exacto por la derecha).
  \item[b)] $Z'_p(P_\bullet)=\operatorname{Coker}(P_{p+1}\longrightarrow P_p)$
    es un $A$-módulo por la derecha plano.
  \item[c)] Existe un complejo $P'_\bullet$ de $A$-módulos por la derecha planos
    tal que la diferencial
    \[
      d_{p+1}:P'_{p+1}\longrightarrow P'_p
    \]
    es nula, y un isomorfismo de funtores homológicos de
    $H_\bullet(P_\bullet\otimes_A M)$ sobre
    $H_\bullet(P'_\bullet\otimes_A M)$.
\end{enumerate}
\end{proposition}

\begin{proof}
Por definición se tiene una sucesión exacta funtorial en $M$
\[
  0\longrightarrow T_p(M)
  \longrightarrow Z'_p(P_\bullet\otimes M)
  \longrightarrow P_{p-1}\otimes M
\]
donde
$Z'_p(P_\bullet\otimes M)=\operatorname{Coker}(P_{p+1}\otimes M\longrightarrow P_p\otimes M)=Z'_p(P_\bullet)\otimes M$
en virtud de la exactitud por la derecha del producto tensorial. Para todo homomorfismo
$f:M\to N$, se tiene pues un diagrama conmutativo
\[
\xymatrix@C=30pt@R=38pt{
  0\ar[r]
  &T_p(M)\ar[r]\ar[d]_u
  &Z'_p(P_\bullet)\otimes M\ar[r]\ar[d]_v
  &P_{p-1}\otimes M\ar[d]_w
  \\
  0\ar[r]
  &T_p(N)\ar[r]
  &Z'_p(P_\bullet)\otimes N\ar[r]
  &P_{p-1}\otimes N
}
\tag{7.4.2.1}\label{III.7.4.2.1.fr}
\]
cuyas filas son exactas. Si $f$ es un monomorfismo, lo mismo ocurre
con $w$ puesto que $P_{p-1}$ es plano; si $T_p$ es exacto por la izquierda, $u$ es
también un monomorfismo; se concluye que $v$ es un monomorfismo, lo que
implica que $Z'_p(P_\bullet)$ es plano. Recíprocamente, si es así, $v$
es un monomorfismo para todo monomorfismo $f:M\to N$, luego el diagrama
(7.4.2.1) muestra que $u$ es un monomorfismo, y por tanto $T_p$ (que ya es
semiexacto) es exacto por la izquierda. Así, $a)$ y $b)$ son equivalentes. Es
inmediato que $c)$ implica $a)$, pues si
$d_{p+1}:P_{p+1}\to P_p$ es nula, y
$0\to M'\to M\to M''\to0$ es una sucesión exacta de $A$-módulos, el operador
borde en la sucesión exacta
\[
  H_{p+1}(P_\bullet\otimes M'')
  \xrightarrow{\partial}H_p(P_\bullet\otimes M')
  \longrightarrow H_p(P_\bullet\otimes M)
\]
\oldpage[III]{187}
es nulo por definición (M, IV, 1), luego $T_p$ es exacto por la izquierda. Mostremos
recíprocamente que $b)$ implica $c)$. Si
$Z_{p+1}(P_\bullet)=\operatorname{Ker}(P_{p+1}\longrightarrow P_p)$,
se tiene una sucesión exacta
\[
  0\longrightarrow Z_{p+1}(P_\bullet)\longrightarrow P_{p+1}
  \longrightarrow Z'_p(P_\bullet)\longrightarrow0
\]
en la cual $P_{p+1}$ y $Z'_p(P_\bullet)$ son planos, luego
$Z_{p+1}(P_\bullet)$ es plano ($\mathbf{0}_{\mathrm I}$, 6.1.2). Se tomará
\[
  P'_i=P_i\quad\text{para }i\ne p\text{ y }i\ne p+1,
  \qquad
  P'_p=Z'_p(P_\bullet)
  \qquad\text{y}\qquad
  P'_{p+1}=Z_{p+1}(P_\bullet);
\]
como diferencial $d'_i:P'_i\to P'_{i-1}$, se tomará la del complejo
$P_\bullet$ para $i\ne p$ y $i\ne p+1$, $0$ para $i=p+1$ y para $i=p$
el homomorfismo $Z'_p(P_\bullet)\to P_{p-1}$ deducido de $d_p$ por paso al
cociente. Como los $P_i$ son planos, se tiene
\[
  Z'_i(P_\bullet\otimes M)=Z'_i(P_\bullet)\otimes M,
  \qquad
  Z_i(P_\bullet\otimes M)=Z_i(P_\bullet)\otimes M
  \quad\text{y}\quad
  B_i(P_\bullet\otimes M)=B_i(P_\bullet)\otimes M
\]
(poniendo $B_i(P_\bullet)=\operatorname{Im}(P_{i+1}\to P_i)$); se deducen
inmediatamente para todo $M$ isomorfismos funtoriales
$H_i(P_\bullet\otimes M)\xrightarrow{\sim}H_i(P'_\bullet\otimes M)$
para todo $i$, y la verificación del hecho de que se trata de un isomorfismo de
$\partial$-funtores se deduce sin dificultad de la definición de $\partial$
(M, IV, 1).
\end{proof}

Se observará que las condiciones de (7.4.2) implican también que
$B_p(P_\bullet)$ es \emph{plano}, pues se tiene una sucesión exacta
$0\longrightarrow B_p(P_\bullet)\longrightarrow P_p
\longrightarrow Z'_p(P_\bullet)\longrightarrow0$,
en la cual $P_p$ y $Z'_p(P_\bullet)$ son planos ($\mathbf{0}_{\mathrm I}$, 6.1.2).

\begin{corollary}[7.4.3]
\label{III.7.4.3}
Supongamos que $A$ es un anillo noetheriano regular de dimensión $1$
\textup{(en otros términos, un \emph{producto de anillos de Dedekind}
($\mathbf{0}_{\mathrm{IV}}$, 17.1.3 y 17.3.7), por ejemplo un anillo \emph{principal})}. Entonces,
para que $T_p$ sea exacto por la izquierda, es necesario y suficiente que
$T_p(A)$ sea un $A$-módulo plano. Para que $T_p$ sea exacto, es necesario y
suficiente que $T_p(A)$ y $T_{p-1}(A)$ sean $A$-módulos planos.
\end{corollary}

\begin{proof}
Recordemos que, para un módulo $M$ sobre un anillo de Dedekind, es equivalente
decir que $M$ es \emph{plano} o que es \emph{sin torsión} ($\mathbf{0}_{\mathrm I}$, 6.3.3 y
6.3.4); bajo las hipótesis de (7.4.3), todo submódulo de un $A$-módulo plano es, por tanto, plano.

La segunda aserción de \sref{III.7.4.3} resulta de la primera, puesto que decir
que $T_p$ es exacto significa que $T_p$ y $T_{p-1}$ son exactos por la
izquierda. Para demostrar la primera aserción, observemos que se tiene una
sucesión exacta
\[
  0\longrightarrow H_p(P_\bullet)\longrightarrow Z'_p(P_\bullet)
  \longrightarrow B_{p-1}(P_\bullet)\longrightarrow0
\]
en la cual $B_{p-1}(P_\bullet)$ es un $A$-módulo \emph{plano}, por ser un submódulo
del $A$-módulo plano $P_{p-1}$. Por tanto, es equivalente decir que
$H_p(P_\bullet)$ es plano o que $Z'_p(P_\bullet)$ es plano
($\mathbf{0}_{\mathrm I}$, 6.1.2).
\end{proof}

Las aplicaciones más importantes de (7.4.2) son las siguientes:

\begin{proposition}[7.4.4]
\label{III.7.4.4}
Sean $A$ un anillo noetheriano y $P_\bullet$ un complejo de $A$-módulos
planos; se supone, o bien que los $P_i$ son de tipo finito, o bien que los
$H_i(P_\bullet)$ son $A$-módulos de tipo finito y que existe $i_0$ tal que
$H_i(P_\bullet)=0$ para $i<i_0$. Sea $T$ el funtor homológico definido por
\eqref{III.7.4.1.1.eq}. Entonces, el conjunto $U$ de los
$y\in\operatorname{Spec}(A)$ tales que $(T_p)_y$ \sref{III.7.1.4} sea exacto por la derecha
\textup{(resp. \emph{exacto por la izquierda, exacto})} es abierto en $\operatorname{Spec}(A)$.
\end{proposition}

\begin{proof}
En la segunda hipótesis sobre $P_\bullet$, se puede reemplazar $P_\bullet$ por
un complejo $P'_\bullet$ de
\oldpage[III]{188}
$A$-módulos libres de tipo finito para el cual el funtor
$H_\bullet(P'_\bullet\otimes_A M)$ es isomorfo (como $\partial$-funtor) a
$T_\bullet(M)$ \sref[0]{0.11.9.3}. Por tanto, siempre se puede reducir a la
primera hipótesis, y, en este caso, los $Z'_i(P_\bullet)$ son de
tipo finito; además, se puede limitar a demostrar las afirmaciones relativas a
la exactitud por la izquierda (cf. (7.4.2, $a)$)). Dicho esto, sea $x\in U$; como el
funtor $M\leadsto M_x$ es exacto, se tiene
$(Z'_p(P_\bullet))_x=Z'_p((P_\bullet)_x)$,
y (teniendo en cuenta (7.4.1.3)) la hipótesis implica, en virtud de
(7.4.2, $b)$), que $(Z'_p(P_\bullet))_x$ es un $A_x$-módulo plano,
luego \emph{libre}, pues es de tipo finito y $A_x$ es un anillo local noetheriano
(Bourbaki, \emph{Algèbre commutative}, cap.~II, \S~3,
n\textsuperscript{o}~2, cor.~2 de la prop.~5). Se deduce que existe $f\in A$
tal que $(Z'_p(P_\bullet))_f$ es libre sobre $A_f$ (Bourbaki,
\emph{Algèbre commutative}, cap.~II, \S~5, n\textsuperscript{o}~1, cor.~de la
prop.~2), y, \emph{a fortiori}, $(Z'_p(P_\bullet))_y$ es libre sobre $A_y$ para
todo $y\in D(f)$, lo
que termina la demostración en virtud de (7.4.2, $b)$).
\end{proof}

\begin{corollary}[7.4.5]
\label{III.7.4.5}
Bajo las hipótesis de \sref{III.7.4.4}, supongamos además que $A$ es un
dominio entero.
Entonces, el conjunto $U$ de los $x\in\operatorname{Spec}(A)$ tales que $(T_p)_x$ sea
exacto es abierto no vacío.
\end{corollary}

\begin{proof}
En efecto, basta demostrar que $(T_p)_x$ es exacto para el punto \emph{genérico} $x$
de $\operatorname{Spec}(A)$, lo cual es inmediato, pues $A_x$ es un cuerpo, y, por tanto,
todo funtor aditivo sobre $\boldsymbol{Ab}_{A_x}$ es exacto.
\end{proof}

\begin{proposition}[7.4.6]
Bajo las hipótesis generales de \textup{(7.4.4)}, las condiciones $a)$, $b)$ y $c)$ \textup{de
(7.4.2)} son también equivalentes a:
\begin{enumerate}
  \item[d)] Existe un $A$-módulo $Q$ y un isomorfismo funtorial
    \[
      T_p(M)\xrightarrow{\sim}\operatorname{Hom}_A(Q,M).
      \tag{7.4.6.1}\label{III.7.4.6.1.fr}
    \]
    Además, el $A$-módulo $Q$ queda determinado, salvo isomorfismo único, por
    esta propiedad, y es de tipo finito.
\end{enumerate}
\end{proposition}

\begin{proof}
La unicidad de $Q$ es un caso particular de la unicidad de un objeto
representante de un funtor representable \sref[0]{0.8.1.5}. Es claro que el
segundo miembro de \eqref{III.7.4.6.1.eq} es exacto por la izquierda.
Recíprocamente, para demostrar la existencia de $Q$, cuando $T_p$ es exacto
por la izquierda, primero se puede, como en \sref{III.7.4.4}, reducir al caso
en que los $P_i$ son planos y de tipo finito, luego (puesto que $A$ es
noetheriano) \emph{proyectivos} de tipo finito (Bourbaki,
\emph{Algèbre commutative}, cap.~II, \S~5, n\textsuperscript{o}~2, cor.~del
teor.~1). El \emph{dual} $\check P_i$ de $P_i$ es entonces también un $A$-módulo
proyectivo de tipo finito, $P_i$ es canónicamente isomorfo al dual de
$\check P_i$, y el homomorfismo canónico
$P_i\otimes_A M\longrightarrow\operatorname{Hom}_A(\check P_i,M)$
es biyectivo (Bourbaki, \emph{Alg.}, cap. II, 3.\textsuperscript{a} ed.,
\S 4, n\textsuperscript{o} 2, prop. 2). Se sabe por otra parte (7.4.2, $c)$)
que se puede suponer que $d_{p+1}:P_{p+1}\to P_p$ es nula, de modo que se tiene una sucesión exacta
\[
  0\longrightarrow T_p(M)\xrightarrow{u}P_p\otimes M
  \xrightarrow{v}P_{p-1}\otimes M
\]
donde $v=d_p\otimes1$. Pongamos entonces $Q'=\operatorname{Ker}(d_p)$, de modo que se tiene
la sucesión exacta
$0\longrightarrow Q'\xrightarrow{w}P_p\xrightarrow{d_p}P_{p-1}$,
de donde, por transposición, la sucesión exacta
$\check P_{p-1}\xrightarrow{{}^t d_p}\check P_p
\xrightarrow{{}^t w}\check Q'\longrightarrow0$.
Vamos a ver que
$Q=\check Q'=\operatorname{Coker}({}^t d_p)$
responde a la cuestión. En efecto, se tiene la sucesión exacta
\[
  0\longrightarrow\operatorname{Hom}(Q,M)
  \longrightarrow\operatorname{Hom}(\check P_p,M)
  \xrightarrow{v'}\operatorname{Hom}(\check P_{p-1},M)
\]
\oldpage[III]{189}
donde $v'=\operatorname{Hom}({}^t d_p,1)$; al identificar canónicamente
$P_i\otimes M$ con $\operatorname{Hom}(\check P_i,M)$, $v'$ se identifica, por tanto, con
$v=d_p\otimes1$, y se obtiene así el isomorfismo funtorial
$T_p(M)\xrightarrow{\sim}\operatorname{Hom}_A(Q,M)$
buscado. Además, $Q$, por ser un cociente de $\check P_p$, es de tipo finito.
\end{proof}

\begin{proposition}[7.4.7]
\label{III.7.4.7}
Supongamos verificadas las condiciones generales de \sref{III.7.4.4}. Entonces,
para todo $A$-módulo $M$ de tipo finito:
\begin{enumerate}
  \item[\textup{(i)}] los $T_i(M)$ son $A$-módulos de tipo finito;
  \item[\textup{(ii)}] para todo ideal $\mathfrak m$ de $A$, el homomorfismo canónico
    \[
      (T_i(M))^{\widehat{\,}}
      \longrightarrow
      \varprojlim_n T_i\bigl(M\otimes_A(A/\mathfrak m^{n+1})\bigr)
      \tag{7.4.7.1}\label{III.7.4.7.1.fr}
    \]
    (donde el primer miembro es el completado separado de $T_i(M)$ para la
    topología $\mathfrak m$-ádica) es biyectivo.
\end{enumerate}
\end{proposition}

\begin{proof}
Como en \sref{III.7.4.4}, se reduce primero al caso en que los $P_i$ son de
tipo finito; como $A$ es noetheriano, los submódulos de los
$P_i\otimes_A M$ son de tipo finito, de donde se deduce trivialmente la
aserción (i). En cuanto a la aserción (ii), resulta más generalmente del lema
siguiente:
\end{proof}

\begin{lemma}[7.4.7.2]
\label{III.7.4.7.2}
Sean $A$ un anillo noetheriano y $u:E\to F$ un homomorfismo de $A$-módulos de
tipo finito. Para todo $A$-módulo $M$ de tipo finito, pongamos
$K(M)=\operatorname{Ker}(u\otimes1_M)$,
$C(M)=\operatorname{Coker}(u\otimes1_M)$;
entonces los homomorfismos canónicos
\[
  (K(M))^{\widehat{\,}}\longrightarrow\varprojlim_n K(M_n),
  \qquad
  (C(M))^{\widehat{\,}}\longrightarrow\varprojlim_n C(M_n)
  \tag{7.4.7.3}\label{III.7.4.7.3.fr}
\]
(donde se ha puesto
$M_n=M\otimes_A(A/\mathfrak m^{n+1})=M/\mathfrak m^{n+1}M$) son biyectivos
para todo ideal $\mathfrak m$ de $A$.
\end{lemma}

\begin{proof}
Como $E\otimes M$ y $F\otimes M$ son de tipo finito, y el funtor
$M\leadsto\widehat M$ es \emph{exacto} en la categoría de los $A$-módulos de tipo finito
($\mathbf{0}_{\mathrm I}$, 7.3.3), $(K(M))^{\widehat{\,}}$ y $(C(M))^{\widehat{\,}}$ son,
respectivamente, el núcleo y el conúcleo de
$(u\otimes1)^{\widehat{\,}}:(E\otimes M)^{\widehat{\,}}\longrightarrow(F\otimes M)^{\widehat{\,}}$.
La exactitud por la izquierda del funtor $\varprojlim$ muestra, por tanto, que
$(K(M))^{\widehat{\,}}=\varprojlim K(M_n)$; por otra parte, la exactitud por la
derecha del producto tensorial prueba que
$C(M_n)=C(M)\otimes_A(A/\mathfrak m^{n+1})$, luego
$(C(M))^{\widehat{\,}}=\varprojlim C(M_n)$ por definición.
\end{proof}

\begin{remark}[7.4.8]
\label{III.7.4.8}
Teniendo en cuenta \sref{III.6.10.5} y \sref{III.6.10.6}, se ve que, bajo una
hipótesis suplementaria de planitud, \sref{III.7.4.7} vuelve a dar el hecho de
que (4.1.7.1) es un isomorfismo, es decir, lo esencial del «primer teorema de
comparación» para los morfismos propios; además, el enunciado se aplica ya no
solo a un $\mathscr O_X$-módulo coherente, sino a un \emph{complejo} de tales módulos.
Sería interesante obtener un enunciado que comprendiese a la vez
\sref{III.7.4.7} y (4.1.7.1) como casos particulares. Se observará que, cuando los
$P_i$ son de tipo finito, la demostración de \sref{III.7.4.7} no utiliza el
hecho de que sean módulos planos; convendría examinar si la conclusión de
\sref{III.7.4.7} sigue siendo válida cuando no se supone que los $P_i$ sean
planos ni de tipo finito, sino que se supone que los $H_i(P_\bullet)$ son de
tipo finito para todo $i$ y nulos para $i<i_0$. ¿Es posible entonces reemplazar
$P_\bullet$ por un complejo $P'_\bullet$ de $A$-módulos de tipo finito tal que
los funtores $H_\bullet(P_\bullet\otimes M)$ y
$H_\bullet(P'_\bullet\otimes M)$ (que ya no son homológicos) sigan siendo
isomorfos?
\end{remark}

\oldpage[III]{190}
\subsection{Caso de los anillos locales noetherianos}
\label{subsection:III.7.5}
\label{subsection:III.7.5.fr}

\begin{env}[7.5.1]
\label{III.7.5.1}
Sean $A$ un anillo local noetheriano y $\mathfrak m$ su ideal
maximal. Para todo $A$-módulo $M$, designemos por $\widehat M$ su
completado separado para la topología $\mathfrak m$-ádica, el cual es
isomorfo a
$\varprojlim\bigl(M\otimes_A(A/\mathfrak m^{n+1})\bigr)=\varprojlim\bigl(M/\mathfrak m^{n+1}M\bigr)$.
Sea $T$ un funtor covariante aditivo de $\boldsymbol{Ab}_A$ en
$\boldsymbol{Ab}$. Los homomorfismos canónicos
\sref{III.7.2.3.1}
\[
  T(M)\otimes_A(A/\mathfrak m^{n+1})
  \longrightarrow
  T\bigl(M\otimes_A(A/\mathfrak m^{n+1})\bigr)
\]
forman evidentemente un sistema proyectivo de $A$-homomorfismos. Por
paso al límite dan, por tanto, un $\widehat A$-homomorfismo funtorial
en $M$
\[
  (T(M))^{\widehat{\,}}
  \longrightarrow
  \varprojlim_n T(M_n)
  \tag{7.5.1.1}\label{III.7.5.1.1.fr}
\]
donde se ha puesto
$M_n=M\otimes_A(A/\mathfrak m^{n+1})$, $A_n=A/\mathfrak m^{n+1}$.
\end{env}

\begin{proposition}[7.5.2]
\label{III.7.5.2}
Sean $A$ un anillo local noetheriano de ideal maximal $\mathfrak m$,
$k=A/\mathfrak m$ su cuerpo residual y $T$ un funtor covariante
aditivo de $\mathbf{Ab}_A$ en $\mathbf{Ab}$, semiexacto y que conmuta con
los límites inductivos. Se supone además que, para todo $A$-módulo
de tipo finito $M$, $T(M)$ es un $A$-módulo de tipo finito y que el
homomorfismo canónico \eqref{III.7.5.1.1.eq} es un isomorfismo. En
estas condiciones, las propiedades siguientes son equivalentes:
\begin{enumerate}
  \item[(a)] $T$ es exacto por la derecha;
  \item[(b)] para todo $n$, el funtor $N\leadsto T(N)$ es exacto por
    la derecha en la categoría de los $A_n$-módulos de tipo finito;
    equivalentemente, $T$ es exacto por la derecha en la categoría
    de los $A$-módulos \emph{de longitud finita};
  \item[(c)] el homomorfismo canónico
    $T(\varepsilon):T(A)\longrightarrow T(k)$
    es sobreyectivo;
  \item[(d)] para todo $n$ suficientemente grande, el homomorfismo
    canónico
    $T(A_n)\longrightarrow T(k)$ es sobreyectivo.
\end{enumerate}
\end{proposition}

\begin{proof}
Está claro que \emph{a)} implica \emph{b)}. Mostremos que \emph{b)} implica \emph{a)}, es decir que si $u:M\to N$ un epimorfismo de $A$-módulos, $T(u)$
es un epimorfismo. Como $T$ conmuta con los límites inductivos y el funtor $\varinjlim$ es
exacto en la categoría de módulos (para conjuntos de índices filtrantes),
se puede limitar al caso en que $M$ y $N$ son de tipo
finito.

Los módulos $T(M)$ y $T(N)$ son entonces de tipo finito. Como $A$ es
noetheriano y local, basta probar que
$(T(u))^{\widehat{\,}}:(T(M))^{\widehat{\,}}\longrightarrow(T(N))^{\widehat{\,}}$
es sobreyectivo ($\mathbf{0}_{\mathrm I}$, 7.3.5 y $\mathbf{0}_{\mathrm I}$, 6.4.1). Por hipótesis,
$(T(M))^{\widehat{\,}}$ y $(T(N))^{\widehat{\,}}$ son respectivamente $\varprojlim T(M_n)$
y $\varprojlim T(N_n)$, luego $(T(u))^{\widehat{\,}}$ es el límite proyectivo del sistema proyectivo de
homomorfismos
$T(u\otimes1_{A_n}):T(M_n)\longrightarrow T(N_n)$.
La condición \emph{(b)} significa precisamente que estos homomorfismos son
\emph{sobreyectivos}. Además, $T(M_n)$ es un $A_n$-módulo de tipo finito y
$A_n$ es un anillo \emph{artiniano} por hipótesis; se concluye que $(T(u))^{\widehat{\,}}$ es sobreyectivo (\textbf{0}, 13.1.2 y
13.2.2). Está claro que \emph{a)} implica \emph{c)}, y como $T(\varepsilon)$ se factoriza
en
$T(A)\longrightarrow T(A_n)\longrightarrow T(k)$,
\emph{c)} implica \emph{d)}; finalmente, resulta de (7.2.7) que \emph{b)} y \emph{d)} son equivalentes
puesto que $T$ es
semiexacto en $\boldsymbol{Ab}_{A_n}$. Esto termina la demostración.
\end{proof}

\begin{corollary}[7.5.3]
\label{III.7.5.3}
Bajo las hipótesis generales de \sref{III.7.5.2}, si $T(k)=0$,
entonces $T(M)=0$ para todo $A$-módulo $M$.
\end{corollary}

\begin{proof}
\oldpage[III]{191}
Como $k$ es el único $A$-módulo simple, de
\eqref{III.7.3.5.4.eq} se deduce que $T(E)=0$ para todo $A$-módulo
de longitud finita $E$. Si $M$ es de tipo finito,
$(T(M))^{\widehat{\,}}$ es isomorfo a $\varprojlim T(M_n)$, y como los $M_n$ es de longitud finita, de modo que
$(T(M))^{\widehat{\,}}=0$. Por hipótesis, $T(M)$ es de tipo finito por hipótesis, es isomorfo a un submódulo de $(T(M))^{\widehat{\,}}$
($\mathbf{0}_{\mathrm I}$, 7.3.5), luego se tiene $T(M)=0$. Finalmente, para un $A$-módulo arbitrario $M$, el módulo
$T(M)$ es el límite inductivo de los $T(N_\alpha)$, donde los
$N_\alpha$ recorren los submódulos de tipo finito de $M$. Esto
termina la demostración.
\end{proof}

\begin{proposition}[7.5.4]
\label{III.7.5.4}
Sean $A$ un anillo local noetheriano de ideal maximal $\mathfrak m$,
$k=A/\mathfrak m$ su cuerpo residual y $T_\bullet$ un funtor
homológico de $\mathbf{Ab}_A$ en $\mathbf{Ab}$ que conmuta con los límites
inductivos. Se supone además que, para todo $i$ y todo $A$-módulo
de tipo finito $M$, $T_i(M)$ es de tipo finito y el homomorfismo
canónico
$(T_i(M))^{\widehat{\,}}\longrightarrow\varprojlim T_i(M_n)$
es biyectivo. Para un entero dado $p$, las condiciones siguientes
son equivalentes:
\begin{enumerate}
  \item[(a)] $T_p$ es exacto;
  \item[(b)] $T_p$ es exacto por la derecha y $T_p(A)$ es un
    $A$-módulo libre.
  \item[c)] Los homomorfismos canónicos
    $T_{p+1}(A)\longrightarrow T_{p+1}(k)$ y $T_p(A)\longrightarrow T_p(k)$
    son sobreyectivos.
  \item[d)] Para todo $n$, los homomorfismos canónicos
    $T_{p+1}(A_n)\longrightarrow T_{p+1}(k)$ y $T_p(A_n)\longrightarrow T_p(k)$
    son sobreyectivos.
  \item[e)] Para todo $n$, el funtor $N\leadsto T_p(N)$ es exacto
    en la categoría de los $A_n$-módulos de tipo finito.
\end{enumerate}
\end{proposition}

\begin{proof}
Por \sref{III.7.3.3}, la condición \emph{(a)} equivale a decir que
$T_{p+1}$ y $T_p$ son exactos por la derecha. Como $T_\bullet$ es
un funtor homológico en $\boldsymbol{Ab}_{A_n}$, el mismo razonamiento que
en \sref{III.7.3.1} muestra que \emph{(e)} equivale a decir que $T_p$ y
$T_{p+1}$ son exactos por la derecha en la categoría de los
$A_n$-módulos de tipo finito. Se deduce de \sref{III.7.5.2} que \emph{(a)}
y \emph{(e)} son equivalentes; la misma proposición da también la
equivalencia de \emph{(a)}, \emph{(c)} y \emph{(d)}.

Finalmente, todo módulo plano de tipo finito sobre el anillo local
$A$ es \emph{libre} (Bourbaki, \emph{Algèbre commutative}, cap.~II,
\S~3, n\textsuperscript{o}~2, cor.~2 de la prop.~5). La equivalencia
de \emph{(a)} y \emph{(b)} resulta entonces de \sref{III.7.3.1} y
\sref{III.7.3.3}.
\end{proof}

\begin{corollary}[7.5.5]
\label{III.7.5.5}
Supongamos verificadas las hipótesis generales de
\sref{III.7.5.4}.
\begin{enumerate}
  \item[\textup{(i)}] Si $T_p(k)=0$, entonces $T_p=0$, el funtor $T_{p+1}$
    es exacto por la derecha y $T_{p-1}$ es exacto por la izquierda.
  \item[\textup{(ii)}] Si $T_{p-1}(k)=T_{p+1}(k)=0$, entonces $T_p$ es
    exacto, el homomorfismo canónico
    \[
      T_p(A)\otimes_A M\longrightarrow T_p(M)
    \]
    es biyectivo y $T_p(A)$ es un $A$-módulo libre.
\end{enumerate}
\end{corollary}

\begin{proof}
La aserción (i) se deduce inmediatamente de \sref{III.7.5.3}, pues
$T_p$ es semiexacto; la última aserción resulta de la definición de
funtor homológico. Teniendo en cuenta \sref{III.7.3.3}, (i) da de
inmediato las dos primeras aserciones de (ii), y el hecho de que
$T_p(A)$ sea libre resulta de \sref{III.7.5.4}.
\end{proof}

\begin{corollary}[7.5.6]
Supongamos verificadas las hipótesis generales de \textup{(7.5.4)}, y supongamos además
que $A$ sea un anillo de valuación discreta.
\begin{enumerate}
  \item[\textup{(i)}] Para que $T_p$ sea exacto por la derecha, es necesario y suficiente que
    $T_{p-1}(A)$ sea un $A$-módulo libre.
  \item[\textup{(ii)}] Para que $T_p$ sea exacto, es necesario y suficiente que $T_p(A)$ y
    $T_{p-1}(A)$ sean $A$-módulos libres.
\end{enumerate}
\end{corollary}

\begin{proof}
\oldpage[III]{192}
Está claro que (i) implica (ii) (cf. (7.3.3)). Para probar (i), observemos
que si $f$ es un generador del ideal maximal de $A$ («uniformizante» de
$A$), para que un $A$-módulo de tipo finito $M$ sea libre (o plano, lo que viene
a ser lo mismo), es necesario y suficiente que la homotecia $h_f:x\to fx$ de $M$ sea
inyectiva, pues esto equivale aquí a decir que $M$ es sin torsión
($\mathbf{0}_{\mathrm I}$, 6.3.4). Consideremos entonces la sucesión exacta
$0\longrightarrow A\xrightarrow{h_f}A\longrightarrow k\longrightarrow0$,
que proporciona la sucesión exacta de homología
\[
  T_p(A)
  \longrightarrow T_p(k)
  \longrightarrow T_{p-1}(A)
  \xrightarrow{h_f}T_{p-1}(A)
\]

Se ve que $T_{p-1}(A)$ es libre si y solo si
$T_p(A)\to T_p(k)$ es sobreyectivo; la conclusión resulta entonces de (7.5.2).
\end{proof}

\begin{remark}[7.5.7]
\label{III.7.5.7}
En virtud de \sref{III.7.4.7}, las hipótesis generales de
\sref{III.7.4.4} implican que el funtor homológico $T_\bullet$
definido por \eqref{III.7.4.1.1.eq} satisface las hipótesis generales
de \sref{III.7.5.4}. En este caso, \sref{III.7.5.6} está, por tanto,
contenido en \sref{III.7.4.3}.
\end{remark}

\subsection{Descenso de las propiedades de exactitud. Teorema de
semicontinuidad y criterio de exactitud de Grauert}
\label{subsection:III.7.6}
\label{subsection:III.7.6.fr}
\label{III.7.6}

\begin{proposition}[7.6.1]
\label{III.7.6.1}
Bajo las hipótesis de \sref{III.7.4.1}, sea $B$ una $A$-álgebra
conmutativa. Si $T_p$ es exacto por la derecha (resp.
exacto por la izquierda, exacto), lo mismo ocurre con $T_p^{(B)}$;
el recíproco es cierto cuando $B$ es un $A$-módulo fielmente plano.
\end{proposition}

\begin{proof}
La primera aserción es un caso particular de una aserción trivial de
\sref{III.7.1.3}. Recíprocamente, supongamos primero que $B$ es un
$A$-módulo plano. Entonces, para todo $A$-módulo $M$,
$H_\bullet(P_\bullet\otimes_A(M\otimes_A B))=(H_\bullet(P_\bullet\otimes_A M))\otimes_A B$,
lo cual también puede escribirse, para todo $p$,
\[
  T_p(M)\otimes_A B
  =
  T_p^{(B)}(M_{(B)})
  \tag{7.6.1.1}\label{III.7.6.1.1.fr}
\]
salvo isomorfismo canónico. Supongamos que $T_p^{(B)}$ es exacto por
la derecha (resp. exacto por la izquierda, exacto). Como
$M\leadsto M_{(B)}$ es un funtor exacto, el miembro izquierdo de
\eqref{III.7.6.1.1} es un funtor exacto por la derecha
(resp. exacto por la izquierda, exacto) en $M$. Si ahora
$B$ es fielmente plano sobre $A$, se sigue que $T_p$ tiene la misma propiedad de exactitud ($\mathbf{0}_{\mathrm I}$, 6.4.1).
\end{proof}

\begin{proposition}[7.6.2]
\label{III.7.6.2}
Bajo las hipótesis de \sref{III.7.4.1}, se supone además que $A$ es
un anillo noetheriano reducido y que los $P_i$ son $A$-módulos de
tipo finito. Para que $T_p$ sea exacto por la derecha \textup{(resp.} exacto por la izquierda, exacto), es necesario y suficiente
que, para toda $A$-álgebra $B$ que es un anillo de valuación discreta,
$T_p^{(B)}$ lo sea.
\end{proposition}

\begin{proof}
En virtud de (7.3.1) y (7.3.3), se puede limitar a considerar la exactitud por la
derecha, y naturalmente solo hay que probar la suficiencia de la condición
(7.6.1). En virtud de (7.4.2), basta mostrar que
$Z'_{p-1}(P_\bullet)$ es un
$A$-módulo plano. Como $P_{p-1}$ es de tipo finito,
$Z'_{p-1}(P_\bullet)$ es también de tipo finito; el criterio ($\mathbf{0}$, 10.2.8) muestra
que basta entonces que $Z'_{p-1}(P_\bullet)\otimes_A B$
sea un $B$-módulo plano para toda $A$-álgebra $B$ que sea un anillo
de valuación discreta. Como $P_\bullet$ es un complejo de
$A$-módulos planos, se tiene
\[
  Z'_{p-1}(P_\bullet)\otimes_A B
  =
  Z'_{p-1}(P_\bullet\otimes_A B);
\]

\oldpage[III]{193}
El complejo $P_\bullet\otimes_A B$ está formado por $B$-módulos planos
($\mathbf{0}_{\mathrm I}$, 6.2.1), y para todo $B$-módulo $N$ se tiene
$H_\bullet(P_\bullet\otimes_A N)=H_\bullet((P_\bullet\otimes_A B)\otimes_B N)$,
y, por tanto,
$T_p^{(B)}(N)=H_p((P_\bullet\otimes_A B)\otimes_B N)$;
Aplicando \sref{III.7.4.2} a $T_p^{(B)}$, se ve que la hipótesis de
que $T_p^{(B)}$ es exacto por la derecha equivale a que
$Z'_{p-1}(P_\bullet\otimes_A B)$ sea un $B$-módulo plano.

El criterio precedente nos lleva a examinar más de cerca el caso de
los anillos de valuación discreta.
\end{proof}

\begin{proposition}[7.6.3]
Bajo las condiciones de \textup{(7.4.1)}, supongamos que $A$ es un anillo noetheriano
regular de dimensión 1 (en otros términos, que $A$ es noetheriano y que, para todo
$x\in\operatorname{Spec}(A)$, $A_x$ es un cuerpo o un anillo de valuación
discreta). Entonces, para todo entero $p$ y todo $A$-módulo $M$, se tiene una sucesión
exacta canónica funtorial en $M$
\[
  0
  \longrightarrow T_p(A)\otimes_A M
  \xrightarrow{t_M}T_p(M)
  \longrightarrow \operatorname{Tor}_1^A\bigl(T_{p-1}(A),M\bigr)
  \longrightarrow0.
  \tag{7.6.3.1}\label{III.7.6.3.1.fr}
\]
\end{proposition}

\begin{proof}
En lo que sigue, suprimiremos para simplificar la mención del complejo
$P_\bullet$ en las notaciones homológicas usuales $H_p(P_\bullet)$,
$B_p(P_\bullet)$, $Z_p(P_\bullet)$ y $Z'_p(P_\bullet)$. Se tienen las tres sucesiones
exactas
\begin{gather*}
  0\longrightarrow H_p\longrightarrow Z'_p\longrightarrow B_{p-1}
  \longrightarrow0\\
  0\longrightarrow B_{p-1}\longrightarrow Z_{p-1}\longrightarrow H_{p-1}
  \longrightarrow0\\
  0\longrightarrow Z_{p-1}\longrightarrow P_{p-1}\longrightarrow B_{p-2}
  \longrightarrow0
\end{gather*}

Como $P_{p-1}$ y $P_{p-2}$ son planos, lo mismo ocurre con sus
\emph{submódulos} respectivos $B_{p-1}$, $Z_{p-1}$ y $B_{p-2}$ puesto que hay
identidad entre $A_x$-módulos planos y $A_x$-módulos sin torsión (para todo
$x\in\operatorname{Spec}(A)$); por tensorización con $M$, se tienen pues las sucesiones
exactas
\[
  0=\operatorname{Tor}_1^A(B_{p-1},M)
  \longrightarrow H_p\otimes M
  \longrightarrow Z'_p\otimes M
  \xrightarrow{u}B_{p-1}\otimes M
  \longrightarrow0
  \tag{7.6.3.2}\label{III.7.6.3.2.fr}
\]
\[
  0=\operatorname{Tor}_1^A(Z_{p-1},M)
  \longrightarrow\operatorname{Tor}_1^A(H_{p-1},M)
  \longrightarrow B_{p-1}\otimes M
  \xrightarrow{v}Z_{p-1}\otimes M
  \tag{7.6.3.3}\label{III.7.6.3.3.fr}
\]
\[
  0=\operatorname{Tor}_1^A(B_{p-2},M)
  \longrightarrow Z_{p-1}\otimes M
  \xrightarrow{w}P_{p-1}\otimes M.
  \tag{7.6.3.4}\label{III.7.6.3.4.fr}
\]

Por definición,
$T_p(M)=\operatorname{Ker}(d_p\otimes1)/\operatorname{Im}(d_{p+1}\otimes1)$;
es pues el núcleo del homomorfismo
$(P_p\otimes M)/\operatorname{Im}(d_{p+1}\otimes1)\longrightarrow P_{p-1}\otimes M$
obtenido a partir de $d_p\otimes1$ por paso al cociente, homomorfismo que
se escribe también $Z'_p\otimes M\to P_{p-1}\otimes M$ por definición de
$Z'_p=P_p/B_p$; ahora bien, este homomorfismo puede considerarse como el compuesto
\[
  Z'_p\otimes M
  \xrightarrow{u}B_{p-1}\otimes M
  \xrightarrow{v}Z_{p-1}\otimes M
  \xrightarrow{w}P_{p-1}\otimes M.
\]

Como $w$ es inyectivo según (7.6.3.4), se tiene una sucesión exacta
\[
  0\longrightarrow\operatorname{Ker}u
  \longrightarrow T_p(M)
  \longrightarrow\operatorname{Ker}v
  \longrightarrow0,
\]
que no es otra que (7.6.3.1), teniendo en cuenta (7.6.3.2) y (7.6.3.3) y
que $H_p=T_p(A)$ por definición.
\end{proof}

\begin{remarks}[7.6.4]
\label{III.7.6.4}
\begin{enumerate}
\item[(i)]
El grupo $H_\bullet(P_\bullet\otimes_A M)$ es la homología del
bicomplejo $P_\bullet\otimes_A M$, donde $M$ se considera como un complejo reducido a su
    término de grado 0; es por tanto (6.3.6 y 6.3.2) el resultado de la
    sucesión espectral regular cuyo término $E_2$ es
\[
      E^2_{pq}
      =
      \operatorname{Tor}_p^A\bigl(H_q(P_\bullet),M\bigr)
      =
      \operatorname{Tor}_p^A\bigl(T_q(A),M\bigr).
    \]

\oldpage[III]{194}
La hipótesis sobre $A$ implica que $\operatorname{Tor}_p^A(E,F)=0$ para $p\geq2$
y para $A$-módulos
    arbitrarios ($\mathbf{0}_{\mathrm{IV}}$, 17.2.2); se sabe (M, XV) que esto implica
    la exactitud de la sucesión
    \[
      0\longrightarrow E^2_{0,q}
      \longrightarrow H_q(P_\bullet\otimes_A M)
      \longrightarrow E^2_{1,q-1}
      \longrightarrow0
    \]
que es precisamente \eqref{III.7.6.3.1}.

\medskip
\item[(ii)]
Teniendo en cuenta \sref{III.7.3.1}, la sucesión exacta
\eqref{III.7.6.3.1} recupera \sref{III.7.4.3} como caso particular.
\end{enumerate}
\end{remarks}

\begin{corollary}[7.6.5]
\label{III.7.6.5}
Bajo las hipótesis de \sref{III.7.4.1}, supongamos que $A$ es un
anillo de valuación discreta, con cuerpo de fracciones $K$ y cuerpo
residual $k$, y que los $T_i(A)$ son $A$-módulos de tipo finito.
Entonces
\[
  \operatorname{rang}_k T_p(k)
  \geq
  \operatorname{rang}_k\bigl(T_p(A)\otimes_A k\bigr)
  \geq
  \operatorname{rang}_A T_p(A)
  =
  \operatorname{rang}_K T_p(K).
  \tag{7.6.5.1}\label{III.7.6.5.1.fr}
\]
Además, los dos términos extremos de esta desigualdad son iguales si
y solo si $T_p$ es exacto o, equivalentemente, si y solo si $T_p(A)$
y $T_{p-1}(A)$ son $A$-módulos libres.
\end{corollary}

\begin{proof}
Hagamos, en efecto, $M=k$ en la sucesión exacta \eqref{III.7.6.3.1}. Como los
términos son espacios vectoriales sobre $k$, se obtiene
\[
  \operatorname{rang}_k T_p(k)
  =
  \operatorname{rang}_k\bigl(T_p(A)\otimes_A k\bigr)
  +
  \operatorname{rang}_k\bigl(\operatorname{Tor}_1^A(T_{p-1}(A),k)\bigr).
\]
Por otra parte, como $T_p(A)$ es un módulo de tipo finito sobre el
anillo local íntegro $A$, se tiene (Bourbaki,
\emph{Alg\`ebre commutative}, cap.~II, \S~3,
n\textsuperscript{o}~2, cor.~1 de la prop.~4)
\[
  \operatorname{rang}_k\bigl(T_p(A)\otimes_A k\bigr)
  \geq
  \operatorname{rang}_A T_p(A)
  =
  \operatorname{rang}_K\bigl(T_p(A)\otimes_A K\bigr)
  \tag{7.6.5.2}\label{III.7.6.5.2.fr}
\]
y los dos miembros de \eqref{III.7.6.5.2} son iguales si y solo si
$T_p(A)$ es un $A$-módulo libre (\emph{loc. cit.}, prop. 7). Se observará además que, puesto que $K$ es un $A$-módulo plano, también se tiene, por definición,
$T_p(A)\otimes_A K=H_p(P_\bullet)\otimes_A K=H_p(P_\bullet\otimes_A K)=T_p(K)$.
Esto demuestra \eqref{III.7.6.5.1}. También muestra que la igualdad
solo puede darse si, en primer lugar, $T_p(A)$ es libre y, en segundo
lugar,
$\operatorname{Tor}_1^A(T_{p-1}(A),k)=0$, condición que
equivale, como se sabe ($\mathbf{0}$, 10.1.3), al hecho de que $T_{p-1}(A)$ es un $A$-módulo libre. Finalmente, como los $T_i(A)$ son $A$-módulos de tipo finito,
equivale a decir que son planos o libres (Bourbaki,
\emph{Alg. comm.}, cap. II, \S 3, n\textsuperscript{o} 2, cor. 2 de la prop.
5), y se concluye por (7.4.3).
\end{proof}

\begin{env}[7.6.6]
Manteniéndose las hipótesis de (7.4.1), pondremos, para todo $x\in\operatorname{Spec}(A)$ \[
  d_p(x)=d_p^T(x)=\operatorname{rang}_{\mathbf{k}(x)}T_p(\boldsymbol{k}(x)).
  \tag{7.6.6.1}\label{III.7.6.6.1.fr}
\]
\end{env}

\begin{lemma}[7.6.7]
\label{III.7.6.7}
Sea $\varphi:A\to A'$ un homomorfismo de anillos, y sea
\[
  f={}^a\!\varphi:
  \operatorname{Spec}(A')\longrightarrow\operatorname{Spec}(A)
\]
la aplicación correspondiente \textup{\sref[I]{I.1.2.1}}. Si
$T'_\bullet=T_\bullet^{(A')}$ \textup{\sref{III.7.1.3}}, entonces
\[
  d_p^{T'}=d_p^T\circ f.
  \tag{7.6.7.1}\label{III.7.6.7.1.fr}
\]
\end{lemma}

\oldpage[III]{195}
\begin{proof}
Para $x'\in\operatorname{Spec}(A')$, se tiene, poniendo $x=f(x')$. Entonces
\[
  H_\bullet(P_\bullet\otimes_A\boldsymbol{k}(x'))
  =
  H_\bullet\bigl((P_\bullet\otimes_A\boldsymbol{k}(x))
    \otimes_{\mathbf{k}(x)}\boldsymbol{k}(x')\bigr)
  =
  H_\bullet(P_\bullet\otimes_A\boldsymbol{k}(x))
    \otimes_{\mathbf{k}(x)}\boldsymbol{k}(x'),
\]
porque $\boldsymbol{k}(x')$ es plano sobre $\boldsymbol{k}(x)$. Esto demuestra
\eqref{III.7.6.7.1}.
\end{proof}

\begin{lemma}[7.6.8]
\label{III.7.6.8}
Si $A$ es noetheriano y $P_\bullet$ está formado por $A$-módulos de
tipo finito, la función $x\leadsto d_p^T(x)$ sobre $\operatorname{Spec}(A)$ es construible.
\end{lemma}

\begin{proof}
Hay que probar que para toda parte cerrada irreducible $Y$ de $X=\operatorname{Spec}(A)$, existe un abierto no vacío $U$ de $Y$ en el cual $d_p$ es constante ($\mathbf{0}$, 9.2.2); como
$Y=\operatorname{Spec}(A/\mathfrak a)$, donde $\mathfrak a$ es un ideal de $A$ tal
que $A/\mathfrak a$ sea reducido, se puede, en virtud de (7.6.7), limitar al caso en que $Y=X$ y $A$ es un anillo noetheriano
\emph{íntegro}; la aserción resulta entonces de \sref{III.7.4.5}.
\end{proof}

\begin{theorem}[7.6.9]
\label{III.7.6.9}
Sean $A$ un anillo noetheriano, $P_\bullet$ un complejo de
$A$-módulos planos de tipo finito y
$T_\bullet(M)=H_\bullet(P_\bullet\otimes_A M)$
el funtor homológico definido por $P_\bullet$. Para
$x\in\operatorname{Spec}(A)$, sea $d_p(x)=\operatorname{rang}_{\mathbf{k}(x)}T_p(\boldsymbol{k}(x))$.
Entonces:
\begin{enumerate}
  \item[\textup{(i)}] La función $d_p$ es construible y semicontinua
    superiormente sobre $\operatorname{Spec}(A)$.
  \item[\textup{(ii)}] Si $T_p$ es exacto, $d_p$ es continua (y, por tanto,
    localmente constante) sobre $\operatorname{Spec}(A)$. El recíproco es cierto
    cuando $A$
    es reducido.
\end{enumerate}
\end{theorem}

\begin{proof}
\begin{enumerate}
  \item[(i)] La primera afirmación se ha demostrado en (7.6.8). Para probar
    la segunda, basta pues ($\mathbf{0}$, 9.2.4) mostrar que si $x'\ne x$ es una generalización de $x$ en
$\operatorname{Spec}(A)$, entonces $d_p(x')\leq d_p(x)$. Existe un anillo de
valuación discreta $B$ y un morfismo
$f:\operatorname{Spec}(B)\longrightarrow\operatorname{Spec}(A)$
tales que, si $a$ es el punto cerrado de $\operatorname{Spec}(B)$ y $b$ su punto
genérico, entonces $f(a)=x$ y $f(b)=x'$ (\textbf{II}, 7.1.9). En virtud de la fórmula (7.6.7.1), se ve que se reduce a demostrar la desigualdad $d_p(a)\geq d_p(b)$ sobre $\operatorname{Spec}(B)$; pero esto no es sino la desigualdad (7.6.5.1)\footnote{El principio de demostración de (i) por reducción al caso de un anillo de valuación discreta nos fue comunicado oralmente por Hironaka.}.
\item[(ii)] La primera aserción ya ha sido demostrada (7.3.4). Para demostrar la recíproca, utilicemos el criterio valuativo (7.6.2); teniendo en cuenta la fórmula (7.6.7.1), se reduce pues al caso en que $A$ es un anillo
de valuación discreta. Como $\operatorname{Spec}(A)$ tiene entonces solamente dos
puntos, la hipótesis de que $d_p$ es constante implica que $T_p$ es
exacto, por \sref{III.7.6.5}.
\end{enumerate}
\end{proof}

\begin{corollary}[7.6.10]
\label{III.7.6.10}
Sean $A$ un anillo noetheriano, $\mathfrak p_i$ $(1\leq i\leq r)$
sus ideales primos minimales, y $k_i$ el cuerpo residual de
$A_{\mathfrak p_i}$
$(1\leq i\leq r)$.
\begin{enumerate}
  \item[\textup{(i)}] Para todo $x\in\operatorname{Spec}(A)$, existe un índice $i$ tal que
    \[
      d_p(x)\geq\operatorname{rang}_{k_i}T_p(k_i).
      \tag{7.6.10.1}\label{III.7.6.10.1.fr}
    \]
\end{enumerate}
En particular, si $A$ es un anillo íntegro y $K$ es su cuerpo de
fracciones, entonces
\[
  d_p(x)\geq\operatorname{rang}_K T_p(K)
  \tag{7.6.10.2}\label{III.7.6.10.2.fr}
\]
para todo $x\in\operatorname{Spec}(A)$.

\oldpage[III]{196}
\begin{enumerate}
  \item[\textup{(ii)}] Supongamos además que $A$ es local y reducido, y sea
    $k$ su cuerpo residual. Entonces $T_p$ es exacto si y solo si
    \[
      \operatorname{rang}_k T_p(k)
      =
      \operatorname{rang}_{k_i}T_p(k_i)
      \qquad\text{para }1\leq i\leq r.
      \tag{7.6.10.3}\label{III.7.6.10.3.fr}
    \]
\end{enumerate}
\end{corollary}

\begin{proof}
La aserción \textnormal{(i)} es inmediata: toda vecindad de $x$
contiene uno de los $\mathfrak p_i$, y basta aplicar la definición de
semicontinuidad. Si $A$ es local, la única vecindad en $\operatorname{Spec}(A)$ del ideal
maximal $\mathfrak m$ es $\operatorname{Spec}(A)$; por tanto,
$d_p(x)\leq\operatorname{rang}_k T_p(k)$
para todo $x\in\operatorname{Spec}(A)$. Junto con \textnormal{(i)}, esto muestra que
\eqref{III.7.6.10.3} implica que $d_p(x)$ es constante sobre
$\operatorname{Spec}(A)$ y, por tanto, que $T_p$ es exacto por
\sref{III.7.6.9}[(ii)]. El recíproco se deduce inmediatamente del
mismo resultado.
\end{proof}

\begin{remark}[7.6.11]
Cabe preguntarse si la afirmación de (7.6.9, (i)) no puede reforzarse con
la desigualdad
\[
  \operatorname{rang}_{\mathbf{k}(x)}T_p(\boldsymbol{k}(x))
  \geq
  \operatorname{rang}_{\mathbf{k}(x)}
    \bigl(T_p(A)\otimes_A\boldsymbol{k}(x)\bigr)
  \tag{7.6.11.1}\label{III.7.6.11.1.fr}
\]
para todo $x\in\operatorname{Spec}(A)$, que se verifica efectivamente cuando $A$
es un anillo de valuación discreta y $x$ su ideal maximal (7.6.5).
Limitémonos al caso en que $A$ es un anillo \emph{local} noetheriano, de ideal maximal
$\mathfrak m$ y de cuerpo residual $k$. Entonces las condiciones siguientes son
\emph{equivalentes}:
\begin{enumerate}
  \item[\emph{a)}] Para \emph{todo} complejo $P_\bullet$ de $A$-módulos planos de tipo finito,
    se tiene
    \[
      \operatorname{rang}_k(T_p(k))
      \geq
      \operatorname{rang}_k\bigl(T_p(A)\otimes_A k\bigr)
      \quad\text{para todo entero }p.
      \tag{7.6.11.2}\label{III.7.6.11.2.fr}
    \]
  \item[\emph{b)}] Para todo $A$-módulo $M$ de tipo finito, se tiene
    \[
      \operatorname{rang}_k(M\otimes_A k)
      \geq
      \operatorname{rang}_k(\check M\otimes_A k).
      \tag{7.6.11.3}\label{III.7.6.11.3.fr}
    \]
  \item[\emph{c)}] Para todo $A$-módulo $N$ de tipo finito, se tiene
    \[
      \operatorname{rang}_k\bigl(\operatorname{Tor}_1^A(N,k)\bigr)
      \geq
      \operatorname{rang}_k\bigl(\operatorname{Tor}_2^A(N,k)\bigr).
      \tag{7.6.11.4}\label{III.7.6.11.4.fr}
    \]
\end{enumerate}

Se observará que equivale, por desplazamiento (M, V, 7.2), a decir que se tiene,
para todo $i\geq1$,
\[
  \operatorname{rang}_k\bigl(\operatorname{Tor}_i^A(N,k)\bigr)
  \geq
  \operatorname{rang}_k\bigl(\operatorname{Tor}_{i+1}^A(N,k)\bigr).
  \tag{7.6.11.5}\label{III.7.6.11.5.fr}
\]

Demos brevemente algunas indicaciones sobre la demostración. Para ver que \emph{a)}
implica \emph{b)}, se considera una sucesión exacta
$L_1\xrightarrow{d}L_0\longrightarrow M\longrightarrow0$
donde $L_0$ y $L_1$ son libres de tipo finito, y se aplica \emph{a)} al complejo
$P_1\xrightarrow{{}^t d}P_0$ con $P_0=\check L_1$, $P_1=\check L_0$,
siendo nulos los demás términos; se tiene entonces $T_1(A)=\check M$ y
$T_1(k)=\operatorname{Hom}_A(M,k)=\operatorname{Hom}_k(M/\mathfrak mM,k)$,
en otros términos, $T_1(k)$ es el dual del espacio vectorial $M\otimes_A k$, y tiene
pues el mismo rango que este último. Para probar que \emph{b)} implica \emph{c)}, estableceremos
primero el lema siguiente:

\begin{lemma}[7.6.11.6]
Dado un complejo
$\ldots\longrightarrow0\longrightarrow P_1\xrightarrow{d}P_0\longrightarrow0\longrightarrow\ldots$
de $A$-módulos planos, se tiene una sucesión exacta
\[
  0\longrightarrow\operatorname{Tor}_2^A(Z'_0,k)
  \longrightarrow T_1(A)\otimes_A k
  \longrightarrow T_1(k)
  \longrightarrow\operatorname{Tor}_1^A(Z'_0,k)
  \longrightarrow0.
  \tag{7.6.11.7}\label{III.7.6.11.7.fr}
\]
\end{lemma}

\begin{proof}
\oldpage[III]{197}
En efecto, partiendo de la sucesión exacta
$0\to Z_1\to P_1\to B_0\to0$, se obtiene la sucesión exacta
$0\longrightarrow\operatorname{Tor}_1^A(B_0,k)\longrightarrow Z_1\otimes k\longrightarrow P_1\otimes k\xrightarrow{u}B_0\otimes k\longrightarrow0$.
De la sucesión exacta $0\to B_0\to Z_0\to Z'_0\to0$ se obtiene, puesto que $Z_0=P_0$
es plano,
$\operatorname{Tor}_1^A(B_0,k)=\operatorname{Tor}_2^A(Z'_0,k)$; por definición,
se tiene $Z_1=T_1(A)$; finalmente, se tiene $T_1(k)=\operatorname{Ker}(d\otimes1)$, y
$d\otimes1$ se factoriza en
$P_1\otimes k\xrightarrow{u}B_0\otimes k\xrightarrow{v}Z_0\otimes k=P_0\otimes k$;
se tiene $T_1(k)=u^{-1}(R)$, donde $R=\operatorname{Ker}v$, y como $u$ es
sobreyectivo, $R=u(T_1(k))$; finalmente,
$R=\operatorname{Tor}_1^A(Z'_0,k)$ por definición de $v$, lo que termina
de establecer la sucesión exacta (7.6.11.7).
\end{proof}

Para deducir entonces \emph{c)} de \emph{b)}, se considera una sucesión exacta
$L_1\xrightarrow{d}L_0\longrightarrow N\longrightarrow0$,
donde $L_0$ y $L_1$ son módulos libres de tipo finito; consideremos el funtor
$T$ asociado al complejo formado por $L_1$ y $L_0$; como $L_0$ y $L_1$ son
libres, se identifican con sus biduales; luego si
$M=\operatorname{Coker}({}^t d)$, $T_1(A)=\operatorname{Ker}(d)=\check M$;
por otra parte, $M\otimes_A k=\operatorname{Coker}({}^t d\otimes1_k)$ tiene el mismo rango
sobre $k$ que $\operatorname{Ker}(d\otimes1_k)$. La hipótesis \emph{b)} implica por tanto
que
\[
  \operatorname{rang}_k\bigl(T_1(A)\otimes_A k\bigr)
  \leq
  \operatorname{rang}_k\bigl(T_1(k)\bigr);
\]
como $Z'_0=N$, la desigualdad (7.6.11.4) resulta pues de la sucesión exacta
(7.6.11.7).
Finalmente, para probar que \emph{c)} implica \emph{a)}, apliquemos (7.6.11.6) reemplazando
$P_0$ y $P_1$ por $P_p$ y $P_{p-1}$; la hipótesis \emph{c)} aplicada al módulo
$Z'_p$ da $\operatorname{rang}_kR\geq\operatorname{rang}_kS$, donde
$R=\operatorname{Ker}(P_p\otimes k\longrightarrow P_{p-1}\otimes k)$ y $S=Z_p\otimes k$.
Ahora bien, si se factoriza $d_{p+1}:P_{p+1}\to P_p$ en
$P_{p+1}\xrightarrow{v}Z_p\xrightarrow{j}P_p$,
se tiene
$\operatorname{Im}(d_{p+1}\otimes1)=(j\otimes1)\bigl(\operatorname{Im}(v\otimes1)\bigr)$.
Como
\[
  T_p(A)\otimes_A k
  =(Z_p/B_p)\otimes_A k
  =(Z_p\otimes k)/\operatorname{Im}(v\otimes1),
\]
y
$T_p(k)=R/\operatorname{Im}(d_{p+1}\otimes1)$,
se concluye efectivamente la desigualdad (7.6.11.2).

Dicho esto, supongamos que el anillo local $A$ es \emph{regular} de dimensión $n$;
se sabe entonces [17] que el $A$-módulo
$\operatorname{Tor}_i^A(k,k)$ es isomorfo a la potencia exterior
$\bigwedge\nolimits^i(\mathfrak m/\mathfrak m^2)$;
se ve pues que la condición (7.6.11.4) no se verifica para $N=k$ en cuanto
$n\geq4$.
En cambio, si el anillo local íntegro $A$ es tal que todo $A$-módulo \emph{reflexivo}
de tipo finito es \emph{libre} (lo que ocurre cuando $A$ es un anillo regular
de dimensión 2), la condición (7.6.11.3) se verifica: en efecto, se sabe que el
dual $\check M$ de un $A$-módulo $M$ de tipo finito es reflexivo, luego libre, y
por tanto
$\operatorname{rang}_k(\check M\otimes_A k)=\operatorname{rang}_K(\check M)=\operatorname{rang}_K(M)$
($K$ cuerpo de fracciones de $A$); por otra parte, se sabe que toda base sobre $k$
de $M\otimes_A k$ está formada por imágenes de un sistema de generadores de $M$
(Bourbaki, \emph{Alg. comm.}, cap. II, \S 3, n\textsuperscript{o} 2, cor. 2 de
la prop. 4), luego
$\operatorname{rang}_K(M)\leq\operatorname{rang}_k(M\otimes_A k)$, lo que
prueba nuestra afirmación.
\end{remark}


\subsection{Aplicación a los morfismos propios: I. La propiedad de intercambio.}
\label{subsection:III.7.7.fr}

Los tres números que siguen son, en lo esencial, traducciones, al
lenguaje de los morfismos de preesquemas, de los resultados de los números anteriores.

\begin{env}[7.7.1]
\label{III.7.7.1}
Sea $f:X\to Y$ un morfismo cuasicompacto y separado de preesquemas, y sea
$\mathscr P_\bullet$ un complejo de $\mathscr O_X$-módulos cuasicoherentes cuyo
operador de derivación es de grado $-1$; supongamos además que los
$\mathscr O_X$-módulos $\mathscr P_i$ son \emph{planos} sobre $Y$-\emph{planos} ($\mathbf{0}_{\mathrm I}$, 6.7.1).
\oldpage[III]{198}
Vamos a considerar el $\partial$-funtor
$\mathscr M\leadsto\mathscr T_\bullet(\mathscr M)$ (también denotado por
$\mathscr T_\bullet(\mathscr P_\bullet,\mathscr M)$) de la categoría de los
$\mathscr O_Y$-módulos cuasicoherentes en la categoría de los
$\mathscr O_Y$-módulos cuasicoherentes, en virtud de \sref{III.6.2.3}, definido
por
\[
  \mathscr T_n(\mathscr P_\bullet,\mathscr M)
  =\mathscr T_n(\mathscr M)
  =\mathscr H^{-n}
    \bigl(f,\mathscr P^\bullet\otimes_{\mathscr O_Y}\mathscr M\bigr)
  \qquad\text{para }n\in\mathbf Z,
  \tag{7.7.1.1}\label{III.7.7.1.1.fr}
\]
donde $\mathscr P^\bullet$ es el complejo cuyo término de grado $j$ es
$\mathscr P_{-j}$, de modo que su operador de derivación es entonces de grado
$+1$. El funtor $\mathscr T_\bullet$ así definido es un \emph{funtor
homológico} en $\mathscr M$ \sref{III.6.2.6}.
\end{env}

\begin{env}[7.7.2]
\label{III.7.7.2}
Sea $g:Y'\to Y$ un morfismo, y pongamos
$X'=X_{(Y')}=X\times_Y Y'$ y $f'=f_{(Y')}:X'\longrightarrow Y'$,
que es un morfismo cuasicompacto y separado; sea por otra parte
$\mathscr P'_\bullet=\mathscr P_\bullet\otimes_{\mathscr O_Y}\mathscr O_{Y'}$;
Este es un complejo de $\mathscr O_{X'}$-módulos cuasicoherentes que son \emph{planos}
sobre $Y'$-\emph{planos} en virtud de (\textbf{I}, 9.1.12) y ($\mathbf{0}_{\mathrm I}$, 6.2.1). Pondremos
(con las mismas convenciones sobre los grados)
\[
  \mathscr T^{Y'}_\bullet(\mathscr M')
  =\mathscr H^\bullet
    \bigl(f',\mathscr P'^\bullet\otimes_{\mathscr O_{Y'}}\mathscr M'\bigr)
  =\mathscr H^\bullet
    \bigl(f',\mathscr P^\bullet\otimes_{\mathscr O_Y}\mathscr M'\bigr)
  \tag{7.7.2.1}\label{III.7.7.2.1.fr}
\]
que es un funtor homológico en el $\mathscr O_{Y'}$-módulo cuasicoherente
$\mathscr M'$. Cuando $Y'$ sea un esquema afín de anillo $A'$, se escribirá
$\mathscr T^{A'}_\bullet$ en lugar de $\mathscr T^{Y'}_\bullet$; para todo
$A'$-módulo $M'$ se tiene entonces
$\mathscr T^{A'}_\bullet(\widetilde{M}')=\bigl(\Gamma(Y',\mathscr T^{Y'}_\bullet(\widetilde{M}'))\bigr)^\sim$;
se pondrá
$T^{A'}_\bullet(M')=\Gamma(Y',\mathscr T^{Y'}_\bullet(\widetilde{M}'))$,
que es un funtor homológico de $A'$-módulos con valores en la categoría de los
$A'$-módulos.

Se observará que, si $Y=\operatorname{Spec}(A)$ es también afín, el funtor de
$A'$-módulos $T^{A'}_\bullet$ coincide con el funtor obtenido por extensión de
escalares de $A$ a $A'$ a partir del funtor homológico de $A$-módulos
$T^A_\bullet$ \sref{III.7.1.3}. En efecto, sea $g:Y'\to Y$ el morfismo
correspondiente al homomorfismo de anillos $A\to A'$, y sea $g':X'\to X$ el
morfismo correspondiente, que es afín \sref[II]{II.1.6.2}; si $\mathfrak U$ es
un recubrimiento abierto afín de $X$, entonces
$\mathfrak U'=g'^{-1}(\mathfrak U)$ es un recubrimiento abierto afín de $X'$.
En virtud de \sref{III.6.2.2}, todo se reduce a ver que
$C^\bullet\bigl(\mathfrak U,\mathscr P_\bullet\otimes_{\mathscr O_Y}g_*(\mathscr M')\bigr)=C^\bullet\bigl(\mathfrak U',\mathscr P_\bullet\otimes_{\mathscr O_Y}\mathscr M'\bigr)$,
y, finalmente, que para todo abierto afín $U$ de $X$, poniendo
$U'=g'^{-1}(U)$, se tiene
$\Gamma\bigl(U,\mathscr P_\bullet\otimes_{\mathscr O_Y}g_*(\mathscr M')\bigr)=\Gamma\bigl(U',\mathscr P_\bullet\otimes_{\mathscr O_Y}\mathscr M'\bigr)$,
lo cual es trivial \sref[I]{I.1.3}, \sref[I]{I.3.2}.

En particular, si $U$ es un abierto de $Y$, para todo
$\mathscr O_Y$-módulo cuasicoherente $\mathscr M$ se tiene
\[
  \mathscr T^U_\bullet(\mathscr M|U)
  =\bigl(\mathscr T_\bullet(\mathscr M)\bigr)|U.
  \tag{7.7.2.2}\label{III.7.7.2.2.fr}
\]
\end{env}

\begin{env}[7.7.3]
\label{III.7.7.3}
Para todo $\mathscr O_Y$-módulo cuasicoherente $\mathscr M$ existe un
homomorfismo canónico, funtorial en $\mathscr M$,
\[
  \mathscr T_p(\mathscr O_Y)\otimes_{\mathscr O_Y}\mathscr M
  \longrightarrow
  \mathscr T_p(\mathscr M).
  \tag{7.7.3.1}\label{III.7.7.3.1.fr}
\]
En efecto, si $Y$ es afín, este homomorfismo se definió en
\sref{III.7.2.2}; esta definición se extiende sin dificultad al caso general,
observando que, si $U$ y $V$ son dos abiertos afines de $Y$ tales que
$V\subset U$, el diagrama

\oldpage[III]{199}
\[
\xymatrix@C=18pt@R=32pt{
  {\begin{gathered}
    \bigl(\mathscr T_p(\mathscr O_Y)
      \otimes_{\mathscr O_Y}\mathscr M\bigr)|U\\
    =\mathscr T^U_p(\mathscr O_Y|U)
      \otimes_{\mathscr O_Y|U}(\mathscr M|U)
  \end{gathered}}
  \ar[r]\ar[d]
  &
  {\begin{gathered}
    \mathscr T^U_p(\mathscr M|U)\\
    =\bigl(\mathscr T_p(\mathscr M)\bigr)|U
  \end{gathered}}
  \ar[d]
  \\
  {\begin{gathered}
    \bigl(\mathscr T_p(\mathscr O_Y)
      \otimes_{\mathscr O_Y}\mathscr M\bigr)|V\\
    =\mathscr T^V_p(\mathscr O_Y|V)
      \otimes_{\mathscr O_Y|V}(\mathscr M|V)
  \end{gathered}}
  \ar[r]
  &
  {\begin{gathered}
    \mathscr T^V_p(\mathscr M|V)\\
    =\bigl(\mathscr T_p(\mathscr M)\bigr)|V
  \end{gathered}}
}
\]
es conmutativo por \sref{III.7.2.3.3}.

Para todo morfismo $g:Y'\to Y$ existe un homomorfismo canónico
\[
  \mathscr T_p(\mathscr O_Y)\otimes_{\mathscr O_Y}\mathscr O_{Y'}
  \longrightarrow
  \mathscr T^{Y'}_p(\mathscr O_{Y'})
  \tag{7.7.3.2}\label{III.7.7.3.2.fr}
\]
que no es sino el caso particular de \sref{III.6.7.11.2}, para los objetos
límite, en el caso en que $S=Y$, $v_1=f$, $v_2=1_Y$ y
$\mathscr P^{(2)}_\bullet$ se reduce al único término $\mathscr M$ de grado 0.

Cuando $Y=\operatorname{Spec}(A)$ e $Y'=\operatorname{Spec}(A')$ son afines,
\eqref{III.7.7.3.2} no es sino el homomorfismo de haces correspondiente al
homomorfismo canónico de $A'$-módulos definido en \sref{III.7.2.2},
\[
  T^A_p(A)\otimes_A A'
  \longrightarrow
  T^{A'}_p(A')=T^A_p(A')
\]
como resulta fácilmente de \sref{III.6.7.11}, pues en el caso considerado se
puede tomar $\mathfrak U'^{(i)}=u_i^{-1}(\mathfrak U^{(i)})$.
\end{env}

\begin{env}[7.7.4]
\label{III.7.7.4}
Cuando $f$ es un morfismo \emph{propio}, $Y=\operatorname{Spec}(A)$ es un
esquema afín noetheriano y $\mathscr P_\bullet$ es un complejo acotado
inferiormente de $\mathscr O_X$-módulos coherentes y planos sobre $Y$-planos \emph{acotado inferiormente}, se ha
visto (6.10.5) que se puede escribir, salvo isomorfismo,
$\mathscr T_p(\mathscr M)=\mathscr H_p\bigl(\mathscr L_\bullet\otimes_{\mathscr O_Y}\mathscr M\bigr)$,
con $\mathscr L_\bullet=\widetilde{L}_\bullet$,
donde $L_\bullet$ es un complejo acotado inferiormente de $A$-módulos libres
de tipo finito; por tanto, el funtor $\mathscr T_\bullet$ es del tipo estudiado
en detalle en \sref{III.7.4} y \sref{III.7.6}. Traduciremos ahora los resultados
de este estudio.
\end{env}

\begin{theorem}[7.7.5]
\label{III.7.7.5}
Sean $Y$ un preesquema localmente noetheriano, $(U_\alpha)$ un recubrimiento
de $Y$ formado por abiertos afines, $f:X\to Y$ un morfismo propio y
$\mathscr P_\bullet$ un complejo acotado inferiormente de
$\mathscr O_X$-módulos coherentes y planos sobre $Y$. El funtor homológico
$\mathscr T_\bullet(\mathscr M)$ definido por \eqref{III.7.7.1.1} posee
entonces las propiedades siguientes:

\medskip
\noindent
\textnormal{I) (La propiedad de semicontinuidad).}\footnote{Un caso particular
de este teorema se encuentra ya en la nota [3] de Chow--Igusa. La propiedad de
semicontinuidad fue descubierta, en el marco de los espacios analíticos y bajo
hipótesis bastante particulares, por Kodaira--Spencer (``On the variations of
almost-complex structures'', \emph{Algebraic Geometry and Topology, A Symposium
in honor of S. Lefschetz}, Princeton Series n\textsuperscript{o} 12,
p.~139-150, Princeton, 1957), y la versión general fue demostrada por Grauert
[5].}.
La función
\[
  y\leadsto
  d_p(y)
  =
  \operatorname{rang}_{\mathbf{k}(y)}
    T_p^{\mathbf{k}(y)}(\boldsymbol{k}(y))
  \tag{7.7.5.1}\label{III.7.7.5.1.fr}
\]
es semicontinua superiormente.

\oldpage[III]{200}
\medskip
\noindent
\textnormal{II) (La propiedad de cambio de base).}
Para un entero $p$ dado, las condiciones siguientes son equivalentes:
\begin{enumerate}
  \item[\textup{a)}] $\mathscr T_p$ es exacto por la derecha.

  \item[\textup{a$'$)}] $\mathscr T_p(\mathscr M)$ es isomorfo a un funtor de la forma
    $\mathscr N\otimes_{\mathscr O_Y}\mathscr M$, siendo $\mathscr N$
    necesariamente isomorfo a
    $\mathscr T_p(\mathscr O_Y)=\mathscr H^{-p}(f,\mathscr P^\bullet)$

  \item[\textup{a$''$)}] El homomorfismo canónico funtorial
    \eqref{III.7.7.3.1} es un isomorfismo.

  \item[\textup{b)}] $\mathscr T_{p-1}$ es exacto por la izquierda.

  \item[\textup{b$'$)}] Existe un $\mathscr O_Y$-módulo $\mathscr Q$ (necesariamente
    coherente y determinado salvo isomorfismo único) y un isomorfismo de
    funtores
    \[
      \mathscr T_{p-1}(\mathscr M)
      \xrightarrow{\sim}
      \mathscr{H}\!om_{\mathscr O_Y}(\mathscr Q,\mathscr M).
      \tag{7.7.5.2}\label{III.7.7.5.2.fr}
    \]

  \item[\textup{c)}] Denotando por $A_\alpha$ el anillo del abierto afín $U_\alpha$,
    para todo índice $\alpha$ el funtor de $A_\alpha$-módulos
    $T_p^{A_\alpha}$ es exacto por la derecha.

  \item[\textup{d)}] Para todo morfismo $g:Y'\to Y$, el homomorfismo canónico
    \[
      \mathscr T_p(\mathscr O_Y)
      \otimes_{\mathscr O_Y}\mathscr O_{Y'}
      \longrightarrow
      \mathscr T_p^{Y'}(\mathscr O_{Y'})
      \tag{7.7.5.3}\label{III.7.7.5.3.fr}
    \]
    es un isomorfismo.
\end{enumerate}
\end{theorem}

\begin{proof}
La propiedad de semicontinuidad es local sobre $Y$ y resulta, por tanto, de
\sref{III.7.7.4} y \sref{III.7.6.9}. Es claro que \emph{a$''$)} implica \emph{a$'$)} y que
\emph{a$'$)} implica \emph{a)}. La equivalencia de \emph{a)}, \emph{a$''$)}, \emph{b)} y \emph{b$'$)} se demostró en
\sref{III.7.3.1} y \sref{III.7.4.6}, habida cuenta de
\sref{III.7.7.4}, cuando $Y$ es \emph{afín}.

Para pasar al caso general, probemos primero que \emph{a)} equivale a \emph{c)}, lo cual
probará que la propiedad \emph{a)} es \emph{local} sobre $Y$; la demostración se aplica
igualmente para probar el carácter local de \emph{a$''$)} y \emph{b)}. Como es claro que \emph{c)}
implica \emph{a)}, todo se reduce a probar el recíproco. Basta evidentemente demostrar
que, para todo abierto afín $U$ de $Y$ y toda sucesión exacta
$0\longrightarrow\mathscr F'\longrightarrow\mathscr F\longrightarrow\mathscr F''\longrightarrow0$
de $(\mathscr O_Y|U)$-módulos cuasicoherentes, existe una sucesión exacta
$0\longrightarrow\mathscr G'\longrightarrow\mathscr G\longrightarrow\mathscr G''\longrightarrow0$
de $\mathscr O_Y$-Módulos cuasicoherentes tal que
$\mathscr F'=\mathscr G'|U$, $\mathscr F=\mathscr G|U$, $\mathscr F''=\mathscr G''|U$;
ahora bien, esto resulta inmediatamente de la hipótesis de que $Y$ es
localmente noetheriano y de \sref[I]{I.9.4.2}: en efecto, se prolonga
$\mathscr F$ a un $\mathscr O_Y$-módulo cuasicoherente $\mathscr G$,
$\mathscr F'$ a un sub-$\mathscr O_Y$-módulo $\mathscr G'$ de $\mathscr G$, y
basta tomar $\mathscr G''=\mathscr G/\mathscr G'$.

Para demostrar la equivalencia de \emph{b)} y \emph{b$'$)} en el caso general, observemos que,
cuando $Y$ es afín, se sabe que $\mathscr Q$ está determinado salvo isomorfismo
único; si entonces $U$ es un abierto afín del esquema afín $Y$, se deduce que
existe un isomorfismo funtorial
$\mathscr T_{p-1}^{U}(\mathscr M|U)\xrightarrow{\sim}\mathscr{H}\!om_{\mathscr O_Y|U}(\mathscr Q|U,\mathscr M|U)$.
En el caso general, para todo abierto afín $U$ de $Y$ existe un
$(\mathscr O_Y|U)$-módulo coherente $\mathscr Q_U$ y un isomorfismo funtorial
$\mathscr T_{p-1}^{U}(\mathscr M|U)\to\mathscr{H}\!om_{\mathscr O_Y|U}(\mathscr Q_U,\mathscr M|U)$;
La observación precedente muestra que, si $V$ es un abierto afín contenido en
$U$, se tiene $\mathscr Q_U|V=\mathscr Q_V$; de ahí la existencia y unicidad
del $\mathscr O_Y$-módulo $\mathscr Q$ que verifica \eqref{III.7.7.5.2}.

Queda por demostrar la equivalencia de \emph{a)} y \emph{d)}. Es claro que \emph{d)} es local sobre
$Y$, y se ha visto arriba que lo mismo ocurre con \emph{a)}; además, \emph{d)} es también local
sobre $Y'$. Ahora bien, cuando $Y=\operatorname{Spec}(A)$ e
$Y'=\operatorname{Spec}(A')$, se ha visto que $T_\bullet^{A'}$ es el funtor obtenido
\oldpage[III]{201}
a partir de $T_\bullet^A$ por extensión de escalares a $A'$, y es claro entonces
que \emph{a$'$)} implica que \eqref{III.7.7.5.3} es un isomorfismo.
Recíprocamente, supongamos todavía que $Y=\operatorname{Spec}(A)$ es afín y sea
$A'$ la $A$-álgebra $A\oplus M$, donde $M$ es un $A$-módulo cualquiera y la
multiplicación en $A'$ viene dada por
$(a_1,m_1)(a_2,m_2)=(a_1a_2,a_1m_2+a_2m_1)$;
Entonces
\[
  T_p^{A'}(A')
  =
  T_p(A\oplus M)
  =
  T_p(A)\oplus T_p(M),
\]
La hipótesis de que \eqref{III.7.7.5.3} sea biyectivo implica que lo mismo
ocurre con la aplicación canónica
$T_p(A)\otimes_A M\longrightarrow T_p(M)$,
En otras palabras, \emph{d)} implica \emph{a$''$)}, lo que termina la demostración.
\end{proof}

\begin{theorem}[7.7.6]
\label{III.7.7.6}
Sean $Y$ un preesquema localmente noetheriano, $f:X\to Y$ un morfismo propio y
$\mathscr F$ un $\mathscr O_X$-módulo coherente y plano sobre $Y$. Existe
entonces un $\mathscr O_Y$-módulo coherente $\mathscr Q$ (determinado salvo
isomorfismo único) y un isomorfismo de funtores en el
$\mathscr O_Y$-módulo cuasicoherente $\mathscr M$,
\[
  f_*(\mathscr F\otimes_{\mathscr O_Y}\mathscr M)
  \xrightarrow{\sim}
  \mathscr{H}\!om_{\mathscr O_Y}(\mathscr Q,\mathscr M)
  \tag{7.7.6.1}\label{III.7.7.6.1.fr}
\]
\textup{(de donde un isomorfismo de funtores}
\[
  \Gamma(X,\mathscr F\otimes_{\mathscr O_Y}\mathscr M)
  \xrightarrow{\sim}
  \operatorname{Hom}_{\mathscr O_Y}(\mathscr Q,\mathscr M).)
  \tag{7.7.6.2}\label{III.7.7.6.2.fr}
\]
\end{theorem}

\begin{proof}
En efecto, como $\mathscr M\leadsto\mathscr F\otimes_{\mathscr O_Y}\mathscr M$ es exacto
\sref[0]{0.6.7.4} y $f_*$ es exacto por la izquierda, el funtor
$\mathscr M\leadsto f_*(\mathscr F\otimes_{\mathscr O_Y}\mathscr M)$ es exacto por la izquierda. Basta entonces aplicar la equivalencia de (7.7.5, \emph{b)}) y (7.7.5, \emph{b$'$)}) para $p=1$.
\end{proof}

\begin{corollary}[7.7.7]
\label{III.7.7.7}
Sean $Y$ un preesquema localmente noetheriano, $f:X\to Y$ un morfismo propio,
$\mathscr F$ y $\mathscr F'$ dos $\mathscr O_X$-módulos coherentes y planos
sobre $Y$, y $u:\mathscr F\to\mathscr F'$ un homomorfismo. Consideremos los
dos funtores en el $\mathscr O_Y$-módulo cuasicoherente $\mathscr M$,
\begin{align*}
  \mathscr T(\mathscr M)
  &=
  \operatorname{Ker}
  \bigl(
    f_*(\mathscr F\otimes_{\mathscr O_Y}\mathscr M)
    \longrightarrow
    f_*(\mathscr F'\otimes_{\mathscr O_Y}\mathscr M)
  \bigr)\\
  T(\mathscr M)
  &=
  \Gamma(Y,\mathscr T(\mathscr M))\\
  &=
  \operatorname{Ker}
  \bigl(
    \Gamma(X,\mathscr F\otimes_{\mathscr O_Y}\mathscr M)
    \longrightarrow
    \Gamma(X,\mathscr F'\otimes_{\mathscr O_Y}\mathscr M)
  \bigr).
\end{align*}
Entonces existe un $\mathscr O_Y$-módulo coherente $\mathscr R$ (determinado
salvo isomorfismo único) y existen isomorfismos de funtores
\[
  \mathscr T(\mathscr M)
  \xrightarrow{\sim}
  \mathscr{H}\!om_{\mathscr O_Y}(\mathscr R,\mathscr M)
  \tag{7.7.7.1}\label{III.7.7.7.1.fr}
\]
\[
  T(\mathscr M)
  \xrightarrow{\sim}
  \operatorname{Hom}_{\mathscr O_Y}(\mathscr R,\mathscr M).
  \tag{7.7.7.2}\label{III.7.7.7.2.fr}
\]
\end{corollary}

\begin{proof}
Basta demostrar \eqref{III.7.7.7.2}; en efecto, esto probará
\eqref{III.7.7.7.1} en el caso en que $Y$ es afín, y de ahí se pasará al caso
general razonando como en la demostración de la equivalencia de \emph{b)} y \emph{b$'$)} en
\sref{III.7.7.5}, gracias a la unicidad salvo isomorfismo único de un
representante de un funtor representable \sref[0]{0.8.1.8}. De
\sref{III.7.7.6} resulta que existen dos $\mathscr O_Y$-módulos coherentes
$\mathscr Q$, $\mathscr Q'$ que definen isomorfismos funtoriales
\[
  \Gamma(X,\mathscr F\otimes_{\mathscr O_Y}\mathscr M)
  \xrightarrow{\sim}
  \operatorname{Hom}_{\mathscr O_Y}(\mathscr Q,\mathscr M),
  \qquad
  \Gamma(X,\mathscr F'\otimes_{\mathscr O_Y}\mathscr M)
  \xrightarrow{\sim}
  \operatorname{Hom}_{\mathscr O_Y}(\mathscr Q',\mathscr M).
\]
Ahora bien, $u:\mathscr F\to\mathscr F'$ define canónicamente un morfismo de
funtores
\[
  \Gamma(X,\mathscr F\otimes_{\mathscr O_Y}\mathscr M)
  \longrightarrow
  \Gamma(X,\mathscr F'\otimes_{\mathscr O_Y}\mathscr M);
\]

\oldpage[III]{202}
A este último le corresponde un homomorfismo único
$v:\mathscr Q'\to\mathscr Q$ de $\mathscr O_Y$-módulos tal que el diagrama
\[
\xymatrix@C=42pt@R=28pt{
  \Gamma(X,\mathscr F\otimes_{\mathscr O_Y}\mathscr M)
    \ar[r]\ar[d]_-{\wr}
  &
  \Gamma(X,\mathscr F'\otimes_{\mathscr O_Y}\mathscr M)
    \ar[d]_-{\wr}
  \\
  \operatorname{Hom}_{\mathscr O_Y}(\mathscr Q,\mathscr M)
    \ar[r]
  &
  \operatorname{Hom}_{\mathscr O_Y}(\mathscr Q',\mathscr M)
}
\]
sea conmutativo \sref[0]{0.8.1.4}. Como el funtor contravariante
$\mathscr N\leadsto\operatorname{Hom}_{\mathscr O_Y}(\mathscr N,\mathscr M)$ es exacto por la
izquierda en la categoría de los $\mathscr O_Y$-módulos, basta tomar
$\mathscr R=\operatorname{Coker}(v)$ para obtener el isomorfismo
\eqref{III.7.7.7.2} buscado.
\end{proof}

\begin{corollary}[7.7.8]
Bajo las hipótesis de \textup{(7.7.6)} relativas a $X$, $Y$ y $f$, sean $\mathscr F$,
$\mathscr G$ dos $\mathscr O_X$-Módulos coherentes que verifican las condiciones
siguientes: \textup{(i)} $\mathscr F$ es $Y$-plano; \textup{(ii)} $\mathscr G$ es isomorfo al
conúcleo de un homomorfismo de $\mathscr O_X$-Módulos localmente libres de tipo
finito $\mathscr E_1\to\mathscr E_0$. Consideremos los dos funtores en el
$\mathscr O_Y$-Módulo cuasicoherente $\mathscr M$:
\begin{align*}
  \mathscr T(\mathscr M)
  &=
  f_*
  \bigl(
    \mathscr{H}\!om_{\mathscr O_X}
      (\mathscr G,\mathscr F\otimes_{\mathscr O_Y}\mathscr M)
  \bigr)\\
  T(\mathscr M)
  &=
  \Gamma(Y,\mathscr T(\mathscr M))
  =
  \operatorname{Hom}_{\mathscr O_X}
    (\mathscr G,\mathscr F\otimes_{\mathscr O_Y}\mathscr M).
\end{align*}

Entonces existe un $\mathscr O_Y$-Módulo coherente $\mathscr N$ (determinado salvo
isomorfismo único) e isomorfismos de funtores
\[
  \mathscr T(\mathscr M)
  \xrightarrow{\sim}
  \mathscr{H}\!om_{\mathscr O_Y}(\mathscr N,\mathscr M)
  \tag{7.7.8.1}\label{III.7.7.8.1.fr}
\]
\[
  T(\mathscr M)
  \xrightarrow{\sim}
  \operatorname{Hom}_{\mathscr O_Y}(\mathscr N,\mathscr M).
  \tag{7.7.8.2}\label{III.7.7.8.2.fr}
\]
\end{corollary}

\begin{proof}
En virtud del isomorfismo funtorial ($\mathbf{0}_{\mathrm I}$, 5.4.2.1), se tienen
isomorfismos funtoriales en $\mathscr M$
\begin{align*}
  \mathscr{H}\!om_{\mathscr O_X}
    (\mathscr E_i,\mathscr F\otimes_{\mathscr O_Y}\mathscr M)
  &\xrightarrow{\sim}
  \check{\mathscr E}_i
    \otimes_{\mathscr O_X}
    (\mathscr F\otimes_{\mathscr O_Y}\mathscr M)\\
  &\xrightarrow{\sim}
  (\check{\mathscr E}_i\otimes_{\mathscr O_X}\mathscr F)
    \otimes_{\mathscr O_Y}\mathscr M\\
  &\xrightarrow{\sim}
  \mathscr{H}\!om_{\mathscr O_X}(\mathscr E_i,\mathscr F)
    \otimes_{\mathscr O_Y}\mathscr M
\end{align*}
para $i=0,1$. Pongamos
$\mathscr F_i=\mathscr{H}\!om_{\mathscr O_X}(\mathscr E_i,\mathscr F)$ para
$i=0,1$; son $\mathscr O_X$-Módulos coherentes ($\mathbf{0}_{\mathrm I}$, 5.3.5)
y $Y$-planos ($\mathbf{0}_{\mathrm I}$, 5.4.2); sea
$u=\mathscr{H}\!om(v,1_{\mathscr F}):\mathscr F_0\to\mathscr F_1$. En virtud de
la exactitud por la izquierda del funtor
$\mathscr H\leadsto\mathscr{H}\!om_{\mathscr O_X}
(\mathscr H,\mathscr F\otimes_{\mathscr O_Y}\mathscr M)$, se tienen isomorfismos
funtoriales en $\mathscr M$
\begin{align*}
  \mathscr{H}\!om_{\mathscr O_X}
    (\mathscr G,\mathscr F\otimes_{\mathscr O_Y}\mathscr M)
  &\xrightarrow{\sim}
  \operatorname{Ker}
  \bigl(
    \mathscr{H}\!om_{\mathscr O_X}
      (\mathscr E_0,\mathscr F\otimes_{\mathscr O_Y}\mathscr M)
    \longrightarrow
    \mathscr{H}\!om_{\mathscr O_X}
      (\mathscr E_1,\mathscr F\otimes_{\mathscr O_Y}\mathscr M)
  \bigr)\\
  &\xrightarrow{\sim}
  \operatorname{Ker}
  \bigl(
    \mathscr F_0\otimes_{\mathscr O_Y}\mathscr M
    \longrightarrow
    \mathscr F_1\otimes_{\mathscr O_Y}\mathscr M
  \bigr)
\end{align*}

Puesto que $f_*$ es exacto por la izquierda, se deduce un isomorfismo funtorial
\[
  f_*
  \bigl(
    \mathscr{H}\!om_{\mathscr O_X}
      (\mathscr G,\mathscr F\otimes_{\mathscr O_Y}\mathscr M)
  \bigr)
  \xrightarrow{\sim}
  \operatorname{Ker}
  \bigl(
    f_*(\mathscr F_0\otimes_{\mathscr O_Y}\mathscr M)
    \longrightarrow
    f_*(\mathscr F_1\otimes_{\mathscr O_Y}\mathscr M)
  \bigr)
\]
y basta entonces aplicar (7.7.7).
\end{proof}

\oldpage[III]{203}

\begin{remarks}[7.7.9]
\label{III.7.7.9}
\textnormal{(i)} En \sref{III.7.7.6}, \sref{III.7.7.7} y
\sref{III.7.7.8}, la formación de los $\mathscr O_Y$-módulos $\mathscr Q$,
$\mathscr R$ y $\mathscr N$ \emph{conmuta con los cambios de base}. Por ejemplo,
conservando las notaciones de \sref{III.7.7.2}, en el caso de
\sref{III.7.7.6} se tiene, para todo $\mathscr O_{Y'}$-módulo cuasicoherente
$\mathscr M'$, el isomorfismo
\[
  f'_*(\mathscr F\otimes_{\mathscr O_Y}\mathscr M')
  \xrightarrow{\sim}
  \mathscr{H}\!om_{\mathscr O_{Y'}}(g^*(\mathscr Q),\mathscr M')
\]
pues, en virtud de la observación hecha en \sref{III.7.7.2}, todo se reduce a
ver que se tiene
\[
  \operatorname{Hom}_{\mathscr O_Y}(\mathscr Q,g_*(\mathscr M'))
  =
  \operatorname{Hom}_{\mathscr O_{Y'}}(g^*(\mathscr Q),\mathscr M')
\]
lo que no es sino ($\mathbf{0}_{\mathrm I}$, 4.4.3.1). Del mismo modo, cuando en (7.7.7)
se reemplaza $Y$, $f$, $\mathscr M$, $\mathscr F$ y
$\mathscr F'$ por $Y'$, $f'$, $\mathscr M'$,
$\mathscr F\otimes_{\mathscr O_X}\mathscr O_{X'}$ y
$\mathscr F'\otimes_{\mathscr O_X}\mathscr O_{X'}$, hay que reemplazar
$\mathscr R$ por $g^*(\mathscr R)$. Finalmente, en
\sref{III.7.7.8}, cuando se sustituyen $X$, $Y$, $f$, $\mathscr M$ y
$\mathscr F$ por $X'$, $Y'$, $f'$, $\mathscr M'$ y
$\mathscr F\otimes_{\mathscr O_X}\mathscr O_{X'}$, y $\mathscr G$ por
$\mathscr G'=\mathscr G\otimes_{\mathscr O_Y}\mathscr O_{X'}$,
hay que sustituir $\mathscr N$ por $g^*(\mathscr N)$: esto resulta de que se tiene también una
sucesión exacta
$\mathscr E'_1\longrightarrow\mathscr E'_0\longrightarrow\mathscr G'\longrightarrow0$ con
$\mathscr E'_i=\mathscr E_i\otimes_{\mathscr O_Y}\mathscr O_{X'}$
y de que
$  \check{\mathscr E}'_i
  \otimes_{\mathscr O_{X'}}
  (\mathscr F\otimes_{\mathscr O_Y}\mathscr M')
  =
  \check{\mathscr E}_i
  \otimes_{\mathscr O_X}
  (\mathscr F\otimes_{\mathscr O_Y}\mathscr M')$ ($i=0,1$).

\medskip
\textnormal{(ii)} La condición (ii) del enunciado de (7.7.8) relativa a
$\mathscr G$ se satisface siempre para cualquier
$\mathscr O_X$-módulo coherente $\mathscr G$ cuando existe un
$\mathscr O_X$-módulo invertible $Y$-\emph{amplio}, por ejemplo cuando $Y$ es afín y
$f:X\to Y$ es un morfismo \emph{proyectivo}. Basta observar, teniendo en cuenta
\sref[II]{II.5.5.1}, que existe un $\mathscr O_X$-módulo localmente libre de
tipo finito $\mathscr E_0$ tal que $\mathscr G$ es isomorfo a un cociente de
$\mathscr E_0$ \sref[II]{II.2.7.10}. Como $\mathscr E_0$ y $\mathscr G$ son
coherentes, también lo es el núcleo $\mathscr G_1$ de
$\mathscr E_0\to\mathscr G$; aplicando el mismo resultado a $\mathscr G_1$, se
obtiene una sucesión exacta
$\mathscr E_1\to\mathscr E_0\to\mathscr G\to0$
donde $\mathscr E_0$ y $\mathscr E_1$ son localmente libres de tipo finito.

\medskip
\textnormal{(iii)} Demostraremos en el capítulo V que, en
\sref{III.7.7.8}, la hipótesis restrictiva (ii) es \emph{superflua}.
\end{remarks}

\begin{proposition}[7.7.10]
\label{III.7.7.10}
\textnormal{(Criterios locales para la propiedad de cambio de base).} Bajo las
condiciones generales de \textup{\sref{III.7.7.5}}, sean $y$ un punto de $Y$ y $p$ un
entero. Las propiedades siguientes son equivalentes:
\begin{enumerate}
  \item[\textup{a)}] El funtor $T_p^{\mathscr O_y}$ es exacto por la derecha.

  \item[\textup{b)}] El homomorfismo canónico
    $      T_p^{\mathscr O_y}(\mathscr O_y)
      \longrightarrow
      T_p^{\mathscr O_y}(\boldsymbol{k}(y))$
    es sobreyectivo.

  \item[\textup{c)}] Para todo entero $n$, el homomorfismo canónico
    $      T_p^{\mathscr O_y}
      (\mathscr O_y/\mathfrak m_y^{n+1})
      \longrightarrow
      T_p^{\mathscr O_y}(\boldsymbol{k}(y))$
    es sobreyectivo.
\end{enumerate}
Además, el conjunto de los $y\in Y$ que verifican estas condiciones es el
mayor abierto $U$ de $Y$ tal que $\mathscr T_p^U$ sea exacto por la derecha.
\end{proposition}

\begin{proof}
Habida cuenta de \sref{III.7.7.4}, la equivalencia de \emph{a)}, \emph{b)} y \emph{c)} resulta de
\sref{III.7.4.7} y \sref{III.7.5.2}. El hecho de que el conjunto $U$ donde
$T_p^{\mathscr O_y}$ es exacto por la derecha sea abierto es también
consecuencia de \sref{III.7.4.4}; recíprocamente, si $\mathscr T_p^V$ es exacto
por la derecha, lo mismo ocurre con $T_p^{\mathscr O_y}$ para todo $y\in V$,
por la condición \emph{c)} de \sref{III.7.7.5} y \sref{III.7.6.1}.
\end{proof}

\begin{corollary}[7.7.11]
\label{III.7.7.11}
Si $\mathscr T_p$ es exacto por la derecha \textup{(resp.} por la izquierda\textup{)},
entonces, para todo morfismo $g:Y'\to Y$, $\mathscr T_p^{Y'}$ es exacto por la
derecha \textup{(resp.} por la izquierda\textup{)}. El recíproco es cierto cuando el
morfismo $g$ es fielmente plano.
\end{corollary}

\oldpage[III]{204}
\begin{proof}
La primera aserción es consecuencia inmediata de \sref{III.7.6.1} y del hecho
de que la cuestión es local sobre $Y$ y $Y'$, por las condiciones \emph{c)} y \emph{b)} de
\sref{III.7.7.5}. Para demostrar la segunda aserción basta ver, en virtud de
\sref{III.7.7.10}, que para todo $y\in Y$, $T_p^{\mathscr O_y}$ es exacto por
la derecha (resp. por la izquierda). Pero, por hipótesis, existe
$y'\in Y'$ tal que $g(y')=y$, y $\mathscr O_{y'}$ es un
$\mathscr O_y$-módulo fielmente plano; la conclusión resulta entonces de la
hipótesis y de \sref{III.7.6.1}.
\end{proof}

\begin{remarks}[7.7.12]
\label{III.7.7.12}
\textnormal{(i)} Bajo las hipótesis de \sref{III.7.7.4}, supongamos además que
$\mathscr P_\bullet$ es un complejo \emph{finito}. Resulta entonces de
\sref{III.7.7.1}, puesto que se puede tomar para $\mathfrak U$ un
recubrimiento abierto afín \emph{finito} de $X$, que el bicomplejo
$C^\bullet(\mathfrak U,\mathscr P^\bullet\otimes_{\mathscr O_Y}\mathscr M)$ es también \emph{finito} y, más
precisamente, que existe un conjunto finito $E$ de pares, \emph{independiente
de $\mathscr M$}, tal que
$C^h(\mathfrak U,\mathscr P^k\otimes_{\mathscr O_Y}\mathscr M)=0$ para todos los pares $(h,k)\notin E$.
Se concluye que existe $i_1$ tal que, para $i\geq i_1$, se tiene
$\mathscr T_i(\mathscr M)=0$ para \emph{todo} $\mathscr O_Y$-módulo
cuasicoherente $\mathscr M$. En particular, $\mathscr T_i$ es trivialmente un
funtor exacto en $\mathscr M$ para estos valores de $i$ y, por consiguiente,
por \sref{III.7.4.1}, $Z'_i(L_\bullet)$ es un $A$-módulo \emph{plano} de tipo
finito (y, por tanto, \emph{proyectivo} de tipo finito, puesto que $A$ es
noetheriano) para estos valores de $i$.

Consideremos entonces el complejo $(L'_\bullet)$ de $A$-módulos tal que
$L'_i=L_i$ para $i<i_1$, $L'_{i_1}=Z'_{i_1}(L_\bullet)$ y $L'_i=0$ para
$i>i_1$ y sea $\mathscr L'_\bullet=\widetilde{L}'_\bullet$. Es claro que
$  \mathscr H_i(\mathscr L'_\bullet\otimes_{\mathscr O_Y}\mathscr M)
  =
  \mathscr H_i(\mathscr L_\bullet\otimes_{\mathscr O_Y}\mathscr M)$
para $i<i_1-1$ y también para $i\geq i_1$ (siendo entonces nulos ambos
miembros). Finalmente, como
$\operatorname{Im}(Z'_{i_1}\otimes_A M)=\operatorname{Im}(L_{i_1}\otimes_A M)$
por definición, se tiene asimismo
$\mathscr H_i(\mathscr L'_\bullet\otimes_{\mathscr O_Y}\mathscr M)
=\mathscr H_i(\mathscr L_\bullet\otimes_{\mathscr O_Y}\mathscr M)$
para $i=i_1-1$. Se ve, pues, que se puede suponer en \sref{III.7.7.4} que
$\mathscr L_\bullet$ es también un complejo \emph{finito}, a condición de
exigir solamente que los $\mathscr L_i$ sean $\mathscr O_Y$-módulos
\emph{localmente libres}, asociados a $A$-módulos proyectivos de tipo finito.

Este razonamiento se aplica, en particular, al caso en que
$\mathscr P_\bullet$ se reduce a \emph{un solo} término $\mathscr F\ne0$, de grado 0 (en cuyo caso
$\mathscr T_n(\mathscr M)=\mathrm R^{-n}f_*
(\mathscr F\otimes_{\mathscr O_Y}\mathscr M)$). Se puede entonces suponer que
los $\mathscr L_i$ son nulos para $i>0$; en este caso se utilizarán
preferentemente las notaciones cohomológicas, escribiendo $\mathscr T^{-p}$ en
lugar de $\mathscr T_p$.

\medskip
\textnormal{(ii)} Cuando en el enunciado de \sref{III.7.7.5} ya no se supone
que los $\mathscr P_i$ sean planos sobre $Y$, las conclusiones siguen siendo
válidas a condición de poner esta vez
\[
  \mathscr T_p(\mathscr M)
  =
  \boldsymbol{\mathscr{T}\!or}^Y_p
    (f,1_Y;\mathscr P_\bullet,\mathscr M).
  \tag{7.7.12.1}\label{III.7.7.12.1.fr}
\]
En efecto,
$\boldsymbol{\mathscr{T}\!or}^Y_n(f,1_Y;\mathscr P_\bullet,\mathscr O_Y)$ es entonces un
$\mathscr O_Y$-módulo \emph{coherente} en virtud de \sref{III.6.7.9}. La
demostración de \sref{III.6.10.5} se aplica sin cambio, teniendo en cuenta
\sref{III.6.10.1}, y muestra también que, cuando
$Y=\operatorname{Spec}(A)$ es afín, se tiene
$\mathscr T_p(\mathscr M)=\mathscr H_p(\mathscr L_\bullet\otimes_{\mathscr O_Y}\mathscr M)$,
con $\mathscr L_\bullet=\widetilde{L}_\bullet$,
donde $L_\bullet$ es un complejo de $A$-módulos libres de tipo finito; esto
prueba la aserción.
\end{remarks}

\subsection{Aplicación a los morfismos propios: II. Criterios de planitud cohomológica}
\label{subsection:III.7.8}
\label{subsection:III.7.8.fr}
\label{III.7.8}

\begin{definition}[7.8.1]
\label{III.7.8.1}
Sean $X$ e $Y$ dos preesquemas, sea $f:X\to Y$ un morfismo
cuasicompacto y separado, sea $\mathscr P_\bullet$ un complejo de
$\mathscr O_X$-módulos cuasicoherentes y planos sobre $Y$, y sea
$\mathscr T_\bullet$ el funtor homológico de $\mathscr O_Y$-módulos
cuasicoherentes definido por \eqref{III.7.7.1.1}. Sea $y$ un punto
de $Y$.

Se dice que $\mathscr P_\bullet$ es
\oldpage[III]{205}
homológicamente plano sobre
$Y$ en el punto $y$, en dimensión $p$ \textup{(o} cohomológicamente
plano sobre $Y$ en el punto $y$, en dimensión $\setmin p$) si existe una
vecindad abierta $U$ de $y$ en $Y$ tal que
$\mathscr T_p^U=\mathscr T_U^{-p}$
es exacto.

Se dice que $\mathscr P_\bullet$ es homológicamente plano en
dimensión $p$ sobre $Y$ \textup{(o} cohomológicamente plano en
dimensión $\setmin p$ sobre $Y$) si es homológicamente plano sobre $Y$ en
todo punto $y\in Y$, en dimensión $p$.

\upshape Cuando $\mathscr P_\bullet$ es homológicamente plano sobre $Y$
(resp. sobre $Y$ en el punto $y$) en \emph{toda}
dimensión $p$, se dice simplemente que $\mathscr P_\bullet$ es
\emph{homológicamente plano sobre $Y$} (resp. sobre $Y$
en el punto $y$), o \emph{cohomológicamente plano sobre $Y$}
(resp. sobre $Y$ en el punto $y$).
\end{definition}

\begin{env}[7.8.2]
Por definición, la noción de planitud homológica sobre $Y$.

Si $Y$ es localmente noetheriano, o si $Y$ es un esquema, para que
$\mathscr P_\bullet$ sea homológicamente plano sobre $Y$ en
dimensión $p$, es necesario y suficiente que el funtor
$\mathscr T_p$ sea \emph{exacto}: la demostración se ha hecho en el caso en que $Y$ es localmente noetheriano en el curso de la demostración de (7.7.5); el razonamiento es el mismo (basado en (\textbf{I}, 9.4.2) aplicado a un abierto afín en un esquema cuasicompacto) cuando $Y$ es un esquema.
\end{env}

\begin{proposition}[7.8.3]
\label{III.7.8.3}
Con las notaciones e hipótesis de \textup{\sref{III.7.8.1}}, las condiciones
siguientes son equivalentes:
\begin{enumerate}
  \item[\textup{a)}] $\mathscr P_\bullet$ es homológicamente plano sobre $Y$
    en el punto $y$, en dimensión $p$.
  \item[\textup{b)}] Existe una vecindad abierta $U$ de $y$ en $Y$ tal que
    $\mathscr T_p^U$ y $\mathscr T_{p+1}^U$ son exactos por la
    derecha.
  \item[\textup{c)}] Existe una vecindad abierta $U$ de $y$ en $Y$ tal que
    $\mathscr T_p^U$ y $\mathscr T_{p-1}^U$ son exactos por la
    izquierda.
  \item[\textup{d)}] Existe una vecindad abierta $U$ de $y$ en $Y$ tal que
    $\mathscr T_{p+1}^U$ es exacto por la derecha y
    $\mathscr T_{p-1}^U$ es exacto por la izquierda.
\end{enumerate}
\end{proposition}

\begin{proof}
Teniendo en cuenta la interpretación de $\mathscr T_p$ cuando $Y$
es afín, esto no es más que una traducción de una parte de
\sref{III.7.3.3}.
\end{proof}

\begin{proposition}[7.8.4]
\label{III.7.8.4}
Sean $Y$ un preesquema localmente noetheriano, $f:X\to Y$ un
morfismo propio, $\mathscr P_\bullet$ un complejo acotado
inferiormente de $\mathscr O_X$-módulos coherentes y planos sobre
$Y$, y $\mathscr T_\bullet$ el funtor definido por
\eqref{III.7.7.1.1}. Para todo $y\in Y$, las condiciones siguientes
son equivalentes:
\begin{enumerate}
  \item[\textup{a)}] $\mathscr P_\bullet$ es homológicamente plano sobre $Y$
    en $y$, en dimensión $p$.
  \item[\textup{b)}] El funtor $T_p^{\mathscr O_y}$ es exacto.
  \item[\textup{c)}] Existe un entero $n_0$ tal que, para $n\geq n_0$,
    \[
      \operatorname{long}T_p^{\mathscr O_y}
        (\mathscr O_y/\mathfrak m_y^{n+1})
      =
      \operatorname{long}T_p^{\mathscr O_y}(\boldsymbol{k}(y))
      .
      \operatorname{long}\mathscr O_y/\mathfrak m_y^{n+1}
      \tag{7.8.4.1}\label{III.7.8.4.1.fr}
    \]
(donde se trata de longitudes de $\mathscr O_y$-módulos.
  \item[\textup{d)}] Existe una vecindad abierta $U$ de $y$ tal que
    $(\mathcal H^{-p}(f,\mathscr P^\bullet))|U$
    es isomorfo a un $(\mathscr O_Y|U)$-módulo de la forma
    $(\mathscr O_Y|U)^m$ y, para todo $(\mathscr O_Y|U)$-módulo
    cuasicoherente $\mathscr M$, el homomorfismo canónico
    \[
      ((\mathcal H^{-p}(f,\mathscr P^\bullet))|U)
      \otimes_{\mathscr O_Y|U}\mathscr M
      \longrightarrow
      \mathcal H^{-p}
      \bigl(
        f,(\mathscr P^\bullet|U)
          \otimes_{\mathscr O_Y|U}\mathscr M
      \bigr)
      \tag{7.8.4.2}\label{III.7.8.4.2.fr}
    \]
    es biyectivo.
\end{enumerate}
Cuando se cumplen estas condiciones, se tiene además:
\begin{enumerate}
  \item[\textup{e)}] Existe una vecindad de $y$ en la cual la función
    $z\leadsto d_p(z)$, definida en \eqref{III.7.7.5.1}, es
    constante.
\end{enumerate}
Además, si $Y$ es \textup{reducido} en $y$ \textup{($\mathbf{0}_{\mathrm I}$, 4.1.4)}, e) es
equivalente a las demás condiciones.
\end{proposition}

\oldpage[III]{206}
\begin{proof}
En efecto, la condición \emph{b)} equivale a decir que $T_p^{\mathscr O_y}$ y $T_{p+1}^{\mathscr O_y}$
son exactos por la derecha \sref{III.7.3.3}. La equivalencia de \emph{a)}
y \emph{b)} resulta entonces de \sref{III.7.7.10} y \sref{III.7.8.3}.

Como $\mathscr O_y/\mathfrak m_y^{n+1}$ es artiniano, y que
$T_p^{\mathscr O_y}(\mathscr O_y/\mathfrak m_y^{n+1})$ y
$T_p^{\mathscr O_y}(\boldsymbol{k}(y))$ son
$(\mathscr O_y/\mathfrak m_y^{n+1})$-módulos de tipo finito
\sref{III.7.7.4}, y por tanto de longitud finita, la equivalencia de
\emph{b)} y \emph{c)} resulta nuevamente de \sref{III.7.7.10} y
\sref{III.7.3.5.7}.

Que \emph{a)} implica \emph{e)}, y que \emph{a)} y \emph{e)} son equivalentes cuando
$\mathscr O_y$ es reducido, resulta de \sref{III.7.6.9}.
Finalmente, \emph{a)} implica que \eqref{III.7.8.4.2} es biyectivo, por la
definición \sref{III.7.8.1} y \sref{III.7.7.5}. También \emph{a)} implica que
$(\mathcal H^{-p}(f,\mathscr P^\bullet))_y$
es un $\mathscr O_y$-módulo plano (7.3.3, \emph{f)}), luego libre (\textbf{0}, 10.1.3), puesto que se trata de un $\mathscr O_y$-módulo libre de
tipo finito \sref[0]{0.10.1.3}, pues es de tipo finito por
\sref{III.7.7.4}. Como $\mathcal H^{-p}(f,\mathscr P^\bullet)$ es
un $\mathscr O_Y$-módulo coherente \sref{III.7.7.4}, es localmente
libre en una vecindad de $y$
($\mathbf{0}_{\mathrm I}$, 5.2.7). Recíprocamente, está claro que \emph{d)} implica \emph{a)} por
definición del funtor $\mathscr T_p^U$ \eqref{III.7.7.2.2}.
\end{proof}

\begin{proposition}[7.8.5]
\label{III.7.8.5}
Bajo las hipótesis de \textup{\sref{III.7.8.4}}, las condiciones siguientes
son equivalentes:
\begin{enumerate}
  \item[\textup{a)}] $\mathscr P_\bullet$ es homológicamente plano sobre $Y$
    en toda dimensión $i\leq p$.
  \item[\textup{b)}] Para $i\leq p+1$, los funtores $\mathscr T_i$ son
    exactos por la derecha.
  \item[\textup{c)}] Para $i\geq\setmin p$, los $\mathscr O_Y$-módulos
    $\mathcal H^i(f,\mathscr P^\bullet)$ son localmente libres.
\end{enumerate}
\end{proposition}

\begin{proof}
La equivalencia de \emph{a)} y \emph{b)} es inmediata por \sref{III.7.8.3}, y \emph{a)}
implica \emph{c)} por \sref{III.7.8.4}.

Recíprocamente, supongamos que se cumple \emph{c)}. Recordemos también que
$\mathscr L_i=0$ para $i\leq i_0$ \sref{III.7.7.4}, y por tanto
$\mathscr T_i=0$ para $i\leq i_0$. Todo punto $y\in Y$ tiene, pues,
una vecindad afín $U=\operatorname{Spec}(A)$ tal que $T_i^A(A)$ es un $A$-módulo
libre para $i\leq p$. De \sref{III.7.3.7} se deduce que
$\mathscr T_i^U=T_i^A$ es exacto para $i\leq p$.
\end{proof}

Aplicaremos principalmente los criterios de planitud cohomológica al
caso en que el complejo $\mathscr P_\bullet$ está concentrado en \emph{un solo} $\mathscr O_X$-módulo coherente $\mathscr F$, plano
sobre $Y$, tomado como $\mathscr P_0$. Recordemos que en este caso
$\mathscr T_p(\mathscr M)=\mathrm R^{-p}f_*
(\mathscr F\otimes_{\mathscr O_Y}\mathscr M)$.

\begin{proposition}[7.8.6]
Sean $Y$ un preesquema localmente noetheriano, $f:X\to Y$ un morfismo propio
y plano, $y$ un punto de $Y$; designemos por $X_y$ la fibra
$f^{-1}(y)=X\otimes_Y\boldsymbol{k}(y)$. Supongamos que
$\Gamma(X_y,\mathscr O_{X_y})=R$ es una $\boldsymbol{k}(y)$-álgebra separable
\textup{(Bourbaki, \emph{Alg.}, cap. VIII, § 7, n\textsuperscript{o} 5)}, en otros términos,
compuesta por un número finito de extensiones separables de grado finito de $\boldsymbol{k}(y)$.
Entonces $\mathscr O_X$ es cohomológicamente plano sobre $Y$ en el punto $y$ en
dimensión 0.
\end{proposition}

\begin{proof}
En virtud de (7.8.4), se puede limitar al caso en que $Y$ es el espectro del anillo
local $A=\mathscr O_y$; la hipótesis de que $f$ es plano implica
$\mathscr T_{-1}=0$, luego se ve ya que $T_0^A$ es exacto por la izquierda y todo
se reduce a ver que es exacto por la derecha; en virtud de (7.7.10 \emph{c)}), se está incluso
reducido al caso en que $A=\mathscr O_y$ es \emph{artiniano}. Sea $k'$ una extensión
finita de $\boldsymbol{k}(y)$ que sea un cuerpo de descomposición de $R$, de modo que
$R\otimes_{\mathbf{k}(y)}k'$ es producto directo de un número finito de cuerpos
isomorfos a $k'$. Se sabe que existe un homomorfismo local de $A$ en un
anillo local $A'$, que hace de $A'$ una $A$-álgebra \emph{libre} finita sobre $A$,
y tal que el cuerpo residual de $A'$ es isomorfo a $k'$ (\textbf{0},
10.3.2). En virtud de (7.6.1), se está reducido a probar que $T_0^{A'}$ es exacto
por la derecha, en otros términos, se puede suponer que $R$ es producto directo de $m$
cuerpos isomorfos a $\boldsymbol{k}(y)$. Observemos ahora el lema elemental
siguiente:

\begin{lemma}[7.8.6.1]
Sea $Z$ un espacio anillado en anillos locales; para que $Z$ sea conexo,
\oldpage[III]{207}
es necesario y suficiente que el anillo $\Gamma(Z,\mathscr O_Z)$ no sea un producto
de dos anillos no reducidos a 0.
\end{lemma}

\begin{proof}
Está claro en efecto que si $Z$ es unión de dos abiertos no vacíos disjuntos,
$\Gamma(Z,\mathscr O_Z)$ es isomorfo al producto de los dos anillos
$\Gamma(Z_1,\mathscr O_Z)$ y $\Gamma(Z_2,\mathscr O_Z)$ no reducidos a 0.
Recíprocamente, decir que $\Gamma(Z,\mathscr O_Z)$ es tal producto equivale a
decir que hay en $\Gamma(Z,\mathscr O_Z)$ un idempotente $s$ distinto de 0 y
de 1; para todo $z\in Z$, $s_z$ es entonces un idempotente en $\mathscr O_z$,
luego igual a 0 o 1. Pero está claro que el conjunto de los $z$ tales que $s_z=0$
es abierto; por otra parte, si $s_z=1$, se tiene por definición $s(z)\ne0$, luego
el conjunto de los $z$ donde $s_z=1$ es también abierto ($\mathbf{0}_{\mathrm I}$, 5.5.2); de donde
la conclusión.
\end{proof}

Resulta de este lema que $X_y$ tiene exactamente $m$ componentes conexas $X'_i$
y que $\Gamma(X'_i,\mathscr O_{X'_i})=\boldsymbol{k}(y)$ para todo $i$. Como $A$ se ha
supuesto local y artiniano, su espectro se reduce a un punto, luego $X$ y $X_y$
tienen el mismo espacio subyacente; $X$ tiene pues $m$ componentes conexas $X_i$ tales
que $X'_i=X_i\otimes_Y\boldsymbol{k}(y)$. Se está así finalmente reducido al caso en que
$R=\boldsymbol{k}(y)$; en virtud de (7.7.10, \emph{b)}), se está reducido a probar que
el homomorfismo canónico
$\Gamma(X,\mathscr O_X)\to\Gamma(X_y,\mathscr O_{X_y})$
es sobreyectivo; pero esto es trivial, pues el compuesto
\[
  \Gamma(Y,\mathscr O_Y)=A
  \longrightarrow
  \Gamma(X,\mathscr O_X)
  \longrightarrow
  \Gamma(X_y,\mathscr O_{X_y})=\boldsymbol{k}(y)
\]
ya es sobreyectivo.
\end{proof}

\begin{corollary}[7.8.7]
\label{III.7.8.7}
Bajo las hipótesis de \textup{\sref{III.7.8.6}}, existe una vecindad abierta
$U$ de $y$ tal que:
\begin{enumerate}
  \item[\textup{(i)}] $f_*(\mathscr O_X)|U$ es isomorfo a un
    $(\mathscr O_Y|U)$-módulo de la forma $(\mathscr O_Y|U)^m$;
  \item[\textup{(ii)}] para todo $z\in U$, el homomorfismo canónico
    \[
      (f_*(\mathscr O_X))_z
      \otimes_{\mathscr O_z}\boldsymbol{k}(z)
      \longrightarrow
      \Gamma(X_z,\mathscr O_{X_z})
    \]
    es biyectivo.
\end{enumerate}
\end{corollary}

\begin{proof}
La aserción (i) resulta de \sref{III.7.8.6} y \sref{III.7.8.4}.
La aserción (ii) resulta de la exactitud de $\mathscr T_0^U$, para
una elección adecuada de $U$, y de \eqref{III.7.7.5.3}.
\end{proof}

\begin{corollary}[7.8.8]
\label{III.7.8.8}
Supongamos que se cumplen las hipótesis de \textup{\sref{III.7.8.6}} y,
además, que
$\Gamma(X_y,\mathscr O_{X_y})=\boldsymbol{k}(y)$.
Entonces existe una vecindad abierta $U$ de $y$ tal que el
homomorfismo canónico
$\mathscr O_Y|U\to f_*(\mathscr O_X)|U$
es biyectivo.
\end{corollary}

\begin{proof}
En efecto, de la aserción (ii) de \sref{III.7.8.7} se deduce que el
entero $m$ que figura en (7.8.7,
(i)) es necesariamente igual a 1.
\end{proof}

\begin{corollary}[7.8.9]
\label{III.7.8.9}
Bajo las hipótesis de \textup{\sref{III.7.8.6}}, existen una vecindad abierta
$U$ de $y$, un $\mathscr O_U$-módulo coherente $\mathscr Q$,
determinado salvo isomorfismo único, y un isomorfismo de funtores en
el $\mathscr O_U$-módulo cuasicoherente $\mathscr M$,
\[
  \mathrm R^1f_*(f^*(\mathscr M))
  \xrightarrow{\sim}
  \mathscr{H}\!om_{\mathscr O_U}(\mathscr Q,\mathscr M).
  \tag{7.8.9.1}\label{III.7.8.9.1.fr}
\]
\end{corollary}

\begin{proof}
La hipótesis implica que $\mathscr T_0^U$ es exacto para un $U$
apropiado; basta pues aplicar la equivalencia de (7.7.5, \emph{a)}) y
(7.7.5, \emph{b$'$)}) al caso $p=0$ y tomando como
$\mathscr P_\bullet$ el complejo
reducido a su término de grado 0 igual a $\mathscr O_X$.
\end{proof}

\begin{remarks}[7.8.10]
\label{III.7.8.10}
\textnormal{(i)}
Bajo las hipótesis de \sref{III.7.8.6}, consideremos la
factorización de Stein de $f$ \sref{III.4.3.3},
\[
\xymatrix{
  X \ar[r]^{f'} & Y' \ar[r]^{g} & Y
}
\]
\oldpage[III]{208}
Aquí
$Y'=\operatorname{Spec}(f_*(\mathscr O_X))$;
El morfismo finito $g$ es tal que
$g_*(\mathscr O_{Y'})=f_*(\mathscr O_X)$
es localmente libre cerca de $y$, y su fibra sobre $y$ es el
espectro de una álgebra separable sobre $\boldsymbol{k}(y)$
\sref[II]{II.1.5.1}. Deduciremos en el capítulo~IV que existe una
vecindad abierta $U$ de $y$ en $Y$ tal que, para la restricción
$g^{-1}(U)\to U$ de $g$, \emph{toda} fibra $g^{-1}(z)$ (donde $z\in U$) sea espectro de un álgebra separable sobre $\boldsymbol{k}(z)$ (es lo que llamaremos un \emph{revestimiento étale} de $U$); resultará entonces de (7.8.7, (ii)) que la hipótesis hecha sobre el punto $y$ en \sref{III.7.8.6} se cumple también
en todo punto de una vecindad de $y$.

\medskip
\textnormal{(ii)}
Veremos en el capítulo~V que, incluso si $X$ es proyectivo sobre $Y$
(e incluso si además es ``simple'' sobre $Y$, propiedad que se
definirá en el capítulo~IV), el $\mathscr O_U$-módulo $\mathscr Q$
de \sref{III.7.8.9} no es necesariamente localmente libre. En otras
palabras, bajo estas hipótesis $\mathscr O_X$ (bajo estas
condiciones) no es necesariamente cohomológicamente plano en \emph{dimensión
1} sobre $Y$ en el
punto $y$. En el capítulo~V interpretaremos $\mathscr Q$ como el
haz de las 1-diferenciales del esquema de Picard de $X$ relativo a
$Y$, a lo largo de la sección unidad.
\end{remarks}

\subsection{Aplicación a los morfismos propios: III. Invariancia de la
característica de Euler-Poincaré y del polinomio de Hilbert.}
\label{subsection:III.7.9}
\label{subsection:III.7.9.fr}
\label{III.7.9}
\label{III.7.9.fr}

\begin{env}[7.9.1]
\label{III.7.9.1}
\label{III.7.9.1.fr}
Sean $A$ un anillo y $M$ un $A$-módulo proyectivo de tipo finito.
Recordemos (Bourbaki, \emph{Alg.~comm.}, cap.~II, \S5, n.\textsuperscript{o}~2)
que esto equivale a decir que el $\mathscr O_X$-módulo $\widetilde M$ sobre
$X=\operatorname{Spec}(A)$ es localmente libre de tipo finito.
Para todo $\mathfrak p\in\operatorname{Spec}(A)$, se llama \emph{rango de $M$ en
$\mathfrak p$}, y se denota por $\operatorname{rang}_{\mathfrak p}(M)$,
al rango del $A_{\mathfrak p}$-módulo libre $M_{\mathfrak p}$ (o,
equivalentemente, al rango en $\mathfrak p$ del $\mathscr O_X$-módulo
localmente libre $\widetilde M$). Por tanto,
\[
  \operatorname{rang}_{\mathfrak p}M
  =\operatorname{rang}_{\mathfrak p}(M_{\mathfrak p})
  =\operatorname{rang}_{\mathbf{k}(\mathfrak p)}
  \bigl(M\otimes_A\boldsymbol{k}(\mathfrak p)\bigr).
  \tag{7.9.1.1}\label{III.7.9.1.1.fr}
\]
\end{env}

\begin{proposition}[7.9.2]
\label{III.7.9.2}
\label{III.7.9.2.fr}
Sea $P_\bullet$ un complejo finito de $A$-módulos proyectivos de
tipo finito y, para todo $A$-módulo $M$, sea $T_\bullet(M)=H_\bullet(P_\bullet\otimes_A M)$.
Entonces, para todo $\mathfrak p\in\operatorname{Spec}(A)$, se tiene
\[
  \sum_i(\setmin 1)^i\operatorname{rang}_{\mathbf{k}(\mathfrak p)}
  T_i\bigl(\boldsymbol{k}(\mathfrak p)\bigr)
  =\sum_i(\setmin 1)^i\operatorname{rang}_{\mathfrak p}(P_i).
  \tag{7.9.2.1}\label{III.7.9.2.1.fr}
\]
\end{proposition}

\begin{proof}
En efecto, por definición,
$T_i(\boldsymbol{k}(\mathfrak p))=H_i(P_\bullet\otimes_A\boldsymbol{k}(\mathfrak p))$.
Teniendo en cuenta \eqref{III.7.9.1.1}, la fórmula
\eqref{III.7.9.2.1} no es sino la invariancia de la característica
de Euler-Poincaré de un complejo finito de espacios vectoriales de
dimensión finita al pasar a la homología
\sref[0]{0.11.10.2}.
\end{proof}

\begin{corollary}[7.9.3]
\label{III.7.9.3}
\label{III.7.9.3.fr}
La función
\[
  \mathfrak p\leadsto
  \sum_i(-1)^i\operatorname{rang}_{\mathbf{k}(\mathfrak p)}
  T_i\bigl(\boldsymbol{k}(\mathfrak p)\bigr)
\]
es localmente constante en $\operatorname{Spec}(A)$.
\end{corollary}

\begin{theorem}[7.9.4]
\label{III.7.9.4}
\label{III.7.9.4.fr}
Sean $Y$ un preesquema localmente noetheriano,
$f:X\to Y$ un morfismo propio y $\mathscr P_\bullet$ un complejo
finito de $\mathscr O_X$-módulos coherentes y planos sobre $Y$-planos. Si se pone
$\mathscr T_\bullet(\mathscr M)
  =\mathscr H^{-\bullet}
  \bigl(f,\mathscr P^\bullet\otimes_{\mathscr O_Y}\mathscr M\bigr)$ \textup{(cf. (7.7.1.1))},
la función
\oldpage[III]{209}
\[
  y\leadsto
  \sum_i(\setmin 1)^i\operatorname{rang}_{\mathbf{k}(y)}
  T_i\bigl(\boldsymbol{k}(y)\bigr)
  \tag{7.9.4.1}\label{III.7.9.4.1.fr}
\]
es localmente constante en $Y$.
\end{theorem}

\begin{proof}
Se puede limitar al caso en que $Y=\operatorname{Spec}(A)$ es afín y $A$ es
noetheriano. Como el complejo $\mathscr P_\bullet$ es finito,
\sref{III.7.7.12}{(i)} muestra que
$\mathscr T_p(\mathscr M)=\mathscr H_p(\mathscr L_\bullet\otimes_{\mathscr O_Y}\mathscr M)$,
donde $\mathscr L_\bullet=\widetilde{L}_\bullet$ y
$L_\bullet$ es un complejo finito de $A$-módulos proyectivos de
tipo finito. El teorema se deduce entonces de
\sref{III.7.9.3}.
\end{proof}

\begin{env}[7.9.5]
\label{III.7.9.5}
\label{III.7.9.5.fr}
Bajo las hipótesis de \sref{III.7.9.4}, la función
\eqref{III.7.9.4.1} es \emph{constante} cuando $Y$ es
\emph{conexo}. Cuando $Y$ es conexo y no vacío, se designa
el valor único (entero) de (7.9.4.1) por $\operatorname{EP}(f,\mathscr P_\bullet)$ o $\operatorname{EP}(Y,\mathscr P_\bullet)$, o simplemente
$\operatorname{EP}(\mathscr P_\bullet)$
si no puede resultar confusión, y se dice que este entero es la
\emph{característica de Euler-Poincaré} de $\mathscr P_\bullet$
relativa a $f$ (o a $Y$). En el caso general, se denotará también por $\operatorname{EP}(f,\mathscr P_\bullet;y)$ o $\operatorname{EP}(Y,\mathscr P_\bullet;y)$ o
$\operatorname{EP}(\mathscr P_\bullet;y)$
el segundo miembro de (7.9.4.1).
\end{env}

\begin{env}[7.9.6]
\label{III.7.9.6}
\label{III.7.9.6.fr}
Bajo las hipótesis de \sref{III.7.9.4} relativas a
$X$, $Y$ y $f$, sea
\[
  0\longrightarrow\mathscr P'_\bullet
  \xrightarrow{u}\mathscr P_\bullet
  \xrightarrow{v}\mathscr P''_\bullet
  \longrightarrow0
\]
una sucesión exacta de complejos finitos de $\mathscr O_X$-módulos
coherentes y planos sobre $Y$, donde los homomorfismos $u$ y $v$
son de grados pares $2d$ y $2d'$, respectivamente.
Como $\mathscr T_\bullet$ es un funtor homológico
\sref{III.7.7.1}, se tiene una sucesión exacta de homología
\[
  \longrightarrow
  \mathscr T_i(\mathscr P'_\bullet,\boldsymbol{k}(y))
  \longrightarrow
  \mathscr T_{i+2d}(\mathscr P_\bullet,\boldsymbol{k}(y))
  \longrightarrow
  \mathscr T_{i+2d+2d'}(\mathscr P''_\bullet,\boldsymbol{k}(y))
  \longrightarrow
  \mathscr T_{i-1}(\mathscr P'_\bullet,\boldsymbol{k}(y))
  \longrightarrow\ldots
\]
que, además, solo tiene un número finito de términos. Al escribir
que la característica de Euler-Poincaré de este complejo es nula
\sref[0]{0.11.10.1}, se obtiene inmediatamente
\[
  \operatorname{EP}(\mathscr P_\bullet;y)
  =\operatorname{EP}(\mathscr P'_\bullet;y)
  +\operatorname{EP}(\mathscr P''_\bullet;y)
  \tag{7.9.6.1}\label{III.7.9.6.1.fr}
\]
para todo $y\in Y$. Ahora bien, si por ejemplo $\mathscr P_\bullet=(\mathscr P_i)$ y que $\mathscr P_i=0$ para
$i<0$. Se tiene entonces la sucesión exacta de complejos
\[
\xymatrix@C=2.4em{
  \ldots \ar[r] & 0 \ar[r]\ar[d] & 0 \ar[r]\ar[d]
    & 0 \ar[r]\ar[d] & 0 \ar[r]\ar[d] & \ldots \\
  \ldots \ar[r] & 0 \ar[r]\ar[d] & 0 \ar[r]\ar[d]
    & \mathscr P_1 \ar[r]\ar[d] & \mathscr P_2 \ar[r]\ar[d] & \ldots \\
  \ldots \ar[r] & 0 \ar[r]\ar[d] & \mathscr P_0 \ar[r]\ar[d]
    & \mathscr P_1 \ar[r]\ar[d] & \mathscr P_2 \ar[r]\ar[d] & \ldots \\
  \ldots \ar[r] & 0 \ar[r]\ar[d] & \mathscr P_0 \ar[r]\ar[d]
    & 0 \ar[r]\ar[d] & 0 \ar[r]\ar[d] & \ldots \\
  \ldots \ar[r] & 0 \ar[r] & 0 \ar[r]
    & 0 \ar[r] & 0 \ar[r] & \ldots
}
\]
cuyas flechas verticales no nulas son los automorfismos identidad.
Aplicando \eqref{III.7.9.6.1} a esta sucesión exacta y razonando por
inducción sobre la longitud de $\mathscr P_\bullet$, se obtiene
\[
  \operatorname{EP}(\mathscr P_\bullet;y)
  =\sum(\setmin 1)^i\operatorname{EP}(\mathscr P_i;y)
  \tag{7.9.6.2}\label{III.7.9.6.2.fr}
\]

\oldpage[III]{210}
Aquí, para todo $\mathscr O_X$-módulo coherente $\mathscr F$ plano
sobre $Y$, la notación $\operatorname{EP}(\mathscr F;y)$ (o
$\operatorname{EP}(f,\mathscr F;y)$, o
$\operatorname{EP}(Y,\mathscr F;y)$) significa
$\operatorname{EP}(\mathscr Q_\bullet;y)$ correspondiente al complejo $\mathscr Q_\bullet$ cuyo único término $\ne0$ es de grado $0$ e igual a
$\mathscr F$. Se ve pues que se puede reducir al estudio de las características
de Euler-Poincaré de complejos reducidos a un solo término.
\end{env}

\begin{proposition}[7.9.7]
\label{III.7.9.7}
\label{III.7.9.7.fr}
Bajo las hipótesis de \textup{\sref{III.7.9.4}}, sean $Y'$ un preesquema
localmente noetheriano y $g:Y'\to Y$ un morfismo, $X'=X\times_Y Y'$, $f'=f_{(Y')}:X'\to Y'$,
$\mathscr P'_\bullet$ el complejo finito
$\mathscr P_\bullet\otimes_{\mathscr O_Y}\mathscr O_{Y'}$
de $\mathscr O_{X'}$-módulos así obtenido.
Entonces $\mathscr P'_\bullet$ está formado por
$\mathscr O_{X'}$-módulos coherentes y planos sobre $Y'$, y, para
todo $y'\in Y'$, se tiene
\[
  \operatorname{EP}(\mathscr P'_\bullet;y')
  =\operatorname{EP}(\mathscr P_\bullet;g(y')).
  \tag{7.9.7.1}\label{III.7.9.7.1.fr}
\]
\end{proposition}

\begin{proof}
Los $\mathscr O_{X'}$-módulos $\mathscr P'_i$, al ser imágenes
inversas de los $\mathscr P_i$ por la proyección $X'\to X$, son
coherentes. Son planos sobre $Y'$-planos en
virtud de ($\mathbf{0}_{\mathrm I}$, 6.2.1) y (I, 4.14.5), al ser la cuestión local sobre
$X$, $Y$ e $Y'$. Por último, $f'$ es propio por
\sref[II]{II.5.4.2}, de modo que el primer miembro de
\eqref{III.7.9.7.1} está definido. La fórmula
\eqref{III.7.9.7.1} se deduce entonces de
\textup{(6.10.4.2)}, \sref{III.7.7.2} y
\sref{III.7.6.7}, reduciéndose, como siempre se puede hacer, al
caso en que $Y$ e $Y'$ son afines.
\end{proof}

\begin{proposition}[7.9.8]
\label{III.7.9.8}
\label{III.7.9.8.fr}
Supongamos que se cumplen las hipótesis de \textup{\sref{III.7.9.4}} y,
además, que existe un entero $i_0$ tal que $T_i(\boldsymbol{k}(y))=0$ para $i\ne i_0$ y todo $y\in Y$.
Entonces $\mathscr T_{i_0}(\mathscr O_Y)=\mathscr H^{-i_0}(f,\mathscr P_\bullet)$
es un $\mathscr O_Y$-módulo localmente libre, cuyo rango en
$y\in Y$ es
$(\setmin 1)^{i_0}\operatorname{EP}(f,\mathscr P_\bullet;y)$.
\end{proposition}

\begin{proof}
Observemos primero que las hipótesis de (7.4.4) son satisfechas por los $T_j^{\mathscr O_y}$ cumplen las hipótesis de
\sref{III.7.4.4}, por lo que \sref{III.7.4.7} les es aplicable.
La hipótesis y \sref{III.7.5.3} implican que
$T_i^{\mathscr O_y}$ es nulo para $i\ne i_0$.
En consecuencia, $\mathscr T_{i_0}$ es exacto por
\sref{III.7.3.3}, y \sref{III.7.8.4} muestra que
$\mathscr H^{-i_0}(f,\mathscr P_\bullet)$ es localmente libre.
Su rango en $y\in Y$ es
\[
  \operatorname{rang}_{\mathbf{k}(y)}
  T_{i_0}(\boldsymbol{k}(y))
  =(\setmin 1)^{i_0}\operatorname{EP}(f,\mathscr P_\bullet;y)
\]
por definición, ya que $T_i(\boldsymbol{k}(y))=0$ para
$i\ne i_0$.
\end{proof}

\begin{corollary}[7.9.9]
\label{III.7.9.9}
\label{III.7.9.9.fr}
Sean $Y$ un preesquema localmente noetheriano,
$f:X\to Y$ un morfismo propio y $\mathscr F$ un
$\mathscr O_X$-módulo coherente y plano sobre $Y$-plano; se supone que existe un entero $i_0$ tal que
$H^i\bigl(f^{-1}(y),\mathscr F\otimes_{\mathscr O_Y}\boldsymbol{k}(y)\bigr)=0$
para todo $i\ne i_0$ y todo $y\in Y$. Entonces
$\mathrm R^{i_0}f_*(\mathscr F)$ es un $\mathscr O_Y$-módulo localmente
libre, cuyo rango en $y$ es
$(\setmin 1)^{i_0}\operatorname{EP}(f,\mathscr F;y)$.
\end{corollary}

En particular:

\begin{corollary}[7.9.10]
\label{III.7.9.10}
\label{III.7.9.10.fr}
Bajo las hipótesis preliminares de \textup{\sref{III.7.9.9}} relativas a
$X$, $Y$ y $\mathscr F$, supongamos que
$\mathrm R^if_*(\mathscr F)=0$ para todo $i>0$. Entonces
$f_*(\mathscr F)$ es un $\mathscr O_Y$-módulo localmente libre,
cuyo rango en $y$ es $\operatorname{EP}(f,\mathscr F;y)$.
\end{corollary}

\begin{proof}
En virtud de \sref{III.7.9.9}, basta demostrar el lema siguiente.
\end{proof}

\begin{lemma}[7.9.10.1]
\label{III.7.9.10.1}
\label{III.7.9.10.1.fr}
Bajo las hipótesis de \textup{\sref{III.7.9.10}}, se tiene
$H^i\bigl(f^{-1}(y),\mathscr F\otimes_{\mathscr O_Y}\boldsymbol{k}(y)\bigr)=0$
para todo $i>0$ y todo $y\in Y$.
\end{lemma}

\begin{proof}
En efecto, se puede limitar al caso en que $Y=\operatorname{Spec}(A)$ es afín.
Con las notaciones de \sref{III.7.9.4}, y estando
$\mathscr P_\bullet$ reducido a su término
de grado $0$ igual a $\mathscr F$, se tiene en efecto
$\mathscr T_p(\mathscr O_Y)=0$ para $p<0$. Se deduce de
\sref{III.7.3.7} que $\mathscr T_p$ es exacto para $p<0$, y el
lema resulta entonces de la equivalencia de las condiciones \emph{a)} y
\emph{d)} de \sref{III.7.7.5}.
\end{proof}

\oldpage[III]{211}
\begin{proposition}[7.9.11]
\label{III.7.9.11.fr}
Siendo las hipótesis las de \textup{(7.9.4)}, sea $\mathscr L$ un
$\mathscr O_X$-Módulo invertible muy amplio para $Y$, y pongamos
$\mathscr P_\bullet(n)=\mathscr P_\bullet\otimes_{\mathscr O_X}
\mathscr L^{\otimes n}$ para todo $n\in\mathbf Z$. Entonces, para todo $y\in Y$,
la función
\[
  n\leadsto\operatorname{EP}(f,\mathscr P_\bullet(n);y)
  \tag{7.9.11.1}\label{III.7.9.11.1.fr}
\]
es un polinomio con coeficientes en $\mathbf Q$, que es el mismo para todos los
puntos de una misma componente conexa de $Y$.
\end{proposition}

\begin{proof}
Está claro que $\mathscr P_\bullet(n)$ es un complejo de
$\mathscr O_X$-Módulos $Y$-\emph{planos}. En virtud de (7.9.6.2), se puede limitar al caso
en que $\mathscr P_\bullet$ se reduce a un solo término $\mathscr F\ne0$ de grado 0;
además, como se trata de cuestiones locales sobre $Y$, se puede suponer $Y$
afín y $f$ proyectivo (\textbf{II}, 5.5.3); pongamos $X_y=f^{-1}(y)$, y sea
$\mathscr L_y=\mathscr L\otimes_{\mathscr O_Y}\boldsymbol{k}(y)$, que es un
$\mathscr O_{X_y}$-Módulo muy amplio (\textbf{II}, 4.4.10); en virtud de (7.7.2), se tiene,
para el funtor $\mathscr T_\bullet$ relativo al complejo
$\mathscr P_\bullet(n)$, $T_i(\boldsymbol{k}(y))=H^{-i}(X_y,\mathscr F_y\otimes\mathscr L_y^{\otimes n})$
(donde $\mathscr F_y=\mathscr F\otimes_{\mathscr O_y}\boldsymbol{k}(y)$);
de donde resulta que $\operatorname{EP}(f,\mathscr F(n);y)$ no es otra que la
característica de Euler-Poincaré $\chi_{\mathbf{k}(y)}(\mathscr F_y(n))$ definida en
(2.5.1); el hecho de que (7.9.11.1) sea un polinomio resulta entonces de (2.5.3);
además, para cada $n$, su valor es constante en una componente conexa de
$Y$ (7.9.4), lo que termina la demostración.
\end{proof}

Designaremos por $\operatorname{PH}(f,\mathscr P_\bullet;y)$ o
$\operatorname{PH}(\mathscr P_\bullet;y)$ el polinomio (7.9.11.1), con coeficientes
racionales, y diremos que es el \emph{polinomio de Hilbert} en $y$
relativo a $\mathscr P_\bullet$, $f$ y $\mathscr L$ (o simplemente el
\emph{polinomio de Hilbert en $y$ de $\mathscr P_\bullet$}, o \emph{de} $f$, si no
resulta confusión); cuando $Y$ es conexo no vacío, se suprime la
mención de $y$ en la notación y la terminología. El invariante así obtenido
desempeñará un papel esencial en el cap. V, en la teoría de los «módulos» de los
haces coherentes cocientes de un haz coherente dado.

\begin{env}[7.9.12]
\label{III.7.9.12.fr}
Con las notaciones de (7.9.6) y (7.9.11), se tiene
\[
  \operatorname{PH}(\mathscr P_\bullet;y)
  =\operatorname{PH}(\mathscr P'_\bullet;y)
  +\operatorname{PH}(\mathscr P''_\bullet;y)
  \tag{7.9.12.1}\label{III.7.9.12.1.fr}
\]
y, en particular,
\[
  \operatorname{PH}(\mathscr P_\bullet;y)
  =\sum_i(\setmin 1)^i\operatorname{PH}(\mathscr P_i;y) ;
  \tag{7.9.12.2}\label{III.7.9.12.2.fr}
\]
Estas fórmulas se deducen inmediatamente de
\eqref{III.7.9.6.1} y \eqref{III.7.9.6.2}. Asimismo, con las
notaciones e hipótesis de \sref{III.7.9.7}, se tiene
\[
  \operatorname{PH}(\mathscr P'_\bullet;y')
  =\operatorname{PH}(\mathscr P_\bullet;g(y')).
  \tag{7.9.12.3}\label{III.7.9.12.3.fr}
\]
\end{env}

La fórmula \eqref{III.7.9.12.2} reduce el estudio de los
polinomios de Hilbert de un complejo al de los polinomios de
Hilbert de un solo $\mathscr O_X$-módulo plano sobre $Y$.
Estos últimos admiten la siguiente interpretación notable,
independiente de consideraciones homológicas.

\begin{corollary}[7.9.13]
\label{III.7.9.13}
\label{III.7.9.13.fr}
Sean $Y$ un preesquema noetheriano, $f:X\to Y$ un morfismo propio,
$\mathscr L$ un $\mathscr O_X$-módulo inversible muy amplio sobre
$Y$ y $\mathscr F$ un $\mathscr O_X$-módulo coherente y plano
sobre $Y$. Existe un entero $n_0$ tal que, para $n\geq n_0$,
$f_*(\mathscr F(n))$ es un $\mathscr O_Y$-módulo localmente libre
cuyo rango en $y\in Y$ es
$\operatorname{PH}(f,\mathscr F;y)(n)$.
\end{corollary}

\begin{proof}
Como $f$ es proyectivo \sref[II]{II.5.5.3}, existe un entero $n_0$
tal que, para $n\geq n_0$
\oldpage[III]{212}
se tiene $\mathrm R^if_*(\mathscr F(n))=0$ para todo $i>0$ (2.2.1); la conclusión
resulta pues de (7.9.10).
\end{proof}

El criterio de planitud siguiente será importante en la teoría de los «módulos»
de los haces coherentes del cap. V:

\begin{proposition}[7.9.14]
\label{III.7.9.14}
\label{III.7.9.14.fr}
Sean $Y$ un preesquema noetheriano, $f:X\to Y$ un morfismo
proyectivo y $\mathscr L$ un $\mathscr O_X$-módulo inversible
amplio para $f$, y pongamos
$\mathscr F(n)=\mathscr F\otimes\mathscr L^{\otimes n}$ para todo
$\mathscr O_X$-módulo $\mathscr F$ y
todo $n\in\mathbf Z$. Para que un
$\mathscr O_X$-módulo coherente $\mathscr F$ sea plano
sobre $Y$, es necesario y suficiente que exista un entero $n_0$
tal que, para todo $n\geq n_0$, $f_*(\mathscr F(n))$ sea un
$\mathscr O_Y$-Módulo localmente libre.
\end{proposition}

\begin{proof}
La necesidad de la condición se demuestra como en (7.9.13) (aplicándose el resultado de
(2.2.1) a un haz amplio $\mathscr L$, puesto que $f$ es
proyectivo). Para demostrar el recíproco, se puede limitar al caso en que $Y$ es
afín de anillo $A$; en virtud de la hipótesis y de (2.2.2, (i)), los
$A$-módulos
$\Gamma(X,\mathscr F(n))$ son de tipo finito y proyectivos (Bourbaki,
\emph{Alg. comm.}, cap. II, \S~5, n\textsuperscript{o}~2, teorema 1). Sea $S$
el anillo graduado $\bigoplus_{n\geq0}\Gamma(X,\mathscr L^{\otimes n})$;
se sabe que $X$ se identifica
canónicamente con $\operatorname{Proj}(S)$ por
\sref[II]{II.4.5.2}{\emph{(b)}} y \sref[II]{II.5.4.4}. Sea
$M=\bigoplus_{n\geq n_0}\Gamma(X,\mathscr F(n))$;
Sustituyendo, si es preciso, $\mathscr L$ por una potencia
$\mathscr L^{\otimes d}$, se puede suponer que $S$ está generado por un número finito de elementos de grado 1
(2.3.5.1), y resulta entonces de (\textbf{II}, 2.7.5 y 2.7.2) que $\mathscr F$ se
identifica con $\mathscr P\!roj_0(M)$.

Para todo elemento homogéneo $g\in S$ de
grado $>0$, se tiene pues $\Gamma(X_g,\mathscr F)=M_{(g)}$;
Ahora bien, $M$, al ser una suma directa de $A$-módulos
proyectivos, es un $A$-módulo plano. Por consiguiente, $M_g$ ($\mathbf{0}_{\mathrm I}$, 6.3.2), y por tanto también de $M_{(g)}$, que es un
factor directo de $M_g$ ($\mathbf{0}_{\mathrm I}$, 6.1.2). Se concluye (I, 4.14.5) que
$\mathscr F$ es plano sobre $Y$ en todo
punto de $X_g$. Como los abiertos $X_g$ recubren $X$, queda
demostrada la proposición.
\end{proof}

\bigskip
\noindent\emph{(Continuará.)}

\clearpage
\oldpage[III]{213}

\section*{ÍNDICE DE NOTACIONES}

\begin{flushleft}
$\mathbf{Tor}_n^A(P_\bullet,Q_\bullet)$: 6.3.1.

\medskip
$\boldsymbol{\mathscr{T}\!or}_n^{\mathscr O_S}(\mathscr P_\bullet,\mathscr Q_\bullet)$,
$\boldsymbol{\mathscr{T}\!or}_n^S(\mathscr P_\bullet,\mathscr Q_\bullet)$: 6.4.1 y 6.5.1.

\medskip
$\mathscr{T}\!or_n^S(\mathscr F,\mathscr G)$: 6.5.1.

\medskip
$\boldsymbol{\mathscr{T}\!or}_q^S(\mathscr P_\bullet^{(1)},\mathscr P_\bullet^{(2)},\ldots,
\mathscr P_\bullet^{(m)})$: 6.5.15.

\medskip
$\mathbf H^{-n}(\mathfrak U^{(i)},\mathscr P_\bullet^{(i)})$,
$\mathbf H^{-n}(X^{(i)},\mathscr P_\bullet^{(i)})$: 6.6.1.

\medskip
$\mathbf{Tor}_n^A(L_{\bullet\bullet}^{(1)},L_{\bullet\bullet}^{(2)})$,
$\mathbf{Tor}_n^S(\mathfrak U^{(1)},\mathfrak U^{(2)};
\mathscr P_\bullet^{(1)},\mathscr P_\bullet^{(2)})$: 6.6.2.

\medskip
$\mathscr{T}\!or_n^S(\mathfrak U^{(1)},\mathfrak U^{(2)};
\mathscr F^{(1)},\mathscr F^{(2)})$,
$\mathscr{T}\!or_n^S(\mathfrak U^{(1)},\mathfrak U^{(2)};
\mathscr P_\bullet^{(1)},\mathscr P_\bullet^{(2)})$: 6.6.2.

\medskip
$\boldsymbol{\mathscr{T}\!or}_n^S(f_1,f_2;\mathscr P_\bullet^{(1)},
\mathscr P_\bullet^{(2)})$: 6.6.8, 6.7.1 y 6.7.3.

\medskip
$\boldsymbol{\mathscr{T}\!or}_n^S((f_i)_{i\in I};(\mathscr P_\bullet^{(i)})_{i\in I})$,
$\boldsymbol{\mathscr{T}\!or}_n^S(f_1,\ldots,f_m;\mathscr P_\bullet^{(1)},\ldots,
\mathscr P_\bullet^{(m)})$: 6.7.11.

\medskip
$\boldsymbol{Ab}$, $\boldsymbol{Ab}_A$: 7.1.1.

\medskip
$T_{(B)}$, $T^{(B)}$, $T\otimes_A B$
($T$ funtor de $\boldsymbol{Ab}_A$ en $\boldsymbol{Ab}$): 7.1.3.

\medskip
$t_M$: 7.2.2.

\medskip
$T_{\mathfrak p}$: 7.1.4.

\medskip
$\mathscr T_\bullet^{Y'}(\mathscr M')$,
$\mathscr T_\bullet^U(\mathscr M|U)$: 7.7.2.

\medskip
$\operatorname{EP}(f,\mathscr P_\bullet;y)$,
$\operatorname{EP}(Y,\mathscr P_\bullet;y)$,
$\operatorname{EP}(\mathscr P_\bullet;y)$: 7.9.5.

\medskip
$\operatorname{PH}(f,\mathscr P_\bullet;y)$,
$\operatorname{PH}(\mathscr P_\bullet;y)$: 7.9.12.
\end{flushleft}

\clearpage
\oldpage[III]{214}
\section*{ÍNDICE TERMINOLÓGICO}

\begin{flushleft}
Característica de Euler-Poincaré de un complejo de módulos: 7.9.5.

\medskip
Cohomológicamente plano en un punto $y\in Y$ en dimensión $p$,
cohomológicamente plano sobre $Y$ en dimensión $p$, cohomológicamente plano
sobre $Y$ en el punto $y$, cohomológicamente plano sobre $Y$ (complejo de
$\mathscr O_X$-módulos planos sobre $Y$): 7.8.1.

\medskip
Complejos homotópicamente equivalentes: 6.1.4.

\medskip
Espacio anillado de dimensión cohomológica $\leq n$: 6.5.5.

\medskip
Extensión de escalares en un funtor covariante aditivo: 7.1.3.

\medskip
Haz de anillos de dimensión cohomológica $\leq n$: 6.5.5.

\medskip
Fórmula de Künneth: 6.7.8.

\medskip
Homológicamente plano en un punto $y\in Y$ en dimensión $p$,
homológicamente plano sobre $Y$ en dimensión $p$, homológicamente plano sobre
$Y$ en el punto $y$, homológicamente plano sobre $Y$ (complejo de
$\mathscr O_X$-módulos planos sobre $Y$): 7.8.1.

\medskip
Equivalencia homotópica: 6.1.4.

\medskip
Hipertor de dos complejos de $A$-módulos: 6.3.1.

\medskip
Hipertor local de dos complejos de módulos: 6.4.1.

\medskip
Hipertor global de dos complejos de módulos: 6.6.2, 6.6.8 y 6.7.3.

\medskip
Polinomio de Hilbert relativo a un complejo de módulos: 7.9.12.

\medskip
Sucesiones espectrales de Künneth: 6.7.3.
\end{flushleft}

\clearpage
\oldpage[III]{215}
\section*{TABLA DE MATERIAS}

\begin{flushleft}
\textsc{Capítulo III.} --- \textbf{Estudio cohomológico de los haces coherentes}
\textit{(continuación)} \dotfill 5

\medskip
\S~6. Funtores «Tor» locales y globales; fórmula de Künneth \dotfill 5

6.1. Introducción \dotfill 5

6.2. Hipercohomología de los complejos de módulos sobre un preesquema \dotfill 6

6.3. Hipertor de dos complejos de módulos \dotfill 9

6.4. Funtores hipertor locales de complejos de módulos cuasicoherentes: caso
de los esquemas afines \dotfill 14

6.5. Funtores hipertor locales de complejos de módulos cuasicoherentes: caso
general \dotfill 16

6.6. Funtores de hipertor global para complejos de módulos cuasicoherentes y
sucesiones espectrales de Künneth: caso de base afín \dotfill 21

6.7. Funtores de hipertor global de complejos de módulos cuasicoherentes y
sucesiones espectrales de Künneth: caso general \dotfill 25

6.8. Sucesiones espectrales de asociatividad de los hipertor globales \dotfill 32

6.9. Las sucesiones espectrales de cambio de base en los hipertor globales
\dotfill 34

6.10. Estructura local de ciertos funtores cohomológicos \dotfill 39

\medskip
\S~7. Estudio del cambio de base en los funtores homológicos covariantes de
módulos \dotfill 43

7.1. Funtores de $A$-módulos \dotfill 43

7.2. Caracterización del funtor producto tensorial \dotfill 44

7.3. Criterios de exactitud de los funtores homológicos de módulos \dotfill 48

7.4. Criterios de exactitud para los funtores
$\mathrm H_\bullet(P_\bullet\otimes_A M)$ \dotfill 53

7.5. Caso de los anillos locales noetherianos \dotfill 58

7.6. Descenso de las propiedades de exactitud. Teorema de semicontinuidad y
criterio de exactitud de Grauert \dotfill 60

7.7. Aplicación a los morfismos propios: I. La propiedad de cambio de base
\dotfill 65

7.8. Aplicación a los morfismos propios: II. Criterios de planitud
cohomológica \dotfill 72

7.9. Aplicación a los morfismos propios: III. Invariancia de la
característica de Euler-Poincaré y del polinomio de Hilbert \dotfill 76

\begin{flushright}
\emph{Recibido el 15 de diciembre de 1962.}
\end{flushright}
\end{flushleft}

% La página impresa 216 está en blanco.
\clearpage
\oldpage[III]{216}
\thispagestyle{empty}
\null
\clearpage

\oldpage[III]{217}
\section*{ERRATAS Y ADENDAS}
\begin{center}
(Lista 2)

\medskip
A) \emph{Errores tipográficos}
\end{center}

\noindent\textbf{($\mathbf{0}_{\mathrm I}$, 3.2.1)} En la línea 2 de la p. 27, sustituir
$X$ por $U$.

\noindent\textbf{($\mathbf{0}_{\mathrm I}$, 3.2.6)} En la primera línea contando desde
abajo de la p. 27 y en la línea 1 de la p. 28, sustituir (3 veces)
$\mathfrak H_\lambda$ por $\mathscr H_\lambda$.

\noindent\textbf{($\mathbf{0}_{\mathrm I}$, 3.4.5)} Añadir un paréntesis delante de
3.4.5).

\noindent\textbf{($\mathbf{0}_{\mathrm I}$, 4.2.1)} En la línea 10 de la p. 39,
sustituir $\psi_*(A)$ por $\psi_*(\mathscr A)$.

\noindent\textbf{($\mathbf{0}_{\mathrm I}$, 5.1.3)} En la línea 16 de la p. 45,
sustituir $\mathscr F|V$ por $\mathscr F|U$.

\noindent\textbf{($\mathbf{0}_{\mathrm I}$, 5.3.9)} En la línea 12 de la p. 48, antes
de «tal que», añadir: sobre una vecindad abierta $U$ de $x$.

\noindent\textbf{($\mathbf{0}_{\mathrm I}$, 7.6.9)} En la tercera línea contando desde
abajo de la p. 73, sustituir $J_\lambda$ por $\mathfrak J_\lambda$.

\noindent\textbf{(I, 1.1.15)} En la línea 3 de la p. 83, sustituir
$\mathfrak a$ por $\mathfrak a\ne A$.

\noindent\textbf{(I, 1.4.1)} En la segunda línea contando desde abajo de la
p. 90, sustituir $\widetilde N$ por $N$.

\noindent\textbf{(I, 2.3.1)} En la tercera línea contando desde abajo de la
p. 101, sustituir (\textbf 0, 4.1.6) por (\textbf 0, 4.1.7).

\noindent\textbf{(I, 2.3.2)} En la línea 11 de la p. 101, sustituir $B$ por
$B_s$ y $C$ por $C_t$.

\noindent\textbf{(I, 2.5.5)} En la línea 19 de la p. 104, sustituir
$S$-morfismo por $S$-preesquema.

\noindent\textbf{(I, 3.3.7)} En la línea 17 de la p. 109, sustituir
$Y_{(S)}$ por $Y_{(S')}$.

\noindent\textbf{(I, 3.4.5)} En la línea 9 de la p. 113, sustituir
$\mathrm s$ por $s$.

\noindent\textbf{(I, 3.4.8)} En la línea 4 de la p. 114, sustituir $p(x)$ por
$f(x)$.

\noindent\textbf{(I, 3.5.1)} En la línea 3 de la p. 115, sustituir
$1_X\times g$ por $1_{X'}\times g$.

\noindent\textbf{(I, 3.7.2)} En la línea 8 de la p. 119, sustituir el primer
$X$ por $X'$.

\noindent\textbf{(I, 4.2.2)} En la línea 8 de la p. 123, en la flecha vertical
izquierda del diagrama, sustituir $\alpha_{\psi(y)}$ por $\rho_{\psi(y)}$.

\noindent\textbf{(I, 4.4.3)} En la línea 16 de la p. 126, sustituir «isomorfa»
por «isomorfas».

\noindent\textbf{(I, 4.5.5)} En la línea 17 contando desde abajo de la p. 127,
sustituir (4.2.4) por (4.2.5).

\noindent\textbf{(I, 5.1.4)} En la línea 14 contando desde abajo de la p. 128,
sustituir (2.1.7) por (2.1.8).

\noindent\textbf{(I, 5.3.13)} En la línea 20 de la p. 134, sustituir (4.2.4)
por (4.2.5).

\noindent\textbf{(I, 6.4.2)} En la primera línea contando desde abajo de la
p. 147, sustituir recubrimiento por recubrimiento finito.

\noindent\textbf{(I, 9.1.13)} En la línea 15 de la p. 171, sustituir
$p^{-1}(\mathscr F)$ por $p^*(\mathscr F)$.

\noindent\textbf{(I, 9.5.11)} En la línea 7 de la p. 179, sustituir $Y$ por
$Y'$.

\noindent\textbf{(I, 9.6.5)} En la línea 17 contando desde abajo de la p. 180,
sustituir (\textbf{0}, 4.1.4) por (\textbf{0}, 4.1.3).

\noindent\textbf{(I, 10.12.2)} En la línea 14 contando desde abajo de la p.
206, sustituir $(X_n)_{(S_n)}$ por $(X_n)_{(S_m)}$.

\clearpage
\oldpage[III]{218}
\noindent\textbf{(I), Índice terminológico:} En la línea 6 contando desde
abajo de la p. 219, sustituir \textbf{0}, 4.1.4 por \textbf{0}, 4.1.5. En las líneas 16, 17 y 18
contando desde abajo de la p. 222, sustituir \textbf{I}, 2.1.7 por \textbf{I}, 2.1.8; en la línea
19 contando desde abajo de la p. 222, sustituir \textbf{I}, 2.1.8 por \textbf{I}, 2.1.7.

\noindent\textbf{(II, 1.5.2)} En la línea 17 de la p. 12, insertar una $f$
junto a la flecha vertical izquierda del diagrama.

\noindent\textbf{(II, 4.2.7)} En la línea 16 contando desde abajo de la p. 75,
sustituir (III, 2.1.14) por (\textbf{III}, 2.1.13).

\noindent\textbf{(II, 5.2.1)} En la línea 16 contando desde abajo de la p. 97,
sustituir $X\setmin U$ por $U$. En la línea 8 contando desde abajo de la p. 97,
sustituir $X^f$ por $X_f$.

\noindent\textbf{(II, 5.3.6)} En la línea 16 de la p. 100, sustituir (4.6.18)
por (4.6.17).

\noindent\textbf{(II, 6.2.7)} En la línea 9 contando desde abajo de la p. 116,
sustituir cap. V por cap. IV.

\noindent\textbf{(II, 6.6.5)} En la línea 15 contando desde abajo de la p.
132, sustituir cap. V por cap. IV.

\noindent\textbf{(II, 7.4.12)} En la línea 1 de la p. 151, sustituir cap. V
por cap. III.

\noindent\textbf{(II, 8.1.4)} En la línea 15 de la p. 154, sustituir (\textbf{III},
2.3.8) por (\textbf{III}, 2.3.7).

\noindent\textbf{(II, 8.10.5)} En la línea 17 contando desde abajo de la p.
187, sustituir $g_{(j)x}$ por $g_{j(x)}$.

\noindent\textbf{($\mathbf{0}_{\mathrm{III}}$, 10.2.7)} En la línea 1 de la p. 20,
sustituir $\mathfrak u$-ádico por $\mathfrak n$-ádico. En la línea 13 contando desde abajo
de la p. 20, sustituir $k$ por $k_i$.

\noindent\textbf{($\mathbf{0}_{\mathrm{III}}$, 11.1.1)} En la línea 18 de la p. 23,
sustituir $\operatorname{Im}(d_r^{p+r,q-r+1})$ por
$\operatorname{Im}(d_r^{p-r,q+r-1})$.

\noindent\textbf{($\mathbf{0}_{\mathrm{III}}$, 11.8.6)} En la línea 14 contando desde
abajo de la p. 45, sustituir las flechas $\to$ por $\rightsquigarrow$ (4 veces).

\noindent\textbf{($\mathbf{0}_{\mathrm{III}}$, 12.3.4)} En la línea 4 de la p. 61,
sustituir $\xrightarrow{\sim}$ por $\to$.

\noindent\textbf{($\mathbf{0}_{\mathrm{III}}$, 13.2.4)} En la segunda línea contando
desde abajo de la p. 67, sustituir $\varinjlim$ por $\varprojlim$ (2 veces).

\noindent\textbf{(III, 1.1.7)} En la línea 15 de la p. 85, sustituir $A^r$ por
$\wedge(A^r)$ (2 veces).

\noindent\textbf{(III, 1.2.4)} En la línea 7 de la p. 87, sustituir
$\Gamma(U,\mathscr F)$ por $\Gamma(U,\mathscr O_X)$ y sustituir $\mathscr G$
por $\mathscr F$. En la línea 8 de la p. 87, sustituir
$\mathrm H^p(\mathfrak U,\mathscr G)$ por $\mathrm H^p(\mathfrak U,\mathscr F)$.

\noindent\textbf{(III, 2.5.3.1)} En la línea 3 de la p. 111, sustituir $n>0$
por $n\geq0$. En la línea 8 de la p. 111, sustituir $\setmin r\leq n\leq0$ por
$\setmin r\leq n<0$.

\noindent\textbf{(III, 3.1.2)} En la línea 13 de la p. 116, sustituir
$\mathscr G$ por $\mathscr G\in\boldsymbol{K}'$.

\noindent\textbf{(III, 4.3.6)} En la línea 7 de la p. 133, sustituir
$K'=K\otimes_A\widehat A$ por $K'\supset K\otimes_A\widehat A$. En la línea 16
de la p. 133, sustituir $R(X)$ por $R(Y)$.

\noindent\textbf{(III, 4.4.12)} En la línea 10 de la p. 138, sustituir cap. V
por cap. IV.

\noindent\textbf{(III, 4.6.6)} En la línea 19 de la p. 142, sustituir «cap. V»
por «en un parágrafo posterior».

\begin{center}
B) \emph{Modificaciones del texto}\footnote{Para facilitar la búsqueda de las
referencias, las modificaciones pertenecientes a la lista de erratas insertada
en el capítulo N serán designadas en lo sucesivo por $\mathbf{Err}_{\mathrm N}$
seguido de un número.}
\end{center}

\noindent\textbf{($\mathbf{Err}_{\mathrm{III}}$, 1)} En
($\mathbf{0}_{\mathrm I}$, 5.1.3), en la línea 18 de la p. 45, sustituir «suma directa»
por «suma directa finita».

\noindent\textbf{($\mathbf{Err}_{\mathrm{III}}$, 2)} En
($\mathbf{0}_{\mathrm I}$, 5.4.3), en las líneas 8 a 4 contando desde abajo de la p. 49,
sustituir desde «En el caso general\ldots{}» por el texto siguiente:

En el caso general en que $X$ es un espacio anillado tal que $\mathscr O_x$ es
un anillo \emph{local} para
\oldpage[III]{219}
todo $x\in X$, si $\mathscr L$ es un $\mathscr O_X$-módulo \emph{de tipo finito}
tal que existe un $\mathscr O_X$-módulo $\mathscr F$ para el cual
$\mathscr L\otimes_{\mathscr O_X}\mathscr F$ es \emph{isomorfo a}
$\mathscr O_X$, el razonamiento anterior muestra que, para todo $x\in X$,
$\mathscr L_x$ es isomorfo a $\mathscr O_x$. De ello se deduce que $\mathscr L$
es entonces \emph{invertible}. En efecto, para todo $x\in X$, sea $U$ una
vecindad abierta de $x$ tal que $\mathscr L|U$ esté engendrado por $n$ secciones
$s_i\ (1\leq i\leq n)$ sobre $U$ (5.2.1). Se puede suponer, por ejemplo, que
$s=s_1$ es tal que $s_x\ne0$ y, como $\mathscr L_x$ es isomorfo a
$\mathscr O_x$, existe para cada $i$ una sección $t_i$ de $\mathscr O_X$ sobre
una vecindad abierta $V_i\subset U$ de $x$ tal que
$(t_i)_x s_x=(s_i)_x$. Por consiguiente, existe una vecindad abierta
$V\subset\bigcap_{i=1}^nV_i$ de $x$ tal que $t_i s=s_i$ en $V$; dicho de otro
modo, $\mathscr L|V$ está engendrado por \emph{una sola} sección $s|V$. Además,
si $z$ es una sección, sobre un abierto $W\subset V$, del núcleo del
homomorfismo $\mathscr O_X|V\to\mathscr L|V$ definido por $s$ (5.1.1), $z_y$
anula $\mathscr L_y$ para todo $y\in W$, luego $z_y=0$ por hipótesis y, por
consiguiente, $z=0$, lo que termina de probar nuestra afirmación. Además, la
consideración del producto tensorial
$\mathscr L^{-1}\otimes\mathscr L\otimes\mathscr F$ muestra de inmediato que
$\mathscr F$ es isomorfo a $\mathscr L^{-1}$.

\noindent\textbf{($\mathbf{Err}_{\mathrm{III}}$, 3)} En
($\mathbf{0}_{\mathrm I}$, 7.2.4), es preciso imponer la condición de que las potencias
$\mathfrak J^n$ sean ideales \emph{cerrados} en el anillo admisible $A$ para
que la proposición sea correcta. La demostración dada es incorrecta, y el
enunciado rectificado resulta de Bourbaki, \emph{Top. gén.}, cap. III,
3.\textsuperscript{a} ed., \S~3, n.\textsuperscript{o}~5, cor. 1 de la prop. 9.
Asimismo, es preciso suponer que los $\mathfrak J^n$ son \emph{cerrados} en $A$
en los enunciados ($\mathbf{0}_{\mathrm I}$, 7.2.5) y ($\mathbf{0}_{\mathrm I}$, 7.2.6).

\noindent\textbf{($\mathbf{Err}_{\mathrm{III}}$, 4)} Después de la
proposición (\textbf I, 2.2.5), añadir: con las notaciones de (\textbf I, 2.2.5), si
$\mathfrak J$ es un ideal de $A$, se denota por
$\mathfrak J\mathscr F$ el sub-$\mathscr O_X$-módulo
$\widetilde{\mathfrak J}\mathscr F$ definido en (\textbf{0}, 4.3.5).

\noindent\textbf{($\mathbf{Err}_{\mathrm{III}}$, 5)} En (\textbf I, 3.4.5), en la
primera línea contando desde abajo de la p. 112, sustituir «\emph{un cuerpo}
$K$» por «\emph{un cuerpo algebraicamente cerrado} $K$». En la línea 1 de la
p. 113, sustituir «\emph{extensión}» por «\emph{extensión algebraicamente
cerrada}». En las líneas 1 a 4 de la p. 113, sustituir desde «$K$ será
llamado\ldots{}» por el texto siguiente:

Para todo punto de $X$ con valores en un cuerpo $K$, $K$ se llamará el
\emph{cuerpo de valores} del punto correspondiente y, si $x$ es la localidad
de este punto, se dice también que este último está \emph{localizado en} $x$.
Se define así una aplicación $X(K)\to X$ que hace corresponder a un punto con
valores en $K$ su localidad.

En las líneas 7 y 16 de la p. 113, suprimir «geométrico»; en las líneas 10 y 13
de la p. 113, suprimir «geométricos». En las líneas 10 y 11 contando desde
abajo de la p. 115, suprimir «\emph{geométrico}».

\noindent\textbf{($\mathbf{Err}_{\mathrm{III}}$, 6)} Al final de la
demostración de (\textbf I, 3.6.1), en la línea 12 contando desde abajo de la p. 117,
añadir: Si se pone
$X'=X\times_Y\operatorname{Spec}(\mathscr O_y/\mathfrak a_y)$, entonces, para
todo punto $x\in X'$, identificado por $p$ con un punto de $X$, se tiene
$\mathscr O_{X',x}=\mathscr O_{X,x}/\mathfrak a_y\mathscr O_{X,x}$. Puesto que
la cuestión es local en $X$ y $Y$, se puede suponer que
$X=\operatorname{Spec}(B)$, $Y=\operatorname{Spec}(A)$, y se tiene
$(B/\mathfrak a_yB)_x=B_x/\mathfrak a_yB_x$ por planitud (\textbf{0}, 1.3.2).

\noindent\textbf{($\mathbf{Err}_{\mathrm{III}}$, 7)} En (\textbf I, 5.1.1), en la
tercera línea contando desde abajo de la p. 127, sustituir
«$\mathscr O_X$-\emph{módulo cuasicoherente}» por «\emph{ideal cuasicoherente de}
$\mathscr B$».

\noindent\textbf{($\mathbf{Err}_{\mathrm{III}}$, 8)} Después de (\textbf I, 5.1.10),
en la línea 16 de la p. 131, añadir:

\emph{Observación (5.1.11).} --- Una ligera adaptación del razonamiento hecho
en (5.1.9) muestra que la conclusión sigue siendo válida \emph{sin suponer a
priori que $X$ sea un preesquema}, sino

\oldpage[III]{220}
suponiendo solamente que $(X,\mathscr O_X)$ sea un espacio anillado en anillos
locales, que $\mathscr J$ sea un ideal de $\mathscr O_X$ tal que
$\mathscr J^n=0$, que el espacio anillado
$X_0=(X,\mathscr O_X/\mathscr J)$ sea un esquema afín y, por último, que los
ideales $\mathscr J^k/\mathscr J^{k+1}$ sean $\mathscr O_{X_0}$-módulos
cuasicoherentes. Basta, en efecto, utilizar (1.8.1) (cf.
\textbf{($\mathbf{Err}_{\mathrm{II}}$)}) en lugar de (2.2.4) en el razonamiento.

\noindent\textbf{($\mathbf{Err}_{\mathrm{III}}$, 9)} En (\textbf I, 5.3.1), en la
línea 11 de la p. 132, insertar, después de «o $\Delta_X$», «o $\Delta_\varphi$
si $\varphi:X\to S$ es el morfismo estructural».

\noindent\textbf{($\mathbf{Err}_{\mathrm{III}}$, 10)} En (\textbf I, 5.3.9), la
demostración es insuficiente, pues no prueba que $\Delta_X(X)$ sea localmente
cerrado en $X\times_SX$. Para obtener una demostración correcta, basta utilizar
(4.2.4, $a$)): para todo $x\in X$ y toda vecindad afín $U$ de $x$ en $X$,
$U\times_SU$ es una vecindad afín de $\Delta_X(x)$; teniendo en cuenta (5.3.16)
(cuya demostración solo utiliza la definición (5.3.1.1)), se reduce a demostrar
(5.3.9) cuando $S=\operatorname{Spec}(B)$ y $X=\operatorname{Spec}(A)$ son
esquemas afines; entonces resulta claro, en virtud de (5.3.1.1), que $\Delta_X$
corresponde al homomorfismo canónico $A\otimes_BA\to A$ que transforma
$x\otimes y$ en $xy$; como este homomorfismo es sobreyectivo, $\Delta_X$ es en
este caso una inmersión cerrada (4.2.3), lo que termina la demostración.

\noindent\textbf{($\mathbf{Err}_{\mathrm{III}}$, 11)} En (\textbf I, 7.1.15) y
(\textbf I, 7.1.16), en las líneas 7 y 15 contando desde abajo de la p. 158, suprimir
«\emph{geométricos}».

\noindent\textbf{($\mathbf{Err}_{\mathrm{III}}$, 12)} Sustituir (\textbf I, 7.4.7)
por: Se extienden (por abuso de lenguaje) las definiciones de (7.4.1) al caso
en que $X$ es un preesquema \emph{reducido} tal que todo punto admite una
vecindad abierta que solo tiene un número \emph{finito} de componentes
irreducibles; se deduce entonces de (7.3.4) y (7.4.6) que, para un
$\mathscr O_X$-módulo cuasicoherente de tipo finito $\mathscr F$, decir que
$\mathscr F$ es un \emph{haz de torsión} equivale a decir que
$\operatorname{Supp}(\mathscr F)$ \emph{no contiene ninguna componente
irreducible de $X$}.

\noindent\textbf{($\mathbf{Err}_{\mathrm{III}}$, 13)} En (\textbf I, 10.11.7), en
las líneas 17 a 23 de la p. 205, el final de la demostración desde «Queda por
probar\ldots{}» es innecesario en virtud de (10.11.6), pues los haces
considerados son coherentes por (\textbf{0}, 5.3.5).

\noindent\textbf{($\mathbf{Err}_{\mathrm{III}}$, 14)} En (\textbf{II}, 1.7.8),
después de la línea 10 contando desde abajo de la p. 16, añadir: cuando
$\mathscr E=\mathscr O_S^n$, se escribe también $\mathbf V_S^n$ en lugar de
$\mathbf V(\mathscr E)$; si además $S=\operatorname{Spec}(A)$ es afín, se
escribe $\mathbf V_A^n$ en lugar de $\mathbf V_S^n$; se tiene
$\mathbf V_A^n=\operatorname{Spec}(A[T_1,\ldots,T_n])$, donde los $T_i$ son
indeterminadas.

\noindent\textbf{($\mathbf{Err}_{\mathrm{III}}$, 15)} En (\textbf{II}, 1.7.10), en
la línea 10 de la p. 17, suprimir «\emph{geométricos}»; en las líneas 12 y
16-17 de la p. 17, suprimir «geométrico».

\noindent\textbf{($\mathbf{Err}_{\mathrm{III}}$, 16)} En (\textbf{II}, 1.7.12),
añadir: Se pone asimismo:
\[
  \mathbf V(\mathscr O_S^n)=\mathbf V_S^n=S[T_1,\ldots,T_n].
\]

\noindent\textbf{($\mathbf{Err}_{\mathrm{III}}$, 17)} En (\textbf{II}, 4.2.6), en la
línea 7 de la p. 75, suprimir «\emph{geométricos}»; en la línea 8 de la p. 75,
suprimir «\emph{geométrico}».

\noindent\textbf{($\mathbf{Err}_{\mathrm{III}}$, 18)} En (\textbf{II}, 4.6.13), en la
línea 7 de la p. 92, es preciso especificar el razonamiento, pues, con las
notaciones de (4.4.10), hay que probar aquí que la inmersión $\Gamma_f$ es
\emph{cuasicompacta}, a fin de aplicar (i \emph{bis}). En virtud de (4.6.4), para
probar que $\mathscr L$ es amplio para $f$, se puede limitar al caso
en que $Y$ es afín. Observemos, por otra parte, que en cada una de las hipótesis
de (v), $f$ es cuasicompacto (\textbf I, 6.6.4). Por otra parte, si $g$ es
\oldpage[III]{221}
separado, $\Gamma_f$ es una inmersión cerrada (\textbf I, 5.4.3), luego cuasicompacta
(\textbf I, 6.6.4). Si, por el contrario, $X$ es localmente noetheriano, como $Y$ es
afín y $f$ cuasicompacto, el espacio subyacente a $X$ es cuasicompacto, luego
noetheriano, y se puede aplicar de nuevo (\textbf I, 6.6.4) para probar que $\Gamma_f$
es cuasicompacta.

\noindent\textbf{($\mathbf{Err}_{\mathrm{III}}$, 19)} El enunciado de (\textbf{II},
5.2.2) es incorrecto, pues las condiciones $b)$, $c)$ y $c')$ no son locales
en $Y$ cuando no se sabe si, para un abierto $U$ de $Y$, un
$(\mathscr O_X|f^{-1}(U))$-módulo cuasicoherente es la restricción a
$f^{-1}(U)$ de un $\mathscr O_X$-módulo cuasicoherente. Por tanto, es preciso
sustituir en la línea 16 de la p. 98 «\emph{Sea $f:X\to Y$ un morfismo separado
y cuasicompacto}» por: «\emph{Sean $X$, $Y$ dos preesquemas tales que $X$ sea
un esquema o que el espacio subyacente a $X$ sea localmente noetheriano, y sea
$f:X\to Y$ un morfismo cuasicompacto.}» En la demostración, se observará que
las hipótesis implican que $f$ es \emph{separado} cuando se supone que $X$ es
un esquema, en virtud de (\textbf I, 5.5.5 y 5.5.8); por tanto, se puede aplicar
(\textbf I, 9.2.2) en ambos casos.

\noindent\textbf{($\mathbf{Err}_{\mathrm{III}}$, 20)} En (\textbf{II}, 6.2.3),
añadir, después de la línea 18 contando desde abajo de la p. 115:

\emph{Se dice que un morfismo $f:X\to Y$ es cuasifinito en un punto $x\in X$
si existen una vecindad abierta afín $V$ de $y=f(x)$ y una vecindad abierta afín
$U$ de $x$ tales que $f(U)\subset V$ y que el morfismo $U\to V$, restricción de
$f$, sea cuasifinito. Se dice que un morfismo $f:X\to Y$ es localmente
cuasifinito si es cuasifinito en todo punto de $X$.}

\noindent\textbf{($\mathbf{Err}_{\mathrm{III}}$, 21)} En (\textbf{II}, 6.4.3), la
demostración es incorrecta, pues los elementos de un sub-$A$-módulo de tipo
finito de $K$ no son necesariamente enteros sobre $A$. Sustituir las 9 últimas
líneas de la p. 121 por el texto siguiente: Como se puede suponer que $E$ es
sin torsión, luego fiel, $E$ es también un $A[u]$-módulo fiel ($A$ está
inmerso en el anillo de endomorfismos de $E$). Como $E$ es un $A$-módulo de
tipo finito, se deduce que $u$ es \emph{entero} sobre $A$ (Bourbaki,
\emph{Alg. comm.}, cap. V, \S~1, n.\textsuperscript{o}~1, lema 1). Se concluye
inmediatamente que los valores propios de $u\otimes1$ (en una clausura
algebraica de $K$) son elementos \emph{enteros} sobre $A$, y por tanto también
lo son los $\sigma_i(u)$.

\noindent\textbf{($\mathbf{Err}_{\mathrm{III}}$, 22)} En
($\mathbf{0}_{\mathrm{III}}$, 9.1.1), en la línea 18 contando desde abajo de la p. 12,
suprimir «abierta».

\noindent\textbf{($\mathbf{Err}_{\mathrm{III}}$, 23)} En
($\mathbf{0}_{\mathrm{III}}$, 11.7.3), en la línea 10 de la p. 43, después de
$\boldsymbol{C}''$, añadir: tal que, para todo objeto proyectivo $P$ de
$\boldsymbol{C}$ (resp. todo objeto proyectivo $P'$ de $\boldsymbol{C}'$),
el funtor $A'\rightsquigarrow T(P,A')$ (resp.
$A\rightsquigarrow T(A,P')$) sea exacto en $\boldsymbol{C}'$ (resp. en
$\boldsymbol{C}$).

\noindent\textbf{($\mathbf{Err}_{\mathrm{III}}$, 24)} El enunciado de
($\mathbf{0}_{\mathrm{III}}$, 13.7.7) es inexacto y debe modificarse como sigue: en la
línea 6 de la p. 78, suprimir «\emph{y}
$(\mathrm R^{n+1}T(A_k))_{k\in\mathbf Z}$»; en la línea 7 de la p. 78, sustituir
«$\mathrm R'^nT(\mathbf A)$ \emph{es un} $S$-\emph{módulo de tipo finito}» por
«$\mathrm R'^nT(\mathbf A)$ \emph{y} $\mathrm R'^{n+1}T(\mathbf A)$ \emph{son} $S$-\emph{módulos de tipo finito}»;
en las líneas 12-13 de la p. 78, suprimir «\emph{y} $p+q=n+1$». En la demostración,
en la línea 23 de la p. 78, suprimir «y para $n+1$».

\noindent\textbf{($\mathbf{Err}_{\mathrm{III}}$, 25)} En (\textbf{III}, 1.4.15), en
la línea 6 contando desde abajo de la p. 92, sustituir «\emph{de tipo finito}»
por «\emph{cuasicompacto}».

\noindent\textbf{($\mathbf{Err}_{\mathrm{III}}$, 26)} En (\textbf{III}, 2.2.4), en la
línea 1 de la p. 102, sustituir «\emph{que los soportes de $\mathscr F$ y de
$\mathscr H$ sean propios sobre $Y$}» por «\emph{que el soporte de
$\operatorname{Im}(\mathscr F\to\mathscr G)$ sea propio sobre $Y$}». En la
demostración, sustituir el texto de las líneas 10 a 16, desde «Dicho
esto\ldots{}», por:

\oldpage[III]{222}
Pongamos
$\mathscr G_1=\operatorname{Im}(\mathscr F\to\mathscr G)
=\operatorname{Ker}(\mathscr G\to\mathscr H)$, que es coherente. Como se tiene
la sucesión exacta
$0\to\mathscr G_1(n)\to\mathscr G(n)\to\mathscr H(n)$ y el funtor $f_*$ es
exacto por la izquierda, la sucesión
$0\to f_*(\mathscr G_1(n))\to f_*(\mathscr G(n))\to f_*(\mathscr H(n))$ es
exacta para todo $n$, y basta por tanto mostrar que, para $n$ suficientemente
grande, la sucesión $f_*(\mathscr F(n))\to f_*(\mathscr G_1(n))\to0$ es exacta.
Dicho de otro modo, se puede limitar al caso en que $\mathscr H=0$ y en que el
soporte de $\mathscr G$ es propio sobre $Y$, luego \emph{cerrado en $Z$}
(\textbf{II}, 5.4.10); por consiguiente, $\mathscr G'=i_*(\mathscr G)$ es un
$\mathscr O_Z$-módulo coherente tal que $\mathscr G'|X=\mathscr G$. Se sabe
(\textbf I, 9.4.3) que existe un $\mathscr O_Z$-módulo coherente $\mathscr F'$ tal que
$\mathscr F'|X=\mathscr F$. Además, como $X$ es abierto en $Z$ y
$\operatorname{Supp}(\mathscr G)\subset X$ es cerrado en $Z$, es inmediato que
se define un homomorfismo sobreyectivo $u':\mathscr F'\to\mathscr G'$ de haces
tal que $u'|X$ sea el homomorfismo sobreyectivo dado
$u:\mathscr F\to\mathscr G$, tomando $u'|U=0$ para todo abierto $U$ de $Z$ que
no interseque $\operatorname{Supp}(\mathscr G)$ y $u'|U=u|U$ para todo abierto
$U\subset X$. Dicho esto, se ve, como al comienzo de la demostración de
(2.2.2), que se puede limitar al caso en que $Y$ es afín, y se trata por tanto
de mostrar que, para $n$ suficientemente grande, el homomorfismo
$\Gamma(u):\Gamma(X,\mathscr F(n))\to\Gamma(X,\mathscr G(n))$ es sobreyectivo.
Ahora bien, se tiene el diagrama conmutativo
\[
\xymatrix@C=35mm@R=12mm{
\Gamma(Z,\mathscr F'(n)) \ar[r]^{\Gamma(u')} \ar[d]_v &
\Gamma(Z,\mathscr G'(n)) \ar[d]^w \\
\Gamma(X,\mathscr F(n)) \ar[r]_{\Gamma(u)} & \Gamma(X,\mathscr G(n))
}
\]
donde $v$ y $w$ son los homomorfismos de restricción. La definición de
$\mathscr G'$ muestra además que $w$ es \emph{biyectivo}; por otra parte, en
virtud de (2.2.3), $\Gamma(u')$ es \emph{sobreyectivo} para $n$ suficientemente
grande, luego lo mismo vale para $\Gamma(u)$.

\noindent\textbf{($\mathbf{Err}_{\mathrm{III}}$, 27)} En (\textbf{III}, 2.2.5),
después de la línea 7 contando desde abajo de la p. 102, añadir:

(iii) Bajo las hipótesis de (2.2.4) relativas a $X$, $Y$, $f$ y $\mathscr L$,
es inmediato que, si $\mathscr H$ es un $\mathscr O_X$-módulo coherente cuyo
soporte es \emph{propio} sobre $Y$, un razonamiento análogo al de (2.2.4)
muestra que existe un entero $N$ tal que, para $n\geq N$, se tiene
$\mathrm R^if_*(\mathscr H(n))=0$ para todo $i>0$. Se deduce, por la sucesión
exacta de cohomología, que, si $u:\mathscr F\to\mathscr G$ es un homomorfismo
de $\mathscr O_X$-módulos coherentes tal que $\operatorname{Ker}(u)$ y
$\operatorname{Coker}(u)$ tienen soportes \emph{propios} sobre $Y$, entonces
existe $N$ tal que, para $n\geq N$, el homomorfismo correspondiente
$\mathrm R^if_*(\mathscr F(n))\to\mathrm R^if_*(\mathscr G(n))$ sea
\emph{biyectivo} para todo $i>0$.

\noindent\textbf{($\mathbf{Err}_{\mathrm{III}}$, 28)} En (\textbf{III}, 4.4.9), en la
línea 17 contando desde abajo de la p. 137, sustituir «\emph{Sean $Y$ un
preesquema íntegro localmente noetheriano}» por «\emph{Sean $X$ e $Y$ dos
preesquemas íntegros localmente noetherianos}»; en la línea 16 contando desde
abajo de la p. 137, sustituir «\emph{de tipo finito}» por «\emph{localmente de
tipo finito}». La demostración permanece esencialmente inalterada, pues
$f^{-1}(y)$ es localmente de tipo finito sobre $\boldsymbol{k}(y)$ y, por tanto, sigue
siendo discreto.

\noindent\textbf{($\mathbf{Err}_{\mathrm{III}}$, 29)} En (\textbf I, 9.3.4), en la
línea 19 contando desde abajo de la p. 173, sustituir «\emph{noetheriano}» por
«\emph{cuasicompacto}»; en las líneas 18 y 19 contando desde abajo de la p.
173, sustituir «\emph{coherente}» por «\emph{cuasicoherente de tipo finito}».
En la demostración, se reduce al caso en que

\oldpage[III]{223}
$X=\operatorname{Spec}(A)$, $\mathscr F=\widetilde M$, donde $M$ es un
$A$-módulo de tipo finito, $\mathscr J=\widetilde{\mathfrak J}$, donde
$\mathfrak J$ es un ideal de tipo finito de $A$; el resto del razonamiento
permanece entonces inalterado.

\noindent\textbf{($\mathbf{Err}_{\mathrm{III}}$, 30)} En (\textbf I, 9.3.5),
sustituir las líneas 1 a 7 contando desde abajo de la p. 173 por el texto
siguiente:

\emph{Proposición (9.3.5). --- Sean $X$ un preesquema y $\mathscr F$ un
$\mathscr O_X$-módulo cuasicoherente de tipo finito. Entonces existen un
subpreesquema cerrado $Y$ de $X$, cuyo espacio subyacente es igual a
$\operatorname{Supp}(\mathscr F)$, y un $\mathscr O_Y$-módulo cuasicoherente de
tipo finito $\mathscr G$ tales que, si $j:Y\to X$ es la inyección canónica,
$\mathscr F$ sea isomorfo a $j_*(\mathscr G)$.}

Bastará mostrar que el ideal $\mathscr J$ de $\mathscr O_X$,
\emph{anulador} de $\mathscr F$, es \emph{cuasicoherente}; se tomará entonces
como $Y$ el subpreesquema cerrado de $X$ definido por $\mathscr J$ (\textbf I, 4.1.2),
y, como $\mathscr J\mathscr F=0$, $\mathscr F$ es un
$(\mathscr O_X/\mathscr J)$-módulo y se resolverá la cuestión tomando
$\mathscr G=j^*(\mathscr F)$. Para ver que $\mathscr J$ es cuasicoherente, se
puede (puesto que la cuestión es local) limitar al caso en que
$X=\operatorname{Spec}(A)$, $\mathscr F=\widetilde M$, donde $M$ es un
$A$-módulo engendrado por un número finito de elementos
$x_i\ (1\leq i\leq r)$; el ideal $\mathscr J$ es entonces la intersección de
los anuladores de los $x_i$. Pero el anulador de $x_i$ es el núcleo del
homomorfismo $\mathscr O_X\to\mathscr F$ correspondiente al homomorfismo
$s\leadsto sx_i$ de $A$ en $M$; es, por tanto, un ideal cuasicoherente (I,
4.1.1), y toda intersección finita de tales ideales es asimismo un ideal
cuasicoherente (\textbf I, 1.3.10).

\begin{center}\rule{25mm}{0.4pt}\end{center}


\end{document}
